% !TeX TXS-program:compile = txs:///arara
% arara: pdflatex: {shell: no, synctex: no, interaction: batchmode}
% arara: pdflatex: {shell: no, synctex: no, interaction: batchmode}

\documentclass{article}
\usepackage[margin=0.625in]{geometry}
\usepackage{materialdesign-icons}
\usepackage{xcolor}
\usepackage{longtable}
\usepackage{booktabs}
\usepackage{hyperref}
\usepackage{fancyvrb}
\usepackage{shortvrb}
\MakeShortVerb{\|}
\setlength{\parindent}{0pt}

\newenvironment{materialdesigncase}
{%
  \scriptsize%
	\begin{longtable}{cllr}%icon + name + macro + page
		\cmidrule[\heavyrulewidth]{1-4}% \toprule
		\bfseries Icon & \bfseries Name & \bfseries Direct command & \bfseries PDF page \\
		\cmidrule{1-4}
		\endhead
}%
{%
		\cmidrule[\heavyrulewidth]{1-4}
	\end{longtable}%
}

\NewDocumentCommand{\showcasematerialdesignoutlined}{mm}{%name + macro + page
	\textcolor{violet}{\Large\materialdesignicon[outlined]{#1}} & \itshape #1 & \ttfamily \textbackslash materialdesignicon[outlined]\{#1\} & \sffamily #2 \\
}

\NewDocumentCommand{\showcasematerialdesignfilled}{mm}{%name + macro + page
	\textcolor{violet}{\Large\materialdesignicon[filled]{#1}} & \itshape #1 & \ttfamily \textbackslash materialdesignicon[filled]\{#1\} & \sffamily #2 \\
}

\NewDocumentCommand{\showcasematerialdesignround}{mm}{%name + macro + page
  \textcolor{violet}{\Large\materialdesignicon[round]{#1}} & \itshape #1 & \ttfamily \textbackslash materialdesignicon[round]\{#1\} & \sffamily #2 \\
}

\NewDocumentCommand{\showcasematerialdesigntwotone}{mm}{%name + macro + page
  \textcolor{violet}{\Large\materialdesignicon[twotone]{#1}} & \itshape #1 & \ttfamily \textbackslash materialdesignicon[twotone]\{#1\} & \sffamily #2 \\
}

\NewDocumentCommand{\showcasematerialdesignsharp}{mm}{%name + macro + page
  \textcolor{violet}{\Large\materialdesignicon[sharp]{#1}} & \itshape #1 & \ttfamily \textbackslash materialdesignicon[sharp]\{#1\} & \sffamily #2 \\
}

\begin{document}

\title{The \textsf{materialdesign-icons} package (0.14.15) -- 07/07/2026}
\author{%
	Cédric Pierquet\\%
	\url{https://github.com/marella/material-design-icons} (Source version 0.14.15)\\%
	\url{https://github.com/cpierquet/latex-packages/tree/main/materialdesign} (\LaTeX{} package)
}
\date{cpierquet -- at -- outlook . fr}\maketitle

This package provides \LaTeX{} support for the Material Design Icons project :

\url{https://github.com/marella/material-design-icons}.

\medskip

\hspace*{5mm}\fbox{\begin{minipage}{0.75\linewidth}{\emph{A set of 2122 free Apache-2.0 licensed high-quality SVG/PDF icons for you to use in your web projects.}}\end{minipage}}

\medskip

To use \textsf{materialdesign-icons} in your document, include the package with |\usepackage{materialdesign-icons}|.

An icon can be accessed using the icon name. A list of all included icons with their respective \textit{page} can be found at the end of this document.

\medskip

The individual icons are provided as pages of a PDF and are ultimately loaded using the |\includegraphics| command from the |graphicx| package.

\section{Example}

\begin{Verbatim}[frame=single]
...
\usepackage{materialdesign}
...
\begin{document}
...
Inline\materialdesign{24mp}version,
and colored\textcolor{red}{\materialdesign{24mp}}
...
\end{document}
\end{Verbatim}

Result: Inline\materialdesignicon{24mp}version, and colored\textcolor{red}{\materialdesignicon{24mp}}

\pagebreak

\section{Usage}

\subsection{Generic macro}

\begin{Verbatim}[frame=single]
\materialdesignicon%
  [%
    filled=TF,outlined=TF,round=TF,sharp=TF,twotone=TF,%
    height=...,(d)strut=...%
  ]%
  <includegraphics options>{name}
\end{Verbatim}

The boolean \texttt{[filled/outlined/round/sharp/twotone]} activate different version of icon. And key \texttt{[height=...]} can be:

\begin{itemize}
	\item \texttt{dauto} (by default): inline insertion with height/depth;
	\item \texttt{auto}: inline insertion with height/nodepth;
	\item \texttt{a length}: normal insertion with fixed height;
	\item \texttt{manuel height/depth}: normal insertion with fixed height/depth.
\end{itemize}

The keys \texttt{[(d)strut=...]} can be used to set the \textit{reference} box for dimension calculations.

\medskip

By default, icons are given with a small border, so the given height is the height of the \textit{full} box, not the hight of the icons.

\medskip

Example: Inline{\setlength\fboxsep{0pt}\fbox{\materialdesignicon{24mp}}}version or inline {\setlength\fboxsep{0pt}\fbox{\materialdesignicon{24mp}}} version

\pagebreak

\subsection{Examples}

\begin{Verbatim}[frame=single]
%normal version: inline with automatic height/depth (from qX characters)
\materialdesignicon[height=dauto]{24mp} %or \materialdesignicon{24mp}
\end{Verbatim}

{\Huge q\materialdesignicon{24mp}X}

\begin{Verbatim}[frame=single]
%normal version: inline with automatic height/depth (from specified characters)
(q\materialdesignicon[height=dauto,dstrut=(qB)]{24mp}B)
\end{Verbatim}

{\Huge (q\materialdesignicon[height=dauto,dstrut=(qB)]{24mp}B)}

\begin{Verbatim}[frame=single]
%normal version: inline with automatic height (from X character) default version
\materialdesignicon[height=auto]{24mp}
\end{Verbatim}

{\Huge X\materialdesignicon[height=auto]{24mp}X}

\begin{Verbatim}[frame=single]
%normal version: inline with manual height/depth
\materialdesignicon[height={0.8em/-0.4em}]{24mp}
\end{Verbatim}

{\Huge qb\materialdesignicon[height={0.8em/-0.4em}]{24mp}Xz}

\begin{Verbatim}[frame=single]
%fixed height
\materialdesignicon[height=2in]{24mp}
\end{Verbatim}

\materialdesignicon[height=1.5in]{24mp}

\subsection{Used packages}

\texttt{ifthen}, \texttt{calc}, \texttt{graphicx}, \texttt{xstring} and \texttt{simplekv} are loaded and used by the package.

\subsection{Bugs}

For bug reports and feature requests, report on the GitHub repository \url{https://github.com/cpierquet/latex-packages/issues}.

\newpage

\section{List of icons}

Icons are colored in \textcolor{violet}{violet}.

\subsection{Outlined version (2122 icons)}

The \textsf{PDF} file is \texttt{materialdesign-outlined-all.pdf}

\begin{materialdesigncase}
\showcasematerialdesignoutlined{1k}{1}
\showcasematerialdesignoutlined{1k-plus}{2}
\showcasematerialdesignoutlined{1x-mobiledata}{3}
\showcasematerialdesignoutlined{2k}{4}
\showcasematerialdesignoutlined{2k-plus}{5}
\showcasematerialdesignoutlined{2mp}{6}
\showcasematerialdesignoutlined{3d-rotation}{7}
\showcasematerialdesignoutlined{3g-mobiledata}{8}
\showcasematerialdesignoutlined{3k}{9}
\showcasematerialdesignoutlined{3k-plus}{10}
\showcasematerialdesignoutlined{3mp}{11}
\showcasematerialdesignoutlined{3p}{12}
\showcasematerialdesignoutlined{4g-mobiledata}{13}
\showcasematerialdesignoutlined{4g-plus-mobiledata}{14}
\showcasematerialdesignoutlined{4k}{15}
\showcasematerialdesignoutlined{4k-plus}{16}
\showcasematerialdesignoutlined{4mp}{17}
\showcasematerialdesignoutlined{5g}{18}
\showcasematerialdesignoutlined{5k}{19}
\showcasematerialdesignoutlined{5k-plus}{20}
\showcasematerialdesignoutlined{5mp}{21}
\showcasematerialdesignoutlined{6k}{22}
\showcasematerialdesignoutlined{6k-plus}{23}
\showcasematerialdesignoutlined{6mp}{24}
\showcasematerialdesignoutlined{6-ft-apart}{25}
\showcasematerialdesignoutlined{7k}{26}
\showcasematerialdesignoutlined{7k-plus}{27}
\showcasematerialdesignoutlined{7mp}{28}
\showcasematerialdesignoutlined{8k}{29}
\showcasematerialdesignoutlined{8k-plus}{30}
\showcasematerialdesignoutlined{8mp}{31}
\showcasematerialdesignoutlined{9k}{32}
\showcasematerialdesignoutlined{9k-plus}{33}
\showcasematerialdesignoutlined{9mp}{34}
\showcasematerialdesignoutlined{10k}{35}
\showcasematerialdesignoutlined{10mp}{36}
\showcasematerialdesignoutlined{11mp}{37}
\showcasematerialdesignoutlined{12mp}{38}
\showcasematerialdesignoutlined{13mp}{39}
\showcasematerialdesignoutlined{14mp}{40}
\showcasematerialdesignoutlined{15mp}{41}
\showcasematerialdesignoutlined{16mp}{42}
\showcasematerialdesignoutlined{17mp}{43}
\showcasematerialdesignoutlined{18mp}{44}
\showcasematerialdesignoutlined{18-up-rating}{45}
\showcasematerialdesignoutlined{19mp}{46}
\showcasematerialdesignoutlined{20mp}{47}
\showcasematerialdesignoutlined{21mp}{48}
\showcasematerialdesignoutlined{22mp}{49}
\showcasematerialdesignoutlined{23mp}{50}
\showcasematerialdesignoutlined{24mp}{51}
\showcasematerialdesignoutlined{30fps}{52}
\showcasematerialdesignoutlined{30fps-select}{53}
\showcasematerialdesignoutlined{60fps}{54}
\showcasematerialdesignoutlined{60fps-select}{55}
\showcasematerialdesignoutlined{123}{56}
\showcasematerialdesignoutlined{360}{57}
\showcasematerialdesignoutlined{abc}{58}
\showcasematerialdesignoutlined{accessibility}{59}
\showcasematerialdesignoutlined{accessibility-new}{60}
\showcasematerialdesignoutlined{accessible}{61}
\showcasematerialdesignoutlined{accessible-forward}{62}
\showcasematerialdesignoutlined{access-alarm}{63}
\showcasematerialdesignoutlined{access-alarms}{64}
\showcasematerialdesignoutlined{access-time}{65}
\showcasematerialdesignoutlined{access-time-filled}{66}
\showcasematerialdesignoutlined{account-balance}{67}
\showcasematerialdesignoutlined{account-balance-wallet}{68}
\showcasematerialdesignoutlined{account-box}{69}
\showcasematerialdesignoutlined{account-circle}{70}
\showcasematerialdesignoutlined{account-tree}{71}
\showcasematerialdesignoutlined{ac-unit}{72}
\showcasematerialdesignoutlined{adb}{73}
\showcasematerialdesignoutlined{add}{74}
\showcasematerialdesignoutlined{addchart}{75}
\showcasematerialdesignoutlined{add-alarm}{76}
\showcasematerialdesignoutlined{add-alert}{77}
\showcasematerialdesignoutlined{add-a-photo}{78}
\showcasematerialdesignoutlined{add-box}{79}
\showcasematerialdesignoutlined{add-business}{80}
\showcasematerialdesignoutlined{add-card}{81}
\showcasematerialdesignoutlined{add-chart}{82}
\showcasematerialdesignoutlined{add-circle}{83}
\showcasematerialdesignoutlined{add-circle-outline}{84}
\showcasematerialdesignoutlined{add-comment}{85}
\showcasematerialdesignoutlined{add-home}{86}
\showcasematerialdesignoutlined{add-home-work}{87}
\showcasematerialdesignoutlined{add-ic-call}{88}
\showcasematerialdesignoutlined{add-link}{89}
\showcasematerialdesignoutlined{add-location}{90}
\showcasematerialdesignoutlined{add-location-alt}{91}
\showcasematerialdesignoutlined{add-moderator}{92}
\showcasematerialdesignoutlined{add-photo-alternate}{93}
\showcasematerialdesignoutlined{add-reaction}{94}
\showcasematerialdesignoutlined{add-road}{95}
\showcasematerialdesignoutlined{add-shopping-cart}{96}
\showcasematerialdesignoutlined{add-task}{97}
\showcasematerialdesignoutlined{add-to-drive}{98}
\showcasematerialdesignoutlined{add-to-home-screen}{99}
\showcasematerialdesignoutlined{add-to-photos}{100}
\showcasematerialdesignoutlined{add-to-queue}{101}
\showcasematerialdesignoutlined{adf-scanner}{102}
\showcasematerialdesignoutlined{adjust}{103}
\showcasematerialdesignoutlined{admin-panel-settings}{104}
\showcasematerialdesignoutlined{ads-click}{105}
\showcasematerialdesignoutlined{ad-units}{106}
\showcasematerialdesignoutlined{agriculture}{107}
\showcasematerialdesignoutlined{air}{108}
\showcasematerialdesignoutlined{airlines}{109}
\showcasematerialdesignoutlined{airline-seat-flat}{110}
\showcasematerialdesignoutlined{airline-seat-flat-angled}{111}
\showcasematerialdesignoutlined{airline-seat-individual-suite}{112}
\showcasematerialdesignoutlined{airline-seat-legroom-extra}{113}
\showcasematerialdesignoutlined{airline-seat-legroom-normal}{114}
\showcasematerialdesignoutlined{airline-seat-legroom-reduced}{115}
\showcasematerialdesignoutlined{airline-seat-recline-extra}{116}
\showcasematerialdesignoutlined{airline-seat-recline-normal}{117}
\showcasematerialdesignoutlined{airline-stops}{118}
\showcasematerialdesignoutlined{airplanemode-active}{119}
\showcasematerialdesignoutlined{airplanemode-inactive}{120}
\showcasematerialdesignoutlined{airplane-ticket}{121}
\showcasematerialdesignoutlined{airplay}{122}
\showcasematerialdesignoutlined{airport-shuttle}{123}
\showcasematerialdesignoutlined{alarm}{124}
\showcasematerialdesignoutlined{alarm-add}{125}
\showcasematerialdesignoutlined{alarm-off}{126}
\showcasematerialdesignoutlined{alarm-on}{127}
\showcasematerialdesignoutlined{album}{128}
\showcasematerialdesignoutlined{align-horizontal-center}{129}
\showcasematerialdesignoutlined{align-horizontal-left}{130}
\showcasematerialdesignoutlined{align-horizontal-right}{131}
\showcasematerialdesignoutlined{align-vertical-bottom}{132}
\showcasematerialdesignoutlined{align-vertical-center}{133}
\showcasematerialdesignoutlined{align-vertical-top}{134}
\showcasematerialdesignoutlined{all-inbox}{135}
\showcasematerialdesignoutlined{all-inclusive}{136}
\showcasematerialdesignoutlined{all-out}{137}
\showcasematerialdesignoutlined{alternate-email}{138}
\showcasematerialdesignoutlined{alt-route}{139}
\showcasematerialdesignoutlined{analytics}{140}
\showcasematerialdesignoutlined{anchor}{141}
\showcasematerialdesignoutlined{android}{142}
\showcasematerialdesignoutlined{animation}{143}
\showcasematerialdesignoutlined{announcement}{144}
\showcasematerialdesignoutlined{aod}{145}
\showcasematerialdesignoutlined{apartment}{146}
\showcasematerialdesignoutlined{api}{147}
\showcasematerialdesignoutlined{approval}{148}
\showcasematerialdesignoutlined{apps}{149}
\showcasematerialdesignoutlined{apps-outage}{150}
\showcasematerialdesignoutlined{app-blocking}{151}
\showcasematerialdesignoutlined{app-registration}{152}
\showcasematerialdesignoutlined{app-settings-alt}{153}
\showcasematerialdesignoutlined{app-shortcut}{154}
\showcasematerialdesignoutlined{architecture}{155}
\showcasematerialdesignoutlined{archive}{156}
\showcasematerialdesignoutlined{area-chart}{157}
\showcasematerialdesignoutlined{arrow-back}{158}
\showcasematerialdesignoutlined{arrow-back-ios}{159}
\showcasematerialdesignoutlined{arrow-back-ios-new}{160}
\showcasematerialdesignoutlined{arrow-circle-down}{161}
\showcasematerialdesignoutlined{arrow-circle-left}{162}
\showcasematerialdesignoutlined{arrow-circle-right}{163}
\showcasematerialdesignoutlined{arrow-circle-up}{164}
\showcasematerialdesignoutlined{arrow-downward}{165}
\showcasematerialdesignoutlined{arrow-drop-down}{166}
\showcasematerialdesignoutlined{arrow-drop-down-circle}{167}
\showcasematerialdesignoutlined{arrow-drop-up}{168}
\showcasematerialdesignoutlined{arrow-forward}{169}
\showcasematerialdesignoutlined{arrow-forward-ios}{170}
\showcasematerialdesignoutlined{arrow-left}{171}
\showcasematerialdesignoutlined{arrow-outward}{172}
\showcasematerialdesignoutlined{arrow-right}{173}
\showcasematerialdesignoutlined{arrow-right-alt}{174}
\showcasematerialdesignoutlined{arrow-upward}{175}
\showcasematerialdesignoutlined{article}{176}
\showcasematerialdesignoutlined{art-track}{177}
\showcasematerialdesignoutlined{aspect-ratio}{178}
\showcasematerialdesignoutlined{assessment}{179}
\showcasematerialdesignoutlined{assignment}{180}
\showcasematerialdesignoutlined{assignment-ind}{181}
\showcasematerialdesignoutlined{assignment-late}{182}
\showcasematerialdesignoutlined{assignment-return}{183}
\showcasematerialdesignoutlined{assignment-returned}{184}
\showcasematerialdesignoutlined{assignment-turned-in}{185}
\showcasematerialdesignoutlined{assistant}{186}
\showcasematerialdesignoutlined{assistant-direction}{187}
\showcasematerialdesignoutlined{assistant-photo}{188}
\showcasematerialdesignoutlined{assist-walker}{189}
\showcasematerialdesignoutlined{assured-workload}{190}
\showcasematerialdesignoutlined{atm}{191}
\showcasematerialdesignoutlined{attachment}{192}
\showcasematerialdesignoutlined{attach-email}{193}
\showcasematerialdesignoutlined{attach-file}{194}
\showcasematerialdesignoutlined{attach-money}{195}
\showcasematerialdesignoutlined{attractions}{196}
\showcasematerialdesignoutlined{attribution}{197}
\showcasematerialdesignoutlined{audiotrack}{198}
\showcasematerialdesignoutlined{audio-file}{199}
\showcasematerialdesignoutlined{autofps-select}{200}
\showcasematerialdesignoutlined{autorenew}{201}
\showcasematerialdesignoutlined{auto-awesome}{202}
\showcasematerialdesignoutlined{auto-awesome-mosaic}{203}
\showcasematerialdesignoutlined{auto-awesome-motion}{204}
\showcasematerialdesignoutlined{auto-delete}{205}
\showcasematerialdesignoutlined{auto-fix-high}{206}
\showcasematerialdesignoutlined{auto-fix-normal}{207}
\showcasematerialdesignoutlined{auto-fix-off}{208}
\showcasematerialdesignoutlined{auto-graph}{209}
\showcasematerialdesignoutlined{auto-mode}{210}
\showcasematerialdesignoutlined{auto-stories}{211}
\showcasematerialdesignoutlined{av-timer}{212}
\showcasematerialdesignoutlined{baby-changing-station}{213}
\showcasematerialdesignoutlined{backpack}{214}
\showcasematerialdesignoutlined{backspace}{215}
\showcasematerialdesignoutlined{backup}{216}
\showcasematerialdesignoutlined{backup-table}{217}
\showcasematerialdesignoutlined{back-hand}{218}
\showcasematerialdesignoutlined{badge}{219}
\showcasematerialdesignoutlined{bakery-dining}{220}
\showcasematerialdesignoutlined{balance}{221}
\showcasematerialdesignoutlined{balcony}{222}
\showcasematerialdesignoutlined{ballot}{223}
\showcasematerialdesignoutlined{bar-chart}{224}
\showcasematerialdesignoutlined{batch-prediction}{225}
\showcasematerialdesignoutlined{bathroom}{226}
\showcasematerialdesignoutlined{bathtub}{227}
\showcasematerialdesignoutlined{battery-0-bar}{228}
\showcasematerialdesignoutlined{battery-1-bar}{229}
\showcasematerialdesignoutlined{battery-2-bar}{230}
\showcasematerialdesignoutlined{battery-3-bar}{231}
\showcasematerialdesignoutlined{battery-4-bar}{232}
\showcasematerialdesignoutlined{battery-5-bar}{233}
\showcasematerialdesignoutlined{battery-6-bar}{234}
\showcasematerialdesignoutlined{battery-alert}{235}
\showcasematerialdesignoutlined{battery-charging-full}{236}
\showcasematerialdesignoutlined{battery-full}{237}
\showcasematerialdesignoutlined{battery-saver}{238}
\showcasematerialdesignoutlined{battery-std}{239}
\showcasematerialdesignoutlined{battery-unknown}{240}
\showcasematerialdesignoutlined{beach-access}{241}
\showcasematerialdesignoutlined{bed}{242}
\showcasematerialdesignoutlined{bedroom-baby}{243}
\showcasematerialdesignoutlined{bedroom-child}{244}
\showcasematerialdesignoutlined{bedroom-parent}{245}
\showcasematerialdesignoutlined{bedtime}{246}
\showcasematerialdesignoutlined{bedtime-off}{247}
\showcasematerialdesignoutlined{beenhere}{248}
\showcasematerialdesignoutlined{bento}{249}
\showcasematerialdesignoutlined{bike-scooter}{250}
\showcasematerialdesignoutlined{biotech}{251}
\showcasematerialdesignoutlined{blender}{252}
\showcasematerialdesignoutlined{blind}{253}
\showcasematerialdesignoutlined{blinds}{254}
\showcasematerialdesignoutlined{blinds-closed}{255}
\showcasematerialdesignoutlined{block}{256}
\showcasematerialdesignoutlined{bloodtype}{257}
\showcasematerialdesignoutlined{bluetooth}{258}
\showcasematerialdesignoutlined{bluetooth-audio}{259}
\showcasematerialdesignoutlined{bluetooth-connected}{260}
\showcasematerialdesignoutlined{bluetooth-disabled}{261}
\showcasematerialdesignoutlined{bluetooth-drive}{262}
\showcasematerialdesignoutlined{bluetooth-searching}{263}
\showcasematerialdesignoutlined{blur-circular}{264}
\showcasematerialdesignoutlined{blur-linear}{265}
\showcasematerialdesignoutlined{blur-off}{266}
\showcasematerialdesignoutlined{blur-on}{267}
\showcasematerialdesignoutlined{bolt}{268}
\showcasematerialdesignoutlined{book}{269}
\showcasematerialdesignoutlined{bookmark}{270}
\showcasematerialdesignoutlined{bookmarks}{271}
\showcasematerialdesignoutlined{bookmark-add}{272}
\showcasematerialdesignoutlined{bookmark-added}{273}
\showcasematerialdesignoutlined{bookmark-border}{274}
\showcasematerialdesignoutlined{bookmark-remove}{275}
\showcasematerialdesignoutlined{book-online}{276}
\showcasematerialdesignoutlined{border-all}{277}
\showcasematerialdesignoutlined{border-bottom}{278}
\showcasematerialdesignoutlined{border-clear}{279}
\showcasematerialdesignoutlined{border-color}{280}
\showcasematerialdesignoutlined{border-horizontal}{281}
\showcasematerialdesignoutlined{border-inner}{282}
\showcasematerialdesignoutlined{border-left}{283}
\showcasematerialdesignoutlined{border-outer}{284}
\showcasematerialdesignoutlined{border-right}{285}
\showcasematerialdesignoutlined{border-style}{286}
\showcasematerialdesignoutlined{border-top}{287}
\showcasematerialdesignoutlined{border-vertical}{288}
\showcasematerialdesignoutlined{boy}{289}
\showcasematerialdesignoutlined{branding-watermark}{290}
\showcasematerialdesignoutlined{breakfast-dining}{291}
\showcasematerialdesignoutlined{brightness-1}{292}
\showcasematerialdesignoutlined{brightness-2}{293}
\showcasematerialdesignoutlined{brightness-3}{294}
\showcasematerialdesignoutlined{brightness-4}{295}
\showcasematerialdesignoutlined{brightness-5}{296}
\showcasematerialdesignoutlined{brightness-6}{297}
\showcasematerialdesignoutlined{brightness-7}{298}
\showcasematerialdesignoutlined{brightness-auto}{299}
\showcasematerialdesignoutlined{brightness-high}{300}
\showcasematerialdesignoutlined{brightness-low}{301}
\showcasematerialdesignoutlined{brightness-medium}{302}
\showcasematerialdesignoutlined{broadcast-on-home}{303}
\showcasematerialdesignoutlined{broadcast-on-personal}{304}
\showcasematerialdesignoutlined{broken-image}{305}
\showcasematerialdesignoutlined{browser-not-supported}{306}
\showcasematerialdesignoutlined{browser-updated}{307}
\showcasematerialdesignoutlined{browse-gallery}{308}
\showcasematerialdesignoutlined{brunch-dining}{309}
\showcasematerialdesignoutlined{brush}{310}
\showcasematerialdesignoutlined{bubble-chart}{311}
\showcasematerialdesignoutlined{bug-report}{312}
\showcasematerialdesignoutlined{build}{313}
\showcasematerialdesignoutlined{build-circle}{314}
\showcasematerialdesignoutlined{bungalow}{315}
\showcasematerialdesignoutlined{burst-mode}{316}
\showcasematerialdesignoutlined{business}{317}
\showcasematerialdesignoutlined{business-center}{318}
\showcasematerialdesignoutlined{bus-alert}{319}
\showcasematerialdesignoutlined{cabin}{320}
\showcasematerialdesignoutlined{cable}{321}
\showcasematerialdesignoutlined{cached}{322}
\showcasematerialdesignoutlined{cake}{323}
\showcasematerialdesignoutlined{calculate}{324}
\showcasematerialdesignoutlined{calendar-month}{325}
\showcasematerialdesignoutlined{calendar-today}{326}
\showcasematerialdesignoutlined{calendar-view-day}{327}
\showcasematerialdesignoutlined{calendar-view-month}{328}
\showcasematerialdesignoutlined{calendar-view-week}{329}
\showcasematerialdesignoutlined{call}{330}
\showcasematerialdesignoutlined{call-end}{331}
\showcasematerialdesignoutlined{call-made}{332}
\showcasematerialdesignoutlined{call-merge}{333}
\showcasematerialdesignoutlined{call-missed}{334}
\showcasematerialdesignoutlined{call-missed-outgoing}{335}
\showcasematerialdesignoutlined{call-received}{336}
\showcasematerialdesignoutlined{call-split}{337}
\showcasematerialdesignoutlined{call-to-action}{338}
\showcasematerialdesignoutlined{camera}{339}
\showcasematerialdesignoutlined{cameraswitch}{340}
\showcasematerialdesignoutlined{camera-alt}{341}
\showcasematerialdesignoutlined{camera-enhance}{342}
\showcasematerialdesignoutlined{camera-front}{343}
\showcasematerialdesignoutlined{camera-indoor}{344}
\showcasematerialdesignoutlined{camera-outdoor}{345}
\showcasematerialdesignoutlined{camera-rear}{346}
\showcasematerialdesignoutlined{camera-roll}{347}
\showcasematerialdesignoutlined{campaign}{348}
\showcasematerialdesignoutlined{cancel}{349}
\showcasematerialdesignoutlined{cancel-presentation}{350}
\showcasematerialdesignoutlined{cancel-schedule-send}{351}
\showcasematerialdesignoutlined{candlestick-chart}{352}
\showcasematerialdesignoutlined{card-giftcard}{353}
\showcasematerialdesignoutlined{card-membership}{354}
\showcasematerialdesignoutlined{card-travel}{355}
\showcasematerialdesignoutlined{carpenter}{356}
\showcasematerialdesignoutlined{car-crash}{357}
\showcasematerialdesignoutlined{car-rental}{358}
\showcasematerialdesignoutlined{car-repair}{359}
\showcasematerialdesignoutlined{cases}{360}
\showcasematerialdesignoutlined{casino}{361}
\showcasematerialdesignoutlined{cast}{362}
\showcasematerialdesignoutlined{castle}{363}
\showcasematerialdesignoutlined{cast-connected}{364}
\showcasematerialdesignoutlined{cast-for-education}{365}
\showcasematerialdesignoutlined{catching-pokemon}{366}
\showcasematerialdesignoutlined{category}{367}
\showcasematerialdesignoutlined{celebration}{368}
\showcasematerialdesignoutlined{cell-tower}{369}
\showcasematerialdesignoutlined{cell-wifi}{370}
\showcasematerialdesignoutlined{center-focus-strong}{371}
\showcasematerialdesignoutlined{center-focus-weak}{372}
\showcasematerialdesignoutlined{chair}{373}
\showcasematerialdesignoutlined{chair-alt}{374}
\showcasematerialdesignoutlined{chalet}{375}
\showcasematerialdesignoutlined{change-circle}{376}
\showcasematerialdesignoutlined{change-history}{377}
\showcasematerialdesignoutlined{charging-station}{378}
\showcasematerialdesignoutlined{chat}{379}
\showcasematerialdesignoutlined{chat-bubble}{380}
\showcasematerialdesignoutlined{chat-bubble-outline}{381}
\showcasematerialdesignoutlined{check}{382}
\showcasematerialdesignoutlined{checklist}{383}
\showcasematerialdesignoutlined{checklist-rtl}{384}
\showcasematerialdesignoutlined{checkroom}{385}
\showcasematerialdesignoutlined{check-box}{386}
\showcasematerialdesignoutlined{check-box-outline-blank}{387}
\showcasematerialdesignoutlined{check-circle}{388}
\showcasematerialdesignoutlined{check-circle-outline}{389}
\showcasematerialdesignoutlined{chevron-left}{390}
\showcasematerialdesignoutlined{chevron-right}{391}
\showcasematerialdesignoutlined{child-care}{392}
\showcasematerialdesignoutlined{child-friendly}{393}
\showcasematerialdesignoutlined{chrome-reader-mode}{394}
\showcasematerialdesignoutlined{church}{395}
\showcasematerialdesignoutlined{circle}{396}
\showcasematerialdesignoutlined{circle-notifications}{397}
\showcasematerialdesignoutlined{class}{398}
\showcasematerialdesignoutlined{cleaning-services}{399}
\showcasematerialdesignoutlined{clean-hands}{400}
\showcasematerialdesignoutlined{clear}{401}
\showcasematerialdesignoutlined{clear-all}{402}
\showcasematerialdesignoutlined{close}{403}
\showcasematerialdesignoutlined{closed-caption}{404}
\showcasematerialdesignoutlined{closed-caption-disabled}{405}
\showcasematerialdesignoutlined{closed-caption-off}{406}
\showcasematerialdesignoutlined{close-fullscreen}{407}
\showcasematerialdesignoutlined{cloud}{408}
\showcasematerialdesignoutlined{cloud-circle}{409}
\showcasematerialdesignoutlined{cloud-done}{410}
\showcasematerialdesignoutlined{cloud-download}{411}
\showcasematerialdesignoutlined{cloud-off}{412}
\showcasematerialdesignoutlined{cloud-queue}{413}
\showcasematerialdesignoutlined{cloud-sync}{414}
\showcasematerialdesignoutlined{cloud-upload}{415}
\showcasematerialdesignoutlined{co2}{416}
\showcasematerialdesignoutlined{code}{417}
\showcasematerialdesignoutlined{code-off}{418}
\showcasematerialdesignoutlined{coffee}{419}
\showcasematerialdesignoutlined{coffee-maker}{420}
\showcasematerialdesignoutlined{collections}{421}
\showcasematerialdesignoutlined{collections-bookmark}{422}
\showcasematerialdesignoutlined{colorize}{423}
\showcasematerialdesignoutlined{color-lens}{424}
\showcasematerialdesignoutlined{comment}{425}
\showcasematerialdesignoutlined{comments-disabled}{426}
\showcasematerialdesignoutlined{comment-bank}{427}
\showcasematerialdesignoutlined{commit}{428}
\showcasematerialdesignoutlined{commute}{429}
\showcasematerialdesignoutlined{compare}{430}
\showcasematerialdesignoutlined{compare-arrows}{431}
\showcasematerialdesignoutlined{compass-calibration}{432}
\showcasematerialdesignoutlined{compost}{433}
\showcasematerialdesignoutlined{compress}{434}
\showcasematerialdesignoutlined{computer}{435}
\showcasematerialdesignoutlined{confirmation-number}{436}
\showcasematerialdesignoutlined{connected-tv}{437}
\showcasematerialdesignoutlined{connecting-airports}{438}
\showcasematerialdesignoutlined{connect-without-contact}{439}
\showcasematerialdesignoutlined{construction}{440}
\showcasematerialdesignoutlined{contactless}{441}
\showcasematerialdesignoutlined{contacts}{442}
\showcasematerialdesignoutlined{contact-emergency}{443}
\showcasematerialdesignoutlined{contact-mail}{444}
\showcasematerialdesignoutlined{contact-page}{445}
\showcasematerialdesignoutlined{contact-phone}{446}
\showcasematerialdesignoutlined{contact-support}{447}
\showcasematerialdesignoutlined{content-copy}{448}
\showcasematerialdesignoutlined{content-cut}{449}
\showcasematerialdesignoutlined{content-paste}{450}
\showcasematerialdesignoutlined{content-paste-go}{451}
\showcasematerialdesignoutlined{content-paste-off}{452}
\showcasematerialdesignoutlined{content-paste-search}{453}
\showcasematerialdesignoutlined{contrast}{454}
\showcasematerialdesignoutlined{control-camera}{455}
\showcasematerialdesignoutlined{control-point}{456}
\showcasematerialdesignoutlined{control-point-duplicate}{457}
\showcasematerialdesignoutlined{cookie}{458}
\showcasematerialdesignoutlined{copyright}{459}
\showcasematerialdesignoutlined{copy-all}{460}
\showcasematerialdesignoutlined{coronavirus}{461}
\showcasematerialdesignoutlined{corporate-fare}{462}
\showcasematerialdesignoutlined{cottage}{463}
\showcasematerialdesignoutlined{countertops}{464}
\showcasematerialdesignoutlined{co-present}{465}
\showcasematerialdesignoutlined{create}{466}
\showcasematerialdesignoutlined{create-new-folder}{467}
\showcasematerialdesignoutlined{credit-card}{468}
\showcasematerialdesignoutlined{credit-card-off}{469}
\showcasematerialdesignoutlined{credit-score}{470}
\showcasematerialdesignoutlined{crib}{471}
\showcasematerialdesignoutlined{crisis-alert}{472}
\showcasematerialdesignoutlined{crop}{473}
\showcasematerialdesignoutlined{crop-3-2}{474}
\showcasematerialdesignoutlined{crop-5-4}{475}
\showcasematerialdesignoutlined{crop-7-5}{476}
\showcasematerialdesignoutlined{crop-16-9}{477}
\showcasematerialdesignoutlined{crop-din}{478}
\showcasematerialdesignoutlined{crop-free}{479}
\showcasematerialdesignoutlined{crop-landscape}{480}
\showcasematerialdesignoutlined{crop-original}{481}
\showcasematerialdesignoutlined{crop-portrait}{482}
\showcasematerialdesignoutlined{crop-rotate}{483}
\showcasematerialdesignoutlined{crop-square}{484}
\showcasematerialdesignoutlined{cruelty-free}{485}
\showcasematerialdesignoutlined{css}{486}
\showcasematerialdesignoutlined{currency-bitcoin}{487}
\showcasematerialdesignoutlined{currency-exchange}{488}
\showcasematerialdesignoutlined{currency-franc}{489}
\showcasematerialdesignoutlined{currency-lira}{490}
\showcasematerialdesignoutlined{currency-pound}{491}
\showcasematerialdesignoutlined{currency-ruble}{492}
\showcasematerialdesignoutlined{currency-rupee}{493}
\showcasematerialdesignoutlined{currency-yen}{494}
\showcasematerialdesignoutlined{currency-yuan}{495}
\showcasematerialdesignoutlined{curtains}{496}
\showcasematerialdesignoutlined{curtains-closed}{497}
\showcasematerialdesignoutlined{cyclone}{498}
\showcasematerialdesignoutlined{dangerous}{499}
\showcasematerialdesignoutlined{dark-mode}{500}
\showcasematerialdesignoutlined{dashboard}{501}
\showcasematerialdesignoutlined{dashboard-customize}{502}
\showcasematerialdesignoutlined{dataset}{503}
\showcasematerialdesignoutlined{dataset-linked}{504}
\showcasematerialdesignoutlined{data-array}{505}
\showcasematerialdesignoutlined{data-exploration}{506}
\showcasematerialdesignoutlined{data-object}{507}
\showcasematerialdesignoutlined{data-saver-off}{508}
\showcasematerialdesignoutlined{data-saver-on}{509}
\showcasematerialdesignoutlined{data-thresholding}{510}
\showcasematerialdesignoutlined{data-usage}{511}
\showcasematerialdesignoutlined{date-range}{512}
\showcasematerialdesignoutlined{deblur}{513}
\showcasematerialdesignoutlined{deck}{514}
\showcasematerialdesignoutlined{dehaze}{515}
\showcasematerialdesignoutlined{delete}{516}
\showcasematerialdesignoutlined{delete-forever}{517}
\showcasematerialdesignoutlined{delete-outline}{518}
\showcasematerialdesignoutlined{delete-sweep}{519}
\showcasematerialdesignoutlined{delivery-dining}{520}
\showcasematerialdesignoutlined{density-large}{521}
\showcasematerialdesignoutlined{density-medium}{522}
\showcasematerialdesignoutlined{density-small}{523}
\showcasematerialdesignoutlined{departure-board}{524}
\showcasematerialdesignoutlined{description}{525}
\showcasematerialdesignoutlined{deselect}{526}
\showcasematerialdesignoutlined{design-services}{527}
\showcasematerialdesignoutlined{desk}{528}
\showcasematerialdesignoutlined{desktop-access-disabled}{529}
\showcasematerialdesignoutlined{desktop-mac}{530}
\showcasematerialdesignoutlined{desktop-windows}{531}
\showcasematerialdesignoutlined{details}{532}
\showcasematerialdesignoutlined{developer-board}{533}
\showcasematerialdesignoutlined{developer-board-off}{534}
\showcasematerialdesignoutlined{developer-mode}{535}
\showcasematerialdesignoutlined{devices}{536}
\showcasematerialdesignoutlined{devices-fold}{537}
\showcasematerialdesignoutlined{devices-other}{538}
\showcasematerialdesignoutlined{device-hub}{539}
\showcasematerialdesignoutlined{device-thermostat}{540}
\showcasematerialdesignoutlined{device-unknown}{541}
\showcasematerialdesignoutlined{dialer-sip}{542}
\showcasematerialdesignoutlined{dialpad}{543}
\showcasematerialdesignoutlined{diamond}{544}
\showcasematerialdesignoutlined{difference}{545}
\showcasematerialdesignoutlined{dining}{546}
\showcasematerialdesignoutlined{dinner-dining}{547}
\showcasematerialdesignoutlined{directions}{548}
\showcasematerialdesignoutlined{directions-bike}{549}
\showcasematerialdesignoutlined{directions-boat}{550}
\showcasematerialdesignoutlined{directions-boat-filled}{551}
\showcasematerialdesignoutlined{directions-bus}{552}
\showcasematerialdesignoutlined{directions-bus-filled}{553}
\showcasematerialdesignoutlined{directions-car}{554}
\showcasematerialdesignoutlined{directions-car-filled}{555}
\showcasematerialdesignoutlined{directions-off}{556}
\showcasematerialdesignoutlined{directions-railway}{557}
\showcasematerialdesignoutlined{directions-railway-filled}{558}
\showcasematerialdesignoutlined{directions-run}{559}
\showcasematerialdesignoutlined{directions-subway}{560}
\showcasematerialdesignoutlined{directions-subway-filled}{561}
\showcasematerialdesignoutlined{directions-transit}{562}
\showcasematerialdesignoutlined{directions-transit-filled}{563}
\showcasematerialdesignoutlined{directions-walk}{564}
\showcasematerialdesignoutlined{dirty-lens}{565}
\showcasematerialdesignoutlined{disabled-by-default}{566}
\showcasematerialdesignoutlined{disabled-visible}{567}
\showcasematerialdesignoutlined{discount}{568}
\showcasematerialdesignoutlined{disc-full}{569}
\showcasematerialdesignoutlined{display-settings}{570}
\showcasematerialdesignoutlined{diversity-1}{571}
\showcasematerialdesignoutlined{diversity-2}{572}
\showcasematerialdesignoutlined{diversity-3}{573}
\showcasematerialdesignoutlined{dns}{574}
\showcasematerialdesignoutlined{dock}{575}
\showcasematerialdesignoutlined{document-scanner}{576}
\showcasematerialdesignoutlined{domain}{577}
\showcasematerialdesignoutlined{domain-add}{578}
\showcasematerialdesignoutlined{domain-disabled}{579}
\showcasematerialdesignoutlined{domain-verification}{580}
\showcasematerialdesignoutlined{done}{581}
\showcasematerialdesignoutlined{done-all}{582}
\showcasematerialdesignoutlined{done-outline}{583}
\showcasematerialdesignoutlined{donut-large}{584}
\showcasematerialdesignoutlined{donut-small}{585}
\showcasematerialdesignoutlined{doorbell}{586}
\showcasematerialdesignoutlined{door-back}{587}
\showcasematerialdesignoutlined{door-front}{588}
\showcasematerialdesignoutlined{door-sliding}{589}
\showcasematerialdesignoutlined{double-arrow}{590}
\showcasematerialdesignoutlined{downhill-skiing}{591}
\showcasematerialdesignoutlined{download}{592}
\showcasematerialdesignoutlined{downloading}{593}
\showcasematerialdesignoutlined{download-done}{594}
\showcasematerialdesignoutlined{download-for-offline}{595}
\showcasematerialdesignoutlined{do-disturb}{596}
\showcasematerialdesignoutlined{do-disturb-alt}{597}
\showcasematerialdesignoutlined{do-disturb-off}{598}
\showcasematerialdesignoutlined{do-disturb-on}{599}
\showcasematerialdesignoutlined{do-not-disturb}{600}
\showcasematerialdesignoutlined{do-not-disturb-alt}{601}
\showcasematerialdesignoutlined{do-not-disturb-off}{602}
\showcasematerialdesignoutlined{do-not-disturb-on}{603}
\showcasematerialdesignoutlined{do-not-disturb-on-total-silence}{604}
\showcasematerialdesignoutlined{do-not-step}{605}
\showcasematerialdesignoutlined{do-not-touch}{606}
\showcasematerialdesignoutlined{drafts}{607}
\showcasematerialdesignoutlined{drag-handle}{608}
\showcasematerialdesignoutlined{drag-indicator}{609}
\showcasematerialdesignoutlined{draw}{610}
\showcasematerialdesignoutlined{drive-eta}{611}
\showcasematerialdesignoutlined{drive-file-move}{612}
\showcasematerialdesignoutlined{drive-file-move-rtl}{613}
\showcasematerialdesignoutlined{drive-file-rename-outline}{614}
\showcasematerialdesignoutlined{drive-folder-upload}{615}
\showcasematerialdesignoutlined{dry}{616}
\showcasematerialdesignoutlined{dry-cleaning}{617}
\showcasematerialdesignoutlined{duo}{618}
\showcasematerialdesignoutlined{dvr}{619}
\showcasematerialdesignoutlined{dynamic-feed}{620}
\showcasematerialdesignoutlined{dynamic-form}{621}
\showcasematerialdesignoutlined{earbuds}{622}
\showcasematerialdesignoutlined{earbuds-battery}{623}
\showcasematerialdesignoutlined{east}{624}
\showcasematerialdesignoutlined{edgesensor-high}{625}
\showcasematerialdesignoutlined{edgesensor-low}{626}
\showcasematerialdesignoutlined{edit}{627}
\showcasematerialdesignoutlined{edit-attributes}{628}
\showcasematerialdesignoutlined{edit-calendar}{629}
\showcasematerialdesignoutlined{edit-location}{630}
\showcasematerialdesignoutlined{edit-location-alt}{631}
\showcasematerialdesignoutlined{edit-note}{632}
\showcasematerialdesignoutlined{edit-notifications}{633}
\showcasematerialdesignoutlined{edit-off}{634}
\showcasematerialdesignoutlined{edit-road}{635}
\showcasematerialdesignoutlined{egg}{636}
\showcasematerialdesignoutlined{egg-alt}{637}
\showcasematerialdesignoutlined{eject}{638}
\showcasematerialdesignoutlined{elderly}{639}
\showcasematerialdesignoutlined{elderly-woman}{640}
\showcasematerialdesignoutlined{electrical-services}{641}
\showcasematerialdesignoutlined{electric-bike}{642}
\showcasematerialdesignoutlined{electric-bolt}{643}
\showcasematerialdesignoutlined{electric-car}{644}
\showcasematerialdesignoutlined{electric-meter}{645}
\showcasematerialdesignoutlined{electric-moped}{646}
\showcasematerialdesignoutlined{electric-rickshaw}{647}
\showcasematerialdesignoutlined{electric-scooter}{648}
\showcasematerialdesignoutlined{elevator}{649}
\showcasematerialdesignoutlined{email}{650}
\showcasematerialdesignoutlined{emergency}{651}
\showcasematerialdesignoutlined{emergency-recording}{652}
\showcasematerialdesignoutlined{emergency-share}{653}
\showcasematerialdesignoutlined{emoji-emotions}{654}
\showcasematerialdesignoutlined{emoji-events}{655}
\showcasematerialdesignoutlined{emoji-food-beverage}{656}
\showcasematerialdesignoutlined{emoji-nature}{657}
\showcasematerialdesignoutlined{emoji-objects}{658}
\showcasematerialdesignoutlined{emoji-people}{659}
\showcasematerialdesignoutlined{emoji-symbols}{660}
\showcasematerialdesignoutlined{emoji-transportation}{661}
\showcasematerialdesignoutlined{energy-savings-leaf}{662}
\showcasematerialdesignoutlined{engineering}{663}
\showcasematerialdesignoutlined{enhanced-encryption}{664}
\showcasematerialdesignoutlined{equalizer}{665}
\showcasematerialdesignoutlined{error}{666}
\showcasematerialdesignoutlined{error-outline}{667}
\showcasematerialdesignoutlined{escalator}{668}
\showcasematerialdesignoutlined{escalator-warning}{669}
\showcasematerialdesignoutlined{euro}{670}
\showcasematerialdesignoutlined{euro-symbol}{671}
\showcasematerialdesignoutlined{event}{672}
\showcasematerialdesignoutlined{event-available}{673}
\showcasematerialdesignoutlined{event-busy}{674}
\showcasematerialdesignoutlined{event-note}{675}
\showcasematerialdesignoutlined{event-repeat}{676}
\showcasematerialdesignoutlined{event-seat}{677}
\showcasematerialdesignoutlined{ev-station}{678}
\showcasematerialdesignoutlined{exit-to-app}{679}
\showcasematerialdesignoutlined{expand}{680}
\showcasematerialdesignoutlined{expand-circle-down}{681}
\showcasematerialdesignoutlined{expand-less}{682}
\showcasematerialdesignoutlined{expand-more}{683}
\showcasematerialdesignoutlined{explicit}{684}
\showcasematerialdesignoutlined{explore}{685}
\showcasematerialdesignoutlined{explore-off}{686}
\showcasematerialdesignoutlined{exposure}{687}
\showcasematerialdesignoutlined{exposure-neg-1}{688}
\showcasematerialdesignoutlined{exposure-neg-2}{689}
\showcasematerialdesignoutlined{exposure-plus-1}{690}
\showcasematerialdesignoutlined{exposure-plus-2}{691}
\showcasematerialdesignoutlined{exposure-zero}{692}
\showcasematerialdesignoutlined{extension}{693}
\showcasematerialdesignoutlined{extension-off}{694}
\showcasematerialdesignoutlined{e-mobiledata}{695}
\showcasematerialdesignoutlined{face}{696}
\showcasematerialdesignoutlined{face-2}{697}
\showcasematerialdesignoutlined{face-3}{698}
\showcasematerialdesignoutlined{face-4}{699}
\showcasematerialdesignoutlined{face-5}{700}
\showcasematerialdesignoutlined{face-6}{701}
\showcasematerialdesignoutlined{face-retouching-natural}{702}
\showcasematerialdesignoutlined{face-retouching-off}{703}
\showcasematerialdesignoutlined{factory}{704}
\showcasematerialdesignoutlined{fact-check}{705}
\showcasematerialdesignoutlined{family-restroom}{706}
\showcasematerialdesignoutlined{fastfood}{707}
\showcasematerialdesignoutlined{fast-forward}{708}
\showcasematerialdesignoutlined{fast-rewind}{709}
\showcasematerialdesignoutlined{favorite}{710}
\showcasematerialdesignoutlined{favorite-border}{711}
\showcasematerialdesignoutlined{fax}{712}
\showcasematerialdesignoutlined{featured-play-list}{713}
\showcasematerialdesignoutlined{featured-video}{714}
\showcasematerialdesignoutlined{feed}{715}
\showcasematerialdesignoutlined{feedback}{716}
\showcasematerialdesignoutlined{female}{717}
\showcasematerialdesignoutlined{fence}{718}
\showcasematerialdesignoutlined{festival}{719}
\showcasematerialdesignoutlined{fiber-dvr}{720}
\showcasematerialdesignoutlined{fiber-manual-record}{721}
\showcasematerialdesignoutlined{fiber-new}{722}
\showcasematerialdesignoutlined{fiber-pin}{723}
\showcasematerialdesignoutlined{fiber-smart-record}{724}
\showcasematerialdesignoutlined{file-copy}{725}
\showcasematerialdesignoutlined{file-download}{726}
\showcasematerialdesignoutlined{file-download-done}{727}
\showcasematerialdesignoutlined{file-download-off}{728}
\showcasematerialdesignoutlined{file-open}{729}
\showcasematerialdesignoutlined{file-present}{730}
\showcasematerialdesignoutlined{file-upload}{731}
\showcasematerialdesignoutlined{filter}{732}
\showcasematerialdesignoutlined{filter-1}{733}
\showcasematerialdesignoutlined{filter-2}{734}
\showcasematerialdesignoutlined{filter-3}{735}
\showcasematerialdesignoutlined{filter-4}{736}
\showcasematerialdesignoutlined{filter-5}{737}
\showcasematerialdesignoutlined{filter-6}{738}
\showcasematerialdesignoutlined{filter-7}{739}
\showcasematerialdesignoutlined{filter-8}{740}
\showcasematerialdesignoutlined{filter-9}{741}
\showcasematerialdesignoutlined{filter-9-plus}{742}
\showcasematerialdesignoutlined{filter-alt}{743}
\showcasematerialdesignoutlined{filter-alt-off}{744}
\showcasematerialdesignoutlined{filter-b-and-w}{745}
\showcasematerialdesignoutlined{filter-center-focus}{746}
\showcasematerialdesignoutlined{filter-drama}{747}
\showcasematerialdesignoutlined{filter-frames}{748}
\showcasematerialdesignoutlined{filter-hdr}{749}
\showcasematerialdesignoutlined{filter-list}{750}
\showcasematerialdesignoutlined{filter-list-off}{751}
\showcasematerialdesignoutlined{filter-none}{752}
\showcasematerialdesignoutlined{filter-tilt-shift}{753}
\showcasematerialdesignoutlined{filter-vintage}{754}
\showcasematerialdesignoutlined{find-in-page}{755}
\showcasematerialdesignoutlined{find-replace}{756}
\showcasematerialdesignoutlined{fingerprint}{757}
\showcasematerialdesignoutlined{fireplace}{758}
\showcasematerialdesignoutlined{fire-extinguisher}{759}
\showcasematerialdesignoutlined{fire-hydrant-alt}{760}
\showcasematerialdesignoutlined{fire-truck}{761}
\showcasematerialdesignoutlined{first-page}{762}
\showcasematerialdesignoutlined{fitbit}{763}
\showcasematerialdesignoutlined{fitness-center}{764}
\showcasematerialdesignoutlined{fit-screen}{765}
\showcasematerialdesignoutlined{flag}{766}
\showcasematerialdesignoutlined{flag-circle}{767}
\showcasematerialdesignoutlined{flaky}{768}
\showcasematerialdesignoutlined{flare}{769}
\showcasematerialdesignoutlined{flashlight-off}{770}
\showcasematerialdesignoutlined{flashlight-on}{771}
\showcasematerialdesignoutlined{flash-auto}{772}
\showcasematerialdesignoutlined{flash-off}{773}
\showcasematerialdesignoutlined{flash-on}{774}
\showcasematerialdesignoutlined{flatware}{775}
\showcasematerialdesignoutlined{flight}{776}
\showcasematerialdesignoutlined{flight-class}{777}
\showcasematerialdesignoutlined{flight-land}{778}
\showcasematerialdesignoutlined{flight-takeoff}{779}
\showcasematerialdesignoutlined{flip}{780}
\showcasematerialdesignoutlined{flip-camera-android}{781}
\showcasematerialdesignoutlined{flip-camera-ios}{782}
\showcasematerialdesignoutlined{flip-to-back}{783}
\showcasematerialdesignoutlined{flip-to-front}{784}
\showcasematerialdesignoutlined{flood}{785}
\showcasematerialdesignoutlined{fluorescent}{786}
\showcasematerialdesignoutlined{flutter-dash}{787}
\showcasematerialdesignoutlined{fmd-bad}{788}
\showcasematerialdesignoutlined{fmd-good}{789}
\showcasematerialdesignoutlined{folder}{790}
\showcasematerialdesignoutlined{folder-copy}{791}
\showcasematerialdesignoutlined{folder-delete}{792}
\showcasematerialdesignoutlined{folder-off}{793}
\showcasematerialdesignoutlined{folder-open}{794}
\showcasematerialdesignoutlined{folder-shared}{795}
\showcasematerialdesignoutlined{folder-special}{796}
\showcasematerialdesignoutlined{folder-zip}{797}
\showcasematerialdesignoutlined{follow-the-signs}{798}
\showcasematerialdesignoutlined{font-download}{799}
\showcasematerialdesignoutlined{font-download-off}{800}
\showcasematerialdesignoutlined{food-bank}{801}
\showcasematerialdesignoutlined{forest}{802}
\showcasematerialdesignoutlined{fork-left}{803}
\showcasematerialdesignoutlined{fork-right}{804}
\showcasematerialdesignoutlined{format-align-center}{805}
\showcasematerialdesignoutlined{format-align-justify}{806}
\showcasematerialdesignoutlined{format-align-left}{807}
\showcasematerialdesignoutlined{format-align-right}{808}
\showcasematerialdesignoutlined{format-bold}{809}
\showcasematerialdesignoutlined{format-clear}{810}
\showcasematerialdesignoutlined{format-color-fill}{811}
\showcasematerialdesignoutlined{format-color-reset}{812}
\showcasematerialdesignoutlined{format-color-text}{813}
\showcasematerialdesignoutlined{format-indent-decrease}{814}
\showcasematerialdesignoutlined{format-indent-increase}{815}
\showcasematerialdesignoutlined{format-italic}{816}
\showcasematerialdesignoutlined{format-line-spacing}{817}
\showcasematerialdesignoutlined{format-list-bulleted}{818}
\showcasematerialdesignoutlined{format-list-numbered}{819}
\showcasematerialdesignoutlined{format-list-numbered-rtl}{820}
\showcasematerialdesignoutlined{format-overline}{821}
\showcasematerialdesignoutlined{format-paint}{822}
\showcasematerialdesignoutlined{format-quote}{823}
\showcasematerialdesignoutlined{format-shapes}{824}
\showcasematerialdesignoutlined{format-size}{825}
\showcasematerialdesignoutlined{format-strikethrough}{826}
\showcasematerialdesignoutlined{format-textdirection-l-to-r}{827}
\showcasematerialdesignoutlined{format-textdirection-r-to-l}{828}
\showcasematerialdesignoutlined{format-underlined}{829}
\showcasematerialdesignoutlined{fort}{830}
\showcasematerialdesignoutlined{forum}{831}
\showcasematerialdesignoutlined{forward}{832}
\showcasematerialdesignoutlined{forward-5}{833}
\showcasematerialdesignoutlined{forward-10}{834}
\showcasematerialdesignoutlined{forward-30}{835}
\showcasematerialdesignoutlined{forward-to-inbox}{836}
\showcasematerialdesignoutlined{foundation}{837}
\showcasematerialdesignoutlined{free-breakfast}{838}
\showcasematerialdesignoutlined{free-cancellation}{839}
\showcasematerialdesignoutlined{front-hand}{840}
\showcasematerialdesignoutlined{fullscreen}{841}
\showcasematerialdesignoutlined{fullscreen-exit}{842}
\showcasematerialdesignoutlined{functions}{843}
\showcasematerialdesignoutlined{gamepad}{844}
\showcasematerialdesignoutlined{games}{845}
\showcasematerialdesignoutlined{garage}{846}
\showcasematerialdesignoutlined{gas-meter}{847}
\showcasematerialdesignoutlined{gavel}{848}
\showcasematerialdesignoutlined{generating-tokens}{849}
\showcasematerialdesignoutlined{gesture}{850}
\showcasematerialdesignoutlined{get-app}{851}
\showcasematerialdesignoutlined{gif}{852}
\showcasematerialdesignoutlined{gif-box}{853}
\showcasematerialdesignoutlined{girl}{854}
\showcasematerialdesignoutlined{gite}{855}
\showcasematerialdesignoutlined{golf-course}{856}
\showcasematerialdesignoutlined{gpp-bad}{857}
\showcasematerialdesignoutlined{gpp-good}{858}
\showcasematerialdesignoutlined{gpp-maybe}{859}
\showcasematerialdesignoutlined{gps-fixed}{860}
\showcasematerialdesignoutlined{gps-not-fixed}{861}
\showcasematerialdesignoutlined{gps-off}{862}
\showcasematerialdesignoutlined{grade}{863}
\showcasematerialdesignoutlined{gradient}{864}
\showcasematerialdesignoutlined{grading}{865}
\showcasematerialdesignoutlined{grain}{866}
\showcasematerialdesignoutlined{graphic-eq}{867}
\showcasematerialdesignoutlined{grass}{868}
\showcasematerialdesignoutlined{grid-3x3}{869}
\showcasematerialdesignoutlined{grid-4x4}{870}
\showcasematerialdesignoutlined{grid-goldenratio}{871}
\showcasematerialdesignoutlined{grid-off}{872}
\showcasematerialdesignoutlined{grid-on}{873}
\showcasematerialdesignoutlined{grid-view}{874}
\showcasematerialdesignoutlined{group}{875}
\showcasematerialdesignoutlined{groups}{876}
\showcasematerialdesignoutlined{groups-2}{877}
\showcasematerialdesignoutlined{groups-3}{878}
\showcasematerialdesignoutlined{group-add}{879}
\showcasematerialdesignoutlined{group-off}{880}
\showcasematerialdesignoutlined{group-remove}{881}
\showcasematerialdesignoutlined{group-work}{882}
\showcasematerialdesignoutlined{g-mobiledata}{883}
\showcasematerialdesignoutlined{g-translate}{884}
\showcasematerialdesignoutlined{hail}{885}
\showcasematerialdesignoutlined{handshake}{886}
\showcasematerialdesignoutlined{handyman}{887}
\showcasematerialdesignoutlined{hardware}{888}
\showcasematerialdesignoutlined{hd}{889}
\showcasematerialdesignoutlined{hdr-auto}{890}
\showcasematerialdesignoutlined{hdr-auto-select}{891}
\showcasematerialdesignoutlined{hdr-enhanced-select}{892}
\showcasematerialdesignoutlined{hdr-off}{893}
\showcasematerialdesignoutlined{hdr-off-select}{894}
\showcasematerialdesignoutlined{hdr-on}{895}
\showcasematerialdesignoutlined{hdr-on-select}{896}
\showcasematerialdesignoutlined{hdr-plus}{897}
\showcasematerialdesignoutlined{hdr-strong}{898}
\showcasematerialdesignoutlined{hdr-weak}{899}
\showcasematerialdesignoutlined{headphones}{900}
\showcasematerialdesignoutlined{headphones-battery}{901}
\showcasematerialdesignoutlined{headset}{902}
\showcasematerialdesignoutlined{headset-mic}{903}
\showcasematerialdesignoutlined{headset-off}{904}
\showcasematerialdesignoutlined{healing}{905}
\showcasematerialdesignoutlined{health-and-safety}{906}
\showcasematerialdesignoutlined{hearing}{907}
\showcasematerialdesignoutlined{hearing-disabled}{908}
\showcasematerialdesignoutlined{heart-broken}{909}
\showcasematerialdesignoutlined{heat-pump}{910}
\showcasematerialdesignoutlined{height}{911}
\showcasematerialdesignoutlined{help}{912}
\showcasematerialdesignoutlined{help-center}{913}
\showcasematerialdesignoutlined{help-outline}{914}
\showcasematerialdesignoutlined{hevc}{915}
\showcasematerialdesignoutlined{hexagon}{916}
\showcasematerialdesignoutlined{hide-image}{917}
\showcasematerialdesignoutlined{hide-source}{918}
\showcasematerialdesignoutlined{highlight}{919}
\showcasematerialdesignoutlined{highlight-alt}{920}
\showcasematerialdesignoutlined{highlight-off}{921}
\showcasematerialdesignoutlined{high-quality}{922}
\showcasematerialdesignoutlined{hiking}{923}
\showcasematerialdesignoutlined{history}{924}
\showcasematerialdesignoutlined{history-edu}{925}
\showcasematerialdesignoutlined{history-toggle-off}{926}
\showcasematerialdesignoutlined{hive}{927}
\showcasematerialdesignoutlined{hls}{928}
\showcasematerialdesignoutlined{hls-off}{929}
\showcasematerialdesignoutlined{holiday-village}{930}
\showcasematerialdesignoutlined{home}{931}
\showcasematerialdesignoutlined{home-max}{932}
\showcasematerialdesignoutlined{home-mini}{933}
\showcasematerialdesignoutlined{home-repair-service}{934}
\showcasematerialdesignoutlined{home-work}{935}
\showcasematerialdesignoutlined{horizontal-distribute}{936}
\showcasematerialdesignoutlined{horizontal-rule}{937}
\showcasematerialdesignoutlined{horizontal-split}{938}
\showcasematerialdesignoutlined{hotel}{939}
\showcasematerialdesignoutlined{hotel-class}{940}
\showcasematerialdesignoutlined{hot-tub}{941}
\showcasematerialdesignoutlined{hourglass-bottom}{942}
\showcasematerialdesignoutlined{hourglass-disabled}{943}
\showcasematerialdesignoutlined{hourglass-empty}{944}
\showcasematerialdesignoutlined{hourglass-full}{945}
\showcasematerialdesignoutlined{hourglass-top}{946}
\showcasematerialdesignoutlined{house}{947}
\showcasematerialdesignoutlined{houseboat}{948}
\showcasematerialdesignoutlined{house-siding}{949}
\showcasematerialdesignoutlined{how-to-reg}{950}
\showcasematerialdesignoutlined{how-to-vote}{951}
\showcasematerialdesignoutlined{html}{952}
\showcasematerialdesignoutlined{http}{953}
\showcasematerialdesignoutlined{https}{954}
\showcasematerialdesignoutlined{hub}{955}
\showcasematerialdesignoutlined{hvac}{956}
\showcasematerialdesignoutlined{h-mobiledata}{957}
\showcasematerialdesignoutlined{h-plus-mobiledata}{958}
\showcasematerialdesignoutlined{icecream}{959}
\showcasematerialdesignoutlined{ice-skating}{960}
\showcasematerialdesignoutlined{image}{961}
\showcasematerialdesignoutlined{imagesearch-roller}{962}
\showcasematerialdesignoutlined{image-aspect-ratio}{963}
\showcasematerialdesignoutlined{image-not-supported}{964}
\showcasematerialdesignoutlined{image-search}{965}
\showcasematerialdesignoutlined{important-devices}{966}
\showcasematerialdesignoutlined{import-contacts}{967}
\showcasematerialdesignoutlined{import-export}{968}
\showcasematerialdesignoutlined{inbox}{969}
\showcasematerialdesignoutlined{incomplete-circle}{970}
\showcasematerialdesignoutlined{indeterminate-check-box}{971}
\showcasematerialdesignoutlined{info}{972}
\showcasematerialdesignoutlined{input}{973}
\showcasematerialdesignoutlined{insert-chart}{974}
\showcasematerialdesignoutlined{insert-chart-outlined}{975}
\showcasematerialdesignoutlined{insert-comment}{976}
\showcasematerialdesignoutlined{insert-drive-file}{977}
\showcasematerialdesignoutlined{insert-emoticon}{978}
\showcasematerialdesignoutlined{insert-invitation}{979}
\showcasematerialdesignoutlined{insert-link}{980}
\showcasematerialdesignoutlined{insert-page-break}{981}
\showcasematerialdesignoutlined{insert-photo}{982}
\showcasematerialdesignoutlined{insights}{983}
\showcasematerialdesignoutlined{install-desktop}{984}
\showcasematerialdesignoutlined{install-mobile}{985}
\showcasematerialdesignoutlined{integration-instructions}{986}
\showcasematerialdesignoutlined{interests}{987}
\showcasematerialdesignoutlined{interpreter-mode}{988}
\showcasematerialdesignoutlined{inventory}{989}
\showcasematerialdesignoutlined{inventory-2}{990}
\showcasematerialdesignoutlined{invert-colors}{991}
\showcasematerialdesignoutlined{invert-colors-off}{992}
\showcasematerialdesignoutlined{ios-share}{993}
\showcasematerialdesignoutlined{iron}{994}
\showcasematerialdesignoutlined{iso}{995}
\showcasematerialdesignoutlined{javascript}{996}
\showcasematerialdesignoutlined{join-full}{997}
\showcasematerialdesignoutlined{join-inner}{998}
\showcasematerialdesignoutlined{join-left}{999}
\showcasematerialdesignoutlined{join-right}{1000}
\showcasematerialdesignoutlined{kayaking}{1001}
\showcasematerialdesignoutlined{kebab-dining}{1002}
\showcasematerialdesignoutlined{key}{1003}
\showcasematerialdesignoutlined{keyboard}{1004}
\showcasematerialdesignoutlined{keyboard-alt}{1005}
\showcasematerialdesignoutlined{keyboard-arrow-down}{1006}
\showcasematerialdesignoutlined{keyboard-arrow-left}{1007}
\showcasematerialdesignoutlined{keyboard-arrow-right}{1008}
\showcasematerialdesignoutlined{keyboard-arrow-up}{1009}
\showcasematerialdesignoutlined{keyboard-backspace}{1010}
\showcasematerialdesignoutlined{keyboard-capslock}{1011}
\showcasematerialdesignoutlined{keyboard-command-key}{1012}
\showcasematerialdesignoutlined{keyboard-control-key}{1013}
\showcasematerialdesignoutlined{keyboard-double-arrow-down}{1014}
\showcasematerialdesignoutlined{keyboard-double-arrow-left}{1015}
\showcasematerialdesignoutlined{keyboard-double-arrow-right}{1016}
\showcasematerialdesignoutlined{keyboard-double-arrow-up}{1017}
\showcasematerialdesignoutlined{keyboard-hide}{1018}
\showcasematerialdesignoutlined{keyboard-option-key}{1019}
\showcasematerialdesignoutlined{keyboard-return}{1020}
\showcasematerialdesignoutlined{keyboard-tab}{1021}
\showcasematerialdesignoutlined{keyboard-voice}{1022}
\showcasematerialdesignoutlined{key-off}{1023}
\showcasematerialdesignoutlined{king-bed}{1024}
\showcasematerialdesignoutlined{kitchen}{1025}
\showcasematerialdesignoutlined{kitesurfing}{1026}
\showcasematerialdesignoutlined{label}{1027}
\showcasematerialdesignoutlined{label-important}{1028}
\showcasematerialdesignoutlined{label-off}{1029}
\showcasematerialdesignoutlined{lan}{1030}
\showcasematerialdesignoutlined{landscape}{1031}
\showcasematerialdesignoutlined{landslide}{1032}
\showcasematerialdesignoutlined{language}{1033}
\showcasematerialdesignoutlined{laptop}{1034}
\showcasematerialdesignoutlined{laptop-chromebook}{1035}
\showcasematerialdesignoutlined{laptop-mac}{1036}
\showcasematerialdesignoutlined{laptop-windows}{1037}
\showcasematerialdesignoutlined{last-page}{1038}
\showcasematerialdesignoutlined{launch}{1039}
\showcasematerialdesignoutlined{layers}{1040}
\showcasematerialdesignoutlined{layers-clear}{1041}
\showcasematerialdesignoutlined{leaderboard}{1042}
\showcasematerialdesignoutlined{leak-add}{1043}
\showcasematerialdesignoutlined{leak-remove}{1044}
\showcasematerialdesignoutlined{legend-toggle}{1045}
\showcasematerialdesignoutlined{lens}{1046}
\showcasematerialdesignoutlined{lens-blur}{1047}
\showcasematerialdesignoutlined{library-add}{1048}
\showcasematerialdesignoutlined{library-add-check}{1049}
\showcasematerialdesignoutlined{library-books}{1050}
\showcasematerialdesignoutlined{library-music}{1051}
\showcasematerialdesignoutlined{light}{1052}
\showcasematerialdesignoutlined{lightbulb}{1053}
\showcasematerialdesignoutlined{lightbulb-circle}{1054}
\showcasematerialdesignoutlined{light-mode}{1055}
\showcasematerialdesignoutlined{linear-scale}{1056}
\showcasematerialdesignoutlined{line-axis}{1057}
\showcasematerialdesignoutlined{line-style}{1058}
\showcasematerialdesignoutlined{line-weight}{1059}
\showcasematerialdesignoutlined{link}{1060}
\showcasematerialdesignoutlined{linked-camera}{1061}
\showcasematerialdesignoutlined{link-off}{1062}
\showcasematerialdesignoutlined{liquor}{1063}
\showcasematerialdesignoutlined{list}{1064}
\showcasematerialdesignoutlined{list-alt}{1065}
\showcasematerialdesignoutlined{live-help}{1066}
\showcasematerialdesignoutlined{live-tv}{1067}
\showcasematerialdesignoutlined{living}{1068}
\showcasematerialdesignoutlined{local-activity}{1069}
\showcasematerialdesignoutlined{local-airport}{1070}
\showcasematerialdesignoutlined{local-atm}{1071}
\showcasematerialdesignoutlined{local-bar}{1072}
\showcasematerialdesignoutlined{local-cafe}{1073}
\showcasematerialdesignoutlined{local-car-wash}{1074}
\showcasematerialdesignoutlined{local-convenience-store}{1075}
\showcasematerialdesignoutlined{local-dining}{1076}
\showcasematerialdesignoutlined{local-drink}{1077}
\showcasematerialdesignoutlined{local-fire-department}{1078}
\showcasematerialdesignoutlined{local-florist}{1079}
\showcasematerialdesignoutlined{local-gas-station}{1080}
\showcasematerialdesignoutlined{local-grocery-store}{1081}
\showcasematerialdesignoutlined{local-hospital}{1082}
\showcasematerialdesignoutlined{local-hotel}{1083}
\showcasematerialdesignoutlined{local-laundry-service}{1084}
\showcasematerialdesignoutlined{local-library}{1085}
\showcasematerialdesignoutlined{local-mall}{1086}
\showcasematerialdesignoutlined{local-movies}{1087}
\showcasematerialdesignoutlined{local-offer}{1088}
\showcasematerialdesignoutlined{local-parking}{1089}
\showcasematerialdesignoutlined{local-pharmacy}{1090}
\showcasematerialdesignoutlined{local-phone}{1091}
\showcasematerialdesignoutlined{local-pizza}{1092}
\showcasematerialdesignoutlined{local-play}{1093}
\showcasematerialdesignoutlined{local-police}{1094}
\showcasematerialdesignoutlined{local-post-office}{1095}
\showcasematerialdesignoutlined{local-printshop}{1096}
\showcasematerialdesignoutlined{local-see}{1097}
\showcasematerialdesignoutlined{local-shipping}{1098}
\showcasematerialdesignoutlined{local-taxi}{1099}
\showcasematerialdesignoutlined{location-city}{1100}
\showcasematerialdesignoutlined{location-disabled}{1101}
\showcasematerialdesignoutlined{location-off}{1102}
\showcasematerialdesignoutlined{location-on}{1103}
\showcasematerialdesignoutlined{location-searching}{1104}
\showcasematerialdesignoutlined{lock}{1105}
\showcasematerialdesignoutlined{lock-clock}{1106}
\showcasematerialdesignoutlined{lock-open}{1107}
\showcasematerialdesignoutlined{lock-person}{1108}
\showcasematerialdesignoutlined{lock-reset}{1109}
\showcasematerialdesignoutlined{login}{1110}
\showcasematerialdesignoutlined{logout}{1111}
\showcasematerialdesignoutlined{logo-dev}{1112}
\showcasematerialdesignoutlined{looks}{1113}
\showcasematerialdesignoutlined{looks-3}{1114}
\showcasematerialdesignoutlined{looks-4}{1115}
\showcasematerialdesignoutlined{looks-5}{1116}
\showcasematerialdesignoutlined{looks-6}{1117}
\showcasematerialdesignoutlined{looks-one}{1118}
\showcasematerialdesignoutlined{looks-two}{1119}
\showcasematerialdesignoutlined{loop}{1120}
\showcasematerialdesignoutlined{loupe}{1121}
\showcasematerialdesignoutlined{low-priority}{1122}
\showcasematerialdesignoutlined{loyalty}{1123}
\showcasematerialdesignoutlined{lte-mobiledata}{1124}
\showcasematerialdesignoutlined{lte-plus-mobiledata}{1125}
\showcasematerialdesignoutlined{luggage}{1126}
\showcasematerialdesignoutlined{lunch-dining}{1127}
\showcasematerialdesignoutlined{lyrics}{1128}
\showcasematerialdesignoutlined{macro-off}{1129}
\showcasematerialdesignoutlined{mail}{1130}
\showcasematerialdesignoutlined{mail-lock}{1131}
\showcasematerialdesignoutlined{mail-outline}{1132}
\showcasematerialdesignoutlined{male}{1133}
\showcasematerialdesignoutlined{man}{1134}
\showcasematerialdesignoutlined{manage-accounts}{1135}
\showcasematerialdesignoutlined{manage-history}{1136}
\showcasematerialdesignoutlined{manage-search}{1137}
\showcasematerialdesignoutlined{man-2}{1138}
\showcasematerialdesignoutlined{man-3}{1139}
\showcasematerialdesignoutlined{man-4}{1140}
\showcasematerialdesignoutlined{map}{1141}
\showcasematerialdesignoutlined{maps-home-work}{1142}
\showcasematerialdesignoutlined{maps-ugc}{1143}
\showcasematerialdesignoutlined{margin}{1144}
\showcasematerialdesignoutlined{markunread}{1145}
\showcasematerialdesignoutlined{markunread-mailbox}{1146}
\showcasematerialdesignoutlined{mark-as-unread}{1147}
\showcasematerialdesignoutlined{mark-chat-read}{1148}
\showcasematerialdesignoutlined{mark-chat-unread}{1149}
\showcasematerialdesignoutlined{mark-email-read}{1150}
\showcasematerialdesignoutlined{mark-email-unread}{1151}
\showcasematerialdesignoutlined{mark-unread-chat-alt}{1152}
\showcasematerialdesignoutlined{masks}{1153}
\showcasematerialdesignoutlined{maximize}{1154}
\showcasematerialdesignoutlined{mediation}{1155}
\showcasematerialdesignoutlined{media-bluetooth-off}{1156}
\showcasematerialdesignoutlined{media-bluetooth-on}{1157}
\showcasematerialdesignoutlined{medical-information}{1158}
\showcasematerialdesignoutlined{medical-services}{1159}
\showcasematerialdesignoutlined{medication}{1160}
\showcasematerialdesignoutlined{medication-liquid}{1161}
\showcasematerialdesignoutlined{meeting-room}{1162}
\showcasematerialdesignoutlined{memory}{1163}
\showcasematerialdesignoutlined{menu}{1164}
\showcasematerialdesignoutlined{menu-book}{1165}
\showcasematerialdesignoutlined{menu-open}{1166}
\showcasematerialdesignoutlined{merge}{1167}
\showcasematerialdesignoutlined{merge-type}{1168}
\showcasematerialdesignoutlined{message}{1169}
\showcasematerialdesignoutlined{mic}{1170}
\showcasematerialdesignoutlined{microwave}{1171}
\showcasematerialdesignoutlined{mic-external-off}{1172}
\showcasematerialdesignoutlined{mic-external-on}{1173}
\showcasematerialdesignoutlined{mic-none}{1174}
\showcasematerialdesignoutlined{mic-off}{1175}
\showcasematerialdesignoutlined{military-tech}{1176}
\showcasematerialdesignoutlined{minimize}{1177}
\showcasematerialdesignoutlined{minor-crash}{1178}
\showcasematerialdesignoutlined{miscellaneous-services}{1179}
\showcasematerialdesignoutlined{missed-video-call}{1180}
\showcasematerialdesignoutlined{mms}{1181}
\showcasematerialdesignoutlined{mobiledata-off}{1182}
\showcasematerialdesignoutlined{mobile-friendly}{1183}
\showcasematerialdesignoutlined{mobile-off}{1184}
\showcasematerialdesignoutlined{mobile-screen-share}{1185}
\showcasematerialdesignoutlined{mode}{1186}
\showcasematerialdesignoutlined{model-training}{1187}
\showcasematerialdesignoutlined{mode-comment}{1188}
\showcasematerialdesignoutlined{mode-edit}{1189}
\showcasematerialdesignoutlined{mode-edit-outline}{1190}
\showcasematerialdesignoutlined{mode-fan-off}{1191}
\showcasematerialdesignoutlined{mode-night}{1192}
\showcasematerialdesignoutlined{mode-of-travel}{1193}
\showcasematerialdesignoutlined{mode-standby}{1194}
\showcasematerialdesignoutlined{monetization-on}{1195}
\showcasematerialdesignoutlined{money}{1196}
\showcasematerialdesignoutlined{money-off}{1197}
\showcasematerialdesignoutlined{money-off-csred}{1198}
\showcasematerialdesignoutlined{monitor}{1199}
\showcasematerialdesignoutlined{monitor-heart}{1200}
\showcasematerialdesignoutlined{monitor-weight}{1201}
\showcasematerialdesignoutlined{monochrome-photos}{1202}
\showcasematerialdesignoutlined{mood}{1203}
\showcasematerialdesignoutlined{mood-bad}{1204}
\showcasematerialdesignoutlined{moped}{1205}
\showcasematerialdesignoutlined{more}{1206}
\showcasematerialdesignoutlined{more-horiz}{1207}
\showcasematerialdesignoutlined{more-time}{1208}
\showcasematerialdesignoutlined{more-vert}{1209}
\showcasematerialdesignoutlined{mosque}{1210}
\showcasematerialdesignoutlined{motion-photos-auto}{1211}
\showcasematerialdesignoutlined{motion-photos-off}{1212}
\showcasematerialdesignoutlined{motion-photos-on}{1213}
\showcasematerialdesignoutlined{motion-photos-pause}{1214}
\showcasematerialdesignoutlined{motion-photos-paused}{1215}
\showcasematerialdesignoutlined{mouse}{1216}
\showcasematerialdesignoutlined{move-down}{1217}
\showcasematerialdesignoutlined{move-to-inbox}{1218}
\showcasematerialdesignoutlined{move-up}{1219}
\showcasematerialdesignoutlined{movie}{1220}
\showcasematerialdesignoutlined{movie-creation}{1221}
\showcasematerialdesignoutlined{movie-filter}{1222}
\showcasematerialdesignoutlined{moving}{1223}
\showcasematerialdesignoutlined{mp}{1224}
\showcasematerialdesignoutlined{multiline-chart}{1225}
\showcasematerialdesignoutlined{multiple-stop}{1226}
\showcasematerialdesignoutlined{museum}{1227}
\showcasematerialdesignoutlined{music-note}{1228}
\showcasematerialdesignoutlined{music-off}{1229}
\showcasematerialdesignoutlined{music-video}{1230}
\showcasematerialdesignoutlined{my-location}{1231}
\showcasematerialdesignoutlined{nat}{1232}
\showcasematerialdesignoutlined{nature}{1233}
\showcasematerialdesignoutlined{nature-people}{1234}
\showcasematerialdesignoutlined{navigate-before}{1235}
\showcasematerialdesignoutlined{navigate-next}{1236}
\showcasematerialdesignoutlined{navigation}{1237}
\showcasematerialdesignoutlined{nearby-error}{1238}
\showcasematerialdesignoutlined{nearby-off}{1239}
\showcasematerialdesignoutlined{near-me}{1240}
\showcasematerialdesignoutlined{near-me-disabled}{1241}
\showcasematerialdesignoutlined{nest-cam-wired-stand}{1242}
\showcasematerialdesignoutlined{network-cell}{1243}
\showcasematerialdesignoutlined{network-check}{1244}
\showcasematerialdesignoutlined{network-locked}{1245}
\showcasematerialdesignoutlined{network-ping}{1246}
\showcasematerialdesignoutlined{network-wifi}{1247}
\showcasematerialdesignoutlined{network-wifi-1-bar}{1248}
\showcasematerialdesignoutlined{network-wifi-2-bar}{1249}
\showcasematerialdesignoutlined{network-wifi-3-bar}{1250}
\showcasematerialdesignoutlined{newspaper}{1251}
\showcasematerialdesignoutlined{new-label}{1252}
\showcasematerialdesignoutlined{new-releases}{1253}
\showcasematerialdesignoutlined{next-plan}{1254}
\showcasematerialdesignoutlined{next-week}{1255}
\showcasematerialdesignoutlined{nfc}{1256}
\showcasematerialdesignoutlined{nightlife}{1257}
\showcasematerialdesignoutlined{nightlight}{1258}
\showcasematerialdesignoutlined{nightlight-round}{1259}
\showcasematerialdesignoutlined{nights-stay}{1260}
\showcasematerialdesignoutlined{night-shelter}{1261}
\showcasematerialdesignoutlined{noise-aware}{1262}
\showcasematerialdesignoutlined{noise-control-off}{1263}
\showcasematerialdesignoutlined{nordic-walking}{1264}
\showcasematerialdesignoutlined{north}{1265}
\showcasematerialdesignoutlined{north-east}{1266}
\showcasematerialdesignoutlined{north-west}{1267}
\showcasematerialdesignoutlined{note}{1268}
\showcasematerialdesignoutlined{notes}{1269}
\showcasematerialdesignoutlined{note-add}{1270}
\showcasematerialdesignoutlined{note-alt}{1271}
\showcasematerialdesignoutlined{notifications}{1272}
\showcasematerialdesignoutlined{notifications-active}{1273}
\showcasematerialdesignoutlined{notifications-none}{1274}
\showcasematerialdesignoutlined{notifications-off}{1275}
\showcasematerialdesignoutlined{notifications-paused}{1276}
\showcasematerialdesignoutlined{notification-add}{1277}
\showcasematerialdesignoutlined{notification-important}{1278}
\showcasematerialdesignoutlined{not-accessible}{1279}
\showcasematerialdesignoutlined{not-interested}{1280}
\showcasematerialdesignoutlined{not-listed-location}{1281}
\showcasematerialdesignoutlined{not-started}{1282}
\showcasematerialdesignoutlined{no-accounts}{1283}
\showcasematerialdesignoutlined{no-adult-content}{1284}
\showcasematerialdesignoutlined{no-backpack}{1285}
\showcasematerialdesignoutlined{no-cell}{1286}
\showcasematerialdesignoutlined{no-crash}{1287}
\showcasematerialdesignoutlined{no-drinks}{1288}
\showcasematerialdesignoutlined{no-encryption}{1289}
\showcasematerialdesignoutlined{no-encryption-gmailerrorred}{1290}
\showcasematerialdesignoutlined{no-flash}{1291}
\showcasematerialdesignoutlined{no-food}{1292}
\showcasematerialdesignoutlined{no-luggage}{1293}
\showcasematerialdesignoutlined{no-meals}{1294}
\showcasematerialdesignoutlined{no-meeting-room}{1295}
\showcasematerialdesignoutlined{no-photography}{1296}
\showcasematerialdesignoutlined{no-sim}{1297}
\showcasematerialdesignoutlined{no-stroller}{1298}
\showcasematerialdesignoutlined{no-transfer}{1299}
\showcasematerialdesignoutlined{numbers}{1300}
\showcasematerialdesignoutlined{offline-bolt}{1301}
\showcasematerialdesignoutlined{offline-pin}{1302}
\showcasematerialdesignoutlined{offline-share}{1303}
\showcasematerialdesignoutlined{oil-barrel}{1304}
\showcasematerialdesignoutlined{ondemand-video}{1305}
\showcasematerialdesignoutlined{online-prediction}{1306}
\showcasematerialdesignoutlined{on-device-training}{1307}
\showcasematerialdesignoutlined{opacity}{1308}
\showcasematerialdesignoutlined{open-in-browser}{1309}
\showcasematerialdesignoutlined{open-in-full}{1310}
\showcasematerialdesignoutlined{open-in-new}{1311}
\showcasematerialdesignoutlined{open-in-new-off}{1312}
\showcasematerialdesignoutlined{open-with}{1313}
\showcasematerialdesignoutlined{other-houses}{1314}
\showcasematerialdesignoutlined{outbound}{1315}
\showcasematerialdesignoutlined{outbox}{1316}
\showcasematerialdesignoutlined{outdoor-grill}{1317}
\showcasematerialdesignoutlined{outlet}{1318}
\showcasematerialdesignoutlined{outlined-flag}{1319}
\showcasematerialdesignoutlined{output}{1320}
\showcasematerialdesignoutlined{padding}{1321}
\showcasematerialdesignoutlined{pages}{1322}
\showcasematerialdesignoutlined{pageview}{1323}
\showcasematerialdesignoutlined{paid}{1324}
\showcasematerialdesignoutlined{palette}{1325}
\showcasematerialdesignoutlined{panorama}{1326}
\showcasematerialdesignoutlined{panorama-fish-eye}{1327}
\showcasematerialdesignoutlined{panorama-horizontal}{1328}
\showcasematerialdesignoutlined{panorama-horizontal-select}{1329}
\showcasematerialdesignoutlined{panorama-photosphere}{1330}
\showcasematerialdesignoutlined{panorama-photosphere-select}{1331}
\showcasematerialdesignoutlined{panorama-vertical}{1332}
\showcasematerialdesignoutlined{panorama-vertical-select}{1333}
\showcasematerialdesignoutlined{panorama-wide-angle}{1334}
\showcasematerialdesignoutlined{panorama-wide-angle-select}{1335}
\showcasematerialdesignoutlined{pan-tool}{1336}
\showcasematerialdesignoutlined{pan-tool-alt}{1337}
\showcasematerialdesignoutlined{paragliding}{1338}
\showcasematerialdesignoutlined{park}{1339}
\showcasematerialdesignoutlined{party-mode}{1340}
\showcasematerialdesignoutlined{password}{1341}
\showcasematerialdesignoutlined{pattern}{1342}
\showcasematerialdesignoutlined{pause}{1343}
\showcasematerialdesignoutlined{pause-circle}{1344}
\showcasematerialdesignoutlined{pause-circle-filled}{1345}
\showcasematerialdesignoutlined{pause-circle-outline}{1346}
\showcasematerialdesignoutlined{pause-presentation}{1347}
\showcasematerialdesignoutlined{payment}{1348}
\showcasematerialdesignoutlined{payments}{1349}
\showcasematerialdesignoutlined{pedal-bike}{1350}
\showcasematerialdesignoutlined{pending}{1351}
\showcasematerialdesignoutlined{pending-actions}{1352}
\showcasematerialdesignoutlined{pentagon}{1353}
\showcasematerialdesignoutlined{people}{1354}
\showcasematerialdesignoutlined{people-alt}{1355}
\showcasematerialdesignoutlined{people-outline}{1356}
\showcasematerialdesignoutlined{percent}{1357}
\showcasematerialdesignoutlined{perm-camera-mic}{1358}
\showcasematerialdesignoutlined{perm-contact-calendar}{1359}
\showcasematerialdesignoutlined{perm-data-setting}{1360}
\showcasematerialdesignoutlined{perm-device-information}{1361}
\showcasematerialdesignoutlined{perm-identity}{1362}
\showcasematerialdesignoutlined{perm-media}{1363}
\showcasematerialdesignoutlined{perm-phone-msg}{1364}
\showcasematerialdesignoutlined{perm-scan-wifi}{1365}
\showcasematerialdesignoutlined{person}{1366}
\showcasematerialdesignoutlined{personal-injury}{1367}
\showcasematerialdesignoutlined{personal-video}{1368}
\showcasematerialdesignoutlined{person-2}{1369}
\showcasematerialdesignoutlined{person-3}{1370}
\showcasematerialdesignoutlined{person-4}{1371}
\showcasematerialdesignoutlined{person-add}{1372}
\showcasematerialdesignoutlined{person-add-alt}{1373}
\showcasematerialdesignoutlined{person-add-alt-1}{1374}
\showcasematerialdesignoutlined{person-add-disabled}{1375}
\showcasematerialdesignoutlined{person-off}{1376}
\showcasematerialdesignoutlined{person-outline}{1377}
\showcasematerialdesignoutlined{person-pin}{1378}
\showcasematerialdesignoutlined{person-pin-circle}{1379}
\showcasematerialdesignoutlined{person-remove}{1380}
\showcasematerialdesignoutlined{person-remove-alt-1}{1381}
\showcasematerialdesignoutlined{person-search}{1382}
\showcasematerialdesignoutlined{pest-control}{1383}
\showcasematerialdesignoutlined{pest-control-rodent}{1384}
\showcasematerialdesignoutlined{pets}{1385}
\showcasematerialdesignoutlined{phishing}{1386}
\showcasematerialdesignoutlined{phone}{1387}
\showcasematerialdesignoutlined{phonelink}{1388}
\showcasematerialdesignoutlined{phonelink-erase}{1389}
\showcasematerialdesignoutlined{phonelink-lock}{1390}
\showcasematerialdesignoutlined{phonelink-off}{1391}
\showcasematerialdesignoutlined{phonelink-ring}{1392}
\showcasematerialdesignoutlined{phonelink-setup}{1393}
\showcasematerialdesignoutlined{phone-android}{1394}
\showcasematerialdesignoutlined{phone-bluetooth-speaker}{1395}
\showcasematerialdesignoutlined{phone-callback}{1396}
\showcasematerialdesignoutlined{phone-disabled}{1397}
\showcasematerialdesignoutlined{phone-enabled}{1398}
\showcasematerialdesignoutlined{phone-forwarded}{1399}
\showcasematerialdesignoutlined{phone-iphone}{1400}
\showcasematerialdesignoutlined{phone-locked}{1401}
\showcasematerialdesignoutlined{phone-missed}{1402}
\showcasematerialdesignoutlined{phone-paused}{1403}
\showcasematerialdesignoutlined{photo}{1404}
\showcasematerialdesignoutlined{photo-album}{1405}
\showcasematerialdesignoutlined{photo-camera}{1406}
\showcasematerialdesignoutlined{photo-camera-back}{1407}
\showcasematerialdesignoutlined{photo-camera-front}{1408}
\showcasematerialdesignoutlined{photo-filter}{1409}
\showcasematerialdesignoutlined{photo-library}{1410}
\showcasematerialdesignoutlined{photo-size-select-actual}{1411}
\showcasematerialdesignoutlined{photo-size-select-large}{1412}
\showcasematerialdesignoutlined{photo-size-select-small}{1413}
\showcasematerialdesignoutlined{php}{1414}
\showcasematerialdesignoutlined{piano}{1415}
\showcasematerialdesignoutlined{piano-off}{1416}
\showcasematerialdesignoutlined{picture-as-pdf}{1417}
\showcasematerialdesignoutlined{picture-in-picture}{1418}
\showcasematerialdesignoutlined{picture-in-picture-alt}{1419}
\showcasematerialdesignoutlined{pie-chart}{1420}
\showcasematerialdesignoutlined{pie-chart-outline}{1421}
\showcasematerialdesignoutlined{pin}{1422}
\showcasematerialdesignoutlined{pinch}{1423}
\showcasematerialdesignoutlined{pin-drop}{1424}
\showcasematerialdesignoutlined{pin-end}{1425}
\showcasematerialdesignoutlined{pin-invoke}{1426}
\showcasematerialdesignoutlined{pivot-table-chart}{1427}
\showcasematerialdesignoutlined{pix}{1428}
\showcasematerialdesignoutlined{place}{1429}
\showcasematerialdesignoutlined{plagiarism}{1430}
\showcasematerialdesignoutlined{playlist-add}{1431}
\showcasematerialdesignoutlined{playlist-add-check}{1432}
\showcasematerialdesignoutlined{playlist-add-check-circle}{1433}
\showcasematerialdesignoutlined{playlist-add-circle}{1434}
\showcasematerialdesignoutlined{playlist-play}{1435}
\showcasematerialdesignoutlined{playlist-remove}{1436}
\showcasematerialdesignoutlined{play-arrow}{1437}
\showcasematerialdesignoutlined{play-circle}{1438}
\showcasematerialdesignoutlined{play-circle-filled}{1439}
\showcasematerialdesignoutlined{play-circle-outline}{1440}
\showcasematerialdesignoutlined{play-disabled}{1441}
\showcasematerialdesignoutlined{play-for-work}{1442}
\showcasematerialdesignoutlined{play-lesson}{1443}
\showcasematerialdesignoutlined{plumbing}{1444}
\showcasematerialdesignoutlined{plus-one}{1445}
\showcasematerialdesignoutlined{podcasts}{1446}
\showcasematerialdesignoutlined{point-of-sale}{1447}
\showcasematerialdesignoutlined{policy}{1448}
\showcasematerialdesignoutlined{poll}{1449}
\showcasematerialdesignoutlined{polyline}{1450}
\showcasematerialdesignoutlined{polymer}{1451}
\showcasematerialdesignoutlined{pool}{1452}
\showcasematerialdesignoutlined{portable-wifi-off}{1453}
\showcasematerialdesignoutlined{portrait}{1454}
\showcasematerialdesignoutlined{post-add}{1455}
\showcasematerialdesignoutlined{power}{1456}
\showcasematerialdesignoutlined{power-input}{1457}
\showcasematerialdesignoutlined{power-off}{1458}
\showcasematerialdesignoutlined{power-settings-new}{1459}
\showcasematerialdesignoutlined{precision-manufacturing}{1460}
\showcasematerialdesignoutlined{pregnant-woman}{1461}
\showcasematerialdesignoutlined{present-to-all}{1462}
\showcasematerialdesignoutlined{preview}{1463}
\showcasematerialdesignoutlined{price-change}{1464}
\showcasematerialdesignoutlined{price-check}{1465}
\showcasematerialdesignoutlined{print}{1466}
\showcasematerialdesignoutlined{print-disabled}{1467}
\showcasematerialdesignoutlined{priority-high}{1468}
\showcasematerialdesignoutlined{privacy-tip}{1469}
\showcasematerialdesignoutlined{private-connectivity}{1470}
\showcasematerialdesignoutlined{production-quantity-limits}{1471}
\showcasematerialdesignoutlined{propane}{1472}
\showcasematerialdesignoutlined{propane-tank}{1473}
\showcasematerialdesignoutlined{psychology}{1474}
\showcasematerialdesignoutlined{psychology-alt}{1475}
\showcasematerialdesignoutlined{public}{1476}
\showcasematerialdesignoutlined{public-off}{1477}
\showcasematerialdesignoutlined{publish}{1478}
\showcasematerialdesignoutlined{published-with-changes}{1479}
\showcasematerialdesignoutlined{punch-clock}{1480}
\showcasematerialdesignoutlined{push-pin}{1481}
\showcasematerialdesignoutlined{qr-code}{1482}
\showcasematerialdesignoutlined{qr-code-2}{1483}
\showcasematerialdesignoutlined{qr-code-scanner}{1484}
\showcasematerialdesignoutlined{query-builder}{1485}
\showcasematerialdesignoutlined{query-stats}{1486}
\showcasematerialdesignoutlined{question-answer}{1487}
\showcasematerialdesignoutlined{question-mark}{1488}
\showcasematerialdesignoutlined{queue}{1489}
\showcasematerialdesignoutlined{queue-music}{1490}
\showcasematerialdesignoutlined{queue-play-next}{1491}
\showcasematerialdesignoutlined{quickreply}{1492}
\showcasematerialdesignoutlined{quiz}{1493}
\showcasematerialdesignoutlined{radar}{1494}
\showcasematerialdesignoutlined{radio}{1495}
\showcasematerialdesignoutlined{radio-button-checked}{1496}
\showcasematerialdesignoutlined{radio-button-unchecked}{1497}
\showcasematerialdesignoutlined{railway-alert}{1498}
\showcasematerialdesignoutlined{ramen-dining}{1499}
\showcasematerialdesignoutlined{ramp-left}{1500}
\showcasematerialdesignoutlined{ramp-right}{1501}
\showcasematerialdesignoutlined{rate-review}{1502}
\showcasematerialdesignoutlined{raw-off}{1503}
\showcasematerialdesignoutlined{raw-on}{1504}
\showcasematerialdesignoutlined{read-more}{1505}
\showcasematerialdesignoutlined{real-estate-agent}{1506}
\showcasematerialdesignoutlined{receipt}{1507}
\showcasematerialdesignoutlined{receipt-long}{1508}
\showcasematerialdesignoutlined{recent-actors}{1509}
\showcasematerialdesignoutlined{recommend}{1510}
\showcasematerialdesignoutlined{record-voice-over}{1511}
\showcasematerialdesignoutlined{rectangle}{1512}
\showcasematerialdesignoutlined{recycling}{1513}
\showcasematerialdesignoutlined{redeem}{1514}
\showcasematerialdesignoutlined{redo}{1515}
\showcasematerialdesignoutlined{reduce-capacity}{1516}
\showcasematerialdesignoutlined{refresh}{1517}
\showcasematerialdesignoutlined{remember-me}{1518}
\showcasematerialdesignoutlined{remove}{1519}
\showcasematerialdesignoutlined{remove-circle}{1520}
\showcasematerialdesignoutlined{remove-circle-outline}{1521}
\showcasematerialdesignoutlined{remove-done}{1522}
\showcasematerialdesignoutlined{remove-from-queue}{1523}
\showcasematerialdesignoutlined{remove-moderator}{1524}
\showcasematerialdesignoutlined{remove-red-eye}{1525}
\showcasematerialdesignoutlined{remove-road}{1526}
\showcasematerialdesignoutlined{remove-shopping-cart}{1527}
\showcasematerialdesignoutlined{reorder}{1528}
\showcasematerialdesignoutlined{repartition}{1529}
\showcasematerialdesignoutlined{repeat}{1530}
\showcasematerialdesignoutlined{repeat-on}{1531}
\showcasematerialdesignoutlined{repeat-one}{1532}
\showcasematerialdesignoutlined{repeat-one-on}{1533}
\showcasematerialdesignoutlined{replay}{1534}
\showcasematerialdesignoutlined{replay-5}{1535}
\showcasematerialdesignoutlined{replay-10}{1536}
\showcasematerialdesignoutlined{replay-30}{1537}
\showcasematerialdesignoutlined{replay-circle-filled}{1538}
\showcasematerialdesignoutlined{reply}{1539}
\showcasematerialdesignoutlined{reply-all}{1540}
\showcasematerialdesignoutlined{report}{1541}
\showcasematerialdesignoutlined{report-gmailerrorred}{1542}
\showcasematerialdesignoutlined{report-off}{1543}
\showcasematerialdesignoutlined{report-problem}{1544}
\showcasematerialdesignoutlined{request-page}{1545}
\showcasematerialdesignoutlined{request-quote}{1546}
\showcasematerialdesignoutlined{reset-tv}{1547}
\showcasematerialdesignoutlined{restart-alt}{1548}
\showcasematerialdesignoutlined{restaurant}{1549}
\showcasematerialdesignoutlined{restaurant-menu}{1550}
\showcasematerialdesignoutlined{restore}{1551}
\showcasematerialdesignoutlined{restore-from-trash}{1552}
\showcasematerialdesignoutlined{restore-page}{1553}
\showcasematerialdesignoutlined{reviews}{1554}
\showcasematerialdesignoutlined{rice-bowl}{1555}
\showcasematerialdesignoutlined{ring-volume}{1556}
\showcasematerialdesignoutlined{rocket}{1557}
\showcasematerialdesignoutlined{rocket-launch}{1558}
\showcasematerialdesignoutlined{roller-shades}{1559}
\showcasematerialdesignoutlined{roller-shades-closed}{1560}
\showcasematerialdesignoutlined{roller-skating}{1561}
\showcasematerialdesignoutlined{roofing}{1562}
\showcasematerialdesignoutlined{room}{1563}
\showcasematerialdesignoutlined{room-preferences}{1564}
\showcasematerialdesignoutlined{room-service}{1565}
\showcasematerialdesignoutlined{rotate-90-degrees-ccw}{1566}
\showcasematerialdesignoutlined{rotate-90-degrees-cw}{1567}
\showcasematerialdesignoutlined{rotate-left}{1568}
\showcasematerialdesignoutlined{rotate-right}{1569}
\showcasematerialdesignoutlined{roundabout-left}{1570}
\showcasematerialdesignoutlined{roundabout-right}{1571}
\showcasematerialdesignoutlined{rounded-corner}{1572}
\showcasematerialdesignoutlined{route}{1573}
\showcasematerialdesignoutlined{router}{1574}
\showcasematerialdesignoutlined{rowing}{1575}
\showcasematerialdesignoutlined{rss-feed}{1576}
\showcasematerialdesignoutlined{rsvp}{1577}
\showcasematerialdesignoutlined{rtt}{1578}
\showcasematerialdesignoutlined{rule}{1579}
\showcasematerialdesignoutlined{rule-folder}{1580}
\showcasematerialdesignoutlined{running-with-errors}{1581}
\showcasematerialdesignoutlined{run-circle}{1582}
\showcasematerialdesignoutlined{rv-hookup}{1583}
\showcasematerialdesignoutlined{r-mobiledata}{1584}
\showcasematerialdesignoutlined{safety-check}{1585}
\showcasematerialdesignoutlined{safety-divider}{1586}
\showcasematerialdesignoutlined{sailing}{1587}
\showcasematerialdesignoutlined{sanitizer}{1588}
\showcasematerialdesignoutlined{satellite}{1589}
\showcasematerialdesignoutlined{satellite-alt}{1590}
\showcasematerialdesignoutlined{save}{1591}
\showcasematerialdesignoutlined{saved-search}{1592}
\showcasematerialdesignoutlined{save-alt}{1593}
\showcasematerialdesignoutlined{save-as}{1594}
\showcasematerialdesignoutlined{savings}{1595}
\showcasematerialdesignoutlined{scale}{1596}
\showcasematerialdesignoutlined{scanner}{1597}
\showcasematerialdesignoutlined{scatter-plot}{1598}
\showcasematerialdesignoutlined{schedule}{1599}
\showcasematerialdesignoutlined{schedule-send}{1600}
\showcasematerialdesignoutlined{schema}{1601}
\showcasematerialdesignoutlined{school}{1602}
\showcasematerialdesignoutlined{science}{1603}
\showcasematerialdesignoutlined{score}{1604}
\showcasematerialdesignoutlined{scoreboard}{1605}
\showcasematerialdesignoutlined{screenshot}{1606}
\showcasematerialdesignoutlined{screenshot-monitor}{1607}
\showcasematerialdesignoutlined{screen-lock-landscape}{1608}
\showcasematerialdesignoutlined{screen-lock-portrait}{1609}
\showcasematerialdesignoutlined{screen-lock-rotation}{1610}
\showcasematerialdesignoutlined{screen-rotation}{1611}
\showcasematerialdesignoutlined{screen-rotation-alt}{1612}
\showcasematerialdesignoutlined{screen-search-desktop}{1613}
\showcasematerialdesignoutlined{screen-share}{1614}
\showcasematerialdesignoutlined{scuba-diving}{1615}
\showcasematerialdesignoutlined{sd}{1616}
\showcasematerialdesignoutlined{sd-card}{1617}
\showcasematerialdesignoutlined{sd-card-alert}{1618}
\showcasematerialdesignoutlined{sd-storage}{1619}
\showcasematerialdesignoutlined{search}{1620}
\showcasematerialdesignoutlined{search-off}{1621}
\showcasematerialdesignoutlined{security}{1622}
\showcasematerialdesignoutlined{security-update}{1623}
\showcasematerialdesignoutlined{security-update-good}{1624}
\showcasematerialdesignoutlined{security-update-warning}{1625}
\showcasematerialdesignoutlined{segment}{1626}
\showcasematerialdesignoutlined{select-all}{1627}
\showcasematerialdesignoutlined{self-improvement}{1628}
\showcasematerialdesignoutlined{sell}{1629}
\showcasematerialdesignoutlined{send}{1630}
\showcasematerialdesignoutlined{send-and-archive}{1631}
\showcasematerialdesignoutlined{send-time-extension}{1632}
\showcasematerialdesignoutlined{send-to-mobile}{1633}
\showcasematerialdesignoutlined{sensors}{1634}
\showcasematerialdesignoutlined{sensors-off}{1635}
\showcasematerialdesignoutlined{sensor-door}{1636}
\showcasematerialdesignoutlined{sensor-occupied}{1637}
\showcasematerialdesignoutlined{sensor-window}{1638}
\showcasematerialdesignoutlined{sentiment-dissatisfied}{1639}
\showcasematerialdesignoutlined{sentiment-neutral}{1640}
\showcasematerialdesignoutlined{sentiment-satisfied}{1641}
\showcasematerialdesignoutlined{sentiment-satisfied-alt}{1642}
\showcasematerialdesignoutlined{sentiment-very-dissatisfied}{1643}
\showcasematerialdesignoutlined{sentiment-very-satisfied}{1644}
\showcasematerialdesignoutlined{settings}{1645}
\showcasematerialdesignoutlined{settings-accessibility}{1646}
\showcasematerialdesignoutlined{settings-applications}{1647}
\showcasematerialdesignoutlined{settings-backup-restore}{1648}
\showcasematerialdesignoutlined{settings-bluetooth}{1649}
\showcasematerialdesignoutlined{settings-brightness}{1650}
\showcasematerialdesignoutlined{settings-cell}{1651}
\showcasematerialdesignoutlined{settings-ethernet}{1652}
\showcasematerialdesignoutlined{settings-input-antenna}{1653}
\showcasematerialdesignoutlined{settings-input-component}{1654}
\showcasematerialdesignoutlined{settings-input-composite}{1655}
\showcasematerialdesignoutlined{settings-input-hdmi}{1656}
\showcasematerialdesignoutlined{settings-input-svideo}{1657}
\showcasematerialdesignoutlined{settings-overscan}{1658}
\showcasematerialdesignoutlined{settings-phone}{1659}
\showcasematerialdesignoutlined{settings-power}{1660}
\showcasematerialdesignoutlined{settings-remote}{1661}
\showcasematerialdesignoutlined{settings-suggest}{1662}
\showcasematerialdesignoutlined{settings-system-daydream}{1663}
\showcasematerialdesignoutlined{settings-voice}{1664}
\showcasematerialdesignoutlined{set-meal}{1665}
\showcasematerialdesignoutlined{severe-cold}{1666}
\showcasematerialdesignoutlined{shape-line}{1667}
\showcasematerialdesignoutlined{share}{1668}
\showcasematerialdesignoutlined{share-location}{1669}
\showcasematerialdesignoutlined{shield}{1670}
\showcasematerialdesignoutlined{shield-moon}{1671}
\showcasematerialdesignoutlined{shop}{1672}
\showcasematerialdesignoutlined{shopping-bag}{1673}
\showcasematerialdesignoutlined{shopping-basket}{1674}
\showcasematerialdesignoutlined{shopping-cart}{1675}
\showcasematerialdesignoutlined{shopping-cart-checkout}{1676}
\showcasematerialdesignoutlined{shop-2}{1677}
\showcasematerialdesignoutlined{shop-two}{1678}
\showcasematerialdesignoutlined{shortcut}{1679}
\showcasematerialdesignoutlined{short-text}{1680}
\showcasematerialdesignoutlined{shower}{1681}
\showcasematerialdesignoutlined{show-chart}{1682}
\showcasematerialdesignoutlined{shuffle}{1683}
\showcasematerialdesignoutlined{shuffle-on}{1684}
\showcasematerialdesignoutlined{shutter-speed}{1685}
\showcasematerialdesignoutlined{sick}{1686}
\showcasematerialdesignoutlined{signal-cellular-0-bar}{1687}
\showcasematerialdesignoutlined{signal-cellular-4-bar}{1688}
\showcasematerialdesignoutlined{signal-cellular-alt}{1689}
\showcasematerialdesignoutlined{signal-cellular-alt-1-bar}{1690}
\showcasematerialdesignoutlined{signal-cellular-alt-2-bar}{1691}
\showcasematerialdesignoutlined{signal-cellular-connected-no-internet-0-bar}{1692}
\showcasematerialdesignoutlined{signal-cellular-connected-no-internet-4-bar}{1693}
\showcasematerialdesignoutlined{signal-cellular-nodata}{1694}
\showcasematerialdesignoutlined{signal-cellular-no-sim}{1695}
\showcasematerialdesignoutlined{signal-cellular-null}{1696}
\showcasematerialdesignoutlined{signal-cellular-off}{1697}
\showcasematerialdesignoutlined{signal-wifi-0-bar}{1698}
\showcasematerialdesignoutlined{signal-wifi-4-bar}{1699}
\showcasematerialdesignoutlined{signal-wifi-4-bar-lock}{1700}
\showcasematerialdesignoutlined{signal-wifi-bad}{1701}
\showcasematerialdesignoutlined{signal-wifi-connected-no-internet-4}{1702}
\showcasematerialdesignoutlined{signal-wifi-off}{1703}
\showcasematerialdesignoutlined{signal-wifi-statusbar-4-bar}{1704}
\showcasematerialdesignoutlined{signal-wifi-statusbar-connected-no-internet-4}{1705}
\showcasematerialdesignoutlined{signal-wifi-statusbar-null}{1706}
\showcasematerialdesignoutlined{signpost}{1707}
\showcasematerialdesignoutlined{sign-language}{1708}
\showcasematerialdesignoutlined{sim-card}{1709}
\showcasematerialdesignoutlined{sim-card-alert}{1710}
\showcasematerialdesignoutlined{sim-card-download}{1711}
\showcasematerialdesignoutlined{single-bed}{1712}
\showcasematerialdesignoutlined{sip}{1713}
\showcasematerialdesignoutlined{skateboarding}{1714}
\showcasematerialdesignoutlined{skip-next}{1715}
\showcasematerialdesignoutlined{skip-previous}{1716}
\showcasematerialdesignoutlined{sledding}{1717}
\showcasematerialdesignoutlined{slideshow}{1718}
\showcasematerialdesignoutlined{slow-motion-video}{1719}
\showcasematerialdesignoutlined{smartphone}{1720}
\showcasematerialdesignoutlined{smart-button}{1721}
\showcasematerialdesignoutlined{smart-display}{1722}
\showcasematerialdesignoutlined{smart-screen}{1723}
\showcasematerialdesignoutlined{smart-toy}{1724}
\showcasematerialdesignoutlined{smoke-free}{1725}
\showcasematerialdesignoutlined{smoking-rooms}{1726}
\showcasematerialdesignoutlined{sms}{1727}
\showcasematerialdesignoutlined{sms-failed}{1728}
\showcasematerialdesignoutlined{snippet-folder}{1729}
\showcasematerialdesignoutlined{snooze}{1730}
\showcasematerialdesignoutlined{snowboarding}{1731}
\showcasematerialdesignoutlined{snowmobile}{1732}
\showcasematerialdesignoutlined{snowshoeing}{1733}
\showcasematerialdesignoutlined{soap}{1734}
\showcasematerialdesignoutlined{social-distance}{1735}
\showcasematerialdesignoutlined{solar-power}{1736}
\showcasematerialdesignoutlined{sort}{1737}
\showcasematerialdesignoutlined{sort-by-alpha}{1738}
\showcasematerialdesignoutlined{sos}{1739}
\showcasematerialdesignoutlined{soup-kitchen}{1740}
\showcasematerialdesignoutlined{source}{1741}
\showcasematerialdesignoutlined{south}{1742}
\showcasematerialdesignoutlined{south-america}{1743}
\showcasematerialdesignoutlined{south-east}{1744}
\showcasematerialdesignoutlined{south-west}{1745}
\showcasematerialdesignoutlined{spa}{1746}
\showcasematerialdesignoutlined{space-bar}{1747}
\showcasematerialdesignoutlined{space-dashboard}{1748}
\showcasematerialdesignoutlined{spatial-audio}{1749}
\showcasematerialdesignoutlined{spatial-audio-off}{1750}
\showcasematerialdesignoutlined{spatial-tracking}{1751}
\showcasematerialdesignoutlined{speaker}{1752}
\showcasematerialdesignoutlined{speaker-group}{1753}
\showcasematerialdesignoutlined{speaker-notes}{1754}
\showcasematerialdesignoutlined{speaker-notes-off}{1755}
\showcasematerialdesignoutlined{speaker-phone}{1756}
\showcasematerialdesignoutlined{speed}{1757}
\showcasematerialdesignoutlined{spellcheck}{1758}
\showcasematerialdesignoutlined{splitscreen}{1759}
\showcasematerialdesignoutlined{spoke}{1760}
\showcasematerialdesignoutlined{sports}{1761}
\showcasematerialdesignoutlined{sports-bar}{1762}
\showcasematerialdesignoutlined{sports-baseball}{1763}
\showcasematerialdesignoutlined{sports-basketball}{1764}
\showcasematerialdesignoutlined{sports-cricket}{1765}
\showcasematerialdesignoutlined{sports-esports}{1766}
\showcasematerialdesignoutlined{sports-football}{1767}
\showcasematerialdesignoutlined{sports-golf}{1768}
\showcasematerialdesignoutlined{sports-gymnastics}{1769}
\showcasematerialdesignoutlined{sports-handball}{1770}
\showcasematerialdesignoutlined{sports-hockey}{1771}
\showcasematerialdesignoutlined{sports-kabaddi}{1772}
\showcasematerialdesignoutlined{sports-martial-arts}{1773}
\showcasematerialdesignoutlined{sports-mma}{1774}
\showcasematerialdesignoutlined{sports-motorsports}{1775}
\showcasematerialdesignoutlined{sports-rugby}{1776}
\showcasematerialdesignoutlined{sports-score}{1777}
\showcasematerialdesignoutlined{sports-soccer}{1778}
\showcasematerialdesignoutlined{sports-tennis}{1779}
\showcasematerialdesignoutlined{sports-volleyball}{1780}
\showcasematerialdesignoutlined{square}{1781}
\showcasematerialdesignoutlined{square-foot}{1782}
\showcasematerialdesignoutlined{ssid-chart}{1783}
\showcasematerialdesignoutlined{stacked-bar-chart}{1784}
\showcasematerialdesignoutlined{stacked-line-chart}{1785}
\showcasematerialdesignoutlined{stadium}{1786}
\showcasematerialdesignoutlined{stairs}{1787}
\showcasematerialdesignoutlined{star}{1788}
\showcasematerialdesignoutlined{stars}{1789}
\showcasematerialdesignoutlined{start}{1790}
\showcasematerialdesignoutlined{star-border}{1791}
\showcasematerialdesignoutlined{star-border-purple500}{1792}
\showcasematerialdesignoutlined{star-half}{1793}
\showcasematerialdesignoutlined{star-outline}{1794}
\showcasematerialdesignoutlined{star-purple500}{1795}
\showcasematerialdesignoutlined{star-rate}{1796}
\showcasematerialdesignoutlined{stay-current-landscape}{1797}
\showcasematerialdesignoutlined{stay-current-portrait}{1798}
\showcasematerialdesignoutlined{stay-primary-landscape}{1799}
\showcasematerialdesignoutlined{stay-primary-portrait}{1800}
\showcasematerialdesignoutlined{sticky-note-2}{1801}
\showcasematerialdesignoutlined{stop}{1802}
\showcasematerialdesignoutlined{stop-circle}{1803}
\showcasematerialdesignoutlined{stop-screen-share}{1804}
\showcasematerialdesignoutlined{storage}{1805}
\showcasematerialdesignoutlined{store}{1806}
\showcasematerialdesignoutlined{storefront}{1807}
\showcasematerialdesignoutlined{store-mall-directory}{1808}
\showcasematerialdesignoutlined{storm}{1809}
\showcasematerialdesignoutlined{straight}{1810}
\showcasematerialdesignoutlined{straighten}{1811}
\showcasematerialdesignoutlined{stream}{1812}
\showcasematerialdesignoutlined{streetview}{1813}
\showcasematerialdesignoutlined{strikethrough-s}{1814}
\showcasematerialdesignoutlined{stroller}{1815}
\showcasematerialdesignoutlined{style}{1816}
\showcasematerialdesignoutlined{subdirectory-arrow-left}{1817}
\showcasematerialdesignoutlined{subdirectory-arrow-right}{1818}
\showcasematerialdesignoutlined{subject}{1819}
\showcasematerialdesignoutlined{subscript}{1820}
\showcasematerialdesignoutlined{subscriptions}{1821}
\showcasematerialdesignoutlined{subtitles}{1822}
\showcasematerialdesignoutlined{subtitles-off}{1823}
\showcasematerialdesignoutlined{subway}{1824}
\showcasematerialdesignoutlined{summarize}{1825}
\showcasematerialdesignoutlined{superscript}{1826}
\showcasematerialdesignoutlined{supervised-user-circle}{1827}
\showcasematerialdesignoutlined{supervisor-account}{1828}
\showcasematerialdesignoutlined{support}{1829}
\showcasematerialdesignoutlined{support-agent}{1830}
\showcasematerialdesignoutlined{surfing}{1831}
\showcasematerialdesignoutlined{surround-sound}{1832}
\showcasematerialdesignoutlined{swap-calls}{1833}
\showcasematerialdesignoutlined{swap-horiz}{1834}
\showcasematerialdesignoutlined{swap-horizontal-circle}{1835}
\showcasematerialdesignoutlined{swap-vert}{1836}
\showcasematerialdesignoutlined{swap-vertical-circle}{1837}
\showcasematerialdesignoutlined{swipe}{1838}
\showcasematerialdesignoutlined{swipe-down}{1839}
\showcasematerialdesignoutlined{swipe-down-alt}{1840}
\showcasematerialdesignoutlined{swipe-left}{1841}
\showcasematerialdesignoutlined{swipe-left-alt}{1842}
\showcasematerialdesignoutlined{swipe-right}{1843}
\showcasematerialdesignoutlined{swipe-right-alt}{1844}
\showcasematerialdesignoutlined{swipe-up}{1845}
\showcasematerialdesignoutlined{swipe-up-alt}{1846}
\showcasematerialdesignoutlined{swipe-vertical}{1847}
\showcasematerialdesignoutlined{switch-access-shortcut}{1848}
\showcasematerialdesignoutlined{switch-access-shortcut-add}{1849}
\showcasematerialdesignoutlined{switch-account}{1850}
\showcasematerialdesignoutlined{switch-camera}{1851}
\showcasematerialdesignoutlined{switch-left}{1852}
\showcasematerialdesignoutlined{switch-right}{1853}
\showcasematerialdesignoutlined{switch-video}{1854}
\showcasematerialdesignoutlined{synagogue}{1855}
\showcasematerialdesignoutlined{sync}{1856}
\showcasematerialdesignoutlined{sync-alt}{1857}
\showcasematerialdesignoutlined{sync-disabled}{1858}
\showcasematerialdesignoutlined{sync-lock}{1859}
\showcasematerialdesignoutlined{sync-problem}{1860}
\showcasematerialdesignoutlined{system-security-update}{1861}
\showcasematerialdesignoutlined{system-security-update-good}{1862}
\showcasematerialdesignoutlined{system-security-update-warning}{1863}
\showcasematerialdesignoutlined{system-update}{1864}
\showcasematerialdesignoutlined{system-update-alt}{1865}
\showcasematerialdesignoutlined{tab}{1866}
\showcasematerialdesignoutlined{tablet}{1867}
\showcasematerialdesignoutlined{tablet-android}{1868}
\showcasematerialdesignoutlined{tablet-mac}{1869}
\showcasematerialdesignoutlined{table-bar}{1870}
\showcasematerialdesignoutlined{table-chart}{1871}
\showcasematerialdesignoutlined{table-restaurant}{1872}
\showcasematerialdesignoutlined{table-rows}{1873}
\showcasematerialdesignoutlined{table-view}{1874}
\showcasematerialdesignoutlined{tab-unselected}{1875}
\showcasematerialdesignoutlined{tag}{1876}
\showcasematerialdesignoutlined{tag-faces}{1877}
\showcasematerialdesignoutlined{takeout-dining}{1878}
\showcasematerialdesignoutlined{tapas}{1879}
\showcasematerialdesignoutlined{tap-and-play}{1880}
\showcasematerialdesignoutlined{task}{1881}
\showcasematerialdesignoutlined{task-alt}{1882}
\showcasematerialdesignoutlined{taxi-alert}{1883}
\showcasematerialdesignoutlined{temple-buddhist}{1884}
\showcasematerialdesignoutlined{temple-hindu}{1885}
\showcasematerialdesignoutlined{terminal}{1886}
\showcasematerialdesignoutlined{terrain}{1887}
\showcasematerialdesignoutlined{textsms}{1888}
\showcasematerialdesignoutlined{texture}{1889}
\showcasematerialdesignoutlined{text-decrease}{1890}
\showcasematerialdesignoutlined{text-fields}{1891}
\showcasematerialdesignoutlined{text-format}{1892}
\showcasematerialdesignoutlined{text-increase}{1893}
\showcasematerialdesignoutlined{text-rotate-up}{1894}
\showcasematerialdesignoutlined{text-rotate-vertical}{1895}
\showcasematerialdesignoutlined{text-rotation-angledown}{1896}
\showcasematerialdesignoutlined{text-rotation-angleup}{1897}
\showcasematerialdesignoutlined{text-rotation-down}{1898}
\showcasematerialdesignoutlined{text-rotation-none}{1899}
\showcasematerialdesignoutlined{text-snippet}{1900}
\showcasematerialdesignoutlined{theaters}{1901}
\showcasematerialdesignoutlined{theater-comedy}{1902}
\showcasematerialdesignoutlined{thermostat}{1903}
\showcasematerialdesignoutlined{thermostat-auto}{1904}
\showcasematerialdesignoutlined{thumbs-up-down}{1905}
\showcasematerialdesignoutlined{thumb-down}{1906}
\showcasematerialdesignoutlined{thumb-down-alt}{1907}
\showcasematerialdesignoutlined{thumb-down-off-alt}{1908}
\showcasematerialdesignoutlined{thumb-up}{1909}
\showcasematerialdesignoutlined{thumb-up-alt}{1910}
\showcasematerialdesignoutlined{thumb-up-off-alt}{1911}
\showcasematerialdesignoutlined{thunderstorm}{1912}
\showcasematerialdesignoutlined{timelapse}{1913}
\showcasematerialdesignoutlined{timeline}{1914}
\showcasematerialdesignoutlined{timer}{1915}
\showcasematerialdesignoutlined{timer-3}{1916}
\showcasematerialdesignoutlined{timer-3-select}{1917}
\showcasematerialdesignoutlined{timer-10}{1918}
\showcasematerialdesignoutlined{timer-10-select}{1919}
\showcasematerialdesignoutlined{timer-off}{1920}
\showcasematerialdesignoutlined{time-to-leave}{1921}
\showcasematerialdesignoutlined{tips-and-updates}{1922}
\showcasematerialdesignoutlined{tire-repair}{1923}
\showcasematerialdesignoutlined{title}{1924}
\showcasematerialdesignoutlined{toc}{1925}
\showcasematerialdesignoutlined{today}{1926}
\showcasematerialdesignoutlined{toggle-off}{1927}
\showcasematerialdesignoutlined{toggle-on}{1928}
\showcasematerialdesignoutlined{token}{1929}
\showcasematerialdesignoutlined{toll}{1930}
\showcasematerialdesignoutlined{tonality}{1931}
\showcasematerialdesignoutlined{topic}{1932}
\showcasematerialdesignoutlined{tornado}{1933}
\showcasematerialdesignoutlined{touch-app}{1934}
\showcasematerialdesignoutlined{tour}{1935}
\showcasematerialdesignoutlined{toys}{1936}
\showcasematerialdesignoutlined{track-changes}{1937}
\showcasematerialdesignoutlined{traffic}{1938}
\showcasematerialdesignoutlined{train}{1939}
\showcasematerialdesignoutlined{tram}{1940}
\showcasematerialdesignoutlined{transcribe}{1941}
\showcasematerialdesignoutlined{transfer-within-a-station}{1942}
\showcasematerialdesignoutlined{transform}{1943}
\showcasematerialdesignoutlined{transgender}{1944}
\showcasematerialdesignoutlined{transit-enterexit}{1945}
\showcasematerialdesignoutlined{translate}{1946}
\showcasematerialdesignoutlined{travel-explore}{1947}
\showcasematerialdesignoutlined{trending-down}{1948}
\showcasematerialdesignoutlined{trending-flat}{1949}
\showcasematerialdesignoutlined{trending-up}{1950}
\showcasematerialdesignoutlined{trip-origin}{1951}
\showcasematerialdesignoutlined{troubleshoot}{1952}
\showcasematerialdesignoutlined{try}{1953}
\showcasematerialdesignoutlined{tsunami}{1954}
\showcasematerialdesignoutlined{tty}{1955}
\showcasematerialdesignoutlined{tune}{1956}
\showcasematerialdesignoutlined{tungsten}{1957}
\showcasematerialdesignoutlined{turned-in}{1958}
\showcasematerialdesignoutlined{turned-in-not}{1959}
\showcasematerialdesignoutlined{turn-left}{1960}
\showcasematerialdesignoutlined{turn-right}{1961}
\showcasematerialdesignoutlined{turn-sharp-left}{1962}
\showcasematerialdesignoutlined{turn-sharp-right}{1963}
\showcasematerialdesignoutlined{turn-slight-left}{1964}
\showcasematerialdesignoutlined{turn-slight-right}{1965}
\showcasematerialdesignoutlined{tv}{1966}
\showcasematerialdesignoutlined{tv-off}{1967}
\showcasematerialdesignoutlined{two-wheeler}{1968}
\showcasematerialdesignoutlined{type-specimen}{1969}
\showcasematerialdesignoutlined{umbrella}{1970}
\showcasematerialdesignoutlined{unarchive}{1971}
\showcasematerialdesignoutlined{undo}{1972}
\showcasematerialdesignoutlined{unfold-less}{1973}
\showcasematerialdesignoutlined{unfold-less-double}{1974}
\showcasematerialdesignoutlined{unfold-more}{1975}
\showcasematerialdesignoutlined{unfold-more-double}{1976}
\showcasematerialdesignoutlined{unpublished}{1977}
\showcasematerialdesignoutlined{unsubscribe}{1978}
\showcasematerialdesignoutlined{upcoming}{1979}
\showcasematerialdesignoutlined{update}{1980}
\showcasematerialdesignoutlined{update-disabled}{1981}
\showcasematerialdesignoutlined{upgrade}{1982}
\showcasematerialdesignoutlined{upload}{1983}
\showcasematerialdesignoutlined{upload-file}{1984}
\showcasematerialdesignoutlined{usb}{1985}
\showcasematerialdesignoutlined{usb-off}{1986}
\showcasematerialdesignoutlined{u-turn-left}{1987}
\showcasematerialdesignoutlined{u-turn-right}{1988}
\showcasematerialdesignoutlined{vaccines}{1989}
\showcasematerialdesignoutlined{vape-free}{1990}
\showcasematerialdesignoutlined{vaping-rooms}{1991}
\showcasematerialdesignoutlined{verified}{1992}
\showcasematerialdesignoutlined{verified-user}{1993}
\showcasematerialdesignoutlined{vertical-align-bottom}{1994}
\showcasematerialdesignoutlined{vertical-align-center}{1995}
\showcasematerialdesignoutlined{vertical-align-top}{1996}
\showcasematerialdesignoutlined{vertical-distribute}{1997}
\showcasematerialdesignoutlined{vertical-shades}{1998}
\showcasematerialdesignoutlined{vertical-shades-closed}{1999}
\showcasematerialdesignoutlined{vertical-split}{2000}
\showcasematerialdesignoutlined{vibration}{2001}
\showcasematerialdesignoutlined{videocam}{2002}
\showcasematerialdesignoutlined{videocam-off}{2003}
\showcasematerialdesignoutlined{videogame-asset}{2004}
\showcasematerialdesignoutlined{videogame-asset-off}{2005}
\showcasematerialdesignoutlined{video-call}{2006}
\showcasematerialdesignoutlined{video-camera-back}{2007}
\showcasematerialdesignoutlined{video-camera-front}{2008}
\showcasematerialdesignoutlined{video-chat}{2009}
\showcasematerialdesignoutlined{video-file}{2010}
\showcasematerialdesignoutlined{video-label}{2011}
\showcasematerialdesignoutlined{video-library}{2012}
\showcasematerialdesignoutlined{video-settings}{2013}
\showcasematerialdesignoutlined{video-stable}{2014}
\showcasematerialdesignoutlined{view-agenda}{2015}
\showcasematerialdesignoutlined{view-array}{2016}
\showcasematerialdesignoutlined{view-carousel}{2017}
\showcasematerialdesignoutlined{view-column}{2018}
\showcasematerialdesignoutlined{view-comfy}{2019}
\showcasematerialdesignoutlined{view-comfy-alt}{2020}
\showcasematerialdesignoutlined{view-compact}{2021}
\showcasematerialdesignoutlined{view-compact-alt}{2022}
\showcasematerialdesignoutlined{view-cozy}{2023}
\showcasematerialdesignoutlined{view-day}{2024}
\showcasematerialdesignoutlined{view-headline}{2025}
\showcasematerialdesignoutlined{view-in-ar}{2026}
\showcasematerialdesignoutlined{view-kanban}{2027}
\showcasematerialdesignoutlined{view-list}{2028}
\showcasematerialdesignoutlined{view-module}{2029}
\showcasematerialdesignoutlined{view-quilt}{2030}
\showcasematerialdesignoutlined{view-sidebar}{2031}
\showcasematerialdesignoutlined{view-stream}{2032}
\showcasematerialdesignoutlined{view-timeline}{2033}
\showcasematerialdesignoutlined{view-week}{2034}
\showcasematerialdesignoutlined{vignette}{2035}
\showcasematerialdesignoutlined{villa}{2036}
\showcasematerialdesignoutlined{visibility}{2037}
\showcasematerialdesignoutlined{visibility-off}{2038}
\showcasematerialdesignoutlined{voicemail}{2039}
\showcasematerialdesignoutlined{voice-chat}{2040}
\showcasematerialdesignoutlined{voice-over-off}{2041}
\showcasematerialdesignoutlined{volcano}{2042}
\showcasematerialdesignoutlined{volume-down}{2043}
\showcasematerialdesignoutlined{volume-mute}{2044}
\showcasematerialdesignoutlined{volume-off}{2045}
\showcasematerialdesignoutlined{volume-up}{2046}
\showcasematerialdesignoutlined{volunteer-activism}{2047}
\showcasematerialdesignoutlined{vpn-key}{2048}
\showcasematerialdesignoutlined{vpn-key-off}{2049}
\showcasematerialdesignoutlined{vpn-lock}{2050}
\showcasematerialdesignoutlined{vrpano}{2051}
\showcasematerialdesignoutlined{wallet}{2052}
\showcasematerialdesignoutlined{wallpaper}{2053}
\showcasematerialdesignoutlined{warehouse}{2054}
\showcasematerialdesignoutlined{warning}{2055}
\showcasematerialdesignoutlined{warning-amber}{2056}
\showcasematerialdesignoutlined{wash}{2057}
\showcasematerialdesignoutlined{watch}{2058}
\showcasematerialdesignoutlined{watch-later}{2059}
\showcasematerialdesignoutlined{watch-off}{2060}
\showcasematerialdesignoutlined{water}{2061}
\showcasematerialdesignoutlined{waterfall-chart}{2062}
\showcasematerialdesignoutlined{water-damage}{2063}
\showcasematerialdesignoutlined{water-drop}{2064}
\showcasematerialdesignoutlined{waves}{2065}
\showcasematerialdesignoutlined{waving-hand}{2066}
\showcasematerialdesignoutlined{wb-auto}{2067}
\showcasematerialdesignoutlined{wb-cloudy}{2068}
\showcasematerialdesignoutlined{wb-incandescent}{2069}
\showcasematerialdesignoutlined{wb-iridescent}{2070}
\showcasematerialdesignoutlined{wb-shade}{2071}
\showcasematerialdesignoutlined{wb-sunny}{2072}
\showcasematerialdesignoutlined{wb-twilight}{2073}
\showcasematerialdesignoutlined{wc}{2074}
\showcasematerialdesignoutlined{web}{2075}
\showcasematerialdesignoutlined{webhook}{2076}
\showcasematerialdesignoutlined{web-asset}{2077}
\showcasematerialdesignoutlined{web-asset-off}{2078}
\showcasematerialdesignoutlined{web-stories}{2079}
\showcasematerialdesignoutlined{weekend}{2080}
\showcasematerialdesignoutlined{west}{2081}
\showcasematerialdesignoutlined{whatshot}{2082}
\showcasematerialdesignoutlined{wheelchair-pickup}{2083}
\showcasematerialdesignoutlined{where-to-vote}{2084}
\showcasematerialdesignoutlined{widgets}{2085}
\showcasematerialdesignoutlined{width-full}{2086}
\showcasematerialdesignoutlined{width-normal}{2087}
\showcasematerialdesignoutlined{width-wide}{2088}
\showcasematerialdesignoutlined{wifi}{2089}
\showcasematerialdesignoutlined{wifi-1-bar}{2090}
\showcasematerialdesignoutlined{wifi-2-bar}{2091}
\showcasematerialdesignoutlined{wifi-calling}{2092}
\showcasematerialdesignoutlined{wifi-calling-3}{2093}
\showcasematerialdesignoutlined{wifi-channel}{2094}
\showcasematerialdesignoutlined{wifi-find}{2095}
\showcasematerialdesignoutlined{wifi-lock}{2096}
\showcasematerialdesignoutlined{wifi-off}{2097}
\showcasematerialdesignoutlined{wifi-password}{2098}
\showcasematerialdesignoutlined{wifi-protected-setup}{2099}
\showcasematerialdesignoutlined{wifi-tethering}{2100}
\showcasematerialdesignoutlined{wifi-tethering-error}{2101}
\showcasematerialdesignoutlined{wifi-tethering-off}{2102}
\showcasematerialdesignoutlined{window}{2103}
\showcasematerialdesignoutlined{wind-power}{2104}
\showcasematerialdesignoutlined{wine-bar}{2105}
\showcasematerialdesignoutlined{woman}{2106}
\showcasematerialdesignoutlined{woman-2}{2107}
\showcasematerialdesignoutlined{work}{2108}
\showcasematerialdesignoutlined{workspaces}{2109}
\showcasematerialdesignoutlined{workspace-premium}{2110}
\showcasematerialdesignoutlined{work-history}{2111}
\showcasematerialdesignoutlined{work-off}{2112}
\showcasematerialdesignoutlined{work-outline}{2113}
\showcasematerialdesignoutlined{wrap-text}{2114}
\showcasematerialdesignoutlined{wrong-location}{2115}
\showcasematerialdesignoutlined{wysiwyg}{2116}
\showcasematerialdesignoutlined{yard}{2117}
\showcasematerialdesignoutlined{youtube-searched-for}{2118}
\showcasematerialdesignoutlined{zoom-in}{2119}
\showcasematerialdesignoutlined{zoom-in-map}{2120}
\showcasematerialdesignoutlined{zoom-out}{2121}
\showcasematerialdesignoutlined{zoom-out-map}{2122}
\end{materialdesigncase}

\subsection{Filled version (2122 icons)}

The \textsf{PDF} file is \texttt{materialdesign-filled-all.pdf}

\begin{materialdesigncase}
\showcasematerialdesignfilled{1k}{1}
\showcasematerialdesignfilled{1k-plus}{2}
\showcasematerialdesignfilled{1x-mobiledata}{3}
\showcasematerialdesignfilled{2k}{4}
\showcasematerialdesignfilled{2k-plus}{5}
\showcasematerialdesignfilled{2mp}{6}
\showcasematerialdesignfilled{3d-rotation}{7}
\showcasematerialdesignfilled{3g-mobiledata}{8}
\showcasematerialdesignfilled{3k}{9}
\showcasematerialdesignfilled{3k-plus}{10}
\showcasematerialdesignfilled{3mp}{11}
\showcasematerialdesignfilled{3p}{12}
\showcasematerialdesignfilled{4g-mobiledata}{13}
\showcasematerialdesignfilled{4g-plus-mobiledata}{14}
\showcasematerialdesignfilled{4k}{15}
\showcasematerialdesignfilled{4k-plus}{16}
\showcasematerialdesignfilled{4mp}{17}
\showcasematerialdesignfilled{5g}{18}
\showcasematerialdesignfilled{5k}{19}
\showcasematerialdesignfilled{5k-plus}{20}
\showcasematerialdesignfilled{5mp}{21}
\showcasematerialdesignfilled{6k}{22}
\showcasematerialdesignfilled{6k-plus}{23}
\showcasematerialdesignfilled{6mp}{24}
\showcasematerialdesignfilled{6-ft-apart}{25}
\showcasematerialdesignfilled{7k}{26}
\showcasematerialdesignfilled{7k-plus}{27}
\showcasematerialdesignfilled{7mp}{28}
\showcasematerialdesignfilled{8k}{29}
\showcasematerialdesignfilled{8k-plus}{30}
\showcasematerialdesignfilled{8mp}{31}
\showcasematerialdesignfilled{9k}{32}
\showcasematerialdesignfilled{9k-plus}{33}
\showcasematerialdesignfilled{9mp}{34}
\showcasematerialdesignfilled{10k}{35}
\showcasematerialdesignfilled{10mp}{36}
\showcasematerialdesignfilled{11mp}{37}
\showcasematerialdesignfilled{12mp}{38}
\showcasematerialdesignfilled{13mp}{39}
\showcasematerialdesignfilled{14mp}{40}
\showcasematerialdesignfilled{15mp}{41}
\showcasematerialdesignfilled{16mp}{42}
\showcasematerialdesignfilled{17mp}{43}
\showcasematerialdesignfilled{18mp}{44}
\showcasematerialdesignfilled{18-up-rating}{45}
\showcasematerialdesignfilled{19mp}{46}
\showcasematerialdesignfilled{20mp}{47}
\showcasematerialdesignfilled{21mp}{48}
\showcasematerialdesignfilled{22mp}{49}
\showcasematerialdesignfilled{23mp}{50}
\showcasematerialdesignfilled{24mp}{51}
\showcasematerialdesignfilled{30fps}{52}
\showcasematerialdesignfilled{30fps-select}{53}
\showcasematerialdesignfilled{60fps}{54}
\showcasematerialdesignfilled{60fps-select}{55}
\showcasematerialdesignfilled{123}{56}
\showcasematerialdesignfilled{360}{57}
\showcasematerialdesignfilled{abc}{58}
\showcasematerialdesignfilled{accessibility}{59}
\showcasematerialdesignfilled{accessibility-new}{60}
\showcasematerialdesignfilled{accessible}{61}
\showcasematerialdesignfilled{accessible-forward}{62}
\showcasematerialdesignfilled{access-alarm}{63}
\showcasematerialdesignfilled{access-alarms}{64}
\showcasematerialdesignfilled{access-time}{65}
\showcasematerialdesignfilled{access-time-filled}{66}
\showcasematerialdesignfilled{account-balance}{67}
\showcasematerialdesignfilled{account-balance-wallet}{68}
\showcasematerialdesignfilled{account-box}{69}
\showcasematerialdesignfilled{account-circle}{70}
\showcasematerialdesignfilled{account-tree}{71}
\showcasematerialdesignfilled{ac-unit}{72}
\showcasematerialdesignfilled{adb}{73}
\showcasematerialdesignfilled{add}{74}
\showcasematerialdesignfilled{addchart}{75}
\showcasematerialdesignfilled{add-alarm}{76}
\showcasematerialdesignfilled{add-alert}{77}
\showcasematerialdesignfilled{add-a-photo}{78}
\showcasematerialdesignfilled{add-box}{79}
\showcasematerialdesignfilled{add-business}{80}
\showcasematerialdesignfilled{add-card}{81}
\showcasematerialdesignfilled{add-chart}{82}
\showcasematerialdesignfilled{add-circle}{83}
\showcasematerialdesignfilled{add-circle-outline}{84}
\showcasematerialdesignfilled{add-comment}{85}
\showcasematerialdesignfilled{add-home}{86}
\showcasematerialdesignfilled{add-home-work}{87}
\showcasematerialdesignfilled{add-ic-call}{88}
\showcasematerialdesignfilled{add-link}{89}
\showcasematerialdesignfilled{add-location}{90}
\showcasematerialdesignfilled{add-location-alt}{91}
\showcasematerialdesignfilled{add-moderator}{92}
\showcasematerialdesignfilled{add-photo-alternate}{93}
\showcasematerialdesignfilled{add-reaction}{94}
\showcasematerialdesignfilled{add-road}{95}
\showcasematerialdesignfilled{add-shopping-cart}{96}
\showcasematerialdesignfilled{add-task}{97}
\showcasematerialdesignfilled{add-to-drive}{98}
\showcasematerialdesignfilled{add-to-home-screen}{99}
\showcasematerialdesignfilled{add-to-photos}{100}
\showcasematerialdesignfilled{add-to-queue}{101}
\showcasematerialdesignfilled{adf-scanner}{102}
\showcasematerialdesignfilled{adjust}{103}
\showcasematerialdesignfilled{admin-panel-settings}{104}
\showcasematerialdesignfilled{ads-click}{105}
\showcasematerialdesignfilled{ad-units}{106}
\showcasematerialdesignfilled{agriculture}{107}
\showcasematerialdesignfilled{air}{108}
\showcasematerialdesignfilled{airlines}{109}
\showcasematerialdesignfilled{airline-seat-flat}{110}
\showcasematerialdesignfilled{airline-seat-flat-angled}{111}
\showcasematerialdesignfilled{airline-seat-individual-suite}{112}
\showcasematerialdesignfilled{airline-seat-legroom-extra}{113}
\showcasematerialdesignfilled{airline-seat-legroom-normal}{114}
\showcasematerialdesignfilled{airline-seat-legroom-reduced}{115}
\showcasematerialdesignfilled{airline-seat-recline-extra}{116}
\showcasematerialdesignfilled{airline-seat-recline-normal}{117}
\showcasematerialdesignfilled{airline-stops}{118}
\showcasematerialdesignfilled{airplanemode-active}{119}
\showcasematerialdesignfilled{airplanemode-inactive}{120}
\showcasematerialdesignfilled{airplane-ticket}{121}
\showcasematerialdesignfilled{airplay}{122}
\showcasematerialdesignfilled{airport-shuttle}{123}
\showcasematerialdesignfilled{alarm}{124}
\showcasematerialdesignfilled{alarm-add}{125}
\showcasematerialdesignfilled{alarm-off}{126}
\showcasematerialdesignfilled{alarm-on}{127}
\showcasematerialdesignfilled{album}{128}
\showcasematerialdesignfilled{align-horizontal-center}{129}
\showcasematerialdesignfilled{align-horizontal-left}{130}
\showcasematerialdesignfilled{align-horizontal-right}{131}
\showcasematerialdesignfilled{align-vertical-bottom}{132}
\showcasematerialdesignfilled{align-vertical-center}{133}
\showcasematerialdesignfilled{align-vertical-top}{134}
\showcasematerialdesignfilled{all-inbox}{135}
\showcasematerialdesignfilled{all-inclusive}{136}
\showcasematerialdesignfilled{all-out}{137}
\showcasematerialdesignfilled{alternate-email}{138}
\showcasematerialdesignfilled{alt-route}{139}
\showcasematerialdesignfilled{analytics}{140}
\showcasematerialdesignfilled{anchor}{141}
\showcasematerialdesignfilled{android}{142}
\showcasematerialdesignfilled{animation}{143}
\showcasematerialdesignfilled{announcement}{144}
\showcasematerialdesignfilled{aod}{145}
\showcasematerialdesignfilled{apartment}{146}
\showcasematerialdesignfilled{api}{147}
\showcasematerialdesignfilled{approval}{148}
\showcasematerialdesignfilled{apps}{149}
\showcasematerialdesignfilled{apps-outage}{150}
\showcasematerialdesignfilled{app-blocking}{151}
\showcasematerialdesignfilled{app-registration}{152}
\showcasematerialdesignfilled{app-settings-alt}{153}
\showcasematerialdesignfilled{app-shortcut}{154}
\showcasematerialdesignfilled{architecture}{155}
\showcasematerialdesignfilled{archive}{156}
\showcasematerialdesignfilled{area-chart}{157}
\showcasematerialdesignfilled{arrow-back}{158}
\showcasematerialdesignfilled{arrow-back-ios}{159}
\showcasematerialdesignfilled{arrow-back-ios-new}{160}
\showcasematerialdesignfilled{arrow-circle-down}{161}
\showcasematerialdesignfilled{arrow-circle-left}{162}
\showcasematerialdesignfilled{arrow-circle-right}{163}
\showcasematerialdesignfilled{arrow-circle-up}{164}
\showcasematerialdesignfilled{arrow-downward}{165}
\showcasematerialdesignfilled{arrow-drop-down}{166}
\showcasematerialdesignfilled{arrow-drop-down-circle}{167}
\showcasematerialdesignfilled{arrow-drop-up}{168}
\showcasematerialdesignfilled{arrow-forward}{169}
\showcasematerialdesignfilled{arrow-forward-ios}{170}
\showcasematerialdesignfilled{arrow-left}{171}
\showcasematerialdesignfilled{arrow-outward}{172}
\showcasematerialdesignfilled{arrow-right}{173}
\showcasematerialdesignfilled{arrow-right-alt}{174}
\showcasematerialdesignfilled{arrow-upward}{175}
\showcasematerialdesignfilled{article}{176}
\showcasematerialdesignfilled{art-track}{177}
\showcasematerialdesignfilled{aspect-ratio}{178}
\showcasematerialdesignfilled{assessment}{179}
\showcasematerialdesignfilled{assignment}{180}
\showcasematerialdesignfilled{assignment-ind}{181}
\showcasematerialdesignfilled{assignment-late}{182}
\showcasematerialdesignfilled{assignment-return}{183}
\showcasematerialdesignfilled{assignment-returned}{184}
\showcasematerialdesignfilled{assignment-turned-in}{185}
\showcasematerialdesignfilled{assistant}{186}
\showcasematerialdesignfilled{assistant-direction}{187}
\showcasematerialdesignfilled{assistant-photo}{188}
\showcasematerialdesignfilled{assist-walker}{189}
\showcasematerialdesignfilled{assured-workload}{190}
\showcasematerialdesignfilled{atm}{191}
\showcasematerialdesignfilled{attachment}{192}
\showcasematerialdesignfilled{attach-email}{193}
\showcasematerialdesignfilled{attach-file}{194}
\showcasematerialdesignfilled{attach-money}{195}
\showcasematerialdesignfilled{attractions}{196}
\showcasematerialdesignfilled{attribution}{197}
\showcasematerialdesignfilled{audiotrack}{198}
\showcasematerialdesignfilled{audio-file}{199}
\showcasematerialdesignfilled{autofps-select}{200}
\showcasematerialdesignfilled{autorenew}{201}
\showcasematerialdesignfilled{auto-awesome}{202}
\showcasematerialdesignfilled{auto-awesome-mosaic}{203}
\showcasematerialdesignfilled{auto-awesome-motion}{204}
\showcasematerialdesignfilled{auto-delete}{205}
\showcasematerialdesignfilled{auto-fix-high}{206}
\showcasematerialdesignfilled{auto-fix-normal}{207}
\showcasematerialdesignfilled{auto-fix-off}{208}
\showcasematerialdesignfilled{auto-graph}{209}
\showcasematerialdesignfilled{auto-mode}{210}
\showcasematerialdesignfilled{auto-stories}{211}
\showcasematerialdesignfilled{av-timer}{212}
\showcasematerialdesignfilled{baby-changing-station}{213}
\showcasematerialdesignfilled{backpack}{214}
\showcasematerialdesignfilled{backspace}{215}
\showcasematerialdesignfilled{backup}{216}
\showcasematerialdesignfilled{backup-table}{217}
\showcasematerialdesignfilled{back-hand}{218}
\showcasematerialdesignfilled{badge}{219}
\showcasematerialdesignfilled{bakery-dining}{220}
\showcasematerialdesignfilled{balance}{221}
\showcasematerialdesignfilled{balcony}{222}
\showcasematerialdesignfilled{ballot}{223}
\showcasematerialdesignfilled{bar-chart}{224}
\showcasematerialdesignfilled{batch-prediction}{225}
\showcasematerialdesignfilled{bathroom}{226}
\showcasematerialdesignfilled{bathtub}{227}
\showcasematerialdesignfilled{battery-0-bar}{228}
\showcasematerialdesignfilled{battery-1-bar}{229}
\showcasematerialdesignfilled{battery-2-bar}{230}
\showcasematerialdesignfilled{battery-3-bar}{231}
\showcasematerialdesignfilled{battery-4-bar}{232}
\showcasematerialdesignfilled{battery-5-bar}{233}
\showcasematerialdesignfilled{battery-6-bar}{234}
\showcasematerialdesignfilled{battery-alert}{235}
\showcasematerialdesignfilled{battery-charging-full}{236}
\showcasematerialdesignfilled{battery-full}{237}
\showcasematerialdesignfilled{battery-saver}{238}
\showcasematerialdesignfilled{battery-std}{239}
\showcasematerialdesignfilled{battery-unknown}{240}
\showcasematerialdesignfilled{beach-access}{241}
\showcasematerialdesignfilled{bed}{242}
\showcasematerialdesignfilled{bedroom-baby}{243}
\showcasematerialdesignfilled{bedroom-child}{244}
\showcasematerialdesignfilled{bedroom-parent}{245}
\showcasematerialdesignfilled{bedtime}{246}
\showcasematerialdesignfilled{bedtime-off}{247}
\showcasematerialdesignfilled{beenhere}{248}
\showcasematerialdesignfilled{bento}{249}
\showcasematerialdesignfilled{bike-scooter}{250}
\showcasematerialdesignfilled{biotech}{251}
\showcasematerialdesignfilled{blender}{252}
\showcasematerialdesignfilled{blind}{253}
\showcasematerialdesignfilled{blinds}{254}
\showcasematerialdesignfilled{blinds-closed}{255}
\showcasematerialdesignfilled{block}{256}
\showcasematerialdesignfilled{bloodtype}{257}
\showcasematerialdesignfilled{bluetooth}{258}
\showcasematerialdesignfilled{bluetooth-audio}{259}
\showcasematerialdesignfilled{bluetooth-connected}{260}
\showcasematerialdesignfilled{bluetooth-disabled}{261}
\showcasematerialdesignfilled{bluetooth-drive}{262}
\showcasematerialdesignfilled{bluetooth-searching}{263}
\showcasematerialdesignfilled{blur-circular}{264}
\showcasematerialdesignfilled{blur-linear}{265}
\showcasematerialdesignfilled{blur-off}{266}
\showcasematerialdesignfilled{blur-on}{267}
\showcasematerialdesignfilled{bolt}{268}
\showcasematerialdesignfilled{book}{269}
\showcasematerialdesignfilled{bookmark}{270}
\showcasematerialdesignfilled{bookmarks}{271}
\showcasematerialdesignfilled{bookmark-add}{272}
\showcasematerialdesignfilled{bookmark-added}{273}
\showcasematerialdesignfilled{bookmark-border}{274}
\showcasematerialdesignfilled{bookmark-remove}{275}
\showcasematerialdesignfilled{book-online}{276}
\showcasematerialdesignfilled{border-all}{277}
\showcasematerialdesignfilled{border-bottom}{278}
\showcasematerialdesignfilled{border-clear}{279}
\showcasematerialdesignfilled{border-color}{280}
\showcasematerialdesignfilled{border-horizontal}{281}
\showcasematerialdesignfilled{border-inner}{282}
\showcasematerialdesignfilled{border-left}{283}
\showcasematerialdesignfilled{border-outer}{284}
\showcasematerialdesignfilled{border-right}{285}
\showcasematerialdesignfilled{border-style}{286}
\showcasematerialdesignfilled{border-top}{287}
\showcasematerialdesignfilled{border-vertical}{288}
\showcasematerialdesignfilled{boy}{289}
\showcasematerialdesignfilled{branding-watermark}{290}
\showcasematerialdesignfilled{breakfast-dining}{291}
\showcasematerialdesignfilled{brightness-1}{292}
\showcasematerialdesignfilled{brightness-2}{293}
\showcasematerialdesignfilled{brightness-3}{294}
\showcasematerialdesignfilled{brightness-4}{295}
\showcasematerialdesignfilled{brightness-5}{296}
\showcasematerialdesignfilled{brightness-6}{297}
\showcasematerialdesignfilled{brightness-7}{298}
\showcasematerialdesignfilled{brightness-auto}{299}
\showcasematerialdesignfilled{brightness-high}{300}
\showcasematerialdesignfilled{brightness-low}{301}
\showcasematerialdesignfilled{brightness-medium}{302}
\showcasematerialdesignfilled{broadcast-on-home}{303}
\showcasematerialdesignfilled{broadcast-on-personal}{304}
\showcasematerialdesignfilled{broken-image}{305}
\showcasematerialdesignfilled{browser-not-supported}{306}
\showcasematerialdesignfilled{browser-updated}{307}
\showcasematerialdesignfilled{browse-gallery}{308}
\showcasematerialdesignfilled{brunch-dining}{309}
\showcasematerialdesignfilled{brush}{310}
\showcasematerialdesignfilled{bubble-chart}{311}
\showcasematerialdesignfilled{bug-report}{312}
\showcasematerialdesignfilled{build}{313}
\showcasematerialdesignfilled{build-circle}{314}
\showcasematerialdesignfilled{bungalow}{315}
\showcasematerialdesignfilled{burst-mode}{316}
\showcasematerialdesignfilled{business}{317}
\showcasematerialdesignfilled{business-center}{318}
\showcasematerialdesignfilled{bus-alert}{319}
\showcasematerialdesignfilled{cabin}{320}
\showcasematerialdesignfilled{cable}{321}
\showcasematerialdesignfilled{cached}{322}
\showcasematerialdesignfilled{cake}{323}
\showcasematerialdesignfilled{calculate}{324}
\showcasematerialdesignfilled{calendar-month}{325}
\showcasematerialdesignfilled{calendar-today}{326}
\showcasematerialdesignfilled{calendar-view-day}{327}
\showcasematerialdesignfilled{calendar-view-month}{328}
\showcasematerialdesignfilled{calendar-view-week}{329}
\showcasematerialdesignfilled{call}{330}
\showcasematerialdesignfilled{call-end}{331}
\showcasematerialdesignfilled{call-made}{332}
\showcasematerialdesignfilled{call-merge}{333}
\showcasematerialdesignfilled{call-missed}{334}
\showcasematerialdesignfilled{call-missed-outgoing}{335}
\showcasematerialdesignfilled{call-received}{336}
\showcasematerialdesignfilled{call-split}{337}
\showcasematerialdesignfilled{call-to-action}{338}
\showcasematerialdesignfilled{camera}{339}
\showcasematerialdesignfilled{cameraswitch}{340}
\showcasematerialdesignfilled{camera-alt}{341}
\showcasematerialdesignfilled{camera-enhance}{342}
\showcasematerialdesignfilled{camera-front}{343}
\showcasematerialdesignfilled{camera-indoor}{344}
\showcasematerialdesignfilled{camera-outdoor}{345}
\showcasematerialdesignfilled{camera-rear}{346}
\showcasematerialdesignfilled{camera-roll}{347}
\showcasematerialdesignfilled{campaign}{348}
\showcasematerialdesignfilled{cancel}{349}
\showcasematerialdesignfilled{cancel-presentation}{350}
\showcasematerialdesignfilled{cancel-schedule-send}{351}
\showcasematerialdesignfilled{candlestick-chart}{352}
\showcasematerialdesignfilled{card-giftcard}{353}
\showcasematerialdesignfilled{card-membership}{354}
\showcasematerialdesignfilled{card-travel}{355}
\showcasematerialdesignfilled{carpenter}{356}
\showcasematerialdesignfilled{car-crash}{357}
\showcasematerialdesignfilled{car-rental}{358}
\showcasematerialdesignfilled{car-repair}{359}
\showcasematerialdesignfilled{cases}{360}
\showcasematerialdesignfilled{casino}{361}
\showcasematerialdesignfilled{cast}{362}
\showcasematerialdesignfilled{castle}{363}
\showcasematerialdesignfilled{cast-connected}{364}
\showcasematerialdesignfilled{cast-for-education}{365}
\showcasematerialdesignfilled{catching-pokemon}{366}
\showcasematerialdesignfilled{category}{367}
\showcasematerialdesignfilled{celebration}{368}
\showcasematerialdesignfilled{cell-tower}{369}
\showcasematerialdesignfilled{cell-wifi}{370}
\showcasematerialdesignfilled{center-focus-strong}{371}
\showcasematerialdesignfilled{center-focus-weak}{372}
\showcasematerialdesignfilled{chair}{373}
\showcasematerialdesignfilled{chair-alt}{374}
\showcasematerialdesignfilled{chalet}{375}
\showcasematerialdesignfilled{change-circle}{376}
\showcasematerialdesignfilled{change-history}{377}
\showcasematerialdesignfilled{charging-station}{378}
\showcasematerialdesignfilled{chat}{379}
\showcasematerialdesignfilled{chat-bubble}{380}
\showcasematerialdesignfilled{chat-bubble-outline}{381}
\showcasematerialdesignfilled{check}{382}
\showcasematerialdesignfilled{checklist}{383}
\showcasematerialdesignfilled{checklist-rtl}{384}
\showcasematerialdesignfilled{checkroom}{385}
\showcasematerialdesignfilled{check-box}{386}
\showcasematerialdesignfilled{check-box-outline-blank}{387}
\showcasematerialdesignfilled{check-circle}{388}
\showcasematerialdesignfilled{check-circle-outline}{389}
\showcasematerialdesignfilled{chevron-left}{390}
\showcasematerialdesignfilled{chevron-right}{391}
\showcasematerialdesignfilled{child-care}{392}
\showcasematerialdesignfilled{child-friendly}{393}
\showcasematerialdesignfilled{chrome-reader-mode}{394}
\showcasematerialdesignfilled{church}{395}
\showcasematerialdesignfilled{circle}{396}
\showcasematerialdesignfilled{circle-notifications}{397}
\showcasematerialdesignfilled{class}{398}
\showcasematerialdesignfilled{cleaning-services}{399}
\showcasematerialdesignfilled{clean-hands}{400}
\showcasematerialdesignfilled{clear}{401}
\showcasematerialdesignfilled{clear-all}{402}
\showcasematerialdesignfilled{close}{403}
\showcasematerialdesignfilled{closed-caption}{404}
\showcasematerialdesignfilled{closed-caption-disabled}{405}
\showcasematerialdesignfilled{closed-caption-off}{406}
\showcasematerialdesignfilled{close-fullscreen}{407}
\showcasematerialdesignfilled{cloud}{408}
\showcasematerialdesignfilled{cloud-circle}{409}
\showcasematerialdesignfilled{cloud-done}{410}
\showcasematerialdesignfilled{cloud-download}{411}
\showcasematerialdesignfilled{cloud-off}{412}
\showcasematerialdesignfilled{cloud-queue}{413}
\showcasematerialdesignfilled{cloud-sync}{414}
\showcasematerialdesignfilled{cloud-upload}{415}
\showcasematerialdesignfilled{co2}{416}
\showcasematerialdesignfilled{code}{417}
\showcasematerialdesignfilled{code-off}{418}
\showcasematerialdesignfilled{coffee}{419}
\showcasematerialdesignfilled{coffee-maker}{420}
\showcasematerialdesignfilled{collections}{421}
\showcasematerialdesignfilled{collections-bookmark}{422}
\showcasematerialdesignfilled{colorize}{423}
\showcasematerialdesignfilled{color-lens}{424}
\showcasematerialdesignfilled{comment}{425}
\showcasematerialdesignfilled{comments-disabled}{426}
\showcasematerialdesignfilled{comment-bank}{427}
\showcasematerialdesignfilled{commit}{428}
\showcasematerialdesignfilled{commute}{429}
\showcasematerialdesignfilled{compare}{430}
\showcasematerialdesignfilled{compare-arrows}{431}
\showcasematerialdesignfilled{compass-calibration}{432}
\showcasematerialdesignfilled{compost}{433}
\showcasematerialdesignfilled{compress}{434}
\showcasematerialdesignfilled{computer}{435}
\showcasematerialdesignfilled{confirmation-number}{436}
\showcasematerialdesignfilled{connected-tv}{437}
\showcasematerialdesignfilled{connecting-airports}{438}
\showcasematerialdesignfilled{connect-without-contact}{439}
\showcasematerialdesignfilled{construction}{440}
\showcasematerialdesignfilled{contactless}{441}
\showcasematerialdesignfilled{contacts}{442}
\showcasematerialdesignfilled{contact-emergency}{443}
\showcasematerialdesignfilled{contact-mail}{444}
\showcasematerialdesignfilled{contact-page}{445}
\showcasematerialdesignfilled{contact-phone}{446}
\showcasematerialdesignfilled{contact-support}{447}
\showcasematerialdesignfilled{content-copy}{448}
\showcasematerialdesignfilled{content-cut}{449}
\showcasematerialdesignfilled{content-paste}{450}
\showcasematerialdesignfilled{content-paste-go}{451}
\showcasematerialdesignfilled{content-paste-off}{452}
\showcasematerialdesignfilled{content-paste-search}{453}
\showcasematerialdesignfilled{contrast}{454}
\showcasematerialdesignfilled{control-camera}{455}
\showcasematerialdesignfilled{control-point}{456}
\showcasematerialdesignfilled{control-point-duplicate}{457}
\showcasematerialdesignfilled{cookie}{458}
\showcasematerialdesignfilled{copyright}{459}
\showcasematerialdesignfilled{copy-all}{460}
\showcasematerialdesignfilled{coronavirus}{461}
\showcasematerialdesignfilled{corporate-fare}{462}
\showcasematerialdesignfilled{cottage}{463}
\showcasematerialdesignfilled{countertops}{464}
\showcasematerialdesignfilled{co-present}{465}
\showcasematerialdesignfilled{create}{466}
\showcasematerialdesignfilled{create-new-folder}{467}
\showcasematerialdesignfilled{credit-card}{468}
\showcasematerialdesignfilled{credit-card-off}{469}
\showcasematerialdesignfilled{credit-score}{470}
\showcasematerialdesignfilled{crib}{471}
\showcasematerialdesignfilled{crisis-alert}{472}
\showcasematerialdesignfilled{crop}{473}
\showcasematerialdesignfilled{crop-3-2}{474}
\showcasematerialdesignfilled{crop-5-4}{475}
\showcasematerialdesignfilled{crop-7-5}{476}
\showcasematerialdesignfilled{crop-16-9}{477}
\showcasematerialdesignfilled{crop-din}{478}
\showcasematerialdesignfilled{crop-free}{479}
\showcasematerialdesignfilled{crop-landscape}{480}
\showcasematerialdesignfilled{crop-original}{481}
\showcasematerialdesignfilled{crop-portrait}{482}
\showcasematerialdesignfilled{crop-rotate}{483}
\showcasematerialdesignfilled{crop-square}{484}
\showcasematerialdesignfilled{cruelty-free}{485}
\showcasematerialdesignfilled{css}{486}
\showcasematerialdesignfilled{currency-bitcoin}{487}
\showcasematerialdesignfilled{currency-exchange}{488}
\showcasematerialdesignfilled{currency-franc}{489}
\showcasematerialdesignfilled{currency-lira}{490}
\showcasematerialdesignfilled{currency-pound}{491}
\showcasematerialdesignfilled{currency-ruble}{492}
\showcasematerialdesignfilled{currency-rupee}{493}
\showcasematerialdesignfilled{currency-yen}{494}
\showcasematerialdesignfilled{currency-yuan}{495}
\showcasematerialdesignfilled{curtains}{496}
\showcasematerialdesignfilled{curtains-closed}{497}
\showcasematerialdesignfilled{cyclone}{498}
\showcasematerialdesignfilled{dangerous}{499}
\showcasematerialdesignfilled{dark-mode}{500}
\showcasematerialdesignfilled{dashboard}{501}
\showcasematerialdesignfilled{dashboard-customize}{502}
\showcasematerialdesignfilled{dataset}{503}
\showcasematerialdesignfilled{dataset-linked}{504}
\showcasematerialdesignfilled{data-array}{505}
\showcasematerialdesignfilled{data-exploration}{506}
\showcasematerialdesignfilled{data-object}{507}
\showcasematerialdesignfilled{data-saver-off}{508}
\showcasematerialdesignfilled{data-saver-on}{509}
\showcasematerialdesignfilled{data-thresholding}{510}
\showcasematerialdesignfilled{data-usage}{511}
\showcasematerialdesignfilled{date-range}{512}
\showcasematerialdesignfilled{deblur}{513}
\showcasematerialdesignfilled{deck}{514}
\showcasematerialdesignfilled{dehaze}{515}
\showcasematerialdesignfilled{delete}{516}
\showcasematerialdesignfilled{delete-forever}{517}
\showcasematerialdesignfilled{delete-outline}{518}
\showcasematerialdesignfilled{delete-sweep}{519}
\showcasematerialdesignfilled{delivery-dining}{520}
\showcasematerialdesignfilled{density-large}{521}
\showcasematerialdesignfilled{density-medium}{522}
\showcasematerialdesignfilled{density-small}{523}
\showcasematerialdesignfilled{departure-board}{524}
\showcasematerialdesignfilled{description}{525}
\showcasematerialdesignfilled{deselect}{526}
\showcasematerialdesignfilled{design-services}{527}
\showcasematerialdesignfilled{desk}{528}
\showcasematerialdesignfilled{desktop-access-disabled}{529}
\showcasematerialdesignfilled{desktop-mac}{530}
\showcasematerialdesignfilled{desktop-windows}{531}
\showcasematerialdesignfilled{details}{532}
\showcasematerialdesignfilled{developer-board}{533}
\showcasematerialdesignfilled{developer-board-off}{534}
\showcasematerialdesignfilled{developer-mode}{535}
\showcasematerialdesignfilled{devices}{536}
\showcasematerialdesignfilled{devices-fold}{537}
\showcasematerialdesignfilled{devices-other}{538}
\showcasematerialdesignfilled{device-hub}{539}
\showcasematerialdesignfilled{device-thermostat}{540}
\showcasematerialdesignfilled{device-unknown}{541}
\showcasematerialdesignfilled{dialer-sip}{542}
\showcasematerialdesignfilled{dialpad}{543}
\showcasematerialdesignfilled{diamond}{544}
\showcasematerialdesignfilled{difference}{545}
\showcasematerialdesignfilled{dining}{546}
\showcasematerialdesignfilled{dinner-dining}{547}
\showcasematerialdesignfilled{directions}{548}
\showcasematerialdesignfilled{directions-bike}{549}
\showcasematerialdesignfilled{directions-boat}{550}
\showcasematerialdesignfilled{directions-boat-filled}{551}
\showcasematerialdesignfilled{directions-bus}{552}
\showcasematerialdesignfilled{directions-bus-filled}{553}
\showcasematerialdesignfilled{directions-car}{554}
\showcasematerialdesignfilled{directions-car-filled}{555}
\showcasematerialdesignfilled{directions-off}{556}
\showcasematerialdesignfilled{directions-railway}{557}
\showcasematerialdesignfilled{directions-railway-filled}{558}
\showcasematerialdesignfilled{directions-run}{559}
\showcasematerialdesignfilled{directions-subway}{560}
\showcasematerialdesignfilled{directions-subway-filled}{561}
\showcasematerialdesignfilled{directions-transit}{562}
\showcasematerialdesignfilled{directions-transit-filled}{563}
\showcasematerialdesignfilled{directions-walk}{564}
\showcasematerialdesignfilled{dirty-lens}{565}
\showcasematerialdesignfilled{disabled-by-default}{566}
\showcasematerialdesignfilled{disabled-visible}{567}
\showcasematerialdesignfilled{discount}{568}
\showcasematerialdesignfilled{disc-full}{569}
\showcasematerialdesignfilled{display-settings}{570}
\showcasematerialdesignfilled{diversity-1}{571}
\showcasematerialdesignfilled{diversity-2}{572}
\showcasematerialdesignfilled{diversity-3}{573}
\showcasematerialdesignfilled{dns}{574}
\showcasematerialdesignfilled{dock}{575}
\showcasematerialdesignfilled{document-scanner}{576}
\showcasematerialdesignfilled{domain}{577}
\showcasematerialdesignfilled{domain-add}{578}
\showcasematerialdesignfilled{domain-disabled}{579}
\showcasematerialdesignfilled{domain-verification}{580}
\showcasematerialdesignfilled{done}{581}
\showcasematerialdesignfilled{done-all}{582}
\showcasematerialdesignfilled{done-outline}{583}
\showcasematerialdesignfilled{donut-large}{584}
\showcasematerialdesignfilled{donut-small}{585}
\showcasematerialdesignfilled{doorbell}{586}
\showcasematerialdesignfilled{door-back}{587}
\showcasematerialdesignfilled{door-front}{588}
\showcasematerialdesignfilled{door-sliding}{589}
\showcasematerialdesignfilled{double-arrow}{590}
\showcasematerialdesignfilled{downhill-skiing}{591}
\showcasematerialdesignfilled{download}{592}
\showcasematerialdesignfilled{downloading}{593}
\showcasematerialdesignfilled{download-done}{594}
\showcasematerialdesignfilled{download-for-offline}{595}
\showcasematerialdesignfilled{do-disturb}{596}
\showcasematerialdesignfilled{do-disturb-alt}{597}
\showcasematerialdesignfilled{do-disturb-off}{598}
\showcasematerialdesignfilled{do-disturb-on}{599}
\showcasematerialdesignfilled{do-not-disturb}{600}
\showcasematerialdesignfilled{do-not-disturb-alt}{601}
\showcasematerialdesignfilled{do-not-disturb-off}{602}
\showcasematerialdesignfilled{do-not-disturb-on}{603}
\showcasematerialdesignfilled{do-not-disturb-on-total-silence}{604}
\showcasematerialdesignfilled{do-not-step}{605}
\showcasematerialdesignfilled{do-not-touch}{606}
\showcasematerialdesignfilled{drafts}{607}
\showcasematerialdesignfilled{drag-handle}{608}
\showcasematerialdesignfilled{drag-indicator}{609}
\showcasematerialdesignfilled{draw}{610}
\showcasematerialdesignfilled{drive-eta}{611}
\showcasematerialdesignfilled{drive-file-move}{612}
\showcasematerialdesignfilled{drive-file-move-rtl}{613}
\showcasematerialdesignfilled{drive-file-rename-outline}{614}
\showcasematerialdesignfilled{drive-folder-upload}{615}
\showcasematerialdesignfilled{dry}{616}
\showcasematerialdesignfilled{dry-cleaning}{617}
\showcasematerialdesignfilled{duo}{618}
\showcasematerialdesignfilled{dvr}{619}
\showcasematerialdesignfilled{dynamic-feed}{620}
\showcasematerialdesignfilled{dynamic-form}{621}
\showcasematerialdesignfilled{earbuds}{622}
\showcasematerialdesignfilled{earbuds-battery}{623}
\showcasematerialdesignfilled{east}{624}
\showcasematerialdesignfilled{edgesensor-high}{625}
\showcasematerialdesignfilled{edgesensor-low}{626}
\showcasematerialdesignfilled{edit}{627}
\showcasematerialdesignfilled{edit-attributes}{628}
\showcasematerialdesignfilled{edit-calendar}{629}
\showcasematerialdesignfilled{edit-location}{630}
\showcasematerialdesignfilled{edit-location-alt}{631}
\showcasematerialdesignfilled{edit-note}{632}
\showcasematerialdesignfilled{edit-notifications}{633}
\showcasematerialdesignfilled{edit-off}{634}
\showcasematerialdesignfilled{edit-road}{635}
\showcasematerialdesignfilled{egg}{636}
\showcasematerialdesignfilled{egg-alt}{637}
\showcasematerialdesignfilled{eject}{638}
\showcasematerialdesignfilled{elderly}{639}
\showcasematerialdesignfilled{elderly-woman}{640}
\showcasematerialdesignfilled{electrical-services}{641}
\showcasematerialdesignfilled{electric-bike}{642}
\showcasematerialdesignfilled{electric-bolt}{643}
\showcasematerialdesignfilled{electric-car}{644}
\showcasematerialdesignfilled{electric-meter}{645}
\showcasematerialdesignfilled{electric-moped}{646}
\showcasematerialdesignfilled{electric-rickshaw}{647}
\showcasematerialdesignfilled{electric-scooter}{648}
\showcasematerialdesignfilled{elevator}{649}
\showcasematerialdesignfilled{email}{650}
\showcasematerialdesignfilled{emergency}{651}
\showcasematerialdesignfilled{emergency-recording}{652}
\showcasematerialdesignfilled{emergency-share}{653}
\showcasematerialdesignfilled{emoji-emotions}{654}
\showcasematerialdesignfilled{emoji-events}{655}
\showcasematerialdesignfilled{emoji-food-beverage}{656}
\showcasematerialdesignfilled{emoji-nature}{657}
\showcasematerialdesignfilled{emoji-objects}{658}
\showcasematerialdesignfilled{emoji-people}{659}
\showcasematerialdesignfilled{emoji-symbols}{660}
\showcasematerialdesignfilled{emoji-transportation}{661}
\showcasematerialdesignfilled{energy-savings-leaf}{662}
\showcasematerialdesignfilled{engineering}{663}
\showcasematerialdesignfilled{enhanced-encryption}{664}
\showcasematerialdesignfilled{equalizer}{665}
\showcasematerialdesignfilled{error}{666}
\showcasematerialdesignfilled{error-outline}{667}
\showcasematerialdesignfilled{escalator}{668}
\showcasematerialdesignfilled{escalator-warning}{669}
\showcasematerialdesignfilled{euro}{670}
\showcasematerialdesignfilled{euro-symbol}{671}
\showcasematerialdesignfilled{event}{672}
\showcasematerialdesignfilled{event-available}{673}
\showcasematerialdesignfilled{event-busy}{674}
\showcasematerialdesignfilled{event-note}{675}
\showcasematerialdesignfilled{event-repeat}{676}
\showcasematerialdesignfilled{event-seat}{677}
\showcasematerialdesignfilled{ev-station}{678}
\showcasematerialdesignfilled{exit-to-app}{679}
\showcasematerialdesignfilled{expand}{680}
\showcasematerialdesignfilled{expand-circle-down}{681}
\showcasematerialdesignfilled{expand-less}{682}
\showcasematerialdesignfilled{expand-more}{683}
\showcasematerialdesignfilled{explicit}{684}
\showcasematerialdesignfilled{explore}{685}
\showcasematerialdesignfilled{explore-off}{686}
\showcasematerialdesignfilled{exposure}{687}
\showcasematerialdesignfilled{exposure-neg-1}{688}
\showcasematerialdesignfilled{exposure-neg-2}{689}
\showcasematerialdesignfilled{exposure-plus-1}{690}
\showcasematerialdesignfilled{exposure-plus-2}{691}
\showcasematerialdesignfilled{exposure-zero}{692}
\showcasematerialdesignfilled{extension}{693}
\showcasematerialdesignfilled{extension-off}{694}
\showcasematerialdesignfilled{e-mobiledata}{695}
\showcasematerialdesignfilled{face}{696}
\showcasematerialdesignfilled{face-2}{697}
\showcasematerialdesignfilled{face-3}{698}
\showcasematerialdesignfilled{face-4}{699}
\showcasematerialdesignfilled{face-5}{700}
\showcasematerialdesignfilled{face-6}{701}
\showcasematerialdesignfilled{face-retouching-natural}{702}
\showcasematerialdesignfilled{face-retouching-off}{703}
\showcasematerialdesignfilled{factory}{704}
\showcasematerialdesignfilled{fact-check}{705}
\showcasematerialdesignfilled{family-restroom}{706}
\showcasematerialdesignfilled{fastfood}{707}
\showcasematerialdesignfilled{fast-forward}{708}
\showcasematerialdesignfilled{fast-rewind}{709}
\showcasematerialdesignfilled{favorite}{710}
\showcasematerialdesignfilled{favorite-border}{711}
\showcasematerialdesignfilled{fax}{712}
\showcasematerialdesignfilled{featured-play-list}{713}
\showcasematerialdesignfilled{featured-video}{714}
\showcasematerialdesignfilled{feed}{715}
\showcasematerialdesignfilled{feedback}{716}
\showcasematerialdesignfilled{female}{717}
\showcasematerialdesignfilled{fence}{718}
\showcasematerialdesignfilled{festival}{719}
\showcasematerialdesignfilled{fiber-dvr}{720}
\showcasematerialdesignfilled{fiber-manual-record}{721}
\showcasematerialdesignfilled{fiber-new}{722}
\showcasematerialdesignfilled{fiber-pin}{723}
\showcasematerialdesignfilled{fiber-smart-record}{724}
\showcasematerialdesignfilled{file-copy}{725}
\showcasematerialdesignfilled{file-download}{726}
\showcasematerialdesignfilled{file-download-done}{727}
\showcasematerialdesignfilled{file-download-off}{728}
\showcasematerialdesignfilled{file-open}{729}
\showcasematerialdesignfilled{file-present}{730}
\showcasematerialdesignfilled{file-upload}{731}
\showcasematerialdesignfilled{filter}{732}
\showcasematerialdesignfilled{filter-1}{733}
\showcasematerialdesignfilled{filter-2}{734}
\showcasematerialdesignfilled{filter-3}{735}
\showcasematerialdesignfilled{filter-4}{736}
\showcasematerialdesignfilled{filter-5}{737}
\showcasematerialdesignfilled{filter-6}{738}
\showcasematerialdesignfilled{filter-7}{739}
\showcasematerialdesignfilled{filter-8}{740}
\showcasematerialdesignfilled{filter-9}{741}
\showcasematerialdesignfilled{filter-9-plus}{742}
\showcasematerialdesignfilled{filter-alt}{743}
\showcasematerialdesignfilled{filter-alt-off}{744}
\showcasematerialdesignfilled{filter-b-and-w}{745}
\showcasematerialdesignfilled{filter-center-focus}{746}
\showcasematerialdesignfilled{filter-drama}{747}
\showcasematerialdesignfilled{filter-frames}{748}
\showcasematerialdesignfilled{filter-hdr}{749}
\showcasematerialdesignfilled{filter-list}{750}
\showcasematerialdesignfilled{filter-list-off}{751}
\showcasematerialdesignfilled{filter-none}{752}
\showcasematerialdesignfilled{filter-tilt-shift}{753}
\showcasematerialdesignfilled{filter-vintage}{754}
\showcasematerialdesignfilled{find-in-page}{755}
\showcasematerialdesignfilled{find-replace}{756}
\showcasematerialdesignfilled{fingerprint}{757}
\showcasematerialdesignfilled{fireplace}{758}
\showcasematerialdesignfilled{fire-extinguisher}{759}
\showcasematerialdesignfilled{fire-hydrant-alt}{760}
\showcasematerialdesignfilled{fire-truck}{761}
\showcasematerialdesignfilled{first-page}{762}
\showcasematerialdesignfilled{fitbit}{763}
\showcasematerialdesignfilled{fitness-center}{764}
\showcasematerialdesignfilled{fit-screen}{765}
\showcasematerialdesignfilled{flag}{766}
\showcasematerialdesignfilled{flag-circle}{767}
\showcasematerialdesignfilled{flaky}{768}
\showcasematerialdesignfilled{flare}{769}
\showcasematerialdesignfilled{flashlight-off}{770}
\showcasematerialdesignfilled{flashlight-on}{771}
\showcasematerialdesignfilled{flash-auto}{772}
\showcasematerialdesignfilled{flash-off}{773}
\showcasematerialdesignfilled{flash-on}{774}
\showcasematerialdesignfilled{flatware}{775}
\showcasematerialdesignfilled{flight}{776}
\showcasematerialdesignfilled{flight-class}{777}
\showcasematerialdesignfilled{flight-land}{778}
\showcasematerialdesignfilled{flight-takeoff}{779}
\showcasematerialdesignfilled{flip}{780}
\showcasematerialdesignfilled{flip-camera-android}{781}
\showcasematerialdesignfilled{flip-camera-ios}{782}
\showcasematerialdesignfilled{flip-to-back}{783}
\showcasematerialdesignfilled{flip-to-front}{784}
\showcasematerialdesignfilled{flood}{785}
\showcasematerialdesignfilled{fluorescent}{786}
\showcasematerialdesignfilled{flutter-dash}{787}
\showcasematerialdesignfilled{fmd-bad}{788}
\showcasematerialdesignfilled{fmd-good}{789}
\showcasematerialdesignfilled{folder}{790}
\showcasematerialdesignfilled{folder-copy}{791}
\showcasematerialdesignfilled{folder-delete}{792}
\showcasematerialdesignfilled{folder-off}{793}
\showcasematerialdesignfilled{folder-open}{794}
\showcasematerialdesignfilled{folder-shared}{795}
\showcasematerialdesignfilled{folder-special}{796}
\showcasematerialdesignfilled{folder-zip}{797}
\showcasematerialdesignfilled{follow-the-signs}{798}
\showcasematerialdesignfilled{font-download}{799}
\showcasematerialdesignfilled{font-download-off}{800}
\showcasematerialdesignfilled{food-bank}{801}
\showcasematerialdesignfilled{forest}{802}
\showcasematerialdesignfilled{fork-left}{803}
\showcasematerialdesignfilled{fork-right}{804}
\showcasematerialdesignfilled{format-align-center}{805}
\showcasematerialdesignfilled{format-align-justify}{806}
\showcasematerialdesignfilled{format-align-left}{807}
\showcasematerialdesignfilled{format-align-right}{808}
\showcasematerialdesignfilled{format-bold}{809}
\showcasematerialdesignfilled{format-clear}{810}
\showcasematerialdesignfilled{format-color-fill}{811}
\showcasematerialdesignfilled{format-color-reset}{812}
\showcasematerialdesignfilled{format-color-text}{813}
\showcasematerialdesignfilled{format-indent-decrease}{814}
\showcasematerialdesignfilled{format-indent-increase}{815}
\showcasematerialdesignfilled{format-italic}{816}
\showcasematerialdesignfilled{format-line-spacing}{817}
\showcasematerialdesignfilled{format-list-bulleted}{818}
\showcasematerialdesignfilled{format-list-numbered}{819}
\showcasematerialdesignfilled{format-list-numbered-rtl}{820}
\showcasematerialdesignfilled{format-overline}{821}
\showcasematerialdesignfilled{format-paint}{822}
\showcasematerialdesignfilled{format-quote}{823}
\showcasematerialdesignfilled{format-shapes}{824}
\showcasematerialdesignfilled{format-size}{825}
\showcasematerialdesignfilled{format-strikethrough}{826}
\showcasematerialdesignfilled{format-textdirection-l-to-r}{827}
\showcasematerialdesignfilled{format-textdirection-r-to-l}{828}
\showcasematerialdesignfilled{format-underlined}{829}
\showcasematerialdesignfilled{fort}{830}
\showcasematerialdesignfilled{forum}{831}
\showcasematerialdesignfilled{forward}{832}
\showcasematerialdesignfilled{forward-5}{833}
\showcasematerialdesignfilled{forward-10}{834}
\showcasematerialdesignfilled{forward-30}{835}
\showcasematerialdesignfilled{forward-to-inbox}{836}
\showcasematerialdesignfilled{foundation}{837}
\showcasematerialdesignfilled{free-breakfast}{838}
\showcasematerialdesignfilled{free-cancellation}{839}
\showcasematerialdesignfilled{front-hand}{840}
\showcasematerialdesignfilled{fullscreen}{841}
\showcasematerialdesignfilled{fullscreen-exit}{842}
\showcasematerialdesignfilled{functions}{843}
\showcasematerialdesignfilled{gamepad}{844}
\showcasematerialdesignfilled{games}{845}
\showcasematerialdesignfilled{garage}{846}
\showcasematerialdesignfilled{gas-meter}{847}
\showcasematerialdesignfilled{gavel}{848}
\showcasematerialdesignfilled{generating-tokens}{849}
\showcasematerialdesignfilled{gesture}{850}
\showcasematerialdesignfilled{get-app}{851}
\showcasematerialdesignfilled{gif}{852}
\showcasematerialdesignfilled{gif-box}{853}
\showcasematerialdesignfilled{girl}{854}
\showcasematerialdesignfilled{gite}{855}
\showcasematerialdesignfilled{golf-course}{856}
\showcasematerialdesignfilled{gpp-bad}{857}
\showcasematerialdesignfilled{gpp-good}{858}
\showcasematerialdesignfilled{gpp-maybe}{859}
\showcasematerialdesignfilled{gps-fixed}{860}
\showcasematerialdesignfilled{gps-not-fixed}{861}
\showcasematerialdesignfilled{gps-off}{862}
\showcasematerialdesignfilled{grade}{863}
\showcasematerialdesignfilled{gradient}{864}
\showcasematerialdesignfilled{grading}{865}
\showcasematerialdesignfilled{grain}{866}
\showcasematerialdesignfilled{graphic-eq}{867}
\showcasematerialdesignfilled{grass}{868}
\showcasematerialdesignfilled{grid-3x3}{869}
\showcasematerialdesignfilled{grid-4x4}{870}
\showcasematerialdesignfilled{grid-goldenratio}{871}
\showcasematerialdesignfilled{grid-off}{872}
\showcasematerialdesignfilled{grid-on}{873}
\showcasematerialdesignfilled{grid-view}{874}
\showcasematerialdesignfilled{group}{875}
\showcasematerialdesignfilled{groups}{876}
\showcasematerialdesignfilled{groups-2}{877}
\showcasematerialdesignfilled{groups-3}{878}
\showcasematerialdesignfilled{group-add}{879}
\showcasematerialdesignfilled{group-off}{880}
\showcasematerialdesignfilled{group-remove}{881}
\showcasematerialdesignfilled{group-work}{882}
\showcasematerialdesignfilled{g-mobiledata}{883}
\showcasematerialdesignfilled{g-translate}{884}
\showcasematerialdesignfilled{hail}{885}
\showcasematerialdesignfilled{handshake}{886}
\showcasematerialdesignfilled{handyman}{887}
\showcasematerialdesignfilled{hardware}{888}
\showcasematerialdesignfilled{hd}{889}
\showcasematerialdesignfilled{hdr-auto}{890}
\showcasematerialdesignfilled{hdr-auto-select}{891}
\showcasematerialdesignfilled{hdr-enhanced-select}{892}
\showcasematerialdesignfilled{hdr-off}{893}
\showcasematerialdesignfilled{hdr-off-select}{894}
\showcasematerialdesignfilled{hdr-on}{895}
\showcasematerialdesignfilled{hdr-on-select}{896}
\showcasematerialdesignfilled{hdr-plus}{897}
\showcasematerialdesignfilled{hdr-strong}{898}
\showcasematerialdesignfilled{hdr-weak}{899}
\showcasematerialdesignfilled{headphones}{900}
\showcasematerialdesignfilled{headphones-battery}{901}
\showcasematerialdesignfilled{headset}{902}
\showcasematerialdesignfilled{headset-mic}{903}
\showcasematerialdesignfilled{headset-off}{904}
\showcasematerialdesignfilled{healing}{905}
\showcasematerialdesignfilled{health-and-safety}{906}
\showcasematerialdesignfilled{hearing}{907}
\showcasematerialdesignfilled{hearing-disabled}{908}
\showcasematerialdesignfilled{heart-broken}{909}
\showcasematerialdesignfilled{heat-pump}{910}
\showcasematerialdesignfilled{height}{911}
\showcasematerialdesignfilled{help}{912}
\showcasematerialdesignfilled{help-center}{913}
\showcasematerialdesignfilled{help-outline}{914}
\showcasematerialdesignfilled{hevc}{915}
\showcasematerialdesignfilled{hexagon}{916}
\showcasematerialdesignfilled{hide-image}{917}
\showcasematerialdesignfilled{hide-source}{918}
\showcasematerialdesignfilled{highlight}{919}
\showcasematerialdesignfilled{highlight-alt}{920}
\showcasematerialdesignfilled{highlight-off}{921}
\showcasematerialdesignfilled{high-quality}{922}
\showcasematerialdesignfilled{hiking}{923}
\showcasematerialdesignfilled{history}{924}
\showcasematerialdesignfilled{history-edu}{925}
\showcasematerialdesignfilled{history-toggle-off}{926}
\showcasematerialdesignfilled{hive}{927}
\showcasematerialdesignfilled{hls}{928}
\showcasematerialdesignfilled{hls-off}{929}
\showcasematerialdesignfilled{holiday-village}{930}
\showcasematerialdesignfilled{home}{931}
\showcasematerialdesignfilled{home-max}{932}
\showcasematerialdesignfilled{home-mini}{933}
\showcasematerialdesignfilled{home-repair-service}{934}
\showcasematerialdesignfilled{home-work}{935}
\showcasematerialdesignfilled{horizontal-distribute}{936}
\showcasematerialdesignfilled{horizontal-rule}{937}
\showcasematerialdesignfilled{horizontal-split}{938}
\showcasematerialdesignfilled{hotel}{939}
\showcasematerialdesignfilled{hotel-class}{940}
\showcasematerialdesignfilled{hot-tub}{941}
\showcasematerialdesignfilled{hourglass-bottom}{942}
\showcasematerialdesignfilled{hourglass-disabled}{943}
\showcasematerialdesignfilled{hourglass-empty}{944}
\showcasematerialdesignfilled{hourglass-full}{945}
\showcasematerialdesignfilled{hourglass-top}{946}
\showcasematerialdesignfilled{house}{947}
\showcasematerialdesignfilled{houseboat}{948}
\showcasematerialdesignfilled{house-siding}{949}
\showcasematerialdesignfilled{how-to-reg}{950}
\showcasematerialdesignfilled{how-to-vote}{951}
\showcasematerialdesignfilled{html}{952}
\showcasematerialdesignfilled{http}{953}
\showcasematerialdesignfilled{https}{954}
\showcasematerialdesignfilled{hub}{955}
\showcasematerialdesignfilled{hvac}{956}
\showcasematerialdesignfilled{h-mobiledata}{957}
\showcasematerialdesignfilled{h-plus-mobiledata}{958}
\showcasematerialdesignfilled{icecream}{959}
\showcasematerialdesignfilled{ice-skating}{960}
\showcasematerialdesignfilled{image}{961}
\showcasematerialdesignfilled{imagesearch-roller}{962}
\showcasematerialdesignfilled{image-aspect-ratio}{963}
\showcasematerialdesignfilled{image-not-supported}{964}
\showcasematerialdesignfilled{image-search}{965}
\showcasematerialdesignfilled{important-devices}{966}
\showcasematerialdesignfilled{import-contacts}{967}
\showcasematerialdesignfilled{import-export}{968}
\showcasematerialdesignfilled{inbox}{969}
\showcasematerialdesignfilled{incomplete-circle}{970}
\showcasematerialdesignfilled{indeterminate-check-box}{971}
\showcasematerialdesignfilled{info}{972}
\showcasematerialdesignfilled{input}{973}
\showcasematerialdesignfilled{insert-chart}{974}
\showcasematerialdesignfilled{insert-chart-outlined}{975}
\showcasematerialdesignfilled{insert-comment}{976}
\showcasematerialdesignfilled{insert-drive-file}{977}
\showcasematerialdesignfilled{insert-emoticon}{978}
\showcasematerialdesignfilled{insert-invitation}{979}
\showcasematerialdesignfilled{insert-link}{980}
\showcasematerialdesignfilled{insert-page-break}{981}
\showcasematerialdesignfilled{insert-photo}{982}
\showcasematerialdesignfilled{insights}{983}
\showcasematerialdesignfilled{install-desktop}{984}
\showcasematerialdesignfilled{install-mobile}{985}
\showcasematerialdesignfilled{integration-instructions}{986}
\showcasematerialdesignfilled{interests}{987}
\showcasematerialdesignfilled{interpreter-mode}{988}
\showcasematerialdesignfilled{inventory}{989}
\showcasematerialdesignfilled{inventory-2}{990}
\showcasematerialdesignfilled{invert-colors}{991}
\showcasematerialdesignfilled{invert-colors-off}{992}
\showcasematerialdesignfilled{ios-share}{993}
\showcasematerialdesignfilled{iron}{994}
\showcasematerialdesignfilled{iso}{995}
\showcasematerialdesignfilled{javascript}{996}
\showcasematerialdesignfilled{join-full}{997}
\showcasematerialdesignfilled{join-inner}{998}
\showcasematerialdesignfilled{join-left}{999}
\showcasematerialdesignfilled{join-right}{1000}
\showcasematerialdesignfilled{kayaking}{1001}
\showcasematerialdesignfilled{kebab-dining}{1002}
\showcasematerialdesignfilled{key}{1003}
\showcasematerialdesignfilled{keyboard}{1004}
\showcasematerialdesignfilled{keyboard-alt}{1005}
\showcasematerialdesignfilled{keyboard-arrow-down}{1006}
\showcasematerialdesignfilled{keyboard-arrow-left}{1007}
\showcasematerialdesignfilled{keyboard-arrow-right}{1008}
\showcasematerialdesignfilled{keyboard-arrow-up}{1009}
\showcasematerialdesignfilled{keyboard-backspace}{1010}
\showcasematerialdesignfilled{keyboard-capslock}{1011}
\showcasematerialdesignfilled{keyboard-command-key}{1012}
\showcasematerialdesignfilled{keyboard-control-key}{1013}
\showcasematerialdesignfilled{keyboard-double-arrow-down}{1014}
\showcasematerialdesignfilled{keyboard-double-arrow-left}{1015}
\showcasematerialdesignfilled{keyboard-double-arrow-right}{1016}
\showcasematerialdesignfilled{keyboard-double-arrow-up}{1017}
\showcasematerialdesignfilled{keyboard-hide}{1018}
\showcasematerialdesignfilled{keyboard-option-key}{1019}
\showcasematerialdesignfilled{keyboard-return}{1020}
\showcasematerialdesignfilled{keyboard-tab}{1021}
\showcasematerialdesignfilled{keyboard-voice}{1022}
\showcasematerialdesignfilled{key-off}{1023}
\showcasematerialdesignfilled{king-bed}{1024}
\showcasematerialdesignfilled{kitchen}{1025}
\showcasematerialdesignfilled{kitesurfing}{1026}
\showcasematerialdesignfilled{label}{1027}
\showcasematerialdesignfilled{label-important}{1028}
\showcasematerialdesignfilled{label-off}{1029}
\showcasematerialdesignfilled{lan}{1030}
\showcasematerialdesignfilled{landscape}{1031}
\showcasematerialdesignfilled{landslide}{1032}
\showcasematerialdesignfilled{language}{1033}
\showcasematerialdesignfilled{laptop}{1034}
\showcasematerialdesignfilled{laptop-chromebook}{1035}
\showcasematerialdesignfilled{laptop-mac}{1036}
\showcasematerialdesignfilled{laptop-windows}{1037}
\showcasematerialdesignfilled{last-page}{1038}
\showcasematerialdesignfilled{launch}{1039}
\showcasematerialdesignfilled{layers}{1040}
\showcasematerialdesignfilled{layers-clear}{1041}
\showcasematerialdesignfilled{leaderboard}{1042}
\showcasematerialdesignfilled{leak-add}{1043}
\showcasematerialdesignfilled{leak-remove}{1044}
\showcasematerialdesignfilled{legend-toggle}{1045}
\showcasematerialdesignfilled{lens}{1046}
\showcasematerialdesignfilled{lens-blur}{1047}
\showcasematerialdesignfilled{library-add}{1048}
\showcasematerialdesignfilled{library-add-check}{1049}
\showcasematerialdesignfilled{library-books}{1050}
\showcasematerialdesignfilled{library-music}{1051}
\showcasematerialdesignfilled{light}{1052}
\showcasematerialdesignfilled{lightbulb}{1053}
\showcasematerialdesignfilled{lightbulb-circle}{1054}
\showcasematerialdesignfilled{light-mode}{1055}
\showcasematerialdesignfilled{linear-scale}{1056}
\showcasematerialdesignfilled{line-axis}{1057}
\showcasematerialdesignfilled{line-style}{1058}
\showcasematerialdesignfilled{line-weight}{1059}
\showcasematerialdesignfilled{link}{1060}
\showcasematerialdesignfilled{linked-camera}{1061}
\showcasematerialdesignfilled{link-off}{1062}
\showcasematerialdesignfilled{liquor}{1063}
\showcasematerialdesignfilled{list}{1064}
\showcasematerialdesignfilled{list-alt}{1065}
\showcasematerialdesignfilled{live-help}{1066}
\showcasematerialdesignfilled{live-tv}{1067}
\showcasematerialdesignfilled{living}{1068}
\showcasematerialdesignfilled{local-activity}{1069}
\showcasematerialdesignfilled{local-airport}{1070}
\showcasematerialdesignfilled{local-atm}{1071}
\showcasematerialdesignfilled{local-bar}{1072}
\showcasematerialdesignfilled{local-cafe}{1073}
\showcasematerialdesignfilled{local-car-wash}{1074}
\showcasematerialdesignfilled{local-convenience-store}{1075}
\showcasematerialdesignfilled{local-dining}{1076}
\showcasematerialdesignfilled{local-drink}{1077}
\showcasematerialdesignfilled{local-fire-department}{1078}
\showcasematerialdesignfilled{local-florist}{1079}
\showcasematerialdesignfilled{local-gas-station}{1080}
\showcasematerialdesignfilled{local-grocery-store}{1081}
\showcasematerialdesignfilled{local-hospital}{1082}
\showcasematerialdesignfilled{local-hotel}{1083}
\showcasematerialdesignfilled{local-laundry-service}{1084}
\showcasematerialdesignfilled{local-library}{1085}
\showcasematerialdesignfilled{local-mall}{1086}
\showcasematerialdesignfilled{local-movies}{1087}
\showcasematerialdesignfilled{local-offer}{1088}
\showcasematerialdesignfilled{local-parking}{1089}
\showcasematerialdesignfilled{local-pharmacy}{1090}
\showcasematerialdesignfilled{local-phone}{1091}
\showcasematerialdesignfilled{local-pizza}{1092}
\showcasematerialdesignfilled{local-play}{1093}
\showcasematerialdesignfilled{local-police}{1094}
\showcasematerialdesignfilled{local-post-office}{1095}
\showcasematerialdesignfilled{local-printshop}{1096}
\showcasematerialdesignfilled{local-see}{1097}
\showcasematerialdesignfilled{local-shipping}{1098}
\showcasematerialdesignfilled{local-taxi}{1099}
\showcasematerialdesignfilled{location-city}{1100}
\showcasematerialdesignfilled{location-disabled}{1101}
\showcasematerialdesignfilled{location-off}{1102}
\showcasematerialdesignfilled{location-on}{1103}
\showcasematerialdesignfilled{location-searching}{1104}
\showcasematerialdesignfilled{lock}{1105}
\showcasematerialdesignfilled{lock-clock}{1106}
\showcasematerialdesignfilled{lock-open}{1107}
\showcasematerialdesignfilled{lock-person}{1108}
\showcasematerialdesignfilled{lock-reset}{1109}
\showcasematerialdesignfilled{login}{1110}
\showcasematerialdesignfilled{logout}{1111}
\showcasematerialdesignfilled{logo-dev}{1112}
\showcasematerialdesignfilled{looks}{1113}
\showcasematerialdesignfilled{looks-3}{1114}
\showcasematerialdesignfilled{looks-4}{1115}
\showcasematerialdesignfilled{looks-5}{1116}
\showcasematerialdesignfilled{looks-6}{1117}
\showcasematerialdesignfilled{looks-one}{1118}
\showcasematerialdesignfilled{looks-two}{1119}
\showcasematerialdesignfilled{loop}{1120}
\showcasematerialdesignfilled{loupe}{1121}
\showcasematerialdesignfilled{low-priority}{1122}
\showcasematerialdesignfilled{loyalty}{1123}
\showcasematerialdesignfilled{lte-mobiledata}{1124}
\showcasematerialdesignfilled{lte-plus-mobiledata}{1125}
\showcasematerialdesignfilled{luggage}{1126}
\showcasematerialdesignfilled{lunch-dining}{1127}
\showcasematerialdesignfilled{lyrics}{1128}
\showcasematerialdesignfilled{macro-off}{1129}
\showcasematerialdesignfilled{mail}{1130}
\showcasematerialdesignfilled{mail-lock}{1131}
\showcasematerialdesignfilled{mail-outline}{1132}
\showcasematerialdesignfilled{male}{1133}
\showcasematerialdesignfilled{man}{1134}
\showcasematerialdesignfilled{manage-accounts}{1135}
\showcasematerialdesignfilled{manage-history}{1136}
\showcasematerialdesignfilled{manage-search}{1137}
\showcasematerialdesignfilled{man-2}{1138}
\showcasematerialdesignfilled{man-3}{1139}
\showcasematerialdesignfilled{man-4}{1140}
\showcasematerialdesignfilled{map}{1141}
\showcasematerialdesignfilled{maps-home-work}{1142}
\showcasematerialdesignfilled{maps-ugc}{1143}
\showcasematerialdesignfilled{margin}{1144}
\showcasematerialdesignfilled{markunread}{1145}
\showcasematerialdesignfilled{markunread-mailbox}{1146}
\showcasematerialdesignfilled{mark-as-unread}{1147}
\showcasematerialdesignfilled{mark-chat-read}{1148}
\showcasematerialdesignfilled{mark-chat-unread}{1149}
\showcasematerialdesignfilled{mark-email-read}{1150}
\showcasematerialdesignfilled{mark-email-unread}{1151}
\showcasematerialdesignfilled{mark-unread-chat-alt}{1152}
\showcasematerialdesignfilled{masks}{1153}
\showcasematerialdesignfilled{maximize}{1154}
\showcasematerialdesignfilled{mediation}{1155}
\showcasematerialdesignfilled{media-bluetooth-off}{1156}
\showcasematerialdesignfilled{media-bluetooth-on}{1157}
\showcasematerialdesignfilled{medical-information}{1158}
\showcasematerialdesignfilled{medical-services}{1159}
\showcasematerialdesignfilled{medication}{1160}
\showcasematerialdesignfilled{medication-liquid}{1161}
\showcasematerialdesignfilled{meeting-room}{1162}
\showcasematerialdesignfilled{memory}{1163}
\showcasematerialdesignfilled{menu}{1164}
\showcasematerialdesignfilled{menu-book}{1165}
\showcasematerialdesignfilled{menu-open}{1166}
\showcasematerialdesignfilled{merge}{1167}
\showcasematerialdesignfilled{merge-type}{1168}
\showcasematerialdesignfilled{message}{1169}
\showcasematerialdesignfilled{mic}{1170}
\showcasematerialdesignfilled{microwave}{1171}
\showcasematerialdesignfilled{mic-external-off}{1172}
\showcasematerialdesignfilled{mic-external-on}{1173}
\showcasematerialdesignfilled{mic-none}{1174}
\showcasematerialdesignfilled{mic-off}{1175}
\showcasematerialdesignfilled{military-tech}{1176}
\showcasematerialdesignfilled{minimize}{1177}
\showcasematerialdesignfilled{minor-crash}{1178}
\showcasematerialdesignfilled{miscellaneous-services}{1179}
\showcasematerialdesignfilled{missed-video-call}{1180}
\showcasematerialdesignfilled{mms}{1181}
\showcasematerialdesignfilled{mobiledata-off}{1182}
\showcasematerialdesignfilled{mobile-friendly}{1183}
\showcasematerialdesignfilled{mobile-off}{1184}
\showcasematerialdesignfilled{mobile-screen-share}{1185}
\showcasematerialdesignfilled{mode}{1186}
\showcasematerialdesignfilled{model-training}{1187}
\showcasematerialdesignfilled{mode-comment}{1188}
\showcasematerialdesignfilled{mode-edit}{1189}
\showcasematerialdesignfilled{mode-edit-outline}{1190}
\showcasematerialdesignfilled{mode-fan-off}{1191}
\showcasematerialdesignfilled{mode-night}{1192}
\showcasematerialdesignfilled{mode-of-travel}{1193}
\showcasematerialdesignfilled{mode-standby}{1194}
\showcasematerialdesignfilled{monetization-on}{1195}
\showcasematerialdesignfilled{money}{1196}
\showcasematerialdesignfilled{money-off}{1197}
\showcasematerialdesignfilled{money-off-csred}{1198}
\showcasematerialdesignfilled{monitor}{1199}
\showcasematerialdesignfilled{monitor-heart}{1200}
\showcasematerialdesignfilled{monitor-weight}{1201}
\showcasematerialdesignfilled{monochrome-photos}{1202}
\showcasematerialdesignfilled{mood}{1203}
\showcasematerialdesignfilled{mood-bad}{1204}
\showcasematerialdesignfilled{moped}{1205}
\showcasematerialdesignfilled{more}{1206}
\showcasematerialdesignfilled{more-horiz}{1207}
\showcasematerialdesignfilled{more-time}{1208}
\showcasematerialdesignfilled{more-vert}{1209}
\showcasematerialdesignfilled{mosque}{1210}
\showcasematerialdesignfilled{motion-photos-auto}{1211}
\showcasematerialdesignfilled{motion-photos-off}{1212}
\showcasematerialdesignfilled{motion-photos-on}{1213}
\showcasematerialdesignfilled{motion-photos-pause}{1214}
\showcasematerialdesignfilled{motion-photos-paused}{1215}
\showcasematerialdesignfilled{mouse}{1216}
\showcasematerialdesignfilled{move-down}{1217}
\showcasematerialdesignfilled{move-to-inbox}{1218}
\showcasematerialdesignfilled{move-up}{1219}
\showcasematerialdesignfilled{movie}{1220}
\showcasematerialdesignfilled{movie-creation}{1221}
\showcasematerialdesignfilled{movie-filter}{1222}
\showcasematerialdesignfilled{moving}{1223}
\showcasematerialdesignfilled{mp}{1224}
\showcasematerialdesignfilled{multiline-chart}{1225}
\showcasematerialdesignfilled{multiple-stop}{1226}
\showcasematerialdesignfilled{museum}{1227}
\showcasematerialdesignfilled{music-note}{1228}
\showcasematerialdesignfilled{music-off}{1229}
\showcasematerialdesignfilled{music-video}{1230}
\showcasematerialdesignfilled{my-location}{1231}
\showcasematerialdesignfilled{nat}{1232}
\showcasematerialdesignfilled{nature}{1233}
\showcasematerialdesignfilled{nature-people}{1234}
\showcasematerialdesignfilled{navigate-before}{1235}
\showcasematerialdesignfilled{navigate-next}{1236}
\showcasematerialdesignfilled{navigation}{1237}
\showcasematerialdesignfilled{nearby-error}{1238}
\showcasematerialdesignfilled{nearby-off}{1239}
\showcasematerialdesignfilled{near-me}{1240}
\showcasematerialdesignfilled{near-me-disabled}{1241}
\showcasematerialdesignfilled{nest-cam-wired-stand}{1242}
\showcasematerialdesignfilled{network-cell}{1243}
\showcasematerialdesignfilled{network-check}{1244}
\showcasematerialdesignfilled{network-locked}{1245}
\showcasematerialdesignfilled{network-ping}{1246}
\showcasematerialdesignfilled{network-wifi}{1247}
\showcasematerialdesignfilled{network-wifi-1-bar}{1248}
\showcasematerialdesignfilled{network-wifi-2-bar}{1249}
\showcasematerialdesignfilled{network-wifi-3-bar}{1250}
\showcasematerialdesignfilled{newspaper}{1251}
\showcasematerialdesignfilled{new-label}{1252}
\showcasematerialdesignfilled{new-releases}{1253}
\showcasematerialdesignfilled{next-plan}{1254}
\showcasematerialdesignfilled{next-week}{1255}
\showcasematerialdesignfilled{nfc}{1256}
\showcasematerialdesignfilled{nightlife}{1257}
\showcasematerialdesignfilled{nightlight}{1258}
\showcasematerialdesignfilled{nightlight-round}{1259}
\showcasematerialdesignfilled{nights-stay}{1260}
\showcasematerialdesignfilled{night-shelter}{1261}
\showcasematerialdesignfilled{noise-aware}{1262}
\showcasematerialdesignfilled{noise-control-off}{1263}
\showcasematerialdesignfilled{nordic-walking}{1264}
\showcasematerialdesignfilled{north}{1265}
\showcasematerialdesignfilled{north-east}{1266}
\showcasematerialdesignfilled{north-west}{1267}
\showcasematerialdesignfilled{note}{1268}
\showcasematerialdesignfilled{notes}{1269}
\showcasematerialdesignfilled{note-add}{1270}
\showcasematerialdesignfilled{note-alt}{1271}
\showcasematerialdesignfilled{notifications}{1272}
\showcasematerialdesignfilled{notifications-active}{1273}
\showcasematerialdesignfilled{notifications-none}{1274}
\showcasematerialdesignfilled{notifications-off}{1275}
\showcasematerialdesignfilled{notifications-paused}{1276}
\showcasematerialdesignfilled{notification-add}{1277}
\showcasematerialdesignfilled{notification-important}{1278}
\showcasematerialdesignfilled{not-accessible}{1279}
\showcasematerialdesignfilled{not-interested}{1280}
\showcasematerialdesignfilled{not-listed-location}{1281}
\showcasematerialdesignfilled{not-started}{1282}
\showcasematerialdesignfilled{no-accounts}{1283}
\showcasematerialdesignfilled{no-adult-content}{1284}
\showcasematerialdesignfilled{no-backpack}{1285}
\showcasematerialdesignfilled{no-cell}{1286}
\showcasematerialdesignfilled{no-crash}{1287}
\showcasematerialdesignfilled{no-drinks}{1288}
\showcasematerialdesignfilled{no-encryption}{1289}
\showcasematerialdesignfilled{no-encryption-gmailerrorred}{1290}
\showcasematerialdesignfilled{no-flash}{1291}
\showcasematerialdesignfilled{no-food}{1292}
\showcasematerialdesignfilled{no-luggage}{1293}
\showcasematerialdesignfilled{no-meals}{1294}
\showcasematerialdesignfilled{no-meeting-room}{1295}
\showcasematerialdesignfilled{no-photography}{1296}
\showcasematerialdesignfilled{no-sim}{1297}
\showcasematerialdesignfilled{no-stroller}{1298}
\showcasematerialdesignfilled{no-transfer}{1299}
\showcasematerialdesignfilled{numbers}{1300}
\showcasematerialdesignfilled{offline-bolt}{1301}
\showcasematerialdesignfilled{offline-pin}{1302}
\showcasematerialdesignfilled{offline-share}{1303}
\showcasematerialdesignfilled{oil-barrel}{1304}
\showcasematerialdesignfilled{ondemand-video}{1305}
\showcasematerialdesignfilled{online-prediction}{1306}
\showcasematerialdesignfilled{on-device-training}{1307}
\showcasematerialdesignfilled{opacity}{1308}
\showcasematerialdesignfilled{open-in-browser}{1309}
\showcasematerialdesignfilled{open-in-full}{1310}
\showcasematerialdesignfilled{open-in-new}{1311}
\showcasematerialdesignfilled{open-in-new-off}{1312}
\showcasematerialdesignfilled{open-with}{1313}
\showcasematerialdesignfilled{other-houses}{1314}
\showcasematerialdesignfilled{outbound}{1315}
\showcasematerialdesignfilled{outbox}{1316}
\showcasematerialdesignfilled{outdoor-grill}{1317}
\showcasematerialdesignfilled{outlet}{1318}
\showcasematerialdesignfilled{outlined-flag}{1319}
\showcasematerialdesignfilled{output}{1320}
\showcasematerialdesignfilled{padding}{1321}
\showcasematerialdesignfilled{pages}{1322}
\showcasematerialdesignfilled{pageview}{1323}
\showcasematerialdesignfilled{paid}{1324}
\showcasematerialdesignfilled{palette}{1325}
\showcasematerialdesignfilled{panorama}{1326}
\showcasematerialdesignfilled{panorama-fish-eye}{1327}
\showcasematerialdesignfilled{panorama-horizontal}{1328}
\showcasematerialdesignfilled{panorama-horizontal-select}{1329}
\showcasematerialdesignfilled{panorama-photosphere}{1330}
\showcasematerialdesignfilled{panorama-photosphere-select}{1331}
\showcasematerialdesignfilled{panorama-vertical}{1332}
\showcasematerialdesignfilled{panorama-vertical-select}{1333}
\showcasematerialdesignfilled{panorama-wide-angle}{1334}
\showcasematerialdesignfilled{panorama-wide-angle-select}{1335}
\showcasematerialdesignfilled{pan-tool}{1336}
\showcasematerialdesignfilled{pan-tool-alt}{1337}
\showcasematerialdesignfilled{paragliding}{1338}
\showcasematerialdesignfilled{park}{1339}
\showcasematerialdesignfilled{party-mode}{1340}
\showcasematerialdesignfilled{password}{1341}
\showcasematerialdesignfilled{pattern}{1342}
\showcasematerialdesignfilled{pause}{1343}
\showcasematerialdesignfilled{pause-circle}{1344}
\showcasematerialdesignfilled{pause-circle-filled}{1345}
\showcasematerialdesignfilled{pause-circle-outline}{1346}
\showcasematerialdesignfilled{pause-presentation}{1347}
\showcasematerialdesignfilled{payment}{1348}
\showcasematerialdesignfilled{payments}{1349}
\showcasematerialdesignfilled{pedal-bike}{1350}
\showcasematerialdesignfilled{pending}{1351}
\showcasematerialdesignfilled{pending-actions}{1352}
\showcasematerialdesignfilled{pentagon}{1353}
\showcasematerialdesignfilled{people}{1354}
\showcasematerialdesignfilled{people-alt}{1355}
\showcasematerialdesignfilled{people-outline}{1356}
\showcasematerialdesignfilled{percent}{1357}
\showcasematerialdesignfilled{perm-camera-mic}{1358}
\showcasematerialdesignfilled{perm-contact-calendar}{1359}
\showcasematerialdesignfilled{perm-data-setting}{1360}
\showcasematerialdesignfilled{perm-device-information}{1361}
\showcasematerialdesignfilled{perm-identity}{1362}
\showcasematerialdesignfilled{perm-media}{1363}
\showcasematerialdesignfilled{perm-phone-msg}{1364}
\showcasematerialdesignfilled{perm-scan-wifi}{1365}
\showcasematerialdesignfilled{person}{1366}
\showcasematerialdesignfilled{personal-injury}{1367}
\showcasematerialdesignfilled{personal-video}{1368}
\showcasematerialdesignfilled{person-2}{1369}
\showcasematerialdesignfilled{person-3}{1370}
\showcasematerialdesignfilled{person-4}{1371}
\showcasematerialdesignfilled{person-add}{1372}
\showcasematerialdesignfilled{person-add-alt}{1373}
\showcasematerialdesignfilled{person-add-alt-1}{1374}
\showcasematerialdesignfilled{person-add-disabled}{1375}
\showcasematerialdesignfilled{person-off}{1376}
\showcasematerialdesignfilled{person-outline}{1377}
\showcasematerialdesignfilled{person-pin}{1378}
\showcasematerialdesignfilled{person-pin-circle}{1379}
\showcasematerialdesignfilled{person-remove}{1380}
\showcasematerialdesignfilled{person-remove-alt-1}{1381}
\showcasematerialdesignfilled{person-search}{1382}
\showcasematerialdesignfilled{pest-control}{1383}
\showcasematerialdesignfilled{pest-control-rodent}{1384}
\showcasematerialdesignfilled{pets}{1385}
\showcasematerialdesignfilled{phishing}{1386}
\showcasematerialdesignfilled{phone}{1387}
\showcasematerialdesignfilled{phonelink}{1388}
\showcasematerialdesignfilled{phonelink-erase}{1389}
\showcasematerialdesignfilled{phonelink-lock}{1390}
\showcasematerialdesignfilled{phonelink-off}{1391}
\showcasematerialdesignfilled{phonelink-ring}{1392}
\showcasematerialdesignfilled{phonelink-setup}{1393}
\showcasematerialdesignfilled{phone-android}{1394}
\showcasematerialdesignfilled{phone-bluetooth-speaker}{1395}
\showcasematerialdesignfilled{phone-callback}{1396}
\showcasematerialdesignfilled{phone-disabled}{1397}
\showcasematerialdesignfilled{phone-enabled}{1398}
\showcasematerialdesignfilled{phone-forwarded}{1399}
\showcasematerialdesignfilled{phone-iphone}{1400}
\showcasematerialdesignfilled{phone-locked}{1401}
\showcasematerialdesignfilled{phone-missed}{1402}
\showcasematerialdesignfilled{phone-paused}{1403}
\showcasematerialdesignfilled{photo}{1404}
\showcasematerialdesignfilled{photo-album}{1405}
\showcasematerialdesignfilled{photo-camera}{1406}
\showcasematerialdesignfilled{photo-camera-back}{1407}
\showcasematerialdesignfilled{photo-camera-front}{1408}
\showcasematerialdesignfilled{photo-filter}{1409}
\showcasematerialdesignfilled{photo-library}{1410}
\showcasematerialdesignfilled{photo-size-select-actual}{1411}
\showcasematerialdesignfilled{photo-size-select-large}{1412}
\showcasematerialdesignfilled{photo-size-select-small}{1413}
\showcasematerialdesignfilled{php}{1414}
\showcasematerialdesignfilled{piano}{1415}
\showcasematerialdesignfilled{piano-off}{1416}
\showcasematerialdesignfilled{picture-as-pdf}{1417}
\showcasematerialdesignfilled{picture-in-picture}{1418}
\showcasematerialdesignfilled{picture-in-picture-alt}{1419}
\showcasematerialdesignfilled{pie-chart}{1420}
\showcasematerialdesignfilled{pie-chart-outline}{1421}
\showcasematerialdesignfilled{pin}{1422}
\showcasematerialdesignfilled{pinch}{1423}
\showcasematerialdesignfilled{pin-drop}{1424}
\showcasematerialdesignfilled{pin-end}{1425}
\showcasematerialdesignfilled{pin-invoke}{1426}
\showcasematerialdesignfilled{pivot-table-chart}{1427}
\showcasematerialdesignfilled{pix}{1428}
\showcasematerialdesignfilled{place}{1429}
\showcasematerialdesignfilled{plagiarism}{1430}
\showcasematerialdesignfilled{playlist-add}{1431}
\showcasematerialdesignfilled{playlist-add-check}{1432}
\showcasematerialdesignfilled{playlist-add-check-circle}{1433}
\showcasematerialdesignfilled{playlist-add-circle}{1434}
\showcasematerialdesignfilled{playlist-play}{1435}
\showcasematerialdesignfilled{playlist-remove}{1436}
\showcasematerialdesignfilled{play-arrow}{1437}
\showcasematerialdesignfilled{play-circle}{1438}
\showcasematerialdesignfilled{play-circle-filled}{1439}
\showcasematerialdesignfilled{play-circle-outline}{1440}
\showcasematerialdesignfilled{play-disabled}{1441}
\showcasematerialdesignfilled{play-for-work}{1442}
\showcasematerialdesignfilled{play-lesson}{1443}
\showcasematerialdesignfilled{plumbing}{1444}
\showcasematerialdesignfilled{plus-one}{1445}
\showcasematerialdesignfilled{podcasts}{1446}
\showcasematerialdesignfilled{point-of-sale}{1447}
\showcasematerialdesignfilled{policy}{1448}
\showcasematerialdesignfilled{poll}{1449}
\showcasematerialdesignfilled{polyline}{1450}
\showcasematerialdesignfilled{polymer}{1451}
\showcasematerialdesignfilled{pool}{1452}
\showcasematerialdesignfilled{portable-wifi-off}{1453}
\showcasematerialdesignfilled{portrait}{1454}
\showcasematerialdesignfilled{post-add}{1455}
\showcasematerialdesignfilled{power}{1456}
\showcasematerialdesignfilled{power-input}{1457}
\showcasematerialdesignfilled{power-off}{1458}
\showcasematerialdesignfilled{power-settings-new}{1459}
\showcasematerialdesignfilled{precision-manufacturing}{1460}
\showcasematerialdesignfilled{pregnant-woman}{1461}
\showcasematerialdesignfilled{present-to-all}{1462}
\showcasematerialdesignfilled{preview}{1463}
\showcasematerialdesignfilled{price-change}{1464}
\showcasematerialdesignfilled{price-check}{1465}
\showcasematerialdesignfilled{print}{1466}
\showcasematerialdesignfilled{print-disabled}{1467}
\showcasematerialdesignfilled{priority-high}{1468}
\showcasematerialdesignfilled{privacy-tip}{1469}
\showcasematerialdesignfilled{private-connectivity}{1470}
\showcasematerialdesignfilled{production-quantity-limits}{1471}
\showcasematerialdesignfilled{propane}{1472}
\showcasematerialdesignfilled{propane-tank}{1473}
\showcasematerialdesignfilled{psychology}{1474}
\showcasematerialdesignfilled{psychology-alt}{1475}
\showcasematerialdesignfilled{public}{1476}
\showcasematerialdesignfilled{public-off}{1477}
\showcasematerialdesignfilled{publish}{1478}
\showcasematerialdesignfilled{published-with-changes}{1479}
\showcasematerialdesignfilled{punch-clock}{1480}
\showcasematerialdesignfilled{push-pin}{1481}
\showcasematerialdesignfilled{qr-code}{1482}
\showcasematerialdesignfilled{qr-code-2}{1483}
\showcasematerialdesignfilled{qr-code-scanner}{1484}
\showcasematerialdesignfilled{query-builder}{1485}
\showcasematerialdesignfilled{query-stats}{1486}
\showcasematerialdesignfilled{question-answer}{1487}
\showcasematerialdesignfilled{question-mark}{1488}
\showcasematerialdesignfilled{queue}{1489}
\showcasematerialdesignfilled{queue-music}{1490}
\showcasematerialdesignfilled{queue-play-next}{1491}
\showcasematerialdesignfilled{quickreply}{1492}
\showcasematerialdesignfilled{quiz}{1493}
\showcasematerialdesignfilled{radar}{1494}
\showcasematerialdesignfilled{radio}{1495}
\showcasematerialdesignfilled{radio-button-checked}{1496}
\showcasematerialdesignfilled{radio-button-unchecked}{1497}
\showcasematerialdesignfilled{railway-alert}{1498}
\showcasematerialdesignfilled{ramen-dining}{1499}
\showcasematerialdesignfilled{ramp-left}{1500}
\showcasematerialdesignfilled{ramp-right}{1501}
\showcasematerialdesignfilled{rate-review}{1502}
\showcasematerialdesignfilled{raw-off}{1503}
\showcasematerialdesignfilled{raw-on}{1504}
\showcasematerialdesignfilled{read-more}{1505}
\showcasematerialdesignfilled{real-estate-agent}{1506}
\showcasematerialdesignfilled{receipt}{1507}
\showcasematerialdesignfilled{receipt-long}{1508}
\showcasematerialdesignfilled{recent-actors}{1509}
\showcasematerialdesignfilled{recommend}{1510}
\showcasematerialdesignfilled{record-voice-over}{1511}
\showcasematerialdesignfilled{rectangle}{1512}
\showcasematerialdesignfilled{recycling}{1513}
\showcasematerialdesignfilled{redeem}{1514}
\showcasematerialdesignfilled{redo}{1515}
\showcasematerialdesignfilled{reduce-capacity}{1516}
\showcasematerialdesignfilled{refresh}{1517}
\showcasematerialdesignfilled{remember-me}{1518}
\showcasematerialdesignfilled{remove}{1519}
\showcasematerialdesignfilled{remove-circle}{1520}
\showcasematerialdesignfilled{remove-circle-outline}{1521}
\showcasematerialdesignfilled{remove-done}{1522}
\showcasematerialdesignfilled{remove-from-queue}{1523}
\showcasematerialdesignfilled{remove-moderator}{1524}
\showcasematerialdesignfilled{remove-red-eye}{1525}
\showcasematerialdesignfilled{remove-road}{1526}
\showcasematerialdesignfilled{remove-shopping-cart}{1527}
\showcasematerialdesignfilled{reorder}{1528}
\showcasematerialdesignfilled{repartition}{1529}
\showcasematerialdesignfilled{repeat}{1530}
\showcasematerialdesignfilled{repeat-on}{1531}
\showcasematerialdesignfilled{repeat-one}{1532}
\showcasematerialdesignfilled{repeat-one-on}{1533}
\showcasematerialdesignfilled{replay}{1534}
\showcasematerialdesignfilled{replay-5}{1535}
\showcasematerialdesignfilled{replay-10}{1536}
\showcasematerialdesignfilled{replay-30}{1537}
\showcasematerialdesignfilled{replay-circle-filled}{1538}
\showcasematerialdesignfilled{reply}{1539}
\showcasematerialdesignfilled{reply-all}{1540}
\showcasematerialdesignfilled{report}{1541}
\showcasematerialdesignfilled{report-gmailerrorred}{1542}
\showcasematerialdesignfilled{report-off}{1543}
\showcasematerialdesignfilled{report-problem}{1544}
\showcasematerialdesignfilled{request-page}{1545}
\showcasematerialdesignfilled{request-quote}{1546}
\showcasematerialdesignfilled{reset-tv}{1547}
\showcasematerialdesignfilled{restart-alt}{1548}
\showcasematerialdesignfilled{restaurant}{1549}
\showcasematerialdesignfilled{restaurant-menu}{1550}
\showcasematerialdesignfilled{restore}{1551}
\showcasematerialdesignfilled{restore-from-trash}{1552}
\showcasematerialdesignfilled{restore-page}{1553}
\showcasematerialdesignfilled{reviews}{1554}
\showcasematerialdesignfilled{rice-bowl}{1555}
\showcasematerialdesignfilled{ring-volume}{1556}
\showcasematerialdesignfilled{rocket}{1557}
\showcasematerialdesignfilled{rocket-launch}{1558}
\showcasematerialdesignfilled{roller-shades}{1559}
\showcasematerialdesignfilled{roller-shades-closed}{1560}
\showcasematerialdesignfilled{roller-skating}{1561}
\showcasematerialdesignfilled{roofing}{1562}
\showcasematerialdesignfilled{room}{1563}
\showcasematerialdesignfilled{room-preferences}{1564}
\showcasematerialdesignfilled{room-service}{1565}
\showcasematerialdesignfilled{rotate-90-degrees-ccw}{1566}
\showcasematerialdesignfilled{rotate-90-degrees-cw}{1567}
\showcasematerialdesignfilled{rotate-left}{1568}
\showcasematerialdesignfilled{rotate-right}{1569}
\showcasematerialdesignfilled{roundabout-left}{1570}
\showcasematerialdesignfilled{roundabout-right}{1571}
\showcasematerialdesignfilled{rounded-corner}{1572}
\showcasematerialdesignfilled{route}{1573}
\showcasematerialdesignfilled{router}{1574}
\showcasematerialdesignfilled{rowing}{1575}
\showcasematerialdesignfilled{rss-feed}{1576}
\showcasematerialdesignfilled{rsvp}{1577}
\showcasematerialdesignfilled{rtt}{1578}
\showcasematerialdesignfilled{rule}{1579}
\showcasematerialdesignfilled{rule-folder}{1580}
\showcasematerialdesignfilled{running-with-errors}{1581}
\showcasematerialdesignfilled{run-circle}{1582}
\showcasematerialdesignfilled{rv-hookup}{1583}
\showcasematerialdesignfilled{r-mobiledata}{1584}
\showcasematerialdesignfilled{safety-check}{1585}
\showcasematerialdesignfilled{safety-divider}{1586}
\showcasematerialdesignfilled{sailing}{1587}
\showcasematerialdesignfilled{sanitizer}{1588}
\showcasematerialdesignfilled{satellite}{1589}
\showcasematerialdesignfilled{satellite-alt}{1590}
\showcasematerialdesignfilled{save}{1591}
\showcasematerialdesignfilled{saved-search}{1592}
\showcasematerialdesignfilled{save-alt}{1593}
\showcasematerialdesignfilled{save-as}{1594}
\showcasematerialdesignfilled{savings}{1595}
\showcasematerialdesignfilled{scale}{1596}
\showcasematerialdesignfilled{scanner}{1597}
\showcasematerialdesignfilled{scatter-plot}{1598}
\showcasematerialdesignfilled{schedule}{1599}
\showcasematerialdesignfilled{schedule-send}{1600}
\showcasematerialdesignfilled{schema}{1601}
\showcasematerialdesignfilled{school}{1602}
\showcasematerialdesignfilled{science}{1603}
\showcasematerialdesignfilled{score}{1604}
\showcasematerialdesignfilled{scoreboard}{1605}
\showcasematerialdesignfilled{screenshot}{1606}
\showcasematerialdesignfilled{screenshot-monitor}{1607}
\showcasematerialdesignfilled{screen-lock-landscape}{1608}
\showcasematerialdesignfilled{screen-lock-portrait}{1609}
\showcasematerialdesignfilled{screen-lock-rotation}{1610}
\showcasematerialdesignfilled{screen-rotation}{1611}
\showcasematerialdesignfilled{screen-rotation-alt}{1612}
\showcasematerialdesignfilled{screen-search-desktop}{1613}
\showcasematerialdesignfilled{screen-share}{1614}
\showcasematerialdesignfilled{scuba-diving}{1615}
\showcasematerialdesignfilled{sd}{1616}
\showcasematerialdesignfilled{sd-card}{1617}
\showcasematerialdesignfilled{sd-card-alert}{1618}
\showcasematerialdesignfilled{sd-storage}{1619}
\showcasematerialdesignfilled{search}{1620}
\showcasematerialdesignfilled{search-off}{1621}
\showcasematerialdesignfilled{security}{1622}
\showcasematerialdesignfilled{security-update}{1623}
\showcasematerialdesignfilled{security-update-good}{1624}
\showcasematerialdesignfilled{security-update-warning}{1625}
\showcasematerialdesignfilled{segment}{1626}
\showcasematerialdesignfilled{select-all}{1627}
\showcasematerialdesignfilled{self-improvement}{1628}
\showcasematerialdesignfilled{sell}{1629}
\showcasematerialdesignfilled{send}{1630}
\showcasematerialdesignfilled{send-and-archive}{1631}
\showcasematerialdesignfilled{send-time-extension}{1632}
\showcasematerialdesignfilled{send-to-mobile}{1633}
\showcasematerialdesignfilled{sensors}{1634}
\showcasematerialdesignfilled{sensors-off}{1635}
\showcasematerialdesignfilled{sensor-door}{1636}
\showcasematerialdesignfilled{sensor-occupied}{1637}
\showcasematerialdesignfilled{sensor-window}{1638}
\showcasematerialdesignfilled{sentiment-dissatisfied}{1639}
\showcasematerialdesignfilled{sentiment-neutral}{1640}
\showcasematerialdesignfilled{sentiment-satisfied}{1641}
\showcasematerialdesignfilled{sentiment-satisfied-alt}{1642}
\showcasematerialdesignfilled{sentiment-very-dissatisfied}{1643}
\showcasematerialdesignfilled{sentiment-very-satisfied}{1644}
\showcasematerialdesignfilled{settings}{1645}
\showcasematerialdesignfilled{settings-accessibility}{1646}
\showcasematerialdesignfilled{settings-applications}{1647}
\showcasematerialdesignfilled{settings-backup-restore}{1648}
\showcasematerialdesignfilled{settings-bluetooth}{1649}
\showcasematerialdesignfilled{settings-brightness}{1650}
\showcasematerialdesignfilled{settings-cell}{1651}
\showcasematerialdesignfilled{settings-ethernet}{1652}
\showcasematerialdesignfilled{settings-input-antenna}{1653}
\showcasematerialdesignfilled{settings-input-component}{1654}
\showcasematerialdesignfilled{settings-input-composite}{1655}
\showcasematerialdesignfilled{settings-input-hdmi}{1656}
\showcasematerialdesignfilled{settings-input-svideo}{1657}
\showcasematerialdesignfilled{settings-overscan}{1658}
\showcasematerialdesignfilled{settings-phone}{1659}
\showcasematerialdesignfilled{settings-power}{1660}
\showcasematerialdesignfilled{settings-remote}{1661}
\showcasematerialdesignfilled{settings-suggest}{1662}
\showcasematerialdesignfilled{settings-system-daydream}{1663}
\showcasematerialdesignfilled{settings-voice}{1664}
\showcasematerialdesignfilled{set-meal}{1665}
\showcasematerialdesignfilled{severe-cold}{1666}
\showcasematerialdesignfilled{shape-line}{1667}
\showcasematerialdesignfilled{share}{1668}
\showcasematerialdesignfilled{share-location}{1669}
\showcasematerialdesignfilled{shield}{1670}
\showcasematerialdesignfilled{shield-moon}{1671}
\showcasematerialdesignfilled{shop}{1672}
\showcasematerialdesignfilled{shopping-bag}{1673}
\showcasematerialdesignfilled{shopping-basket}{1674}
\showcasematerialdesignfilled{shopping-cart}{1675}
\showcasematerialdesignfilled{shopping-cart-checkout}{1676}
\showcasematerialdesignfilled{shop-2}{1677}
\showcasematerialdesignfilled{shop-two}{1678}
\showcasematerialdesignfilled{shortcut}{1679}
\showcasematerialdesignfilled{short-text}{1680}
\showcasematerialdesignfilled{shower}{1681}
\showcasematerialdesignfilled{show-chart}{1682}
\showcasematerialdesignfilled{shuffle}{1683}
\showcasematerialdesignfilled{shuffle-on}{1684}
\showcasematerialdesignfilled{shutter-speed}{1685}
\showcasematerialdesignfilled{sick}{1686}
\showcasematerialdesignfilled{signal-cellular-0-bar}{1687}
\showcasematerialdesignfilled{signal-cellular-4-bar}{1688}
\showcasematerialdesignfilled{signal-cellular-alt}{1689}
\showcasematerialdesignfilled{signal-cellular-alt-1-bar}{1690}
\showcasematerialdesignfilled{signal-cellular-alt-2-bar}{1691}
\showcasematerialdesignfilled{signal-cellular-connected-no-internet-0-bar}{1692}
\showcasematerialdesignfilled{signal-cellular-connected-no-internet-4-bar}{1693}
\showcasematerialdesignfilled{signal-cellular-nodata}{1694}
\showcasematerialdesignfilled{signal-cellular-no-sim}{1695}
\showcasematerialdesignfilled{signal-cellular-null}{1696}
\showcasematerialdesignfilled{signal-cellular-off}{1697}
\showcasematerialdesignfilled{signal-wifi-0-bar}{1698}
\showcasematerialdesignfilled{signal-wifi-4-bar}{1699}
\showcasematerialdesignfilled{signal-wifi-4-bar-lock}{1700}
\showcasematerialdesignfilled{signal-wifi-bad}{1701}
\showcasematerialdesignfilled{signal-wifi-connected-no-internet-4}{1702}
\showcasematerialdesignfilled{signal-wifi-off}{1703}
\showcasematerialdesignfilled{signal-wifi-statusbar-4-bar}{1704}
\showcasematerialdesignfilled{signal-wifi-statusbar-connected-no-internet-4}{1705}
\showcasematerialdesignfilled{signal-wifi-statusbar-null}{1706}
\showcasematerialdesignfilled{signpost}{1707}
\showcasematerialdesignfilled{sign-language}{1708}
\showcasematerialdesignfilled{sim-card}{1709}
\showcasematerialdesignfilled{sim-card-alert}{1710}
\showcasematerialdesignfilled{sim-card-download}{1711}
\showcasematerialdesignfilled{single-bed}{1712}
\showcasematerialdesignfilled{sip}{1713}
\showcasematerialdesignfilled{skateboarding}{1714}
\showcasematerialdesignfilled{skip-next}{1715}
\showcasematerialdesignfilled{skip-previous}{1716}
\showcasematerialdesignfilled{sledding}{1717}
\showcasematerialdesignfilled{slideshow}{1718}
\showcasematerialdesignfilled{slow-motion-video}{1719}
\showcasematerialdesignfilled{smartphone}{1720}
\showcasematerialdesignfilled{smart-button}{1721}
\showcasematerialdesignfilled{smart-display}{1722}
\showcasematerialdesignfilled{smart-screen}{1723}
\showcasematerialdesignfilled{smart-toy}{1724}
\showcasematerialdesignfilled{smoke-free}{1725}
\showcasematerialdesignfilled{smoking-rooms}{1726}
\showcasematerialdesignfilled{sms}{1727}
\showcasematerialdesignfilled{sms-failed}{1728}
\showcasematerialdesignfilled{snippet-folder}{1729}
\showcasematerialdesignfilled{snooze}{1730}
\showcasematerialdesignfilled{snowboarding}{1731}
\showcasematerialdesignfilled{snowmobile}{1732}
\showcasematerialdesignfilled{snowshoeing}{1733}
\showcasematerialdesignfilled{soap}{1734}
\showcasematerialdesignfilled{social-distance}{1735}
\showcasematerialdesignfilled{solar-power}{1736}
\showcasematerialdesignfilled{sort}{1737}
\showcasematerialdesignfilled{sort-by-alpha}{1738}
\showcasematerialdesignfilled{sos}{1739}
\showcasematerialdesignfilled{soup-kitchen}{1740}
\showcasematerialdesignfilled{source}{1741}
\showcasematerialdesignfilled{south}{1742}
\showcasematerialdesignfilled{south-america}{1743}
\showcasematerialdesignfilled{south-east}{1744}
\showcasematerialdesignfilled{south-west}{1745}
\showcasematerialdesignfilled{spa}{1746}
\showcasematerialdesignfilled{space-bar}{1747}
\showcasematerialdesignfilled{space-dashboard}{1748}
\showcasematerialdesignfilled{spatial-audio}{1749}
\showcasematerialdesignfilled{spatial-audio-off}{1750}
\showcasematerialdesignfilled{spatial-tracking}{1751}
\showcasematerialdesignfilled{speaker}{1752}
\showcasematerialdesignfilled{speaker-group}{1753}
\showcasematerialdesignfilled{speaker-notes}{1754}
\showcasematerialdesignfilled{speaker-notes-off}{1755}
\showcasematerialdesignfilled{speaker-phone}{1756}
\showcasematerialdesignfilled{speed}{1757}
\showcasematerialdesignfilled{spellcheck}{1758}
\showcasematerialdesignfilled{splitscreen}{1759}
\showcasematerialdesignfilled{spoke}{1760}
\showcasematerialdesignfilled{sports}{1761}
\showcasematerialdesignfilled{sports-bar}{1762}
\showcasematerialdesignfilled{sports-baseball}{1763}
\showcasematerialdesignfilled{sports-basketball}{1764}
\showcasematerialdesignfilled{sports-cricket}{1765}
\showcasematerialdesignfilled{sports-esports}{1766}
\showcasematerialdesignfilled{sports-football}{1767}
\showcasematerialdesignfilled{sports-golf}{1768}
\showcasematerialdesignfilled{sports-gymnastics}{1769}
\showcasematerialdesignfilled{sports-handball}{1770}
\showcasematerialdesignfilled{sports-hockey}{1771}
\showcasematerialdesignfilled{sports-kabaddi}{1772}
\showcasematerialdesignfilled{sports-martial-arts}{1773}
\showcasematerialdesignfilled{sports-mma}{1774}
\showcasematerialdesignfilled{sports-motorsports}{1775}
\showcasematerialdesignfilled{sports-rugby}{1776}
\showcasematerialdesignfilled{sports-score}{1777}
\showcasematerialdesignfilled{sports-soccer}{1778}
\showcasematerialdesignfilled{sports-tennis}{1779}
\showcasematerialdesignfilled{sports-volleyball}{1780}
\showcasematerialdesignfilled{square}{1781}
\showcasematerialdesignfilled{square-foot}{1782}
\showcasematerialdesignfilled{ssid-chart}{1783}
\showcasematerialdesignfilled{stacked-bar-chart}{1784}
\showcasematerialdesignfilled{stacked-line-chart}{1785}
\showcasematerialdesignfilled{stadium}{1786}
\showcasematerialdesignfilled{stairs}{1787}
\showcasematerialdesignfilled{star}{1788}
\showcasematerialdesignfilled{stars}{1789}
\showcasematerialdesignfilled{start}{1790}
\showcasematerialdesignfilled{star-border}{1791}
\showcasematerialdesignfilled{star-border-purple500}{1792}
\showcasematerialdesignfilled{star-half}{1793}
\showcasematerialdesignfilled{star-outline}{1794}
\showcasematerialdesignfilled{star-purple500}{1795}
\showcasematerialdesignfilled{star-rate}{1796}
\showcasematerialdesignfilled{stay-current-landscape}{1797}
\showcasematerialdesignfilled{stay-current-portrait}{1798}
\showcasematerialdesignfilled{stay-primary-landscape}{1799}
\showcasematerialdesignfilled{stay-primary-portrait}{1800}
\showcasematerialdesignfilled{sticky-note-2}{1801}
\showcasematerialdesignfilled{stop}{1802}
\showcasematerialdesignfilled{stop-circle}{1803}
\showcasematerialdesignfilled{stop-screen-share}{1804}
\showcasematerialdesignfilled{storage}{1805}
\showcasematerialdesignfilled{store}{1806}
\showcasematerialdesignfilled{storefront}{1807}
\showcasematerialdesignfilled{store-mall-directory}{1808}
\showcasematerialdesignfilled{storm}{1809}
\showcasematerialdesignfilled{straight}{1810}
\showcasematerialdesignfilled{straighten}{1811}
\showcasematerialdesignfilled{stream}{1812}
\showcasematerialdesignfilled{streetview}{1813}
\showcasematerialdesignfilled{strikethrough-s}{1814}
\showcasematerialdesignfilled{stroller}{1815}
\showcasematerialdesignfilled{style}{1816}
\showcasematerialdesignfilled{subdirectory-arrow-left}{1817}
\showcasematerialdesignfilled{subdirectory-arrow-right}{1818}
\showcasematerialdesignfilled{subject}{1819}
\showcasematerialdesignfilled{subscript}{1820}
\showcasematerialdesignfilled{subscriptions}{1821}
\showcasematerialdesignfilled{subtitles}{1822}
\showcasematerialdesignfilled{subtitles-off}{1823}
\showcasematerialdesignfilled{subway}{1824}
\showcasematerialdesignfilled{summarize}{1825}
\showcasematerialdesignfilled{superscript}{1826}
\showcasematerialdesignfilled{supervised-user-circle}{1827}
\showcasematerialdesignfilled{supervisor-account}{1828}
\showcasematerialdesignfilled{support}{1829}
\showcasematerialdesignfilled{support-agent}{1830}
\showcasematerialdesignfilled{surfing}{1831}
\showcasematerialdesignfilled{surround-sound}{1832}
\showcasematerialdesignfilled{swap-calls}{1833}
\showcasematerialdesignfilled{swap-horiz}{1834}
\showcasematerialdesignfilled{swap-horizontal-circle}{1835}
\showcasematerialdesignfilled{swap-vert}{1836}
\showcasematerialdesignfilled{swap-vertical-circle}{1837}
\showcasematerialdesignfilled{swipe}{1838}
\showcasematerialdesignfilled{swipe-down}{1839}
\showcasematerialdesignfilled{swipe-down-alt}{1840}
\showcasematerialdesignfilled{swipe-left}{1841}
\showcasematerialdesignfilled{swipe-left-alt}{1842}
\showcasematerialdesignfilled{swipe-right}{1843}
\showcasematerialdesignfilled{swipe-right-alt}{1844}
\showcasematerialdesignfilled{swipe-up}{1845}
\showcasematerialdesignfilled{swipe-up-alt}{1846}
\showcasematerialdesignfilled{swipe-vertical}{1847}
\showcasematerialdesignfilled{switch-access-shortcut}{1848}
\showcasematerialdesignfilled{switch-access-shortcut-add}{1849}
\showcasematerialdesignfilled{switch-account}{1850}
\showcasematerialdesignfilled{switch-camera}{1851}
\showcasematerialdesignfilled{switch-left}{1852}
\showcasematerialdesignfilled{switch-right}{1853}
\showcasematerialdesignfilled{switch-video}{1854}
\showcasematerialdesignfilled{synagogue}{1855}
\showcasematerialdesignfilled{sync}{1856}
\showcasematerialdesignfilled{sync-alt}{1857}
\showcasematerialdesignfilled{sync-disabled}{1858}
\showcasematerialdesignfilled{sync-lock}{1859}
\showcasematerialdesignfilled{sync-problem}{1860}
\showcasematerialdesignfilled{system-security-update}{1861}
\showcasematerialdesignfilled{system-security-update-good}{1862}
\showcasematerialdesignfilled{system-security-update-warning}{1863}
\showcasematerialdesignfilled{system-update}{1864}
\showcasematerialdesignfilled{system-update-alt}{1865}
\showcasematerialdesignfilled{tab}{1866}
\showcasematerialdesignfilled{tablet}{1867}
\showcasematerialdesignfilled{tablet-android}{1868}
\showcasematerialdesignfilled{tablet-mac}{1869}
\showcasematerialdesignfilled{table-bar}{1870}
\showcasematerialdesignfilled{table-chart}{1871}
\showcasematerialdesignfilled{table-restaurant}{1872}
\showcasematerialdesignfilled{table-rows}{1873}
\showcasematerialdesignfilled{table-view}{1874}
\showcasematerialdesignfilled{tab-unselected}{1875}
\showcasematerialdesignfilled{tag}{1876}
\showcasematerialdesignfilled{tag-faces}{1877}
\showcasematerialdesignfilled{takeout-dining}{1878}
\showcasematerialdesignfilled{tapas}{1879}
\showcasematerialdesignfilled{tap-and-play}{1880}
\showcasematerialdesignfilled{task}{1881}
\showcasematerialdesignfilled{task-alt}{1882}
\showcasematerialdesignfilled{taxi-alert}{1883}
\showcasematerialdesignfilled{temple-buddhist}{1884}
\showcasematerialdesignfilled{temple-hindu}{1885}
\showcasematerialdesignfilled{terminal}{1886}
\showcasematerialdesignfilled{terrain}{1887}
\showcasematerialdesignfilled{textsms}{1888}
\showcasematerialdesignfilled{texture}{1889}
\showcasematerialdesignfilled{text-decrease}{1890}
\showcasematerialdesignfilled{text-fields}{1891}
\showcasematerialdesignfilled{text-format}{1892}
\showcasematerialdesignfilled{text-increase}{1893}
\showcasematerialdesignfilled{text-rotate-up}{1894}
\showcasematerialdesignfilled{text-rotate-vertical}{1895}
\showcasematerialdesignfilled{text-rotation-angledown}{1896}
\showcasematerialdesignfilled{text-rotation-angleup}{1897}
\showcasematerialdesignfilled{text-rotation-down}{1898}
\showcasematerialdesignfilled{text-rotation-none}{1899}
\showcasematerialdesignfilled{text-snippet}{1900}
\showcasematerialdesignfilled{theaters}{1901}
\showcasematerialdesignfilled{theater-comedy}{1902}
\showcasematerialdesignfilled{thermostat}{1903}
\showcasematerialdesignfilled{thermostat-auto}{1904}
\showcasematerialdesignfilled{thumbs-up-down}{1905}
\showcasematerialdesignfilled{thumb-down}{1906}
\showcasematerialdesignfilled{thumb-down-alt}{1907}
\showcasematerialdesignfilled{thumb-down-off-alt}{1908}
\showcasematerialdesignfilled{thumb-up}{1909}
\showcasematerialdesignfilled{thumb-up-alt}{1910}
\showcasematerialdesignfilled{thumb-up-off-alt}{1911}
\showcasematerialdesignfilled{thunderstorm}{1912}
\showcasematerialdesignfilled{timelapse}{1913}
\showcasematerialdesignfilled{timeline}{1914}
\showcasematerialdesignfilled{timer}{1915}
\showcasematerialdesignfilled{timer-3}{1916}
\showcasematerialdesignfilled{timer-3-select}{1917}
\showcasematerialdesignfilled{timer-10}{1918}
\showcasematerialdesignfilled{timer-10-select}{1919}
\showcasematerialdesignfilled{timer-off}{1920}
\showcasematerialdesignfilled{time-to-leave}{1921}
\showcasematerialdesignfilled{tips-and-updates}{1922}
\showcasematerialdesignfilled{tire-repair}{1923}
\showcasematerialdesignfilled{title}{1924}
\showcasematerialdesignfilled{toc}{1925}
\showcasematerialdesignfilled{today}{1926}
\showcasematerialdesignfilled{toggle-off}{1927}
\showcasematerialdesignfilled{toggle-on}{1928}
\showcasematerialdesignfilled{token}{1929}
\showcasematerialdesignfilled{toll}{1930}
\showcasematerialdesignfilled{tonality}{1931}
\showcasematerialdesignfilled{topic}{1932}
\showcasematerialdesignfilled{tornado}{1933}
\showcasematerialdesignfilled{touch-app}{1934}
\showcasematerialdesignfilled{tour}{1935}
\showcasematerialdesignfilled{toys}{1936}
\showcasematerialdesignfilled{track-changes}{1937}
\showcasematerialdesignfilled{traffic}{1938}
\showcasematerialdesignfilled{train}{1939}
\showcasematerialdesignfilled{tram}{1940}
\showcasematerialdesignfilled{transcribe}{1941}
\showcasematerialdesignfilled{transfer-within-a-station}{1942}
\showcasematerialdesignfilled{transform}{1943}
\showcasematerialdesignfilled{transgender}{1944}
\showcasematerialdesignfilled{transit-enterexit}{1945}
\showcasematerialdesignfilled{translate}{1946}
\showcasematerialdesignfilled{travel-explore}{1947}
\showcasematerialdesignfilled{trending-down}{1948}
\showcasematerialdesignfilled{trending-flat}{1949}
\showcasematerialdesignfilled{trending-up}{1950}
\showcasematerialdesignfilled{trip-origin}{1951}
\showcasematerialdesignfilled{troubleshoot}{1952}
\showcasematerialdesignfilled{try}{1953}
\showcasematerialdesignfilled{tsunami}{1954}
\showcasematerialdesignfilled{tty}{1955}
\showcasematerialdesignfilled{tune}{1956}
\showcasematerialdesignfilled{tungsten}{1957}
\showcasematerialdesignfilled{turned-in}{1958}
\showcasematerialdesignfilled{turned-in-not}{1959}
\showcasematerialdesignfilled{turn-left}{1960}
\showcasematerialdesignfilled{turn-right}{1961}
\showcasematerialdesignfilled{turn-sharp-left}{1962}
\showcasematerialdesignfilled{turn-sharp-right}{1963}
\showcasematerialdesignfilled{turn-slight-left}{1964}
\showcasematerialdesignfilled{turn-slight-right}{1965}
\showcasematerialdesignfilled{tv}{1966}
\showcasematerialdesignfilled{tv-off}{1967}
\showcasematerialdesignfilled{two-wheeler}{1968}
\showcasematerialdesignfilled{type-specimen}{1969}
\showcasematerialdesignfilled{umbrella}{1970}
\showcasematerialdesignfilled{unarchive}{1971}
\showcasematerialdesignfilled{undo}{1972}
\showcasematerialdesignfilled{unfold-less}{1973}
\showcasematerialdesignfilled{unfold-less-double}{1974}
\showcasematerialdesignfilled{unfold-more}{1975}
\showcasematerialdesignfilled{unfold-more-double}{1976}
\showcasematerialdesignfilled{unpublished}{1977}
\showcasematerialdesignfilled{unsubscribe}{1978}
\showcasematerialdesignfilled{upcoming}{1979}
\showcasematerialdesignfilled{update}{1980}
\showcasematerialdesignfilled{update-disabled}{1981}
\showcasematerialdesignfilled{upgrade}{1982}
\showcasematerialdesignfilled{upload}{1983}
\showcasematerialdesignfilled{upload-file}{1984}
\showcasematerialdesignfilled{usb}{1985}
\showcasematerialdesignfilled{usb-off}{1986}
\showcasematerialdesignfilled{u-turn-left}{1987}
\showcasematerialdesignfilled{u-turn-right}{1988}
\showcasematerialdesignfilled{vaccines}{1989}
\showcasematerialdesignfilled{vape-free}{1990}
\showcasematerialdesignfilled{vaping-rooms}{1991}
\showcasematerialdesignfilled{verified}{1992}
\showcasematerialdesignfilled{verified-user}{1993}
\showcasematerialdesignfilled{vertical-align-bottom}{1994}
\showcasematerialdesignfilled{vertical-align-center}{1995}
\showcasematerialdesignfilled{vertical-align-top}{1996}
\showcasematerialdesignfilled{vertical-distribute}{1997}
\showcasematerialdesignfilled{vertical-shades}{1998}
\showcasematerialdesignfilled{vertical-shades-closed}{1999}
\showcasematerialdesignfilled{vertical-split}{2000}
\showcasematerialdesignfilled{vibration}{2001}
\showcasematerialdesignfilled{videocam}{2002}
\showcasematerialdesignfilled{videocam-off}{2003}
\showcasematerialdesignfilled{videogame-asset}{2004}
\showcasematerialdesignfilled{videogame-asset-off}{2005}
\showcasematerialdesignfilled{video-call}{2006}
\showcasematerialdesignfilled{video-camera-back}{2007}
\showcasematerialdesignfilled{video-camera-front}{2008}
\showcasematerialdesignfilled{video-chat}{2009}
\showcasematerialdesignfilled{video-file}{2010}
\showcasematerialdesignfilled{video-label}{2011}
\showcasematerialdesignfilled{video-library}{2012}
\showcasematerialdesignfilled{video-settings}{2013}
\showcasematerialdesignfilled{video-stable}{2014}
\showcasematerialdesignfilled{view-agenda}{2015}
\showcasematerialdesignfilled{view-array}{2016}
\showcasematerialdesignfilled{view-carousel}{2017}
\showcasematerialdesignfilled{view-column}{2018}
\showcasematerialdesignfilled{view-comfy}{2019}
\showcasematerialdesignfilled{view-comfy-alt}{2020}
\showcasematerialdesignfilled{view-compact}{2021}
\showcasematerialdesignfilled{view-compact-alt}{2022}
\showcasematerialdesignfilled{view-cozy}{2023}
\showcasematerialdesignfilled{view-day}{2024}
\showcasematerialdesignfilled{view-headline}{2025}
\showcasematerialdesignfilled{view-in-ar}{2026}
\showcasematerialdesignfilled{view-kanban}{2027}
\showcasematerialdesignfilled{view-list}{2028}
\showcasematerialdesignfilled{view-module}{2029}
\showcasematerialdesignfilled{view-quilt}{2030}
\showcasematerialdesignfilled{view-sidebar}{2031}
\showcasematerialdesignfilled{view-stream}{2032}
\showcasematerialdesignfilled{view-timeline}{2033}
\showcasematerialdesignfilled{view-week}{2034}
\showcasematerialdesignfilled{vignette}{2035}
\showcasematerialdesignfilled{villa}{2036}
\showcasematerialdesignfilled{visibility}{2037}
\showcasematerialdesignfilled{visibility-off}{2038}
\showcasematerialdesignfilled{voicemail}{2039}
\showcasematerialdesignfilled{voice-chat}{2040}
\showcasematerialdesignfilled{voice-over-off}{2041}
\showcasematerialdesignfilled{volcano}{2042}
\showcasematerialdesignfilled{volume-down}{2043}
\showcasematerialdesignfilled{volume-mute}{2044}
\showcasematerialdesignfilled{volume-off}{2045}
\showcasematerialdesignfilled{volume-up}{2046}
\showcasematerialdesignfilled{volunteer-activism}{2047}
\showcasematerialdesignfilled{vpn-key}{2048}
\showcasematerialdesignfilled{vpn-key-off}{2049}
\showcasematerialdesignfilled{vpn-lock}{2050}
\showcasematerialdesignfilled{vrpano}{2051}
\showcasematerialdesignfilled{wallet}{2052}
\showcasematerialdesignfilled{wallpaper}{2053}
\showcasematerialdesignfilled{warehouse}{2054}
\showcasematerialdesignfilled{warning}{2055}
\showcasematerialdesignfilled{warning-amber}{2056}
\showcasematerialdesignfilled{wash}{2057}
\showcasematerialdesignfilled{watch}{2058}
\showcasematerialdesignfilled{watch-later}{2059}
\showcasematerialdesignfilled{watch-off}{2060}
\showcasematerialdesignfilled{water}{2061}
\showcasematerialdesignfilled{waterfall-chart}{2062}
\showcasematerialdesignfilled{water-damage}{2063}
\showcasematerialdesignfilled{water-drop}{2064}
\showcasematerialdesignfilled{waves}{2065}
\showcasematerialdesignfilled{waving-hand}{2066}
\showcasematerialdesignfilled{wb-auto}{2067}
\showcasematerialdesignfilled{wb-cloudy}{2068}
\showcasematerialdesignfilled{wb-incandescent}{2069}
\showcasematerialdesignfilled{wb-iridescent}{2070}
\showcasematerialdesignfilled{wb-shade}{2071}
\showcasematerialdesignfilled{wb-sunny}{2072}
\showcasematerialdesignfilled{wb-twilight}{2073}
\showcasematerialdesignfilled{wc}{2074}
\showcasematerialdesignfilled{web}{2075}
\showcasematerialdesignfilled{webhook}{2076}
\showcasematerialdesignfilled{web-asset}{2077}
\showcasematerialdesignfilled{web-asset-off}{2078}
\showcasematerialdesignfilled{web-stories}{2079}
\showcasematerialdesignfilled{weekend}{2080}
\showcasematerialdesignfilled{west}{2081}
\showcasematerialdesignfilled{whatshot}{2082}
\showcasematerialdesignfilled{wheelchair-pickup}{2083}
\showcasematerialdesignfilled{where-to-vote}{2084}
\showcasematerialdesignfilled{widgets}{2085}
\showcasematerialdesignfilled{width-full}{2086}
\showcasematerialdesignfilled{width-normal}{2087}
\showcasematerialdesignfilled{width-wide}{2088}
\showcasematerialdesignfilled{wifi}{2089}
\showcasematerialdesignfilled{wifi-1-bar}{2090}
\showcasematerialdesignfilled{wifi-2-bar}{2091}
\showcasematerialdesignfilled{wifi-calling}{2092}
\showcasematerialdesignfilled{wifi-calling-3}{2093}
\showcasematerialdesignfilled{wifi-channel}{2094}
\showcasematerialdesignfilled{wifi-find}{2095}
\showcasematerialdesignfilled{wifi-lock}{2096}
\showcasematerialdesignfilled{wifi-off}{2097}
\showcasematerialdesignfilled{wifi-password}{2098}
\showcasematerialdesignfilled{wifi-protected-setup}{2099}
\showcasematerialdesignfilled{wifi-tethering}{2100}
\showcasematerialdesignfilled{wifi-tethering-error}{2101}
\showcasematerialdesignfilled{wifi-tethering-off}{2102}
\showcasematerialdesignfilled{window}{2103}
\showcasematerialdesignfilled{wind-power}{2104}
\showcasematerialdesignfilled{wine-bar}{2105}
\showcasematerialdesignfilled{woman}{2106}
\showcasematerialdesignfilled{woman-2}{2107}
\showcasematerialdesignfilled{work}{2108}
\showcasematerialdesignfilled{workspaces}{2109}
\showcasematerialdesignfilled{workspace-premium}{2110}
\showcasematerialdesignfilled{work-history}{2111}
\showcasematerialdesignfilled{work-off}{2112}
\showcasematerialdesignfilled{work-outline}{2113}
\showcasematerialdesignfilled{wrap-text}{2114}
\showcasematerialdesignfilled{wrong-location}{2115}
\showcasematerialdesignfilled{wysiwyg}{2116}
\showcasematerialdesignfilled{yard}{2117}
\showcasematerialdesignfilled{youtube-searched-for}{2118}
\showcasematerialdesignfilled{zoom-in}{2119}
\showcasematerialdesignfilled{zoom-in-map}{2120}
\showcasematerialdesignfilled{zoom-out}{2121}
\showcasematerialdesignfilled{zoom-out-map}{2122}
\end{materialdesigncase}

\subsection{Round version (2122 icons)}

The \textsf{PDF} file is \texttt{materialdesign-round-all.pdf}

\begin{materialdesigncase}
\showcasematerialdesignround{1k}{1}
\showcasematerialdesignround{1k-plus}{2}
\showcasematerialdesignround{1x-mobiledata}{3}
\showcasematerialdesignround{2k}{4}
\showcasematerialdesignround{2k-plus}{5}
\showcasematerialdesignround{2mp}{6}
\showcasematerialdesignround{3d-rotation}{7}
\showcasematerialdesignround{3g-mobiledata}{8}
\showcasematerialdesignround{3k}{9}
\showcasematerialdesignround{3k-plus}{10}
\showcasematerialdesignround{3mp}{11}
\showcasematerialdesignround{3p}{12}
\showcasematerialdesignround{4g-mobiledata}{13}
\showcasematerialdesignround{4g-plus-mobiledata}{14}
\showcasematerialdesignround{4k}{15}
\showcasematerialdesignround{4k-plus}{16}
\showcasematerialdesignround{4mp}{17}
\showcasematerialdesignround{5g}{18}
\showcasematerialdesignround{5k}{19}
\showcasematerialdesignround{5k-plus}{20}
\showcasematerialdesignround{5mp}{21}
\showcasematerialdesignround{6k}{22}
\showcasematerialdesignround{6k-plus}{23}
\showcasematerialdesignround{6mp}{24}
\showcasematerialdesignround{6-ft-apart}{25}
\showcasematerialdesignround{7k}{26}
\showcasematerialdesignround{7k-plus}{27}
\showcasematerialdesignround{7mp}{28}
\showcasematerialdesignround{8k}{29}
\showcasematerialdesignround{8k-plus}{30}
\showcasematerialdesignround{8mp}{31}
\showcasematerialdesignround{9k}{32}
\showcasematerialdesignround{9k-plus}{33}
\showcasematerialdesignround{9mp}{34}
\showcasematerialdesignround{10k}{35}
\showcasematerialdesignround{10mp}{36}
\showcasematerialdesignround{11mp}{37}
\showcasematerialdesignround{12mp}{38}
\showcasematerialdesignround{13mp}{39}
\showcasematerialdesignround{14mp}{40}
\showcasematerialdesignround{15mp}{41}
\showcasematerialdesignround{16mp}{42}
\showcasematerialdesignround{17mp}{43}
\showcasematerialdesignround{18mp}{44}
\showcasematerialdesignround{18-up-rating}{45}
\showcasematerialdesignround{19mp}{46}
\showcasematerialdesignround{20mp}{47}
\showcasematerialdesignround{21mp}{48}
\showcasematerialdesignround{22mp}{49}
\showcasematerialdesignround{23mp}{50}
\showcasematerialdesignround{24mp}{51}
\showcasematerialdesignround{30fps}{52}
\showcasematerialdesignround{30fps-select}{53}
\showcasematerialdesignround{60fps}{54}
\showcasematerialdesignround{60fps-select}{55}
\showcasematerialdesignround{123}{56}
\showcasematerialdesignround{360}{57}
\showcasematerialdesignround{abc}{58}
\showcasematerialdesignround{accessibility}{59}
\showcasematerialdesignround{accessibility-new}{60}
\showcasematerialdesignround{accessible}{61}
\showcasematerialdesignround{accessible-forward}{62}
\showcasematerialdesignround{access-alarm}{63}
\showcasematerialdesignround{access-alarms}{64}
\showcasematerialdesignround{access-time}{65}
\showcasematerialdesignround{access-time-filled}{66}
\showcasematerialdesignround{account-balance}{67}
\showcasematerialdesignround{account-balance-wallet}{68}
\showcasematerialdesignround{account-box}{69}
\showcasematerialdesignround{account-circle}{70}
\showcasematerialdesignround{account-tree}{71}
\showcasematerialdesignround{ac-unit}{72}
\showcasematerialdesignround{adb}{73}
\showcasematerialdesignround{add}{74}
\showcasematerialdesignround{addchart}{75}
\showcasematerialdesignround{add-alarm}{76}
\showcasematerialdesignround{add-alert}{77}
\showcasematerialdesignround{add-a-photo}{78}
\showcasematerialdesignround{add-box}{79}
\showcasematerialdesignround{add-business}{80}
\showcasematerialdesignround{add-card}{81}
\showcasematerialdesignround{add-chart}{82}
\showcasematerialdesignround{add-circle}{83}
\showcasematerialdesignround{add-circle-outline}{84}
\showcasematerialdesignround{add-comment}{85}
\showcasematerialdesignround{add-home}{86}
\showcasematerialdesignround{add-home-work}{87}
\showcasematerialdesignround{add-ic-call}{88}
\showcasematerialdesignround{add-link}{89}
\showcasematerialdesignround{add-location}{90}
\showcasematerialdesignround{add-location-alt}{91}
\showcasematerialdesignround{add-moderator}{92}
\showcasematerialdesignround{add-photo-alternate}{93}
\showcasematerialdesignround{add-reaction}{94}
\showcasematerialdesignround{add-road}{95}
\showcasematerialdesignround{add-shopping-cart}{96}
\showcasematerialdesignround{add-task}{97}
\showcasematerialdesignround{add-to-drive}{98}
\showcasematerialdesignround{add-to-home-screen}{99}
\showcasematerialdesignround{add-to-photos}{100}
\showcasematerialdesignround{add-to-queue}{101}
\showcasematerialdesignround{adf-scanner}{102}
\showcasematerialdesignround{adjust}{103}
\showcasematerialdesignround{admin-panel-settings}{104}
\showcasematerialdesignround{ads-click}{105}
\showcasematerialdesignround{ad-units}{106}
\showcasematerialdesignround{agriculture}{107}
\showcasematerialdesignround{air}{108}
\showcasematerialdesignround{airlines}{109}
\showcasematerialdesignround{airline-seat-flat}{110}
\showcasematerialdesignround{airline-seat-flat-angled}{111}
\showcasematerialdesignround{airline-seat-individual-suite}{112}
\showcasematerialdesignround{airline-seat-legroom-extra}{113}
\showcasematerialdesignround{airline-seat-legroom-normal}{114}
\showcasematerialdesignround{airline-seat-legroom-reduced}{115}
\showcasematerialdesignround{airline-seat-recline-extra}{116}
\showcasematerialdesignround{airline-seat-recline-normal}{117}
\showcasematerialdesignround{airline-stops}{118}
\showcasematerialdesignround{airplanemode-active}{119}
\showcasematerialdesignround{airplanemode-inactive}{120}
\showcasematerialdesignround{airplane-ticket}{121}
\showcasematerialdesignround{airplay}{122}
\showcasematerialdesignround{airport-shuttle}{123}
\showcasematerialdesignround{alarm}{124}
\showcasematerialdesignround{alarm-add}{125}
\showcasematerialdesignround{alarm-off}{126}
\showcasematerialdesignround{alarm-on}{127}
\showcasematerialdesignround{album}{128}
\showcasematerialdesignround{align-horizontal-center}{129}
\showcasematerialdesignround{align-horizontal-left}{130}
\showcasematerialdesignround{align-horizontal-right}{131}
\showcasematerialdesignround{align-vertical-bottom}{132}
\showcasematerialdesignround{align-vertical-center}{133}
\showcasematerialdesignround{align-vertical-top}{134}
\showcasematerialdesignround{all-inbox}{135}
\showcasematerialdesignround{all-inclusive}{136}
\showcasematerialdesignround{all-out}{137}
\showcasematerialdesignround{alternate-email}{138}
\showcasematerialdesignround{alt-route}{139}
\showcasematerialdesignround{analytics}{140}
\showcasematerialdesignround{anchor}{141}
\showcasematerialdesignround{android}{142}
\showcasematerialdesignround{animation}{143}
\showcasematerialdesignround{announcement}{144}
\showcasematerialdesignround{aod}{145}
\showcasematerialdesignround{apartment}{146}
\showcasematerialdesignround{api}{147}
\showcasematerialdesignround{approval}{148}
\showcasematerialdesignround{apps}{149}
\showcasematerialdesignround{apps-outage}{150}
\showcasematerialdesignround{app-blocking}{151}
\showcasematerialdesignround{app-registration}{152}
\showcasematerialdesignround{app-settings-alt}{153}
\showcasematerialdesignround{app-shortcut}{154}
\showcasematerialdesignround{architecture}{155}
\showcasematerialdesignround{archive}{156}
\showcasematerialdesignround{area-chart}{157}
\showcasematerialdesignround{arrow-back}{158}
\showcasematerialdesignround{arrow-back-ios}{159}
\showcasematerialdesignround{arrow-back-ios-new}{160}
\showcasematerialdesignround{arrow-circle-down}{161}
\showcasematerialdesignround{arrow-circle-left}{162}
\showcasematerialdesignround{arrow-circle-right}{163}
\showcasematerialdesignround{arrow-circle-up}{164}
\showcasematerialdesignround{arrow-downward}{165}
\showcasematerialdesignround{arrow-drop-down}{166}
\showcasematerialdesignround{arrow-drop-down-circle}{167}
\showcasematerialdesignround{arrow-drop-up}{168}
\showcasematerialdesignround{arrow-forward}{169}
\showcasematerialdesignround{arrow-forward-ios}{170}
\showcasematerialdesignround{arrow-left}{171}
\showcasematerialdesignround{arrow-outward}{172}
\showcasematerialdesignround{arrow-right}{173}
\showcasematerialdesignround{arrow-right-alt}{174}
\showcasematerialdesignround{arrow-upward}{175}
\showcasematerialdesignround{article}{176}
\showcasematerialdesignround{art-track}{177}
\showcasematerialdesignround{aspect-ratio}{178}
\showcasematerialdesignround{assessment}{179}
\showcasematerialdesignround{assignment}{180}
\showcasematerialdesignround{assignment-ind}{181}
\showcasematerialdesignround{assignment-late}{182}
\showcasematerialdesignround{assignment-return}{183}
\showcasematerialdesignround{assignment-returned}{184}
\showcasematerialdesignround{assignment-turned-in}{185}
\showcasematerialdesignround{assistant}{186}
\showcasematerialdesignround{assistant-direction}{187}
\showcasematerialdesignround{assistant-photo}{188}
\showcasematerialdesignround{assist-walker}{189}
\showcasematerialdesignround{assured-workload}{190}
\showcasematerialdesignround{atm}{191}
\showcasematerialdesignround{attachment}{192}
\showcasematerialdesignround{attach-email}{193}
\showcasematerialdesignround{attach-file}{194}
\showcasematerialdesignround{attach-money}{195}
\showcasematerialdesignround{attractions}{196}
\showcasematerialdesignround{attribution}{197}
\showcasematerialdesignround{audiotrack}{198}
\showcasematerialdesignround{audio-file}{199}
\showcasematerialdesignround{autofps-select}{200}
\showcasematerialdesignround{autorenew}{201}
\showcasematerialdesignround{auto-awesome}{202}
\showcasematerialdesignround{auto-awesome-mosaic}{203}
\showcasematerialdesignround{auto-awesome-motion}{204}
\showcasematerialdesignround{auto-delete}{205}
\showcasematerialdesignround{auto-fix-high}{206}
\showcasematerialdesignround{auto-fix-normal}{207}
\showcasematerialdesignround{auto-fix-off}{208}
\showcasematerialdesignround{auto-graph}{209}
\showcasematerialdesignround{auto-mode}{210}
\showcasematerialdesignround{auto-stories}{211}
\showcasematerialdesignround{av-timer}{212}
\showcasematerialdesignround{baby-changing-station}{213}
\showcasematerialdesignround{backpack}{214}
\showcasematerialdesignround{backspace}{215}
\showcasematerialdesignround{backup}{216}
\showcasematerialdesignround{backup-table}{217}
\showcasematerialdesignround{back-hand}{218}
\showcasematerialdesignround{badge}{219}
\showcasematerialdesignround{bakery-dining}{220}
\showcasematerialdesignround{balance}{221}
\showcasematerialdesignround{balcony}{222}
\showcasematerialdesignround{ballot}{223}
\showcasematerialdesignround{bar-chart}{224}
\showcasematerialdesignround{batch-prediction}{225}
\showcasematerialdesignround{bathroom}{226}
\showcasematerialdesignround{bathtub}{227}
\showcasematerialdesignround{battery-0-bar}{228}
\showcasematerialdesignround{battery-1-bar}{229}
\showcasematerialdesignround{battery-2-bar}{230}
\showcasematerialdesignround{battery-3-bar}{231}
\showcasematerialdesignround{battery-4-bar}{232}
\showcasematerialdesignround{battery-5-bar}{233}
\showcasematerialdesignround{battery-6-bar}{234}
\showcasematerialdesignround{battery-alert}{235}
\showcasematerialdesignround{battery-charging-full}{236}
\showcasematerialdesignround{battery-full}{237}
\showcasematerialdesignround{battery-saver}{238}
\showcasematerialdesignround{battery-std}{239}
\showcasematerialdesignround{battery-unknown}{240}
\showcasematerialdesignround{beach-access}{241}
\showcasematerialdesignround{bed}{242}
\showcasematerialdesignround{bedroom-baby}{243}
\showcasematerialdesignround{bedroom-child}{244}
\showcasematerialdesignround{bedroom-parent}{245}
\showcasematerialdesignround{bedtime}{246}
\showcasematerialdesignround{bedtime-off}{247}
\showcasematerialdesignround{beenhere}{248}
\showcasematerialdesignround{bento}{249}
\showcasematerialdesignround{bike-scooter}{250}
\showcasematerialdesignround{biotech}{251}
\showcasematerialdesignround{blender}{252}
\showcasematerialdesignround{blind}{253}
\showcasematerialdesignround{blinds}{254}
\showcasematerialdesignround{blinds-closed}{255}
\showcasematerialdesignround{block}{256}
\showcasematerialdesignround{bloodtype}{257}
\showcasematerialdesignround{bluetooth}{258}
\showcasematerialdesignround{bluetooth-audio}{259}
\showcasematerialdesignround{bluetooth-connected}{260}
\showcasematerialdesignround{bluetooth-disabled}{261}
\showcasematerialdesignround{bluetooth-drive}{262}
\showcasematerialdesignround{bluetooth-searching}{263}
\showcasematerialdesignround{blur-circular}{264}
\showcasematerialdesignround{blur-linear}{265}
\showcasematerialdesignround{blur-off}{266}
\showcasematerialdesignround{blur-on}{267}
\showcasematerialdesignround{bolt}{268}
\showcasematerialdesignround{book}{269}
\showcasematerialdesignround{bookmark}{270}
\showcasematerialdesignround{bookmarks}{271}
\showcasematerialdesignround{bookmark-add}{272}
\showcasematerialdesignround{bookmark-added}{273}
\showcasematerialdesignround{bookmark-border}{274}
\showcasematerialdesignround{bookmark-remove}{275}
\showcasematerialdesignround{book-online}{276}
\showcasematerialdesignround{border-all}{277}
\showcasematerialdesignround{border-bottom}{278}
\showcasematerialdesignround{border-clear}{279}
\showcasematerialdesignround{border-color}{280}
\showcasematerialdesignround{border-horizontal}{281}
\showcasematerialdesignround{border-inner}{282}
\showcasematerialdesignround{border-left}{283}
\showcasematerialdesignround{border-outer}{284}
\showcasematerialdesignround{border-right}{285}
\showcasematerialdesignround{border-style}{286}
\showcasematerialdesignround{border-top}{287}
\showcasematerialdesignround{border-vertical}{288}
\showcasematerialdesignround{boy}{289}
\showcasematerialdesignround{branding-watermark}{290}
\showcasematerialdesignround{breakfast-dining}{291}
\showcasematerialdesignround{brightness-1}{292}
\showcasematerialdesignround{brightness-2}{293}
\showcasematerialdesignround{brightness-3}{294}
\showcasematerialdesignround{brightness-4}{295}
\showcasematerialdesignround{brightness-5}{296}
\showcasematerialdesignround{brightness-6}{297}
\showcasematerialdesignround{brightness-7}{298}
\showcasematerialdesignround{brightness-auto}{299}
\showcasematerialdesignround{brightness-high}{300}
\showcasematerialdesignround{brightness-low}{301}
\showcasematerialdesignround{brightness-medium}{302}
\showcasematerialdesignround{broadcast-on-home}{303}
\showcasematerialdesignround{broadcast-on-personal}{304}
\showcasematerialdesignround{broken-image}{305}
\showcasematerialdesignround{browser-not-supported}{306}
\showcasematerialdesignround{browser-updated}{307}
\showcasematerialdesignround{browse-gallery}{308}
\showcasematerialdesignround{brunch-dining}{309}
\showcasematerialdesignround{brush}{310}
\showcasematerialdesignround{bubble-chart}{311}
\showcasematerialdesignround{bug-report}{312}
\showcasematerialdesignround{build}{313}
\showcasematerialdesignround{build-circle}{314}
\showcasematerialdesignround{bungalow}{315}
\showcasematerialdesignround{burst-mode}{316}
\showcasematerialdesignround{business}{317}
\showcasematerialdesignround{business-center}{318}
\showcasematerialdesignround{bus-alert}{319}
\showcasematerialdesignround{cabin}{320}
\showcasematerialdesignround{cable}{321}
\showcasematerialdesignround{cached}{322}
\showcasematerialdesignround{cake}{323}
\showcasematerialdesignround{calculate}{324}
\showcasematerialdesignround{calendar-month}{325}
\showcasematerialdesignround{calendar-today}{326}
\showcasematerialdesignround{calendar-view-day}{327}
\showcasematerialdesignround{calendar-view-month}{328}
\showcasematerialdesignround{calendar-view-week}{329}
\showcasematerialdesignround{call}{330}
\showcasematerialdesignround{call-end}{331}
\showcasematerialdesignround{call-made}{332}
\showcasematerialdesignround{call-merge}{333}
\showcasematerialdesignround{call-missed}{334}
\showcasematerialdesignround{call-missed-outgoing}{335}
\showcasematerialdesignround{call-received}{336}
\showcasematerialdesignround{call-split}{337}
\showcasematerialdesignround{call-to-action}{338}
\showcasematerialdesignround{camera}{339}
\showcasematerialdesignround{cameraswitch}{340}
\showcasematerialdesignround{camera-alt}{341}
\showcasematerialdesignround{camera-enhance}{342}
\showcasematerialdesignround{camera-front}{343}
\showcasematerialdesignround{camera-indoor}{344}
\showcasematerialdesignround{camera-outdoor}{345}
\showcasematerialdesignround{camera-rear}{346}
\showcasematerialdesignround{camera-roll}{347}
\showcasematerialdesignround{campaign}{348}
\showcasematerialdesignround{cancel}{349}
\showcasematerialdesignround{cancel-presentation}{350}
\showcasematerialdesignround{cancel-schedule-send}{351}
\showcasematerialdesignround{candlestick-chart}{352}
\showcasematerialdesignround{card-giftcard}{353}
\showcasematerialdesignround{card-membership}{354}
\showcasematerialdesignround{card-travel}{355}
\showcasematerialdesignround{carpenter}{356}
\showcasematerialdesignround{car-crash}{357}
\showcasematerialdesignround{car-rental}{358}
\showcasematerialdesignround{car-repair}{359}
\showcasematerialdesignround{cases}{360}
\showcasematerialdesignround{casino}{361}
\showcasematerialdesignround{cast}{362}
\showcasematerialdesignround{castle}{363}
\showcasematerialdesignround{cast-connected}{364}
\showcasematerialdesignround{cast-for-education}{365}
\showcasematerialdesignround{catching-pokemon}{366}
\showcasematerialdesignround{category}{367}
\showcasematerialdesignround{celebration}{368}
\showcasematerialdesignround{cell-tower}{369}
\showcasematerialdesignround{cell-wifi}{370}
\showcasematerialdesignround{center-focus-strong}{371}
\showcasematerialdesignround{center-focus-weak}{372}
\showcasematerialdesignround{chair}{373}
\showcasematerialdesignround{chair-alt}{374}
\showcasematerialdesignround{chalet}{375}
\showcasematerialdesignround{change-circle}{376}
\showcasematerialdesignround{change-history}{377}
\showcasematerialdesignround{charging-station}{378}
\showcasematerialdesignround{chat}{379}
\showcasematerialdesignround{chat-bubble}{380}
\showcasematerialdesignround{chat-bubble-outline}{381}
\showcasematerialdesignround{check}{382}
\showcasematerialdesignround{checklist}{383}
\showcasematerialdesignround{checklist-rtl}{384}
\showcasematerialdesignround{checkroom}{385}
\showcasematerialdesignround{check-box}{386}
\showcasematerialdesignround{check-box-outline-blank}{387}
\showcasematerialdesignround{check-circle}{388}
\showcasematerialdesignround{check-circle-outline}{389}
\showcasematerialdesignround{chevron-left}{390}
\showcasematerialdesignround{chevron-right}{391}
\showcasematerialdesignround{child-care}{392}
\showcasematerialdesignround{child-friendly}{393}
\showcasematerialdesignround{chrome-reader-mode}{394}
\showcasematerialdesignround{church}{395}
\showcasematerialdesignround{circle}{396}
\showcasematerialdesignround{circle-notifications}{397}
\showcasematerialdesignround{class}{398}
\showcasematerialdesignround{cleaning-services}{399}
\showcasematerialdesignround{clean-hands}{400}
\showcasematerialdesignround{clear}{401}
\showcasematerialdesignround{clear-all}{402}
\showcasematerialdesignround{close}{403}
\showcasematerialdesignround{closed-caption}{404}
\showcasematerialdesignround{closed-caption-disabled}{405}
\showcasematerialdesignround{closed-caption-off}{406}
\showcasematerialdesignround{close-fullscreen}{407}
\showcasematerialdesignround{cloud}{408}
\showcasematerialdesignround{cloud-circle}{409}
\showcasematerialdesignround{cloud-done}{410}
\showcasematerialdesignround{cloud-download}{411}
\showcasematerialdesignround{cloud-off}{412}
\showcasematerialdesignround{cloud-queue}{413}
\showcasematerialdesignround{cloud-sync}{414}
\showcasematerialdesignround{cloud-upload}{415}
\showcasematerialdesignround{co2}{416}
\showcasematerialdesignround{code}{417}
\showcasematerialdesignround{code-off}{418}
\showcasematerialdesignround{coffee}{419}
\showcasematerialdesignround{coffee-maker}{420}
\showcasematerialdesignround{collections}{421}
\showcasematerialdesignround{collections-bookmark}{422}
\showcasematerialdesignround{colorize}{423}
\showcasematerialdesignround{color-lens}{424}
\showcasematerialdesignround{comment}{425}
\showcasematerialdesignround{comments-disabled}{426}
\showcasematerialdesignround{comment-bank}{427}
\showcasematerialdesignround{commit}{428}
\showcasematerialdesignround{commute}{429}
\showcasematerialdesignround{compare}{430}
\showcasematerialdesignround{compare-arrows}{431}
\showcasematerialdesignround{compass-calibration}{432}
\showcasematerialdesignround{compost}{433}
\showcasematerialdesignround{compress}{434}
\showcasematerialdesignround{computer}{435}
\showcasematerialdesignround{confirmation-number}{436}
\showcasematerialdesignround{connected-tv}{437}
\showcasematerialdesignround{connecting-airports}{438}
\showcasematerialdesignround{connect-without-contact}{439}
\showcasematerialdesignround{construction}{440}
\showcasematerialdesignround{contactless}{441}
\showcasematerialdesignround{contacts}{442}
\showcasematerialdesignround{contact-emergency}{443}
\showcasematerialdesignround{contact-mail}{444}
\showcasematerialdesignround{contact-page}{445}
\showcasematerialdesignround{contact-phone}{446}
\showcasematerialdesignround{contact-support}{447}
\showcasematerialdesignround{content-copy}{448}
\showcasematerialdesignround{content-cut}{449}
\showcasematerialdesignround{content-paste}{450}
\showcasematerialdesignround{content-paste-go}{451}
\showcasematerialdesignround{content-paste-off}{452}
\showcasematerialdesignround{content-paste-search}{453}
\showcasematerialdesignround{contrast}{454}
\showcasematerialdesignround{control-camera}{455}
\showcasematerialdesignround{control-point}{456}
\showcasematerialdesignround{control-point-duplicate}{457}
\showcasematerialdesignround{cookie}{458}
\showcasematerialdesignround{copyright}{459}
\showcasematerialdesignround{copy-all}{460}
\showcasematerialdesignround{coronavirus}{461}
\showcasematerialdesignround{corporate-fare}{462}
\showcasematerialdesignround{cottage}{463}
\showcasematerialdesignround{countertops}{464}
\showcasematerialdesignround{co-present}{465}
\showcasematerialdesignround{create}{466}
\showcasematerialdesignround{create-new-folder}{467}
\showcasematerialdesignround{credit-card}{468}
\showcasematerialdesignround{credit-card-off}{469}
\showcasematerialdesignround{credit-score}{470}
\showcasematerialdesignround{crib}{471}
\showcasematerialdesignround{crisis-alert}{472}
\showcasematerialdesignround{crop}{473}
\showcasematerialdesignround{crop-3-2}{474}
\showcasematerialdesignround{crop-5-4}{475}
\showcasematerialdesignround{crop-7-5}{476}
\showcasematerialdesignround{crop-16-9}{477}
\showcasematerialdesignround{crop-din}{478}
\showcasematerialdesignround{crop-free}{479}
\showcasematerialdesignround{crop-landscape}{480}
\showcasematerialdesignround{crop-original}{481}
\showcasematerialdesignround{crop-portrait}{482}
\showcasematerialdesignround{crop-rotate}{483}
\showcasematerialdesignround{crop-square}{484}
\showcasematerialdesignround{cruelty-free}{485}
\showcasematerialdesignround{css}{486}
\showcasematerialdesignround{currency-bitcoin}{487}
\showcasematerialdesignround{currency-exchange}{488}
\showcasematerialdesignround{currency-franc}{489}
\showcasematerialdesignround{currency-lira}{490}
\showcasematerialdesignround{currency-pound}{491}
\showcasematerialdesignround{currency-ruble}{492}
\showcasematerialdesignround{currency-rupee}{493}
\showcasematerialdesignround{currency-yen}{494}
\showcasematerialdesignround{currency-yuan}{495}
\showcasematerialdesignround{curtains}{496}
\showcasematerialdesignround{curtains-closed}{497}
\showcasematerialdesignround{cyclone}{498}
\showcasematerialdesignround{dangerous}{499}
\showcasematerialdesignround{dark-mode}{500}
\showcasematerialdesignround{dashboard}{501}
\showcasematerialdesignround{dashboard-customize}{502}
\showcasematerialdesignround{dataset}{503}
\showcasematerialdesignround{dataset-linked}{504}
\showcasematerialdesignround{data-array}{505}
\showcasematerialdesignround{data-exploration}{506}
\showcasematerialdesignround{data-object}{507}
\showcasematerialdesignround{data-saver-off}{508}
\showcasematerialdesignround{data-saver-on}{509}
\showcasematerialdesignround{data-thresholding}{510}
\showcasematerialdesignround{data-usage}{511}
\showcasematerialdesignround{date-range}{512}
\showcasematerialdesignround{deblur}{513}
\showcasematerialdesignround{deck}{514}
\showcasematerialdesignround{dehaze}{515}
\showcasematerialdesignround{delete}{516}
\showcasematerialdesignround{delete-forever}{517}
\showcasematerialdesignround{delete-outline}{518}
\showcasematerialdesignround{delete-sweep}{519}
\showcasematerialdesignround{delivery-dining}{520}
\showcasematerialdesignround{density-large}{521}
\showcasematerialdesignround{density-medium}{522}
\showcasematerialdesignround{density-small}{523}
\showcasematerialdesignround{departure-board}{524}
\showcasematerialdesignround{description}{525}
\showcasematerialdesignround{deselect}{526}
\showcasematerialdesignround{design-services}{527}
\showcasematerialdesignround{desk}{528}
\showcasematerialdesignround{desktop-access-disabled}{529}
\showcasematerialdesignround{desktop-mac}{530}
\showcasematerialdesignround{desktop-windows}{531}
\showcasematerialdesignround{details}{532}
\showcasematerialdesignround{developer-board}{533}
\showcasematerialdesignround{developer-board-off}{534}
\showcasematerialdesignround{developer-mode}{535}
\showcasematerialdesignround{devices}{536}
\showcasematerialdesignround{devices-fold}{537}
\showcasematerialdesignround{devices-other}{538}
\showcasematerialdesignround{device-hub}{539}
\showcasematerialdesignround{device-thermostat}{540}
\showcasematerialdesignround{device-unknown}{541}
\showcasematerialdesignround{dialer-sip}{542}
\showcasematerialdesignround{dialpad}{543}
\showcasematerialdesignround{diamond}{544}
\showcasematerialdesignround{difference}{545}
\showcasematerialdesignround{dining}{546}
\showcasematerialdesignround{dinner-dining}{547}
\showcasematerialdesignround{directions}{548}
\showcasematerialdesignround{directions-bike}{549}
\showcasematerialdesignround{directions-boat}{550}
\showcasematerialdesignround{directions-boat-filled}{551}
\showcasematerialdesignround{directions-bus}{552}
\showcasematerialdesignround{directions-bus-filled}{553}
\showcasematerialdesignround{directions-car}{554}
\showcasematerialdesignround{directions-car-filled}{555}
\showcasematerialdesignround{directions-off}{556}
\showcasematerialdesignround{directions-railway}{557}
\showcasematerialdesignround{directions-railway-filled}{558}
\showcasematerialdesignround{directions-run}{559}
\showcasematerialdesignround{directions-subway}{560}
\showcasematerialdesignround{directions-subway-filled}{561}
\showcasematerialdesignround{directions-transit}{562}
\showcasematerialdesignround{directions-transit-filled}{563}
\showcasematerialdesignround{directions-walk}{564}
\showcasematerialdesignround{dirty-lens}{565}
\showcasematerialdesignround{disabled-by-default}{566}
\showcasematerialdesignround{disabled-visible}{567}
\showcasematerialdesignround{discount}{568}
\showcasematerialdesignround{disc-full}{569}
\showcasematerialdesignround{display-settings}{570}
\showcasematerialdesignround{diversity-1}{571}
\showcasematerialdesignround{diversity-2}{572}
\showcasematerialdesignround{diversity-3}{573}
\showcasematerialdesignround{dns}{574}
\showcasematerialdesignround{dock}{575}
\showcasematerialdesignround{document-scanner}{576}
\showcasematerialdesignround{domain}{577}
\showcasematerialdesignround{domain-add}{578}
\showcasematerialdesignround{domain-disabled}{579}
\showcasematerialdesignround{domain-verification}{580}
\showcasematerialdesignround{done}{581}
\showcasematerialdesignround{done-all}{582}
\showcasematerialdesignround{done-outline}{583}
\showcasematerialdesignround{donut-large}{584}
\showcasematerialdesignround{donut-small}{585}
\showcasematerialdesignround{doorbell}{586}
\showcasematerialdesignround{door-back}{587}
\showcasematerialdesignround{door-front}{588}
\showcasematerialdesignround{door-sliding}{589}
\showcasematerialdesignround{double-arrow}{590}
\showcasematerialdesignround{downhill-skiing}{591}
\showcasematerialdesignround{download}{592}
\showcasematerialdesignround{downloading}{593}
\showcasematerialdesignround{download-done}{594}
\showcasematerialdesignround{download-for-offline}{595}
\showcasematerialdesignround{do-disturb}{596}
\showcasematerialdesignround{do-disturb-alt}{597}
\showcasematerialdesignround{do-disturb-off}{598}
\showcasematerialdesignround{do-disturb-on}{599}
\showcasematerialdesignround{do-not-disturb}{600}
\showcasematerialdesignround{do-not-disturb-alt}{601}
\showcasematerialdesignround{do-not-disturb-off}{602}
\showcasematerialdesignround{do-not-disturb-on}{603}
\showcasematerialdesignround{do-not-disturb-on-total-silence}{604}
\showcasematerialdesignround{do-not-step}{605}
\showcasematerialdesignround{do-not-touch}{606}
\showcasematerialdesignround{drafts}{607}
\showcasematerialdesignround{drag-handle}{608}
\showcasematerialdesignround{drag-indicator}{609}
\showcasematerialdesignround{draw}{610}
\showcasematerialdesignround{drive-eta}{611}
\showcasematerialdesignround{drive-file-move}{612}
\showcasematerialdesignround{drive-file-move-rtl}{613}
\showcasematerialdesignround{drive-file-rename-outline}{614}
\showcasematerialdesignround{drive-folder-upload}{615}
\showcasematerialdesignround{dry}{616}
\showcasematerialdesignround{dry-cleaning}{617}
\showcasematerialdesignround{duo}{618}
\showcasematerialdesignround{dvr}{619}
\showcasematerialdesignround{dynamic-feed}{620}
\showcasematerialdesignround{dynamic-form}{621}
\showcasematerialdesignround{earbuds}{622}
\showcasematerialdesignround{earbuds-battery}{623}
\showcasematerialdesignround{east}{624}
\showcasematerialdesignround{edgesensor-high}{625}
\showcasematerialdesignround{edgesensor-low}{626}
\showcasematerialdesignround{edit}{627}
\showcasematerialdesignround{edit-attributes}{628}
\showcasematerialdesignround{edit-calendar}{629}
\showcasematerialdesignround{edit-location}{630}
\showcasematerialdesignround{edit-location-alt}{631}
\showcasematerialdesignround{edit-note}{632}
\showcasematerialdesignround{edit-notifications}{633}
\showcasematerialdesignround{edit-off}{634}
\showcasematerialdesignround{edit-road}{635}
\showcasematerialdesignround{egg}{636}
\showcasematerialdesignround{egg-alt}{637}
\showcasematerialdesignround{eject}{638}
\showcasematerialdesignround{elderly}{639}
\showcasematerialdesignround{elderly-woman}{640}
\showcasematerialdesignround{electrical-services}{641}
\showcasematerialdesignround{electric-bike}{642}
\showcasematerialdesignround{electric-bolt}{643}
\showcasematerialdesignround{electric-car}{644}
\showcasematerialdesignround{electric-meter}{645}
\showcasematerialdesignround{electric-moped}{646}
\showcasematerialdesignround{electric-rickshaw}{647}
\showcasematerialdesignround{electric-scooter}{648}
\showcasematerialdesignround{elevator}{649}
\showcasematerialdesignround{email}{650}
\showcasematerialdesignround{emergency}{651}
\showcasematerialdesignround{emergency-recording}{652}
\showcasematerialdesignround{emergency-share}{653}
\showcasematerialdesignround{emoji-emotions}{654}
\showcasematerialdesignround{emoji-events}{655}
\showcasematerialdesignround{emoji-food-beverage}{656}
\showcasematerialdesignround{emoji-nature}{657}
\showcasematerialdesignround{emoji-objects}{658}
\showcasematerialdesignround{emoji-people}{659}
\showcasematerialdesignround{emoji-symbols}{660}
\showcasematerialdesignround{emoji-transportation}{661}
\showcasematerialdesignround{energy-savings-leaf}{662}
\showcasematerialdesignround{engineering}{663}
\showcasematerialdesignround{enhanced-encryption}{664}
\showcasematerialdesignround{equalizer}{665}
\showcasematerialdesignround{error}{666}
\showcasematerialdesignround{error-outline}{667}
\showcasematerialdesignround{escalator}{668}
\showcasematerialdesignround{escalator-warning}{669}
\showcasematerialdesignround{euro}{670}
\showcasematerialdesignround{euro-symbol}{671}
\showcasematerialdesignround{event}{672}
\showcasematerialdesignround{event-available}{673}
\showcasematerialdesignround{event-busy}{674}
\showcasematerialdesignround{event-note}{675}
\showcasematerialdesignround{event-repeat}{676}
\showcasematerialdesignround{event-seat}{677}
\showcasematerialdesignround{ev-station}{678}
\showcasematerialdesignround{exit-to-app}{679}
\showcasematerialdesignround{expand}{680}
\showcasematerialdesignround{expand-circle-down}{681}
\showcasematerialdesignround{expand-less}{682}
\showcasematerialdesignround{expand-more}{683}
\showcasematerialdesignround{explicit}{684}
\showcasematerialdesignround{explore}{685}
\showcasematerialdesignround{explore-off}{686}
\showcasematerialdesignround{exposure}{687}
\showcasematerialdesignround{exposure-neg-1}{688}
\showcasematerialdesignround{exposure-neg-2}{689}
\showcasematerialdesignround{exposure-plus-1}{690}
\showcasematerialdesignround{exposure-plus-2}{691}
\showcasematerialdesignround{exposure-zero}{692}
\showcasematerialdesignround{extension}{693}
\showcasematerialdesignround{extension-off}{694}
\showcasematerialdesignround{e-mobiledata}{695}
\showcasematerialdesignround{face}{696}
\showcasematerialdesignround{face-2}{697}
\showcasematerialdesignround{face-3}{698}
\showcasematerialdesignround{face-4}{699}
\showcasematerialdesignround{face-5}{700}
\showcasematerialdesignround{face-6}{701}
\showcasematerialdesignround{face-retouching-natural}{702}
\showcasematerialdesignround{face-retouching-off}{703}
\showcasematerialdesignround{factory}{704}
\showcasematerialdesignround{fact-check}{705}
\showcasematerialdesignround{family-restroom}{706}
\showcasematerialdesignround{fastfood}{707}
\showcasematerialdesignround{fast-forward}{708}
\showcasematerialdesignround{fast-rewind}{709}
\showcasematerialdesignround{favorite}{710}
\showcasematerialdesignround{favorite-border}{711}
\showcasematerialdesignround{fax}{712}
\showcasematerialdesignround{featured-play-list}{713}
\showcasematerialdesignround{featured-video}{714}
\showcasematerialdesignround{feed}{715}
\showcasematerialdesignround{feedback}{716}
\showcasematerialdesignround{female}{717}
\showcasematerialdesignround{fence}{718}
\showcasematerialdesignround{festival}{719}
\showcasematerialdesignround{fiber-dvr}{720}
\showcasematerialdesignround{fiber-manual-record}{721}
\showcasematerialdesignround{fiber-new}{722}
\showcasematerialdesignround{fiber-pin}{723}
\showcasematerialdesignround{fiber-smart-record}{724}
\showcasematerialdesignround{file-copy}{725}
\showcasematerialdesignround{file-download}{726}
\showcasematerialdesignround{file-download-done}{727}
\showcasematerialdesignround{file-download-off}{728}
\showcasematerialdesignround{file-open}{729}
\showcasematerialdesignround{file-present}{730}
\showcasematerialdesignround{file-upload}{731}
\showcasematerialdesignround{filter}{732}
\showcasematerialdesignround{filter-1}{733}
\showcasematerialdesignround{filter-2}{734}
\showcasematerialdesignround{filter-3}{735}
\showcasematerialdesignround{filter-4}{736}
\showcasematerialdesignround{filter-5}{737}
\showcasematerialdesignround{filter-6}{738}
\showcasematerialdesignround{filter-7}{739}
\showcasematerialdesignround{filter-8}{740}
\showcasematerialdesignround{filter-9}{741}
\showcasematerialdesignround{filter-9-plus}{742}
\showcasematerialdesignround{filter-alt}{743}
\showcasematerialdesignround{filter-alt-off}{744}
\showcasematerialdesignround{filter-b-and-w}{745}
\showcasematerialdesignround{filter-center-focus}{746}
\showcasematerialdesignround{filter-drama}{747}
\showcasematerialdesignround{filter-frames}{748}
\showcasematerialdesignround{filter-hdr}{749}
\showcasematerialdesignround{filter-list}{750}
\showcasematerialdesignround{filter-list-off}{751}
\showcasematerialdesignround{filter-none}{752}
\showcasematerialdesignround{filter-tilt-shift}{753}
\showcasematerialdesignround{filter-vintage}{754}
\showcasematerialdesignround{find-in-page}{755}
\showcasematerialdesignround{find-replace}{756}
\showcasematerialdesignround{fingerprint}{757}
\showcasematerialdesignround{fireplace}{758}
\showcasematerialdesignround{fire-extinguisher}{759}
\showcasematerialdesignround{fire-hydrant-alt}{760}
\showcasematerialdesignround{fire-truck}{761}
\showcasematerialdesignround{first-page}{762}
\showcasematerialdesignround{fitbit}{763}
\showcasematerialdesignround{fitness-center}{764}
\showcasematerialdesignround{fit-screen}{765}
\showcasematerialdesignround{flag}{766}
\showcasematerialdesignround{flag-circle}{767}
\showcasematerialdesignround{flaky}{768}
\showcasematerialdesignround{flare}{769}
\showcasematerialdesignround{flashlight-off}{770}
\showcasematerialdesignround{flashlight-on}{771}
\showcasematerialdesignround{flash-auto}{772}
\showcasematerialdesignround{flash-off}{773}
\showcasematerialdesignround{flash-on}{774}
\showcasematerialdesignround{flatware}{775}
\showcasematerialdesignround{flight}{776}
\showcasematerialdesignround{flight-class}{777}
\showcasematerialdesignround{flight-land}{778}
\showcasematerialdesignround{flight-takeoff}{779}
\showcasematerialdesignround{flip}{780}
\showcasematerialdesignround{flip-camera-android}{781}
\showcasematerialdesignround{flip-camera-ios}{782}
\showcasematerialdesignround{flip-to-back}{783}
\showcasematerialdesignround{flip-to-front}{784}
\showcasematerialdesignround{flood}{785}
\showcasematerialdesignround{fluorescent}{786}
\showcasematerialdesignround{flutter-dash}{787}
\showcasematerialdesignround{fmd-bad}{788}
\showcasematerialdesignround{fmd-good}{789}
\showcasematerialdesignround{folder}{790}
\showcasematerialdesignround{folder-copy}{791}
\showcasematerialdesignround{folder-delete}{792}
\showcasematerialdesignround{folder-off}{793}
\showcasematerialdesignround{folder-open}{794}
\showcasematerialdesignround{folder-shared}{795}
\showcasematerialdesignround{folder-special}{796}
\showcasematerialdesignround{folder-zip}{797}
\showcasematerialdesignround{follow-the-signs}{798}
\showcasematerialdesignround{font-download}{799}
\showcasematerialdesignround{font-download-off}{800}
\showcasematerialdesignround{food-bank}{801}
\showcasematerialdesignround{forest}{802}
\showcasematerialdesignround{fork-left}{803}
\showcasematerialdesignround{fork-right}{804}
\showcasematerialdesignround{format-align-center}{805}
\showcasematerialdesignround{format-align-justify}{806}
\showcasematerialdesignround{format-align-left}{807}
\showcasematerialdesignround{format-align-right}{808}
\showcasematerialdesignround{format-bold}{809}
\showcasematerialdesignround{format-clear}{810}
\showcasematerialdesignround{format-color-fill}{811}
\showcasematerialdesignround{format-color-reset}{812}
\showcasematerialdesignround{format-color-text}{813}
\showcasematerialdesignround{format-indent-decrease}{814}
\showcasematerialdesignround{format-indent-increase}{815}
\showcasematerialdesignround{format-italic}{816}
\showcasematerialdesignround{format-line-spacing}{817}
\showcasematerialdesignround{format-list-bulleted}{818}
\showcasematerialdesignround{format-list-numbered}{819}
\showcasematerialdesignround{format-list-numbered-rtl}{820}
\showcasematerialdesignround{format-overline}{821}
\showcasematerialdesignround{format-paint}{822}
\showcasematerialdesignround{format-quote}{823}
\showcasematerialdesignround{format-shapes}{824}
\showcasematerialdesignround{format-size}{825}
\showcasematerialdesignround{format-strikethrough}{826}
\showcasematerialdesignround{format-textdirection-l-to-r}{827}
\showcasematerialdesignround{format-textdirection-r-to-l}{828}
\showcasematerialdesignround{format-underlined}{829}
\showcasematerialdesignround{fort}{830}
\showcasematerialdesignround{forum}{831}
\showcasematerialdesignround{forward}{832}
\showcasematerialdesignround{forward-5}{833}
\showcasematerialdesignround{forward-10}{834}
\showcasematerialdesignround{forward-30}{835}
\showcasematerialdesignround{forward-to-inbox}{836}
\showcasematerialdesignround{foundation}{837}
\showcasematerialdesignround{free-breakfast}{838}
\showcasematerialdesignround{free-cancellation}{839}
\showcasematerialdesignround{front-hand}{840}
\showcasematerialdesignround{fullscreen}{841}
\showcasematerialdesignround{fullscreen-exit}{842}
\showcasematerialdesignround{functions}{843}
\showcasematerialdesignround{gamepad}{844}
\showcasematerialdesignround{games}{845}
\showcasematerialdesignround{garage}{846}
\showcasematerialdesignround{gas-meter}{847}
\showcasematerialdesignround{gavel}{848}
\showcasematerialdesignround{generating-tokens}{849}
\showcasematerialdesignround{gesture}{850}
\showcasematerialdesignround{get-app}{851}
\showcasematerialdesignround{gif}{852}
\showcasematerialdesignround{gif-box}{853}
\showcasematerialdesignround{girl}{854}
\showcasematerialdesignround{gite}{855}
\showcasematerialdesignround{golf-course}{856}
\showcasematerialdesignround{gpp-bad}{857}
\showcasematerialdesignround{gpp-good}{858}
\showcasematerialdesignround{gpp-maybe}{859}
\showcasematerialdesignround{gps-fixed}{860}
\showcasematerialdesignround{gps-not-fixed}{861}
\showcasematerialdesignround{gps-off}{862}
\showcasematerialdesignround{grade}{863}
\showcasematerialdesignround{gradient}{864}
\showcasematerialdesignround{grading}{865}
\showcasematerialdesignround{grain}{866}
\showcasematerialdesignround{graphic-eq}{867}
\showcasematerialdesignround{grass}{868}
\showcasematerialdesignround{grid-3x3}{869}
\showcasematerialdesignround{grid-4x4}{870}
\showcasematerialdesignround{grid-goldenratio}{871}
\showcasematerialdesignround{grid-off}{872}
\showcasematerialdesignround{grid-on}{873}
\showcasematerialdesignround{grid-view}{874}
\showcasematerialdesignround{group}{875}
\showcasematerialdesignround{groups}{876}
\showcasematerialdesignround{groups-2}{877}
\showcasematerialdesignround{groups-3}{878}
\showcasematerialdesignround{group-add}{879}
\showcasematerialdesignround{group-off}{880}
\showcasematerialdesignround{group-remove}{881}
\showcasematerialdesignround{group-work}{882}
\showcasematerialdesignround{g-mobiledata}{883}
\showcasematerialdesignround{g-translate}{884}
\showcasematerialdesignround{hail}{885}
\showcasematerialdesignround{handshake}{886}
\showcasematerialdesignround{handyman}{887}
\showcasematerialdesignround{hardware}{888}
\showcasematerialdesignround{hd}{889}
\showcasematerialdesignround{hdr-auto}{890}
\showcasematerialdesignround{hdr-auto-select}{891}
\showcasematerialdesignround{hdr-enhanced-select}{892}
\showcasematerialdesignround{hdr-off}{893}
\showcasematerialdesignround{hdr-off-select}{894}
\showcasematerialdesignround{hdr-on}{895}
\showcasematerialdesignround{hdr-on-select}{896}
\showcasematerialdesignround{hdr-plus}{897}
\showcasematerialdesignround{hdr-strong}{898}
\showcasematerialdesignround{hdr-weak}{899}
\showcasematerialdesignround{headphones}{900}
\showcasematerialdesignround{headphones-battery}{901}
\showcasematerialdesignround{headset}{902}
\showcasematerialdesignround{headset-mic}{903}
\showcasematerialdesignround{headset-off}{904}
\showcasematerialdesignround{healing}{905}
\showcasematerialdesignround{health-and-safety}{906}
\showcasematerialdesignround{hearing}{907}
\showcasematerialdesignround{hearing-disabled}{908}
\showcasematerialdesignround{heart-broken}{909}
\showcasematerialdesignround{heat-pump}{910}
\showcasematerialdesignround{height}{911}
\showcasematerialdesignround{help}{912}
\showcasematerialdesignround{help-center}{913}
\showcasematerialdesignround{help-outline}{914}
\showcasematerialdesignround{hevc}{915}
\showcasematerialdesignround{hexagon}{916}
\showcasematerialdesignround{hide-image}{917}
\showcasematerialdesignround{hide-source}{918}
\showcasematerialdesignround{highlight}{919}
\showcasematerialdesignround{highlight-alt}{920}
\showcasematerialdesignround{highlight-off}{921}
\showcasematerialdesignround{high-quality}{922}
\showcasematerialdesignround{hiking}{923}
\showcasematerialdesignround{history}{924}
\showcasematerialdesignround{history-edu}{925}
\showcasematerialdesignround{history-toggle-off}{926}
\showcasematerialdesignround{hive}{927}
\showcasematerialdesignround{hls}{928}
\showcasematerialdesignround{hls-off}{929}
\showcasematerialdesignround{holiday-village}{930}
\showcasematerialdesignround{home}{931}
\showcasematerialdesignround{home-max}{932}
\showcasematerialdesignround{home-mini}{933}
\showcasematerialdesignround{home-repair-service}{934}
\showcasematerialdesignround{home-work}{935}
\showcasematerialdesignround{horizontal-distribute}{936}
\showcasematerialdesignround{horizontal-rule}{937}
\showcasematerialdesignround{horizontal-split}{938}
\showcasematerialdesignround{hotel}{939}
\showcasematerialdesignround{hotel-class}{940}
\showcasematerialdesignround{hot-tub}{941}
\showcasematerialdesignround{hourglass-bottom}{942}
\showcasematerialdesignround{hourglass-disabled}{943}
\showcasematerialdesignround{hourglass-empty}{944}
\showcasematerialdesignround{hourglass-full}{945}
\showcasematerialdesignround{hourglass-top}{946}
\showcasematerialdesignround{house}{947}
\showcasematerialdesignround{houseboat}{948}
\showcasematerialdesignround{house-siding}{949}
\showcasematerialdesignround{how-to-reg}{950}
\showcasematerialdesignround{how-to-vote}{951}
\showcasematerialdesignround{html}{952}
\showcasematerialdesignround{http}{953}
\showcasematerialdesignround{https}{954}
\showcasematerialdesignround{hub}{955}
\showcasematerialdesignround{hvac}{956}
\showcasematerialdesignround{h-mobiledata}{957}
\showcasematerialdesignround{h-plus-mobiledata}{958}
\showcasematerialdesignround{icecream}{959}
\showcasematerialdesignround{ice-skating}{960}
\showcasematerialdesignround{image}{961}
\showcasematerialdesignround{imagesearch-roller}{962}
\showcasematerialdesignround{image-aspect-ratio}{963}
\showcasematerialdesignround{image-not-supported}{964}
\showcasematerialdesignround{image-search}{965}
\showcasematerialdesignround{important-devices}{966}
\showcasematerialdesignround{import-contacts}{967}
\showcasematerialdesignround{import-export}{968}
\showcasematerialdesignround{inbox}{969}
\showcasematerialdesignround{incomplete-circle}{970}
\showcasematerialdesignround{indeterminate-check-box}{971}
\showcasematerialdesignround{info}{972}
\showcasematerialdesignround{input}{973}
\showcasematerialdesignround{insert-chart}{974}
\showcasematerialdesignround{insert-chart-outlined}{975}
\showcasematerialdesignround{insert-comment}{976}
\showcasematerialdesignround{insert-drive-file}{977}
\showcasematerialdesignround{insert-emoticon}{978}
\showcasematerialdesignround{insert-invitation}{979}
\showcasematerialdesignround{insert-link}{980}
\showcasematerialdesignround{insert-page-break}{981}
\showcasematerialdesignround{insert-photo}{982}
\showcasematerialdesignround{insights}{983}
\showcasematerialdesignround{install-desktop}{984}
\showcasematerialdesignround{install-mobile}{985}
\showcasematerialdesignround{integration-instructions}{986}
\showcasematerialdesignround{interests}{987}
\showcasematerialdesignround{interpreter-mode}{988}
\showcasematerialdesignround{inventory}{989}
\showcasematerialdesignround{inventory-2}{990}
\showcasematerialdesignround{invert-colors}{991}
\showcasematerialdesignround{invert-colors-off}{992}
\showcasematerialdesignround{ios-share}{993}
\showcasematerialdesignround{iron}{994}
\showcasematerialdesignround{iso}{995}
\showcasematerialdesignround{javascript}{996}
\showcasematerialdesignround{join-full}{997}
\showcasematerialdesignround{join-inner}{998}
\showcasematerialdesignround{join-left}{999}
\showcasematerialdesignround{join-right}{1000}
\showcasematerialdesignround{kayaking}{1001}
\showcasematerialdesignround{kebab-dining}{1002}
\showcasematerialdesignround{key}{1003}
\showcasematerialdesignround{keyboard}{1004}
\showcasematerialdesignround{keyboard-alt}{1005}
\showcasematerialdesignround{keyboard-arrow-down}{1006}
\showcasematerialdesignround{keyboard-arrow-left}{1007}
\showcasematerialdesignround{keyboard-arrow-right}{1008}
\showcasematerialdesignround{keyboard-arrow-up}{1009}
\showcasematerialdesignround{keyboard-backspace}{1010}
\showcasematerialdesignround{keyboard-capslock}{1011}
\showcasematerialdesignround{keyboard-command-key}{1012}
\showcasematerialdesignround{keyboard-control-key}{1013}
\showcasematerialdesignround{keyboard-double-arrow-down}{1014}
\showcasematerialdesignround{keyboard-double-arrow-left}{1015}
\showcasematerialdesignround{keyboard-double-arrow-right}{1016}
\showcasematerialdesignround{keyboard-double-arrow-up}{1017}
\showcasematerialdesignround{keyboard-hide}{1018}
\showcasematerialdesignround{keyboard-option-key}{1019}
\showcasematerialdesignround{keyboard-return}{1020}
\showcasematerialdesignround{keyboard-tab}{1021}
\showcasematerialdesignround{keyboard-voice}{1022}
\showcasematerialdesignround{key-off}{1023}
\showcasematerialdesignround{king-bed}{1024}
\showcasematerialdesignround{kitchen}{1025}
\showcasematerialdesignround{kitesurfing}{1026}
\showcasematerialdesignround{label}{1027}
\showcasematerialdesignround{label-important}{1028}
\showcasematerialdesignround{label-off}{1029}
\showcasematerialdesignround{lan}{1030}
\showcasematerialdesignround{landscape}{1031}
\showcasematerialdesignround{landslide}{1032}
\showcasematerialdesignround{language}{1033}
\showcasematerialdesignround{laptop}{1034}
\showcasematerialdesignround{laptop-chromebook}{1035}
\showcasematerialdesignround{laptop-mac}{1036}
\showcasematerialdesignround{laptop-windows}{1037}
\showcasematerialdesignround{last-page}{1038}
\showcasematerialdesignround{launch}{1039}
\showcasematerialdesignround{layers}{1040}
\showcasematerialdesignround{layers-clear}{1041}
\showcasematerialdesignround{leaderboard}{1042}
\showcasematerialdesignround{leak-add}{1043}
\showcasematerialdesignround{leak-remove}{1044}
\showcasematerialdesignround{legend-toggle}{1045}
\showcasematerialdesignround{lens}{1046}
\showcasematerialdesignround{lens-blur}{1047}
\showcasematerialdesignround{library-add}{1048}
\showcasematerialdesignround{library-add-check}{1049}
\showcasematerialdesignround{library-books}{1050}
\showcasematerialdesignround{library-music}{1051}
\showcasematerialdesignround{light}{1052}
\showcasematerialdesignround{lightbulb}{1053}
\showcasematerialdesignround{lightbulb-circle}{1054}
\showcasematerialdesignround{light-mode}{1055}
\showcasematerialdesignround{linear-scale}{1056}
\showcasematerialdesignround{line-axis}{1057}
\showcasematerialdesignround{line-style}{1058}
\showcasematerialdesignround{line-weight}{1059}
\showcasematerialdesignround{link}{1060}
\showcasematerialdesignround{linked-camera}{1061}
\showcasematerialdesignround{link-off}{1062}
\showcasematerialdesignround{liquor}{1063}
\showcasematerialdesignround{list}{1064}
\showcasematerialdesignround{list-alt}{1065}
\showcasematerialdesignround{live-help}{1066}
\showcasematerialdesignround{live-tv}{1067}
\showcasematerialdesignround{living}{1068}
\showcasematerialdesignround{local-activity}{1069}
\showcasematerialdesignround{local-airport}{1070}
\showcasematerialdesignround{local-atm}{1071}
\showcasematerialdesignround{local-bar}{1072}
\showcasematerialdesignround{local-cafe}{1073}
\showcasematerialdesignround{local-car-wash}{1074}
\showcasematerialdesignround{local-convenience-store}{1075}
\showcasematerialdesignround{local-dining}{1076}
\showcasematerialdesignround{local-drink}{1077}
\showcasematerialdesignround{local-fire-department}{1078}
\showcasematerialdesignround{local-florist}{1079}
\showcasematerialdesignround{local-gas-station}{1080}
\showcasematerialdesignround{local-grocery-store}{1081}
\showcasematerialdesignround{local-hospital}{1082}
\showcasematerialdesignround{local-hotel}{1083}
\showcasematerialdesignround{local-laundry-service}{1084}
\showcasematerialdesignround{local-library}{1085}
\showcasematerialdesignround{local-mall}{1086}
\showcasematerialdesignround{local-movies}{1087}
\showcasematerialdesignround{local-offer}{1088}
\showcasematerialdesignround{local-parking}{1089}
\showcasematerialdesignround{local-pharmacy}{1090}
\showcasematerialdesignround{local-phone}{1091}
\showcasematerialdesignround{local-pizza}{1092}
\showcasematerialdesignround{local-play}{1093}
\showcasematerialdesignround{local-police}{1094}
\showcasematerialdesignround{local-post-office}{1095}
\showcasematerialdesignround{local-printshop}{1096}
\showcasematerialdesignround{local-see}{1097}
\showcasematerialdesignround{local-shipping}{1098}
\showcasematerialdesignround{local-taxi}{1099}
\showcasematerialdesignround{location-city}{1100}
\showcasematerialdesignround{location-disabled}{1101}
\showcasematerialdesignround{location-off}{1102}
\showcasematerialdesignround{location-on}{1103}
\showcasematerialdesignround{location-searching}{1104}
\showcasematerialdesignround{lock}{1105}
\showcasematerialdesignround{lock-clock}{1106}
\showcasematerialdesignround{lock-open}{1107}
\showcasematerialdesignround{lock-person}{1108}
\showcasematerialdesignround{lock-reset}{1109}
\showcasematerialdesignround{login}{1110}
\showcasematerialdesignround{logout}{1111}
\showcasematerialdesignround{logo-dev}{1112}
\showcasematerialdesignround{looks}{1113}
\showcasematerialdesignround{looks-3}{1114}
\showcasematerialdesignround{looks-4}{1115}
\showcasematerialdesignround{looks-5}{1116}
\showcasematerialdesignround{looks-6}{1117}
\showcasematerialdesignround{looks-one}{1118}
\showcasematerialdesignround{looks-two}{1119}
\showcasematerialdesignround{loop}{1120}
\showcasematerialdesignround{loupe}{1121}
\showcasematerialdesignround{low-priority}{1122}
\showcasematerialdesignround{loyalty}{1123}
\showcasematerialdesignround{lte-mobiledata}{1124}
\showcasematerialdesignround{lte-plus-mobiledata}{1125}
\showcasematerialdesignround{luggage}{1126}
\showcasematerialdesignround{lunch-dining}{1127}
\showcasematerialdesignround{lyrics}{1128}
\showcasematerialdesignround{macro-off}{1129}
\showcasematerialdesignround{mail}{1130}
\showcasematerialdesignround{mail-lock}{1131}
\showcasematerialdesignround{mail-outline}{1132}
\showcasematerialdesignround{male}{1133}
\showcasematerialdesignround{man}{1134}
\showcasematerialdesignround{manage-accounts}{1135}
\showcasematerialdesignround{manage-history}{1136}
\showcasematerialdesignround{manage-search}{1137}
\showcasematerialdesignround{man-2}{1138}
\showcasematerialdesignround{man-3}{1139}
\showcasematerialdesignround{man-4}{1140}
\showcasematerialdesignround{map}{1141}
\showcasematerialdesignround{maps-home-work}{1142}
\showcasematerialdesignround{maps-ugc}{1143}
\showcasematerialdesignround{margin}{1144}
\showcasematerialdesignround{markunread}{1145}
\showcasematerialdesignround{markunread-mailbox}{1146}
\showcasematerialdesignround{mark-as-unread}{1147}
\showcasematerialdesignround{mark-chat-read}{1148}
\showcasematerialdesignround{mark-chat-unread}{1149}
\showcasematerialdesignround{mark-email-read}{1150}
\showcasematerialdesignround{mark-email-unread}{1151}
\showcasematerialdesignround{mark-unread-chat-alt}{1152}
\showcasematerialdesignround{masks}{1153}
\showcasematerialdesignround{maximize}{1154}
\showcasematerialdesignround{mediation}{1155}
\showcasematerialdesignround{media-bluetooth-off}{1156}
\showcasematerialdesignround{media-bluetooth-on}{1157}
\showcasematerialdesignround{medical-information}{1158}
\showcasematerialdesignround{medical-services}{1159}
\showcasematerialdesignround{medication}{1160}
\showcasematerialdesignround{medication-liquid}{1161}
\showcasematerialdesignround{meeting-room}{1162}
\showcasematerialdesignround{memory}{1163}
\showcasematerialdesignround{menu}{1164}
\showcasematerialdesignround{menu-book}{1165}
\showcasematerialdesignround{menu-open}{1166}
\showcasematerialdesignround{merge}{1167}
\showcasematerialdesignround{merge-type}{1168}
\showcasematerialdesignround{message}{1169}
\showcasematerialdesignround{mic}{1170}
\showcasematerialdesignround{microwave}{1171}
\showcasematerialdesignround{mic-external-off}{1172}
\showcasematerialdesignround{mic-external-on}{1173}
\showcasematerialdesignround{mic-none}{1174}
\showcasematerialdesignround{mic-off}{1175}
\showcasematerialdesignround{military-tech}{1176}
\showcasematerialdesignround{minimize}{1177}
\showcasematerialdesignround{minor-crash}{1178}
\showcasematerialdesignround{miscellaneous-services}{1179}
\showcasematerialdesignround{missed-video-call}{1180}
\showcasematerialdesignround{mms}{1181}
\showcasematerialdesignround{mobiledata-off}{1182}
\showcasematerialdesignround{mobile-friendly}{1183}
\showcasematerialdesignround{mobile-off}{1184}
\showcasematerialdesignround{mobile-screen-share}{1185}
\showcasematerialdesignround{mode}{1186}
\showcasematerialdesignround{model-training}{1187}
\showcasematerialdesignround{mode-comment}{1188}
\showcasematerialdesignround{mode-edit}{1189}
\showcasematerialdesignround{mode-edit-outline}{1190}
\showcasematerialdesignround{mode-fan-off}{1191}
\showcasematerialdesignround{mode-night}{1192}
\showcasematerialdesignround{mode-of-travel}{1193}
\showcasematerialdesignround{mode-standby}{1194}
\showcasematerialdesignround{monetization-on}{1195}
\showcasematerialdesignround{money}{1196}
\showcasematerialdesignround{money-off}{1197}
\showcasematerialdesignround{money-off-csred}{1198}
\showcasematerialdesignround{monitor}{1199}
\showcasematerialdesignround{monitor-heart}{1200}
\showcasematerialdesignround{monitor-weight}{1201}
\showcasematerialdesignround{monochrome-photos}{1202}
\showcasematerialdesignround{mood}{1203}
\showcasematerialdesignround{mood-bad}{1204}
\showcasematerialdesignround{moped}{1205}
\showcasematerialdesignround{more}{1206}
\showcasematerialdesignround{more-horiz}{1207}
\showcasematerialdesignround{more-time}{1208}
\showcasematerialdesignround{more-vert}{1209}
\showcasematerialdesignround{mosque}{1210}
\showcasematerialdesignround{motion-photos-auto}{1211}
\showcasematerialdesignround{motion-photos-off}{1212}
\showcasematerialdesignround{motion-photos-on}{1213}
\showcasematerialdesignround{motion-photos-pause}{1214}
\showcasematerialdesignround{motion-photos-paused}{1215}
\showcasematerialdesignround{mouse}{1216}
\showcasematerialdesignround{move-down}{1217}
\showcasematerialdesignround{move-to-inbox}{1218}
\showcasematerialdesignround{move-up}{1219}
\showcasematerialdesignround{movie}{1220}
\showcasematerialdesignround{movie-creation}{1221}
\showcasematerialdesignround{movie-filter}{1222}
\showcasematerialdesignround{moving}{1223}
\showcasematerialdesignround{mp}{1224}
\showcasematerialdesignround{multiline-chart}{1225}
\showcasematerialdesignround{multiple-stop}{1226}
\showcasematerialdesignround{museum}{1227}
\showcasematerialdesignround{music-note}{1228}
\showcasematerialdesignround{music-off}{1229}
\showcasematerialdesignround{music-video}{1230}
\showcasematerialdesignround{my-location}{1231}
\showcasematerialdesignround{nat}{1232}
\showcasematerialdesignround{nature}{1233}
\showcasematerialdesignround{nature-people}{1234}
\showcasematerialdesignround{navigate-before}{1235}
\showcasematerialdesignround{navigate-next}{1236}
\showcasematerialdesignround{navigation}{1237}
\showcasematerialdesignround{nearby-error}{1238}
\showcasematerialdesignround{nearby-off}{1239}
\showcasematerialdesignround{near-me}{1240}
\showcasematerialdesignround{near-me-disabled}{1241}
\showcasematerialdesignround{nest-cam-wired-stand}{1242}
\showcasematerialdesignround{network-cell}{1243}
\showcasematerialdesignround{network-check}{1244}
\showcasematerialdesignround{network-locked}{1245}
\showcasematerialdesignround{network-ping}{1246}
\showcasematerialdesignround{network-wifi}{1247}
\showcasematerialdesignround{network-wifi-1-bar}{1248}
\showcasematerialdesignround{network-wifi-2-bar}{1249}
\showcasematerialdesignround{network-wifi-3-bar}{1250}
\showcasematerialdesignround{newspaper}{1251}
\showcasematerialdesignround{new-label}{1252}
\showcasematerialdesignround{new-releases}{1253}
\showcasematerialdesignround{next-plan}{1254}
\showcasematerialdesignround{next-week}{1255}
\showcasematerialdesignround{nfc}{1256}
\showcasematerialdesignround{nightlife}{1257}
\showcasematerialdesignround{nightlight}{1258}
\showcasematerialdesignround{nightlight-round}{1259}
\showcasematerialdesignround{nights-stay}{1260}
\showcasematerialdesignround{night-shelter}{1261}
\showcasematerialdesignround{noise-aware}{1262}
\showcasematerialdesignround{noise-control-off}{1263}
\showcasematerialdesignround{nordic-walking}{1264}
\showcasematerialdesignround{north}{1265}
\showcasematerialdesignround{north-east}{1266}
\showcasematerialdesignround{north-west}{1267}
\showcasematerialdesignround{note}{1268}
\showcasematerialdesignround{notes}{1269}
\showcasematerialdesignround{note-add}{1270}
\showcasematerialdesignround{note-alt}{1271}
\showcasematerialdesignround{notifications}{1272}
\showcasematerialdesignround{notifications-active}{1273}
\showcasematerialdesignround{notifications-none}{1274}
\showcasematerialdesignround{notifications-off}{1275}
\showcasematerialdesignround{notifications-paused}{1276}
\showcasematerialdesignround{notification-add}{1277}
\showcasematerialdesignround{notification-important}{1278}
\showcasematerialdesignround{not-accessible}{1279}
\showcasematerialdesignround{not-interested}{1280}
\showcasematerialdesignround{not-listed-location}{1281}
\showcasematerialdesignround{not-started}{1282}
\showcasematerialdesignround{no-accounts}{1283}
\showcasematerialdesignround{no-adult-content}{1284}
\showcasematerialdesignround{no-backpack}{1285}
\showcasematerialdesignround{no-cell}{1286}
\showcasematerialdesignround{no-crash}{1287}
\showcasematerialdesignround{no-drinks}{1288}
\showcasematerialdesignround{no-encryption}{1289}
\showcasematerialdesignround{no-encryption-gmailerrorred}{1290}
\showcasematerialdesignround{no-flash}{1291}
\showcasematerialdesignround{no-food}{1292}
\showcasematerialdesignround{no-luggage}{1293}
\showcasematerialdesignround{no-meals}{1294}
\showcasematerialdesignround{no-meeting-room}{1295}
\showcasematerialdesignround{no-photography}{1296}
\showcasematerialdesignround{no-sim}{1297}
\showcasematerialdesignround{no-stroller}{1298}
\showcasematerialdesignround{no-transfer}{1299}
\showcasematerialdesignround{numbers}{1300}
\showcasematerialdesignround{offline-bolt}{1301}
\showcasematerialdesignround{offline-pin}{1302}
\showcasematerialdesignround{offline-share}{1303}
\showcasematerialdesignround{oil-barrel}{1304}
\showcasematerialdesignround{ondemand-video}{1305}
\showcasematerialdesignround{online-prediction}{1306}
\showcasematerialdesignround{on-device-training}{1307}
\showcasematerialdesignround{opacity}{1308}
\showcasematerialdesignround{open-in-browser}{1309}
\showcasematerialdesignround{open-in-full}{1310}
\showcasematerialdesignround{open-in-new}{1311}
\showcasematerialdesignround{open-in-new-off}{1312}
\showcasematerialdesignround{open-with}{1313}
\showcasematerialdesignround{other-houses}{1314}
\showcasematerialdesignround{outbound}{1315}
\showcasematerialdesignround{outbox}{1316}
\showcasematerialdesignround{outdoor-grill}{1317}
\showcasematerialdesignround{outlet}{1318}
\showcasematerialdesignround{outlined-flag}{1319}
\showcasematerialdesignround{output}{1320}
\showcasematerialdesignround{padding}{1321}
\showcasematerialdesignround{pages}{1322}
\showcasematerialdesignround{pageview}{1323}
\showcasematerialdesignround{paid}{1324}
\showcasematerialdesignround{palette}{1325}
\showcasematerialdesignround{panorama}{1326}
\showcasematerialdesignround{panorama-fish-eye}{1327}
\showcasematerialdesignround{panorama-horizontal}{1328}
\showcasematerialdesignround{panorama-horizontal-select}{1329}
\showcasematerialdesignround{panorama-photosphere}{1330}
\showcasematerialdesignround{panorama-photosphere-select}{1331}
\showcasematerialdesignround{panorama-vertical}{1332}
\showcasematerialdesignround{panorama-vertical-select}{1333}
\showcasematerialdesignround{panorama-wide-angle}{1334}
\showcasematerialdesignround{panorama-wide-angle-select}{1335}
\showcasematerialdesignround{pan-tool}{1336}
\showcasematerialdesignround{pan-tool-alt}{1337}
\showcasematerialdesignround{paragliding}{1338}
\showcasematerialdesignround{park}{1339}
\showcasematerialdesignround{party-mode}{1340}
\showcasematerialdesignround{password}{1341}
\showcasematerialdesignround{pattern}{1342}
\showcasematerialdesignround{pause}{1343}
\showcasematerialdesignround{pause-circle}{1344}
\showcasematerialdesignround{pause-circle-filled}{1345}
\showcasematerialdesignround{pause-circle-outline}{1346}
\showcasematerialdesignround{pause-presentation}{1347}
\showcasematerialdesignround{payment}{1348}
\showcasematerialdesignround{payments}{1349}
\showcasematerialdesignround{pedal-bike}{1350}
\showcasematerialdesignround{pending}{1351}
\showcasematerialdesignround{pending-actions}{1352}
\showcasematerialdesignround{pentagon}{1353}
\showcasematerialdesignround{people}{1354}
\showcasematerialdesignround{people-alt}{1355}
\showcasematerialdesignround{people-outline}{1356}
\showcasematerialdesignround{percent}{1357}
\showcasematerialdesignround{perm-camera-mic}{1358}
\showcasematerialdesignround{perm-contact-calendar}{1359}
\showcasematerialdesignround{perm-data-setting}{1360}
\showcasematerialdesignround{perm-device-information}{1361}
\showcasematerialdesignround{perm-identity}{1362}
\showcasematerialdesignround{perm-media}{1363}
\showcasematerialdesignround{perm-phone-msg}{1364}
\showcasematerialdesignround{perm-scan-wifi}{1365}
\showcasematerialdesignround{person}{1366}
\showcasematerialdesignround{personal-injury}{1367}
\showcasematerialdesignround{personal-video}{1368}
\showcasematerialdesignround{person-2}{1369}
\showcasematerialdesignround{person-3}{1370}
\showcasematerialdesignround{person-4}{1371}
\showcasematerialdesignround{person-add}{1372}
\showcasematerialdesignround{person-add-alt}{1373}
\showcasematerialdesignround{person-add-alt-1}{1374}
\showcasematerialdesignround{person-add-disabled}{1375}
\showcasematerialdesignround{person-off}{1376}
\showcasematerialdesignround{person-outline}{1377}
\showcasematerialdesignround{person-pin}{1378}
\showcasematerialdesignround{person-pin-circle}{1379}
\showcasematerialdesignround{person-remove}{1380}
\showcasematerialdesignround{person-remove-alt-1}{1381}
\showcasematerialdesignround{person-search}{1382}
\showcasematerialdesignround{pest-control}{1383}
\showcasematerialdesignround{pest-control-rodent}{1384}
\showcasematerialdesignround{pets}{1385}
\showcasematerialdesignround{phishing}{1386}
\showcasematerialdesignround{phone}{1387}
\showcasematerialdesignround{phonelink}{1388}
\showcasematerialdesignround{phonelink-erase}{1389}
\showcasematerialdesignround{phonelink-lock}{1390}
\showcasematerialdesignround{phonelink-off}{1391}
\showcasematerialdesignround{phonelink-ring}{1392}
\showcasematerialdesignround{phonelink-setup}{1393}
\showcasematerialdesignround{phone-android}{1394}
\showcasematerialdesignround{phone-bluetooth-speaker}{1395}
\showcasematerialdesignround{phone-callback}{1396}
\showcasematerialdesignround{phone-disabled}{1397}
\showcasematerialdesignround{phone-enabled}{1398}
\showcasematerialdesignround{phone-forwarded}{1399}
\showcasematerialdesignround{phone-iphone}{1400}
\showcasematerialdesignround{phone-locked}{1401}
\showcasematerialdesignround{phone-missed}{1402}
\showcasematerialdesignround{phone-paused}{1403}
\showcasematerialdesignround{photo}{1404}
\showcasematerialdesignround{photo-album}{1405}
\showcasematerialdesignround{photo-camera}{1406}
\showcasematerialdesignround{photo-camera-back}{1407}
\showcasematerialdesignround{photo-camera-front}{1408}
\showcasematerialdesignround{photo-filter}{1409}
\showcasematerialdesignround{photo-library}{1410}
\showcasematerialdesignround{photo-size-select-actual}{1411}
\showcasematerialdesignround{photo-size-select-large}{1412}
\showcasematerialdesignround{photo-size-select-small}{1413}
\showcasematerialdesignround{php}{1414}
\showcasematerialdesignround{piano}{1415}
\showcasematerialdesignround{piano-off}{1416}
\showcasematerialdesignround{picture-as-pdf}{1417}
\showcasematerialdesignround{picture-in-picture}{1418}
\showcasematerialdesignround{picture-in-picture-alt}{1419}
\showcasematerialdesignround{pie-chart}{1420}
\showcasematerialdesignround{pie-chart-outline}{1421}
\showcasematerialdesignround{pin}{1422}
\showcasematerialdesignround{pinch}{1423}
\showcasematerialdesignround{pin-drop}{1424}
\showcasematerialdesignround{pin-end}{1425}
\showcasematerialdesignround{pin-invoke}{1426}
\showcasematerialdesignround{pivot-table-chart}{1427}
\showcasematerialdesignround{pix}{1428}
\showcasematerialdesignround{place}{1429}
\showcasematerialdesignround{plagiarism}{1430}
\showcasematerialdesignround{playlist-add}{1431}
\showcasematerialdesignround{playlist-add-check}{1432}
\showcasematerialdesignround{playlist-add-check-circle}{1433}
\showcasematerialdesignround{playlist-add-circle}{1434}
\showcasematerialdesignround{playlist-play}{1435}
\showcasematerialdesignround{playlist-remove}{1436}
\showcasematerialdesignround{play-arrow}{1437}
\showcasematerialdesignround{play-circle}{1438}
\showcasematerialdesignround{play-circle-filled}{1439}
\showcasematerialdesignround{play-circle-outline}{1440}
\showcasematerialdesignround{play-disabled}{1441}
\showcasematerialdesignround{play-for-work}{1442}
\showcasematerialdesignround{play-lesson}{1443}
\showcasematerialdesignround{plumbing}{1444}
\showcasematerialdesignround{plus-one}{1445}
\showcasematerialdesignround{podcasts}{1446}
\showcasematerialdesignround{point-of-sale}{1447}
\showcasematerialdesignround{policy}{1448}
\showcasematerialdesignround{poll}{1449}
\showcasematerialdesignround{polyline}{1450}
\showcasematerialdesignround{polymer}{1451}
\showcasematerialdesignround{pool}{1452}
\showcasematerialdesignround{portable-wifi-off}{1453}
\showcasematerialdesignround{portrait}{1454}
\showcasematerialdesignround{post-add}{1455}
\showcasematerialdesignround{power}{1456}
\showcasematerialdesignround{power-input}{1457}
\showcasematerialdesignround{power-off}{1458}
\showcasematerialdesignround{power-settings-new}{1459}
\showcasematerialdesignround{precision-manufacturing}{1460}
\showcasematerialdesignround{pregnant-woman}{1461}
\showcasematerialdesignround{present-to-all}{1462}
\showcasematerialdesignround{preview}{1463}
\showcasematerialdesignround{price-change}{1464}
\showcasematerialdesignround{price-check}{1465}
\showcasematerialdesignround{print}{1466}
\showcasematerialdesignround{print-disabled}{1467}
\showcasematerialdesignround{priority-high}{1468}
\showcasematerialdesignround{privacy-tip}{1469}
\showcasematerialdesignround{private-connectivity}{1470}
\showcasematerialdesignround{production-quantity-limits}{1471}
\showcasematerialdesignround{propane}{1472}
\showcasematerialdesignround{propane-tank}{1473}
\showcasematerialdesignround{psychology}{1474}
\showcasematerialdesignround{psychology-alt}{1475}
\showcasematerialdesignround{public}{1476}
\showcasematerialdesignround{public-off}{1477}
\showcasematerialdesignround{publish}{1478}
\showcasematerialdesignround{published-with-changes}{1479}
\showcasematerialdesignround{punch-clock}{1480}
\showcasematerialdesignround{push-pin}{1481}
\showcasematerialdesignround{qr-code}{1482}
\showcasematerialdesignround{qr-code-2}{1483}
\showcasematerialdesignround{qr-code-scanner}{1484}
\showcasematerialdesignround{query-builder}{1485}
\showcasematerialdesignround{query-stats}{1486}
\showcasematerialdesignround{question-answer}{1487}
\showcasematerialdesignround{question-mark}{1488}
\showcasematerialdesignround{queue}{1489}
\showcasematerialdesignround{queue-music}{1490}
\showcasematerialdesignround{queue-play-next}{1491}
\showcasematerialdesignround{quickreply}{1492}
\showcasematerialdesignround{quiz}{1493}
\showcasematerialdesignround{radar}{1494}
\showcasematerialdesignround{radio}{1495}
\showcasematerialdesignround{radio-button-checked}{1496}
\showcasematerialdesignround{radio-button-unchecked}{1497}
\showcasematerialdesignround{railway-alert}{1498}
\showcasematerialdesignround{ramen-dining}{1499}
\showcasematerialdesignround{ramp-left}{1500}
\showcasematerialdesignround{ramp-right}{1501}
\showcasematerialdesignround{rate-review}{1502}
\showcasematerialdesignround{raw-off}{1503}
\showcasematerialdesignround{raw-on}{1504}
\showcasematerialdesignround{read-more}{1505}
\showcasematerialdesignround{real-estate-agent}{1506}
\showcasematerialdesignround{receipt}{1507}
\showcasematerialdesignround{receipt-long}{1508}
\showcasematerialdesignround{recent-actors}{1509}
\showcasematerialdesignround{recommend}{1510}
\showcasematerialdesignround{record-voice-over}{1511}
\showcasematerialdesignround{rectangle}{1512}
\showcasematerialdesignround{recycling}{1513}
\showcasematerialdesignround{redeem}{1514}
\showcasematerialdesignround{redo}{1515}
\showcasematerialdesignround{reduce-capacity}{1516}
\showcasematerialdesignround{refresh}{1517}
\showcasematerialdesignround{remember-me}{1518}
\showcasematerialdesignround{remove}{1519}
\showcasematerialdesignround{remove-circle}{1520}
\showcasematerialdesignround{remove-circle-outline}{1521}
\showcasematerialdesignround{remove-done}{1522}
\showcasematerialdesignround{remove-from-queue}{1523}
\showcasematerialdesignround{remove-moderator}{1524}
\showcasematerialdesignround{remove-red-eye}{1525}
\showcasematerialdesignround{remove-road}{1526}
\showcasematerialdesignround{remove-shopping-cart}{1527}
\showcasematerialdesignround{reorder}{1528}
\showcasematerialdesignround{repartition}{1529}
\showcasematerialdesignround{repeat}{1530}
\showcasematerialdesignround{repeat-on}{1531}
\showcasematerialdesignround{repeat-one}{1532}
\showcasematerialdesignround{repeat-one-on}{1533}
\showcasematerialdesignround{replay}{1534}
\showcasematerialdesignround{replay-5}{1535}
\showcasematerialdesignround{replay-10}{1536}
\showcasematerialdesignround{replay-30}{1537}
\showcasematerialdesignround{replay-circle-filled}{1538}
\showcasematerialdesignround{reply}{1539}
\showcasematerialdesignround{reply-all}{1540}
\showcasematerialdesignround{report}{1541}
\showcasematerialdesignround{report-gmailerrorred}{1542}
\showcasematerialdesignround{report-off}{1543}
\showcasematerialdesignround{report-problem}{1544}
\showcasematerialdesignround{request-page}{1545}
\showcasematerialdesignround{request-quote}{1546}
\showcasematerialdesignround{reset-tv}{1547}
\showcasematerialdesignround{restart-alt}{1548}
\showcasematerialdesignround{restaurant}{1549}
\showcasematerialdesignround{restaurant-menu}{1550}
\showcasematerialdesignround{restore}{1551}
\showcasematerialdesignround{restore-from-trash}{1552}
\showcasematerialdesignround{restore-page}{1553}
\showcasematerialdesignround{reviews}{1554}
\showcasematerialdesignround{rice-bowl}{1555}
\showcasematerialdesignround{ring-volume}{1556}
\showcasematerialdesignround{rocket}{1557}
\showcasematerialdesignround{rocket-launch}{1558}
\showcasematerialdesignround{roller-shades}{1559}
\showcasematerialdesignround{roller-shades-closed}{1560}
\showcasematerialdesignround{roller-skating}{1561}
\showcasematerialdesignround{roofing}{1562}
\showcasematerialdesignround{room}{1563}
\showcasematerialdesignround{room-preferences}{1564}
\showcasematerialdesignround{room-service}{1565}
\showcasematerialdesignround{rotate-90-degrees-ccw}{1566}
\showcasematerialdesignround{rotate-90-degrees-cw}{1567}
\showcasematerialdesignround{rotate-left}{1568}
\showcasematerialdesignround{rotate-right}{1569}
\showcasematerialdesignround{roundabout-left}{1570}
\showcasematerialdesignround{roundabout-right}{1571}
\showcasematerialdesignround{rounded-corner}{1572}
\showcasematerialdesignround{route}{1573}
\showcasematerialdesignround{router}{1574}
\showcasematerialdesignround{rowing}{1575}
\showcasematerialdesignround{rss-feed}{1576}
\showcasematerialdesignround{rsvp}{1577}
\showcasematerialdesignround{rtt}{1578}
\showcasematerialdesignround{rule}{1579}
\showcasematerialdesignround{rule-folder}{1580}
\showcasematerialdesignround{running-with-errors}{1581}
\showcasematerialdesignround{run-circle}{1582}
\showcasematerialdesignround{rv-hookup}{1583}
\showcasematerialdesignround{r-mobiledata}{1584}
\showcasematerialdesignround{safety-check}{1585}
\showcasematerialdesignround{safety-divider}{1586}
\showcasematerialdesignround{sailing}{1587}
\showcasematerialdesignround{sanitizer}{1588}
\showcasematerialdesignround{satellite}{1589}
\showcasematerialdesignround{satellite-alt}{1590}
\showcasematerialdesignround{save}{1591}
\showcasematerialdesignround{saved-search}{1592}
\showcasematerialdesignround{save-alt}{1593}
\showcasematerialdesignround{save-as}{1594}
\showcasematerialdesignround{savings}{1595}
\showcasematerialdesignround{scale}{1596}
\showcasematerialdesignround{scanner}{1597}
\showcasematerialdesignround{scatter-plot}{1598}
\showcasematerialdesignround{schedule}{1599}
\showcasematerialdesignround{schedule-send}{1600}
\showcasematerialdesignround{schema}{1601}
\showcasematerialdesignround{school}{1602}
\showcasematerialdesignround{science}{1603}
\showcasematerialdesignround{score}{1604}
\showcasematerialdesignround{scoreboard}{1605}
\showcasematerialdesignround{screenshot}{1606}
\showcasematerialdesignround{screenshot-monitor}{1607}
\showcasematerialdesignround{screen-lock-landscape}{1608}
\showcasematerialdesignround{screen-lock-portrait}{1609}
\showcasematerialdesignround{screen-lock-rotation}{1610}
\showcasematerialdesignround{screen-rotation}{1611}
\showcasematerialdesignround{screen-rotation-alt}{1612}
\showcasematerialdesignround{screen-search-desktop}{1613}
\showcasematerialdesignround{screen-share}{1614}
\showcasematerialdesignround{scuba-diving}{1615}
\showcasematerialdesignround{sd}{1616}
\showcasematerialdesignround{sd-card}{1617}
\showcasematerialdesignround{sd-card-alert}{1618}
\showcasematerialdesignround{sd-storage}{1619}
\showcasematerialdesignround{search}{1620}
\showcasematerialdesignround{search-off}{1621}
\showcasematerialdesignround{security}{1622}
\showcasematerialdesignround{security-update}{1623}
\showcasematerialdesignround{security-update-good}{1624}
\showcasematerialdesignround{security-update-warning}{1625}
\showcasematerialdesignround{segment}{1626}
\showcasematerialdesignround{select-all}{1627}
\showcasematerialdesignround{self-improvement}{1628}
\showcasematerialdesignround{sell}{1629}
\showcasematerialdesignround{send}{1630}
\showcasematerialdesignround{send-and-archive}{1631}
\showcasematerialdesignround{send-time-extension}{1632}
\showcasematerialdesignround{send-to-mobile}{1633}
\showcasematerialdesignround{sensors}{1634}
\showcasematerialdesignround{sensors-off}{1635}
\showcasematerialdesignround{sensor-door}{1636}
\showcasematerialdesignround{sensor-occupied}{1637}
\showcasematerialdesignround{sensor-window}{1638}
\showcasematerialdesignround{sentiment-dissatisfied}{1639}
\showcasematerialdesignround{sentiment-neutral}{1640}
\showcasematerialdesignround{sentiment-satisfied}{1641}
\showcasematerialdesignround{sentiment-satisfied-alt}{1642}
\showcasematerialdesignround{sentiment-very-dissatisfied}{1643}
\showcasematerialdesignround{sentiment-very-satisfied}{1644}
\showcasematerialdesignround{settings}{1645}
\showcasematerialdesignround{settings-accessibility}{1646}
\showcasematerialdesignround{settings-applications}{1647}
\showcasematerialdesignround{settings-backup-restore}{1648}
\showcasematerialdesignround{settings-bluetooth}{1649}
\showcasematerialdesignround{settings-brightness}{1650}
\showcasematerialdesignround{settings-cell}{1651}
\showcasematerialdesignround{settings-ethernet}{1652}
\showcasematerialdesignround{settings-input-antenna}{1653}
\showcasematerialdesignround{settings-input-component}{1654}
\showcasematerialdesignround{settings-input-composite}{1655}
\showcasematerialdesignround{settings-input-hdmi}{1656}
\showcasematerialdesignround{settings-input-svideo}{1657}
\showcasematerialdesignround{settings-overscan}{1658}
\showcasematerialdesignround{settings-phone}{1659}
\showcasematerialdesignround{settings-power}{1660}
\showcasematerialdesignround{settings-remote}{1661}
\showcasematerialdesignround{settings-suggest}{1662}
\showcasematerialdesignround{settings-system-daydream}{1663}
\showcasematerialdesignround{settings-voice}{1664}
\showcasematerialdesignround{set-meal}{1665}
\showcasematerialdesignround{severe-cold}{1666}
\showcasematerialdesignround{shape-line}{1667}
\showcasematerialdesignround{share}{1668}
\showcasematerialdesignround{share-location}{1669}
\showcasematerialdesignround{shield}{1670}
\showcasematerialdesignround{shield-moon}{1671}
\showcasematerialdesignround{shop}{1672}
\showcasematerialdesignround{shopping-bag}{1673}
\showcasematerialdesignround{shopping-basket}{1674}
\showcasematerialdesignround{shopping-cart}{1675}
\showcasematerialdesignround{shopping-cart-checkout}{1676}
\showcasematerialdesignround{shop-2}{1677}
\showcasematerialdesignround{shop-two}{1678}
\showcasematerialdesignround{shortcut}{1679}
\showcasematerialdesignround{short-text}{1680}
\showcasematerialdesignround{shower}{1681}
\showcasematerialdesignround{show-chart}{1682}
\showcasematerialdesignround{shuffle}{1683}
\showcasematerialdesignround{shuffle-on}{1684}
\showcasematerialdesignround{shutter-speed}{1685}
\showcasematerialdesignround{sick}{1686}
\showcasematerialdesignround{signal-cellular-0-bar}{1687}
\showcasematerialdesignround{signal-cellular-4-bar}{1688}
\showcasematerialdesignround{signal-cellular-alt}{1689}
\showcasematerialdesignround{signal-cellular-alt-1-bar}{1690}
\showcasematerialdesignround{signal-cellular-alt-2-bar}{1691}
\showcasematerialdesignround{signal-cellular-connected-no-internet-0-bar}{1692}
\showcasematerialdesignround{signal-cellular-connected-no-internet-4-bar}{1693}
\showcasematerialdesignround{signal-cellular-nodata}{1694}
\showcasematerialdesignround{signal-cellular-no-sim}{1695}
\showcasematerialdesignround{signal-cellular-null}{1696}
\showcasematerialdesignround{signal-cellular-off}{1697}
\showcasematerialdesignround{signal-wifi-0-bar}{1698}
\showcasematerialdesignround{signal-wifi-4-bar}{1699}
\showcasematerialdesignround{signal-wifi-4-bar-lock}{1700}
\showcasematerialdesignround{signal-wifi-bad}{1701}
\showcasematerialdesignround{signal-wifi-connected-no-internet-4}{1702}
\showcasematerialdesignround{signal-wifi-off}{1703}
\showcasematerialdesignround{signal-wifi-statusbar-4-bar}{1704}
\showcasematerialdesignround{signal-wifi-statusbar-connected-no-internet-4}{1705}
\showcasematerialdesignround{signal-wifi-statusbar-null}{1706}
\showcasematerialdesignround{signpost}{1707}
\showcasematerialdesignround{sign-language}{1708}
\showcasematerialdesignround{sim-card}{1709}
\showcasematerialdesignround{sim-card-alert}{1710}
\showcasematerialdesignround{sim-card-download}{1711}
\showcasematerialdesignround{single-bed}{1712}
\showcasematerialdesignround{sip}{1713}
\showcasematerialdesignround{skateboarding}{1714}
\showcasematerialdesignround{skip-next}{1715}
\showcasematerialdesignround{skip-previous}{1716}
\showcasematerialdesignround{sledding}{1717}
\showcasematerialdesignround{slideshow}{1718}
\showcasematerialdesignround{slow-motion-video}{1719}
\showcasematerialdesignround{smartphone}{1720}
\showcasematerialdesignround{smart-button}{1721}
\showcasematerialdesignround{smart-display}{1722}
\showcasematerialdesignround{smart-screen}{1723}
\showcasematerialdesignround{smart-toy}{1724}
\showcasematerialdesignround{smoke-free}{1725}
\showcasematerialdesignround{smoking-rooms}{1726}
\showcasematerialdesignround{sms}{1727}
\showcasematerialdesignround{sms-failed}{1728}
\showcasematerialdesignround{snippet-folder}{1729}
\showcasematerialdesignround{snooze}{1730}
\showcasematerialdesignround{snowboarding}{1731}
\showcasematerialdesignround{snowmobile}{1732}
\showcasematerialdesignround{snowshoeing}{1733}
\showcasematerialdesignround{soap}{1734}
\showcasematerialdesignround{social-distance}{1735}
\showcasematerialdesignround{solar-power}{1736}
\showcasematerialdesignround{sort}{1737}
\showcasematerialdesignround{sort-by-alpha}{1738}
\showcasematerialdesignround{sos}{1739}
\showcasematerialdesignround{soup-kitchen}{1740}
\showcasematerialdesignround{source}{1741}
\showcasematerialdesignround{south}{1742}
\showcasematerialdesignround{south-america}{1743}
\showcasematerialdesignround{south-east}{1744}
\showcasematerialdesignround{south-west}{1745}
\showcasematerialdesignround{spa}{1746}
\showcasematerialdesignround{space-bar}{1747}
\showcasematerialdesignround{space-dashboard}{1748}
\showcasematerialdesignround{spatial-audio}{1749}
\showcasematerialdesignround{spatial-audio-off}{1750}
\showcasematerialdesignround{spatial-tracking}{1751}
\showcasematerialdesignround{speaker}{1752}
\showcasematerialdesignround{speaker-group}{1753}
\showcasematerialdesignround{speaker-notes}{1754}
\showcasematerialdesignround{speaker-notes-off}{1755}
\showcasematerialdesignround{speaker-phone}{1756}
\showcasematerialdesignround{speed}{1757}
\showcasematerialdesignround{spellcheck}{1758}
\showcasematerialdesignround{splitscreen}{1759}
\showcasematerialdesignround{spoke}{1760}
\showcasematerialdesignround{sports}{1761}
\showcasematerialdesignround{sports-bar}{1762}
\showcasematerialdesignround{sports-baseball}{1763}
\showcasematerialdesignround{sports-basketball}{1764}
\showcasematerialdesignround{sports-cricket}{1765}
\showcasematerialdesignround{sports-esports}{1766}
\showcasematerialdesignround{sports-football}{1767}
\showcasematerialdesignround{sports-golf}{1768}
\showcasematerialdesignround{sports-gymnastics}{1769}
\showcasematerialdesignround{sports-handball}{1770}
\showcasematerialdesignround{sports-hockey}{1771}
\showcasematerialdesignround{sports-kabaddi}{1772}
\showcasematerialdesignround{sports-martial-arts}{1773}
\showcasematerialdesignround{sports-mma}{1774}
\showcasematerialdesignround{sports-motorsports}{1775}
\showcasematerialdesignround{sports-rugby}{1776}
\showcasematerialdesignround{sports-score}{1777}
\showcasematerialdesignround{sports-soccer}{1778}
\showcasematerialdesignround{sports-tennis}{1779}
\showcasematerialdesignround{sports-volleyball}{1780}
\showcasematerialdesignround{square}{1781}
\showcasematerialdesignround{square-foot}{1782}
\showcasematerialdesignround{ssid-chart}{1783}
\showcasematerialdesignround{stacked-bar-chart}{1784}
\showcasematerialdesignround{stacked-line-chart}{1785}
\showcasematerialdesignround{stadium}{1786}
\showcasematerialdesignround{stairs}{1787}
\showcasematerialdesignround{star}{1788}
\showcasematerialdesignround{stars}{1789}
\showcasematerialdesignround{start}{1790}
\showcasematerialdesignround{star-border}{1791}
\showcasematerialdesignround{star-border-purple500}{1792}
\showcasematerialdesignround{star-half}{1793}
\showcasematerialdesignround{star-outline}{1794}
\showcasematerialdesignround{star-purple500}{1795}
\showcasematerialdesignround{star-rate}{1796}
\showcasematerialdesignround{stay-current-landscape}{1797}
\showcasematerialdesignround{stay-current-portrait}{1798}
\showcasematerialdesignround{stay-primary-landscape}{1799}
\showcasematerialdesignround{stay-primary-portrait}{1800}
\showcasematerialdesignround{sticky-note-2}{1801}
\showcasematerialdesignround{stop}{1802}
\showcasematerialdesignround{stop-circle}{1803}
\showcasematerialdesignround{stop-screen-share}{1804}
\showcasematerialdesignround{storage}{1805}
\showcasematerialdesignround{store}{1806}
\showcasematerialdesignround{storefront}{1807}
\showcasematerialdesignround{store-mall-directory}{1808}
\showcasematerialdesignround{storm}{1809}
\showcasematerialdesignround{straight}{1810}
\showcasematerialdesignround{straighten}{1811}
\showcasematerialdesignround{stream}{1812}
\showcasematerialdesignround{streetview}{1813}
\showcasematerialdesignround{strikethrough-s}{1814}
\showcasematerialdesignround{stroller}{1815}
\showcasematerialdesignround{style}{1816}
\showcasematerialdesignround{subdirectory-arrow-left}{1817}
\showcasematerialdesignround{subdirectory-arrow-right}{1818}
\showcasematerialdesignround{subject}{1819}
\showcasematerialdesignround{subscript}{1820}
\showcasematerialdesignround{subscriptions}{1821}
\showcasematerialdesignround{subtitles}{1822}
\showcasematerialdesignround{subtitles-off}{1823}
\showcasematerialdesignround{subway}{1824}
\showcasematerialdesignround{summarize}{1825}
\showcasematerialdesignround{superscript}{1826}
\showcasematerialdesignround{supervised-user-circle}{1827}
\showcasematerialdesignround{supervisor-account}{1828}
\showcasematerialdesignround{support}{1829}
\showcasematerialdesignround{support-agent}{1830}
\showcasematerialdesignround{surfing}{1831}
\showcasematerialdesignround{surround-sound}{1832}
\showcasematerialdesignround{swap-calls}{1833}
\showcasematerialdesignround{swap-horiz}{1834}
\showcasematerialdesignround{swap-horizontal-circle}{1835}
\showcasematerialdesignround{swap-vert}{1836}
\showcasematerialdesignround{swap-vertical-circle}{1837}
\showcasematerialdesignround{swipe}{1838}
\showcasematerialdesignround{swipe-down}{1839}
\showcasematerialdesignround{swipe-down-alt}{1840}
\showcasematerialdesignround{swipe-left}{1841}
\showcasematerialdesignround{swipe-left-alt}{1842}
\showcasematerialdesignround{swipe-right}{1843}
\showcasematerialdesignround{swipe-right-alt}{1844}
\showcasematerialdesignround{swipe-up}{1845}
\showcasematerialdesignround{swipe-up-alt}{1846}
\showcasematerialdesignround{swipe-vertical}{1847}
\showcasematerialdesignround{switch-access-shortcut}{1848}
\showcasematerialdesignround{switch-access-shortcut-add}{1849}
\showcasematerialdesignround{switch-account}{1850}
\showcasematerialdesignround{switch-camera}{1851}
\showcasematerialdesignround{switch-left}{1852}
\showcasematerialdesignround{switch-right}{1853}
\showcasematerialdesignround{switch-video}{1854}
\showcasematerialdesignround{synagogue}{1855}
\showcasematerialdesignround{sync}{1856}
\showcasematerialdesignround{sync-alt}{1857}
\showcasematerialdesignround{sync-disabled}{1858}
\showcasematerialdesignround{sync-lock}{1859}
\showcasematerialdesignround{sync-problem}{1860}
\showcasematerialdesignround{system-security-update}{1861}
\showcasematerialdesignround{system-security-update-good}{1862}
\showcasematerialdesignround{system-security-update-warning}{1863}
\showcasematerialdesignround{system-update}{1864}
\showcasematerialdesignround{system-update-alt}{1865}
\showcasematerialdesignround{tab}{1866}
\showcasematerialdesignround{tablet}{1867}
\showcasematerialdesignround{tablet-android}{1868}
\showcasematerialdesignround{tablet-mac}{1869}
\showcasematerialdesignround{table-bar}{1870}
\showcasematerialdesignround{table-chart}{1871}
\showcasematerialdesignround{table-restaurant}{1872}
\showcasematerialdesignround{table-rows}{1873}
\showcasematerialdesignround{table-view}{1874}
\showcasematerialdesignround{tab-unselected}{1875}
\showcasematerialdesignround{tag}{1876}
\showcasematerialdesignround{tag-faces}{1877}
\showcasematerialdesignround{takeout-dining}{1878}
\showcasematerialdesignround{tapas}{1879}
\showcasematerialdesignround{tap-and-play}{1880}
\showcasematerialdesignround{task}{1881}
\showcasematerialdesignround{task-alt}{1882}
\showcasematerialdesignround{taxi-alert}{1883}
\showcasematerialdesignround{temple-buddhist}{1884}
\showcasematerialdesignround{temple-hindu}{1885}
\showcasematerialdesignround{terminal}{1886}
\showcasematerialdesignround{terrain}{1887}
\showcasematerialdesignround{textsms}{1888}
\showcasematerialdesignround{texture}{1889}
\showcasematerialdesignround{text-decrease}{1890}
\showcasematerialdesignround{text-fields}{1891}
\showcasematerialdesignround{text-format}{1892}
\showcasematerialdesignround{text-increase}{1893}
\showcasematerialdesignround{text-rotate-up}{1894}
\showcasematerialdesignround{text-rotate-vertical}{1895}
\showcasematerialdesignround{text-rotation-angledown}{1896}
\showcasematerialdesignround{text-rotation-angleup}{1897}
\showcasematerialdesignround{text-rotation-down}{1898}
\showcasematerialdesignround{text-rotation-none}{1899}
\showcasematerialdesignround{text-snippet}{1900}
\showcasematerialdesignround{theaters}{1901}
\showcasematerialdesignround{theater-comedy}{1902}
\showcasematerialdesignround{thermostat}{1903}
\showcasematerialdesignround{thermostat-auto}{1904}
\showcasematerialdesignround{thumbs-up-down}{1905}
\showcasematerialdesignround{thumb-down}{1906}
\showcasematerialdesignround{thumb-down-alt}{1907}
\showcasematerialdesignround{thumb-down-off-alt}{1908}
\showcasematerialdesignround{thumb-up}{1909}
\showcasematerialdesignround{thumb-up-alt}{1910}
\showcasematerialdesignround{thumb-up-off-alt}{1911}
\showcasematerialdesignround{thunderstorm}{1912}
\showcasematerialdesignround{timelapse}{1913}
\showcasematerialdesignround{timeline}{1914}
\showcasematerialdesignround{timer}{1915}
\showcasematerialdesignround{timer-3}{1916}
\showcasematerialdesignround{timer-3-select}{1917}
\showcasematerialdesignround{timer-10}{1918}
\showcasematerialdesignround{timer-10-select}{1919}
\showcasematerialdesignround{timer-off}{1920}
\showcasematerialdesignround{time-to-leave}{1921}
\showcasematerialdesignround{tips-and-updates}{1922}
\showcasematerialdesignround{tire-repair}{1923}
\showcasematerialdesignround{title}{1924}
\showcasematerialdesignround{toc}{1925}
\showcasematerialdesignround{today}{1926}
\showcasematerialdesignround{toggle-off}{1927}
\showcasematerialdesignround{toggle-on}{1928}
\showcasematerialdesignround{token}{1929}
\showcasematerialdesignround{toll}{1930}
\showcasematerialdesignround{tonality}{1931}
\showcasematerialdesignround{topic}{1932}
\showcasematerialdesignround{tornado}{1933}
\showcasematerialdesignround{touch-app}{1934}
\showcasematerialdesignround{tour}{1935}
\showcasematerialdesignround{toys}{1936}
\showcasematerialdesignround{track-changes}{1937}
\showcasematerialdesignround{traffic}{1938}
\showcasematerialdesignround{train}{1939}
\showcasematerialdesignround{tram}{1940}
\showcasematerialdesignround{transcribe}{1941}
\showcasematerialdesignround{transfer-within-a-station}{1942}
\showcasematerialdesignround{transform}{1943}
\showcasematerialdesignround{transgender}{1944}
\showcasematerialdesignround{transit-enterexit}{1945}
\showcasematerialdesignround{translate}{1946}
\showcasematerialdesignround{travel-explore}{1947}
\showcasematerialdesignround{trending-down}{1948}
\showcasematerialdesignround{trending-flat}{1949}
\showcasematerialdesignround{trending-up}{1950}
\showcasematerialdesignround{trip-origin}{1951}
\showcasematerialdesignround{troubleshoot}{1952}
\showcasematerialdesignround{try}{1953}
\showcasematerialdesignround{tsunami}{1954}
\showcasematerialdesignround{tty}{1955}
\showcasematerialdesignround{tune}{1956}
\showcasematerialdesignround{tungsten}{1957}
\showcasematerialdesignround{turned-in}{1958}
\showcasematerialdesignround{turned-in-not}{1959}
\showcasematerialdesignround{turn-left}{1960}
\showcasematerialdesignround{turn-right}{1961}
\showcasematerialdesignround{turn-sharp-left}{1962}
\showcasematerialdesignround{turn-sharp-right}{1963}
\showcasematerialdesignround{turn-slight-left}{1964}
\showcasematerialdesignround{turn-slight-right}{1965}
\showcasematerialdesignround{tv}{1966}
\showcasematerialdesignround{tv-off}{1967}
\showcasematerialdesignround{two-wheeler}{1968}
\showcasematerialdesignround{type-specimen}{1969}
\showcasematerialdesignround{umbrella}{1970}
\showcasematerialdesignround{unarchive}{1971}
\showcasematerialdesignround{undo}{1972}
\showcasematerialdesignround{unfold-less}{1973}
\showcasematerialdesignround{unfold-less-double}{1974}
\showcasematerialdesignround{unfold-more}{1975}
\showcasematerialdesignround{unfold-more-double}{1976}
\showcasematerialdesignround{unpublished}{1977}
\showcasematerialdesignround{unsubscribe}{1978}
\showcasematerialdesignround{upcoming}{1979}
\showcasematerialdesignround{update}{1980}
\showcasematerialdesignround{update-disabled}{1981}
\showcasematerialdesignround{upgrade}{1982}
\showcasematerialdesignround{upload}{1983}
\showcasematerialdesignround{upload-file}{1984}
\showcasematerialdesignround{usb}{1985}
\showcasematerialdesignround{usb-off}{1986}
\showcasematerialdesignround{u-turn-left}{1987}
\showcasematerialdesignround{u-turn-right}{1988}
\showcasematerialdesignround{vaccines}{1989}
\showcasematerialdesignround{vape-free}{1990}
\showcasematerialdesignround{vaping-rooms}{1991}
\showcasematerialdesignround{verified}{1992}
\showcasematerialdesignround{verified-user}{1993}
\showcasematerialdesignround{vertical-align-bottom}{1994}
\showcasematerialdesignround{vertical-align-center}{1995}
\showcasematerialdesignround{vertical-align-top}{1996}
\showcasematerialdesignround{vertical-distribute}{1997}
\showcasematerialdesignround{vertical-shades}{1998}
\showcasematerialdesignround{vertical-shades-closed}{1999}
\showcasematerialdesignround{vertical-split}{2000}
\showcasematerialdesignround{vibration}{2001}
\showcasematerialdesignround{videocam}{2002}
\showcasematerialdesignround{videocam-off}{2003}
\showcasematerialdesignround{videogame-asset}{2004}
\showcasematerialdesignround{videogame-asset-off}{2005}
\showcasematerialdesignround{video-call}{2006}
\showcasematerialdesignround{video-camera-back}{2007}
\showcasematerialdesignround{video-camera-front}{2008}
\showcasematerialdesignround{video-chat}{2009}
\showcasematerialdesignround{video-file}{2010}
\showcasematerialdesignround{video-label}{2011}
\showcasematerialdesignround{video-library}{2012}
\showcasematerialdesignround{video-settings}{2013}
\showcasematerialdesignround{video-stable}{2014}
\showcasematerialdesignround{view-agenda}{2015}
\showcasematerialdesignround{view-array}{2016}
\showcasematerialdesignround{view-carousel}{2017}
\showcasematerialdesignround{view-column}{2018}
\showcasematerialdesignround{view-comfy}{2019}
\showcasematerialdesignround{view-comfy-alt}{2020}
\showcasematerialdesignround{view-compact}{2021}
\showcasematerialdesignround{view-compact-alt}{2022}
\showcasematerialdesignround{view-cozy}{2023}
\showcasematerialdesignround{view-day}{2024}
\showcasematerialdesignround{view-headline}{2025}
\showcasematerialdesignround{view-in-ar}{2026}
\showcasematerialdesignround{view-kanban}{2027}
\showcasematerialdesignround{view-list}{2028}
\showcasematerialdesignround{view-module}{2029}
\showcasematerialdesignround{view-quilt}{2030}
\showcasematerialdesignround{view-sidebar}{2031}
\showcasematerialdesignround{view-stream}{2032}
\showcasematerialdesignround{view-timeline}{2033}
\showcasematerialdesignround{view-week}{2034}
\showcasematerialdesignround{vignette}{2035}
\showcasematerialdesignround{villa}{2036}
\showcasematerialdesignround{visibility}{2037}
\showcasematerialdesignround{visibility-off}{2038}
\showcasematerialdesignround{voicemail}{2039}
\showcasematerialdesignround{voice-chat}{2040}
\showcasematerialdesignround{voice-over-off}{2041}
\showcasematerialdesignround{volcano}{2042}
\showcasematerialdesignround{volume-down}{2043}
\showcasematerialdesignround{volume-mute}{2044}
\showcasematerialdesignround{volume-off}{2045}
\showcasematerialdesignround{volume-up}{2046}
\showcasematerialdesignround{volunteer-activism}{2047}
\showcasematerialdesignround{vpn-key}{2048}
\showcasematerialdesignround{vpn-key-off}{2049}
\showcasematerialdesignround{vpn-lock}{2050}
\showcasematerialdesignround{vrpano}{2051}
\showcasematerialdesignround{wallet}{2052}
\showcasematerialdesignround{wallpaper}{2053}
\showcasematerialdesignround{warehouse}{2054}
\showcasematerialdesignround{warning}{2055}
\showcasematerialdesignround{warning-amber}{2056}
\showcasematerialdesignround{wash}{2057}
\showcasematerialdesignround{watch}{2058}
\showcasematerialdesignround{watch-later}{2059}
\showcasematerialdesignround{watch-off}{2060}
\showcasematerialdesignround{water}{2061}
\showcasematerialdesignround{waterfall-chart}{2062}
\showcasematerialdesignround{water-damage}{2063}
\showcasematerialdesignround{water-drop}{2064}
\showcasematerialdesignround{waves}{2065}
\showcasematerialdesignround{waving-hand}{2066}
\showcasematerialdesignround{wb-auto}{2067}
\showcasematerialdesignround{wb-cloudy}{2068}
\showcasematerialdesignround{wb-incandescent}{2069}
\showcasematerialdesignround{wb-iridescent}{2070}
\showcasematerialdesignround{wb-shade}{2071}
\showcasematerialdesignround{wb-sunny}{2072}
\showcasematerialdesignround{wb-twilight}{2073}
\showcasematerialdesignround{wc}{2074}
\showcasematerialdesignround{web}{2075}
\showcasematerialdesignround{webhook}{2076}
\showcasematerialdesignround{web-asset}{2077}
\showcasematerialdesignround{web-asset-off}{2078}
\showcasematerialdesignround{web-stories}{2079}
\showcasematerialdesignround{weekend}{2080}
\showcasematerialdesignround{west}{2081}
\showcasematerialdesignround{whatshot}{2082}
\showcasematerialdesignround{wheelchair-pickup}{2083}
\showcasematerialdesignround{where-to-vote}{2084}
\showcasematerialdesignround{widgets}{2085}
\showcasematerialdesignround{width-full}{2086}
\showcasematerialdesignround{width-normal}{2087}
\showcasematerialdesignround{width-wide}{2088}
\showcasematerialdesignround{wifi}{2089}
\showcasematerialdesignround{wifi-1-bar}{2090}
\showcasematerialdesignround{wifi-2-bar}{2091}
\showcasematerialdesignround{wifi-calling}{2092}
\showcasematerialdesignround{wifi-calling-3}{2093}
\showcasematerialdesignround{wifi-channel}{2094}
\showcasematerialdesignround{wifi-find}{2095}
\showcasematerialdesignround{wifi-lock}{2096}
\showcasematerialdesignround{wifi-off}{2097}
\showcasematerialdesignround{wifi-password}{2098}
\showcasematerialdesignround{wifi-protected-setup}{2099}
\showcasematerialdesignround{wifi-tethering}{2100}
\showcasematerialdesignround{wifi-tethering-error}{2101}
\showcasematerialdesignround{wifi-tethering-off}{2102}
\showcasematerialdesignround{window}{2103}
\showcasematerialdesignround{wind-power}{2104}
\showcasematerialdesignround{wine-bar}{2105}
\showcasematerialdesignround{woman}{2106}
\showcasematerialdesignround{woman-2}{2107}
\showcasematerialdesignround{work}{2108}
\showcasematerialdesignround{workspaces}{2109}
\showcasematerialdesignround{workspace-premium}{2110}
\showcasematerialdesignround{work-history}{2111}
\showcasematerialdesignround{work-off}{2112}
\showcasematerialdesignround{work-outline}{2113}
\showcasematerialdesignround{wrap-text}{2114}
\showcasematerialdesignround{wrong-location}{2115}
\showcasematerialdesignround{wysiwyg}{2116}
\showcasematerialdesignround{yard}{2117}
\showcasematerialdesignround{youtube-searched-for}{2118}
\showcasematerialdesignround{zoom-in}{2119}
\showcasematerialdesignround{zoom-in-map}{2120}
\showcasematerialdesignround{zoom-out}{2121}
\showcasematerialdesignround{zoom-out-map}{2122}
\end{materialdesigncase}

\subsection{Sharp version (2122 icons)}

The \textsf{PDF} file is \texttt{materialdesign-sharp-all.pdf}

\begin{materialdesigncase}
\showcasematerialdesignsharp{1k}{1}
\showcasematerialdesignsharp{1k-plus}{2}
\showcasematerialdesignsharp{1x-mobiledata}{3}
\showcasematerialdesignsharp{2k}{4}
\showcasematerialdesignsharp{2k-plus}{5}
\showcasematerialdesignsharp{2mp}{6}
\showcasematerialdesignsharp{3d-rotation}{7}
\showcasematerialdesignsharp{3g-mobiledata}{8}
\showcasematerialdesignsharp{3k}{9}
\showcasematerialdesignsharp{3k-plus}{10}
\showcasematerialdesignsharp{3mp}{11}
\showcasematerialdesignsharp{3p}{12}
\showcasematerialdesignsharp{4g-mobiledata}{13}
\showcasematerialdesignsharp{4g-plus-mobiledata}{14}
\showcasematerialdesignsharp{4k}{15}
\showcasematerialdesignsharp{4k-plus}{16}
\showcasematerialdesignsharp{4mp}{17}
\showcasematerialdesignsharp{5g}{18}
\showcasematerialdesignsharp{5k}{19}
\showcasematerialdesignsharp{5k-plus}{20}
\showcasematerialdesignsharp{5mp}{21}
\showcasematerialdesignsharp{6k}{22}
\showcasematerialdesignsharp{6k-plus}{23}
\showcasematerialdesignsharp{6mp}{24}
\showcasematerialdesignsharp{6-ft-apart}{25}
\showcasematerialdesignsharp{7k}{26}
\showcasematerialdesignsharp{7k-plus}{27}
\showcasematerialdesignsharp{7mp}{28}
\showcasematerialdesignsharp{8k}{29}
\showcasematerialdesignsharp{8k-plus}{30}
\showcasematerialdesignsharp{8mp}{31}
\showcasematerialdesignsharp{9k}{32}
\showcasematerialdesignsharp{9k-plus}{33}
\showcasematerialdesignsharp{9mp}{34}
\showcasematerialdesignsharp{10k}{35}
\showcasematerialdesignsharp{10mp}{36}
\showcasematerialdesignsharp{11mp}{37}
\showcasematerialdesignsharp{12mp}{38}
\showcasematerialdesignsharp{13mp}{39}
\showcasematerialdesignsharp{14mp}{40}
\showcasematerialdesignsharp{15mp}{41}
\showcasematerialdesignsharp{16mp}{42}
\showcasematerialdesignsharp{17mp}{43}
\showcasematerialdesignsharp{18mp}{44}
\showcasematerialdesignsharp{18-up-rating}{45}
\showcasematerialdesignsharp{19mp}{46}
\showcasematerialdesignsharp{20mp}{47}
\showcasematerialdesignsharp{21mp}{48}
\showcasematerialdesignsharp{22mp}{49}
\showcasematerialdesignsharp{23mp}{50}
\showcasematerialdesignsharp{24mp}{51}
\showcasematerialdesignsharp{30fps}{52}
\showcasematerialdesignsharp{30fps-select}{53}
\showcasematerialdesignsharp{60fps}{54}
\showcasematerialdesignsharp{60fps-select}{55}
\showcasematerialdesignsharp{123}{56}
\showcasematerialdesignsharp{360}{57}
\showcasematerialdesignsharp{abc}{58}
\showcasematerialdesignsharp{accessibility}{59}
\showcasematerialdesignsharp{accessibility-new}{60}
\showcasematerialdesignsharp{accessible}{61}
\showcasematerialdesignsharp{accessible-forward}{62}
\showcasematerialdesignsharp{access-alarm}{63}
\showcasematerialdesignsharp{access-alarms}{64}
\showcasematerialdesignsharp{access-time}{65}
\showcasematerialdesignsharp{access-time-filled}{66}
\showcasematerialdesignsharp{account-balance}{67}
\showcasematerialdesignsharp{account-balance-wallet}{68}
\showcasematerialdesignsharp{account-box}{69}
\showcasematerialdesignsharp{account-circle}{70}
\showcasematerialdesignsharp{account-tree}{71}
\showcasematerialdesignsharp{ac-unit}{72}
\showcasematerialdesignsharp{adb}{73}
\showcasematerialdesignsharp{add}{74}
\showcasematerialdesignsharp{addchart}{75}
\showcasematerialdesignsharp{add-alarm}{76}
\showcasematerialdesignsharp{add-alert}{77}
\showcasematerialdesignsharp{add-a-photo}{78}
\showcasematerialdesignsharp{add-box}{79}
\showcasematerialdesignsharp{add-business}{80}
\showcasematerialdesignsharp{add-card}{81}
\showcasematerialdesignsharp{add-chart}{82}
\showcasematerialdesignsharp{add-circle}{83}
\showcasematerialdesignsharp{add-circle-outline}{84}
\showcasematerialdesignsharp{add-comment}{85}
\showcasematerialdesignsharp{add-home}{86}
\showcasematerialdesignsharp{add-home-work}{87}
\showcasematerialdesignsharp{add-ic-call}{88}
\showcasematerialdesignsharp{add-link}{89}
\showcasematerialdesignsharp{add-location}{90}
\showcasematerialdesignsharp{add-location-alt}{91}
\showcasematerialdesignsharp{add-moderator}{92}
\showcasematerialdesignsharp{add-photo-alternate}{93}
\showcasematerialdesignsharp{add-reaction}{94}
\showcasematerialdesignsharp{add-road}{95}
\showcasematerialdesignsharp{add-shopping-cart}{96}
\showcasematerialdesignsharp{add-task}{97}
\showcasematerialdesignsharp{add-to-drive}{98}
\showcasematerialdesignsharp{add-to-home-screen}{99}
\showcasematerialdesignsharp{add-to-photos}{100}
\showcasematerialdesignsharp{add-to-queue}{101}
\showcasematerialdesignsharp{adf-scanner}{102}
\showcasematerialdesignsharp{adjust}{103}
\showcasematerialdesignsharp{admin-panel-settings}{104}
\showcasematerialdesignsharp{ads-click}{105}
\showcasematerialdesignsharp{ad-units}{106}
\showcasematerialdesignsharp{agriculture}{107}
\showcasematerialdesignsharp{air}{108}
\showcasematerialdesignsharp{airlines}{109}
\showcasematerialdesignsharp{airline-seat-flat}{110}
\showcasematerialdesignsharp{airline-seat-flat-angled}{111}
\showcasematerialdesignsharp{airline-seat-individual-suite}{112}
\showcasematerialdesignsharp{airline-seat-legroom-extra}{113}
\showcasematerialdesignsharp{airline-seat-legroom-normal}{114}
\showcasematerialdesignsharp{airline-seat-legroom-reduced}{115}
\showcasematerialdesignsharp{airline-seat-recline-extra}{116}
\showcasematerialdesignsharp{airline-seat-recline-normal}{117}
\showcasematerialdesignsharp{airline-stops}{118}
\showcasematerialdesignsharp{airplanemode-active}{119}
\showcasematerialdesignsharp{airplanemode-inactive}{120}
\showcasematerialdesignsharp{airplane-ticket}{121}
\showcasematerialdesignsharp{airplay}{122}
\showcasematerialdesignsharp{airport-shuttle}{123}
\showcasematerialdesignsharp{alarm}{124}
\showcasematerialdesignsharp{alarm-add}{125}
\showcasematerialdesignsharp{alarm-off}{126}
\showcasematerialdesignsharp{alarm-on}{127}
\showcasematerialdesignsharp{album}{128}
\showcasematerialdesignsharp{align-horizontal-center}{129}
\showcasematerialdesignsharp{align-horizontal-left}{130}
\showcasematerialdesignsharp{align-horizontal-right}{131}
\showcasematerialdesignsharp{align-vertical-bottom}{132}
\showcasematerialdesignsharp{align-vertical-center}{133}
\showcasematerialdesignsharp{align-vertical-top}{134}
\showcasematerialdesignsharp{all-inbox}{135}
\showcasematerialdesignsharp{all-inclusive}{136}
\showcasematerialdesignsharp{all-out}{137}
\showcasematerialdesignsharp{alternate-email}{138}
\showcasematerialdesignsharp{alt-route}{139}
\showcasematerialdesignsharp{analytics}{140}
\showcasematerialdesignsharp{anchor}{141}
\showcasematerialdesignsharp{android}{142}
\showcasematerialdesignsharp{animation}{143}
\showcasematerialdesignsharp{announcement}{144}
\showcasematerialdesignsharp{aod}{145}
\showcasematerialdesignsharp{apartment}{146}
\showcasematerialdesignsharp{api}{147}
\showcasematerialdesignsharp{approval}{148}
\showcasematerialdesignsharp{apps}{149}
\showcasematerialdesignsharp{apps-outage}{150}
\showcasematerialdesignsharp{app-blocking}{151}
\showcasematerialdesignsharp{app-registration}{152}
\showcasematerialdesignsharp{app-settings-alt}{153}
\showcasematerialdesignsharp{app-shortcut}{154}
\showcasematerialdesignsharp{architecture}{155}
\showcasematerialdesignsharp{archive}{156}
\showcasematerialdesignsharp{area-chart}{157}
\showcasematerialdesignsharp{arrow-back}{158}
\showcasematerialdesignsharp{arrow-back-ios}{159}
\showcasematerialdesignsharp{arrow-back-ios-new}{160}
\showcasematerialdesignsharp{arrow-circle-down}{161}
\showcasematerialdesignsharp{arrow-circle-left}{162}
\showcasematerialdesignsharp{arrow-circle-right}{163}
\showcasematerialdesignsharp{arrow-circle-up}{164}
\showcasematerialdesignsharp{arrow-downward}{165}
\showcasematerialdesignsharp{arrow-drop-down}{166}
\showcasematerialdesignsharp{arrow-drop-down-circle}{167}
\showcasematerialdesignsharp{arrow-drop-up}{168}
\showcasematerialdesignsharp{arrow-forward}{169}
\showcasematerialdesignsharp{arrow-forward-ios}{170}
\showcasematerialdesignsharp{arrow-left}{171}
\showcasematerialdesignsharp{arrow-outward}{172}
\showcasematerialdesignsharp{arrow-right}{173}
\showcasematerialdesignsharp{arrow-right-alt}{174}
\showcasematerialdesignsharp{arrow-upward}{175}
\showcasematerialdesignsharp{article}{176}
\showcasematerialdesignsharp{art-track}{177}
\showcasematerialdesignsharp{aspect-ratio}{178}
\showcasematerialdesignsharp{assessment}{179}
\showcasematerialdesignsharp{assignment}{180}
\showcasematerialdesignsharp{assignment-ind}{181}
\showcasematerialdesignsharp{assignment-late}{182}
\showcasematerialdesignsharp{assignment-return}{183}
\showcasematerialdesignsharp{assignment-returned}{184}
\showcasematerialdesignsharp{assignment-turned-in}{185}
\showcasematerialdesignsharp{assistant}{186}
\showcasematerialdesignsharp{assistant-direction}{187}
\showcasematerialdesignsharp{assistant-photo}{188}
\showcasematerialdesignsharp{assist-walker}{189}
\showcasematerialdesignsharp{assured-workload}{190}
\showcasematerialdesignsharp{atm}{191}
\showcasematerialdesignsharp{attachment}{192}
\showcasematerialdesignsharp{attach-email}{193}
\showcasematerialdesignsharp{attach-file}{194}
\showcasematerialdesignsharp{attach-money}{195}
\showcasematerialdesignsharp{attractions}{196}
\showcasematerialdesignsharp{attribution}{197}
\showcasematerialdesignsharp{audiotrack}{198}
\showcasematerialdesignsharp{audio-file}{199}
\showcasematerialdesignsharp{autofps-select}{200}
\showcasematerialdesignsharp{autorenew}{201}
\showcasematerialdesignsharp{auto-awesome}{202}
\showcasematerialdesignsharp{auto-awesome-mosaic}{203}
\showcasematerialdesignsharp{auto-awesome-motion}{204}
\showcasematerialdesignsharp{auto-delete}{205}
\showcasematerialdesignsharp{auto-fix-high}{206}
\showcasematerialdesignsharp{auto-fix-normal}{207}
\showcasematerialdesignsharp{auto-fix-off}{208}
\showcasematerialdesignsharp{auto-graph}{209}
\showcasematerialdesignsharp{auto-mode}{210}
\showcasematerialdesignsharp{auto-stories}{211}
\showcasematerialdesignsharp{av-timer}{212}
\showcasematerialdesignsharp{baby-changing-station}{213}
\showcasematerialdesignsharp{backpack}{214}
\showcasematerialdesignsharp{backspace}{215}
\showcasematerialdesignsharp{backup}{216}
\showcasematerialdesignsharp{backup-table}{217}
\showcasematerialdesignsharp{back-hand}{218}
\showcasematerialdesignsharp{badge}{219}
\showcasematerialdesignsharp{bakery-dining}{220}
\showcasematerialdesignsharp{balance}{221}
\showcasematerialdesignsharp{balcony}{222}
\showcasematerialdesignsharp{ballot}{223}
\showcasematerialdesignsharp{bar-chart}{224}
\showcasematerialdesignsharp{batch-prediction}{225}
\showcasematerialdesignsharp{bathroom}{226}
\showcasematerialdesignsharp{bathtub}{227}
\showcasematerialdesignsharp{battery-0-bar}{228}
\showcasematerialdesignsharp{battery-1-bar}{229}
\showcasematerialdesignsharp{battery-2-bar}{230}
\showcasematerialdesignsharp{battery-3-bar}{231}
\showcasematerialdesignsharp{battery-4-bar}{232}
\showcasematerialdesignsharp{battery-5-bar}{233}
\showcasematerialdesignsharp{battery-6-bar}{234}
\showcasematerialdesignsharp{battery-alert}{235}
\showcasematerialdesignsharp{battery-charging-full}{236}
\showcasematerialdesignsharp{battery-full}{237}
\showcasematerialdesignsharp{battery-saver}{238}
\showcasematerialdesignsharp{battery-std}{239}
\showcasematerialdesignsharp{battery-unknown}{240}
\showcasematerialdesignsharp{beach-access}{241}
\showcasematerialdesignsharp{bed}{242}
\showcasematerialdesignsharp{bedroom-baby}{243}
\showcasematerialdesignsharp{bedroom-child}{244}
\showcasematerialdesignsharp{bedroom-parent}{245}
\showcasematerialdesignsharp{bedtime}{246}
\showcasematerialdesignsharp{bedtime-off}{247}
\showcasematerialdesignsharp{beenhere}{248}
\showcasematerialdesignsharp{bento}{249}
\showcasematerialdesignsharp{bike-scooter}{250}
\showcasematerialdesignsharp{biotech}{251}
\showcasematerialdesignsharp{blender}{252}
\showcasematerialdesignsharp{blind}{253}
\showcasematerialdesignsharp{blinds}{254}
\showcasematerialdesignsharp{blinds-closed}{255}
\showcasematerialdesignsharp{block}{256}
\showcasematerialdesignsharp{bloodtype}{257}
\showcasematerialdesignsharp{bluetooth}{258}
\showcasematerialdesignsharp{bluetooth-audio}{259}
\showcasematerialdesignsharp{bluetooth-connected}{260}
\showcasematerialdesignsharp{bluetooth-disabled}{261}
\showcasematerialdesignsharp{bluetooth-drive}{262}
\showcasematerialdesignsharp{bluetooth-searching}{263}
\showcasematerialdesignsharp{blur-circular}{264}
\showcasematerialdesignsharp{blur-linear}{265}
\showcasematerialdesignsharp{blur-off}{266}
\showcasematerialdesignsharp{blur-on}{267}
\showcasematerialdesignsharp{bolt}{268}
\showcasematerialdesignsharp{book}{269}
\showcasematerialdesignsharp{bookmark}{270}
\showcasematerialdesignsharp{bookmarks}{271}
\showcasematerialdesignsharp{bookmark-add}{272}
\showcasematerialdesignsharp{bookmark-added}{273}
\showcasematerialdesignsharp{bookmark-border}{274}
\showcasematerialdesignsharp{bookmark-remove}{275}
\showcasematerialdesignsharp{book-online}{276}
\showcasematerialdesignsharp{border-all}{277}
\showcasematerialdesignsharp{border-bottom}{278}
\showcasematerialdesignsharp{border-clear}{279}
\showcasematerialdesignsharp{border-color}{280}
\showcasematerialdesignsharp{border-horizontal}{281}
\showcasematerialdesignsharp{border-inner}{282}
\showcasematerialdesignsharp{border-left}{283}
\showcasematerialdesignsharp{border-outer}{284}
\showcasematerialdesignsharp{border-right}{285}
\showcasematerialdesignsharp{border-style}{286}
\showcasematerialdesignsharp{border-top}{287}
\showcasematerialdesignsharp{border-vertical}{288}
\showcasematerialdesignsharp{boy}{289}
\showcasematerialdesignsharp{branding-watermark}{290}
\showcasematerialdesignsharp{breakfast-dining}{291}
\showcasematerialdesignsharp{brightness-1}{292}
\showcasematerialdesignsharp{brightness-2}{293}
\showcasematerialdesignsharp{brightness-3}{294}
\showcasematerialdesignsharp{brightness-4}{295}
\showcasematerialdesignsharp{brightness-5}{296}
\showcasematerialdesignsharp{brightness-6}{297}
\showcasematerialdesignsharp{brightness-7}{298}
\showcasematerialdesignsharp{brightness-auto}{299}
\showcasematerialdesignsharp{brightness-high}{300}
\showcasematerialdesignsharp{brightness-low}{301}
\showcasematerialdesignsharp{brightness-medium}{302}
\showcasematerialdesignsharp{broadcast-on-home}{303}
\showcasematerialdesignsharp{broadcast-on-personal}{304}
\showcasematerialdesignsharp{broken-image}{305}
\showcasematerialdesignsharp{browser-not-supported}{306}
\showcasematerialdesignsharp{browser-updated}{307}
\showcasematerialdesignsharp{browse-gallery}{308}
\showcasematerialdesignsharp{brunch-dining}{309}
\showcasematerialdesignsharp{brush}{310}
\showcasematerialdesignsharp{bubble-chart}{311}
\showcasematerialdesignsharp{bug-report}{312}
\showcasematerialdesignsharp{build}{313}
\showcasematerialdesignsharp{build-circle}{314}
\showcasematerialdesignsharp{bungalow}{315}
\showcasematerialdesignsharp{burst-mode}{316}
\showcasematerialdesignsharp{business}{317}
\showcasematerialdesignsharp{business-center}{318}
\showcasematerialdesignsharp{bus-alert}{319}
\showcasematerialdesignsharp{cabin}{320}
\showcasematerialdesignsharp{cable}{321}
\showcasematerialdesignsharp{cached}{322}
\showcasematerialdesignsharp{cake}{323}
\showcasematerialdesignsharp{calculate}{324}
\showcasematerialdesignsharp{calendar-month}{325}
\showcasematerialdesignsharp{calendar-today}{326}
\showcasematerialdesignsharp{calendar-view-day}{327}
\showcasematerialdesignsharp{calendar-view-month}{328}
\showcasematerialdesignsharp{calendar-view-week}{329}
\showcasematerialdesignsharp{call}{330}
\showcasematerialdesignsharp{call-end}{331}
\showcasematerialdesignsharp{call-made}{332}
\showcasematerialdesignsharp{call-merge}{333}
\showcasematerialdesignsharp{call-missed}{334}
\showcasematerialdesignsharp{call-missed-outgoing}{335}
\showcasematerialdesignsharp{call-received}{336}
\showcasematerialdesignsharp{call-split}{337}
\showcasematerialdesignsharp{call-to-action}{338}
\showcasematerialdesignsharp{camera}{339}
\showcasematerialdesignsharp{cameraswitch}{340}
\showcasematerialdesignsharp{camera-alt}{341}
\showcasematerialdesignsharp{camera-enhance}{342}
\showcasematerialdesignsharp{camera-front}{343}
\showcasematerialdesignsharp{camera-indoor}{344}
\showcasematerialdesignsharp{camera-outdoor}{345}
\showcasematerialdesignsharp{camera-rear}{346}
\showcasematerialdesignsharp{camera-roll}{347}
\showcasematerialdesignsharp{campaign}{348}
\showcasematerialdesignsharp{cancel}{349}
\showcasematerialdesignsharp{cancel-presentation}{350}
\showcasematerialdesignsharp{cancel-schedule-send}{351}
\showcasematerialdesignsharp{candlestick-chart}{352}
\showcasematerialdesignsharp{card-giftcard}{353}
\showcasematerialdesignsharp{card-membership}{354}
\showcasematerialdesignsharp{card-travel}{355}
\showcasematerialdesignsharp{carpenter}{356}
\showcasematerialdesignsharp{car-crash}{357}
\showcasematerialdesignsharp{car-rental}{358}
\showcasematerialdesignsharp{car-repair}{359}
\showcasematerialdesignsharp{cases}{360}
\showcasematerialdesignsharp{casino}{361}
\showcasematerialdesignsharp{cast}{362}
\showcasematerialdesignsharp{castle}{363}
\showcasematerialdesignsharp{cast-connected}{364}
\showcasematerialdesignsharp{cast-for-education}{365}
\showcasematerialdesignsharp{catching-pokemon}{366}
\showcasematerialdesignsharp{category}{367}
\showcasematerialdesignsharp{celebration}{368}
\showcasematerialdesignsharp{cell-tower}{369}
\showcasematerialdesignsharp{cell-wifi}{370}
\showcasematerialdesignsharp{center-focus-strong}{371}
\showcasematerialdesignsharp{center-focus-weak}{372}
\showcasematerialdesignsharp{chair}{373}
\showcasematerialdesignsharp{chair-alt}{374}
\showcasematerialdesignsharp{chalet}{375}
\showcasematerialdesignsharp{change-circle}{376}
\showcasematerialdesignsharp{change-history}{377}
\showcasematerialdesignsharp{charging-station}{378}
\showcasematerialdesignsharp{chat}{379}
\showcasematerialdesignsharp{chat-bubble}{380}
\showcasematerialdesignsharp{chat-bubble-outline}{381}
\showcasematerialdesignsharp{check}{382}
\showcasematerialdesignsharp{checklist}{383}
\showcasematerialdesignsharp{checklist-rtl}{384}
\showcasematerialdesignsharp{checkroom}{385}
\showcasematerialdesignsharp{check-box}{386}
\showcasematerialdesignsharp{check-box-outline-blank}{387}
\showcasematerialdesignsharp{check-circle}{388}
\showcasematerialdesignsharp{check-circle-outline}{389}
\showcasematerialdesignsharp{chevron-left}{390}
\showcasematerialdesignsharp{chevron-right}{391}
\showcasematerialdesignsharp{child-care}{392}
\showcasematerialdesignsharp{child-friendly}{393}
\showcasematerialdesignsharp{chrome-reader-mode}{394}
\showcasematerialdesignsharp{church}{395}
\showcasematerialdesignsharp{circle}{396}
\showcasematerialdesignsharp{circle-notifications}{397}
\showcasematerialdesignsharp{class}{398}
\showcasematerialdesignsharp{cleaning-services}{399}
\showcasematerialdesignsharp{clean-hands}{400}
\showcasematerialdesignsharp{clear}{401}
\showcasematerialdesignsharp{clear-all}{402}
\showcasematerialdesignsharp{close}{403}
\showcasematerialdesignsharp{closed-caption}{404}
\showcasematerialdesignsharp{closed-caption-disabled}{405}
\showcasematerialdesignsharp{closed-caption-off}{406}
\showcasematerialdesignsharp{close-fullscreen}{407}
\showcasematerialdesignsharp{cloud}{408}
\showcasematerialdesignsharp{cloud-circle}{409}
\showcasematerialdesignsharp{cloud-done}{410}
\showcasematerialdesignsharp{cloud-download}{411}
\showcasematerialdesignsharp{cloud-off}{412}
\showcasematerialdesignsharp{cloud-queue}{413}
\showcasematerialdesignsharp{cloud-sync}{414}
\showcasematerialdesignsharp{cloud-upload}{415}
\showcasematerialdesignsharp{co2}{416}
\showcasematerialdesignsharp{code}{417}
\showcasematerialdesignsharp{code-off}{418}
\showcasematerialdesignsharp{coffee}{419}
\showcasematerialdesignsharp{coffee-maker}{420}
\showcasematerialdesignsharp{collections}{421}
\showcasematerialdesignsharp{collections-bookmark}{422}
\showcasematerialdesignsharp{colorize}{423}
\showcasematerialdesignsharp{color-lens}{424}
\showcasematerialdesignsharp{comment}{425}
\showcasematerialdesignsharp{comments-disabled}{426}
\showcasematerialdesignsharp{comment-bank}{427}
\showcasematerialdesignsharp{commit}{428}
\showcasematerialdesignsharp{commute}{429}
\showcasematerialdesignsharp{compare}{430}
\showcasematerialdesignsharp{compare-arrows}{431}
\showcasematerialdesignsharp{compass-calibration}{432}
\showcasematerialdesignsharp{compost}{433}
\showcasematerialdesignsharp{compress}{434}
\showcasematerialdesignsharp{computer}{435}
\showcasematerialdesignsharp{confirmation-number}{436}
\showcasematerialdesignsharp{connected-tv}{437}
\showcasematerialdesignsharp{connecting-airports}{438}
\showcasematerialdesignsharp{connect-without-contact}{439}
\showcasematerialdesignsharp{construction}{440}
\showcasematerialdesignsharp{contactless}{441}
\showcasematerialdesignsharp{contacts}{442}
\showcasematerialdesignsharp{contact-emergency}{443}
\showcasematerialdesignsharp{contact-mail}{444}
\showcasematerialdesignsharp{contact-page}{445}
\showcasematerialdesignsharp{contact-phone}{446}
\showcasematerialdesignsharp{contact-support}{447}
\showcasematerialdesignsharp{content-copy}{448}
\showcasematerialdesignsharp{content-cut}{449}
\showcasematerialdesignsharp{content-paste}{450}
\showcasematerialdesignsharp{content-paste-go}{451}
\showcasematerialdesignsharp{content-paste-off}{452}
\showcasematerialdesignsharp{content-paste-search}{453}
\showcasematerialdesignsharp{contrast}{454}
\showcasematerialdesignsharp{control-camera}{455}
\showcasematerialdesignsharp{control-point}{456}
\showcasematerialdesignsharp{control-point-duplicate}{457}
\showcasematerialdesignsharp{cookie}{458}
\showcasematerialdesignsharp{copyright}{459}
\showcasematerialdesignsharp{copy-all}{460}
\showcasematerialdesignsharp{coronavirus}{461}
\showcasematerialdesignsharp{corporate-fare}{462}
\showcasematerialdesignsharp{cottage}{463}
\showcasematerialdesignsharp{countertops}{464}
\showcasematerialdesignsharp{co-present}{465}
\showcasematerialdesignsharp{create}{466}
\showcasematerialdesignsharp{create-new-folder}{467}
\showcasematerialdesignsharp{credit-card}{468}
\showcasematerialdesignsharp{credit-card-off}{469}
\showcasematerialdesignsharp{credit-score}{470}
\showcasematerialdesignsharp{crib}{471}
\showcasematerialdesignsharp{crisis-alert}{472}
\showcasematerialdesignsharp{crop}{473}
\showcasematerialdesignsharp{crop-3-2}{474}
\showcasematerialdesignsharp{crop-5-4}{475}
\showcasematerialdesignsharp{crop-7-5}{476}
\showcasematerialdesignsharp{crop-16-9}{477}
\showcasematerialdesignsharp{crop-din}{478}
\showcasematerialdesignsharp{crop-free}{479}
\showcasematerialdesignsharp{crop-landscape}{480}
\showcasematerialdesignsharp{crop-original}{481}
\showcasematerialdesignsharp{crop-portrait}{482}
\showcasematerialdesignsharp{crop-rotate}{483}
\showcasematerialdesignsharp{crop-square}{484}
\showcasematerialdesignsharp{cruelty-free}{485}
\showcasematerialdesignsharp{css}{486}
\showcasematerialdesignsharp{currency-bitcoin}{487}
\showcasematerialdesignsharp{currency-exchange}{488}
\showcasematerialdesignsharp{currency-franc}{489}
\showcasematerialdesignsharp{currency-lira}{490}
\showcasematerialdesignsharp{currency-pound}{491}
\showcasematerialdesignsharp{currency-ruble}{492}
\showcasematerialdesignsharp{currency-rupee}{493}
\showcasematerialdesignsharp{currency-yen}{494}
\showcasematerialdesignsharp{currency-yuan}{495}
\showcasematerialdesignsharp{curtains}{496}
\showcasematerialdesignsharp{curtains-closed}{497}
\showcasematerialdesignsharp{cyclone}{498}
\showcasematerialdesignsharp{dangerous}{499}
\showcasematerialdesignsharp{dark-mode}{500}
\showcasematerialdesignsharp{dashboard}{501}
\showcasematerialdesignsharp{dashboard-customize}{502}
\showcasematerialdesignsharp{dataset}{503}
\showcasematerialdesignsharp{dataset-linked}{504}
\showcasematerialdesignsharp{data-array}{505}
\showcasematerialdesignsharp{data-exploration}{506}
\showcasematerialdesignsharp{data-object}{507}
\showcasematerialdesignsharp{data-saver-off}{508}
\showcasematerialdesignsharp{data-saver-on}{509}
\showcasematerialdesignsharp{data-thresholding}{510}
\showcasematerialdesignsharp{data-usage}{511}
\showcasematerialdesignsharp{date-range}{512}
\showcasematerialdesignsharp{deblur}{513}
\showcasematerialdesignsharp{deck}{514}
\showcasematerialdesignsharp{dehaze}{515}
\showcasematerialdesignsharp{delete}{516}
\showcasematerialdesignsharp{delete-forever}{517}
\showcasematerialdesignsharp{delete-outline}{518}
\showcasematerialdesignsharp{delete-sweep}{519}
\showcasematerialdesignsharp{delivery-dining}{520}
\showcasematerialdesignsharp{density-large}{521}
\showcasematerialdesignsharp{density-medium}{522}
\showcasematerialdesignsharp{density-small}{523}
\showcasematerialdesignsharp{departure-board}{524}
\showcasematerialdesignsharp{description}{525}
\showcasematerialdesignsharp{deselect}{526}
\showcasematerialdesignsharp{design-services}{527}
\showcasematerialdesignsharp{desk}{528}
\showcasematerialdesignsharp{desktop-access-disabled}{529}
\showcasematerialdesignsharp{desktop-mac}{530}
\showcasematerialdesignsharp{desktop-windows}{531}
\showcasematerialdesignsharp{details}{532}
\showcasematerialdesignsharp{developer-board}{533}
\showcasematerialdesignsharp{developer-board-off}{534}
\showcasematerialdesignsharp{developer-mode}{535}
\showcasematerialdesignsharp{devices}{536}
\showcasematerialdesignsharp{devices-fold}{537}
\showcasematerialdesignsharp{devices-other}{538}
\showcasematerialdesignsharp{device-hub}{539}
\showcasematerialdesignsharp{device-thermostat}{540}
\showcasematerialdesignsharp{device-unknown}{541}
\showcasematerialdesignsharp{dialer-sip}{542}
\showcasematerialdesignsharp{dialpad}{543}
\showcasematerialdesignsharp{diamond}{544}
\showcasematerialdesignsharp{difference}{545}
\showcasematerialdesignsharp{dining}{546}
\showcasematerialdesignsharp{dinner-dining}{547}
\showcasematerialdesignsharp{directions}{548}
\showcasematerialdesignsharp{directions-bike}{549}
\showcasematerialdesignsharp{directions-boat}{550}
\showcasematerialdesignsharp{directions-boat-filled}{551}
\showcasematerialdesignsharp{directions-bus}{552}
\showcasematerialdesignsharp{directions-bus-filled}{553}
\showcasematerialdesignsharp{directions-car}{554}
\showcasematerialdesignsharp{directions-car-filled}{555}
\showcasematerialdesignsharp{directions-off}{556}
\showcasematerialdesignsharp{directions-railway}{557}
\showcasematerialdesignsharp{directions-railway-filled}{558}
\showcasematerialdesignsharp{directions-run}{559}
\showcasematerialdesignsharp{directions-subway}{560}
\showcasematerialdesignsharp{directions-subway-filled}{561}
\showcasematerialdesignsharp{directions-transit}{562}
\showcasematerialdesignsharp{directions-transit-filled}{563}
\showcasematerialdesignsharp{directions-walk}{564}
\showcasematerialdesignsharp{dirty-lens}{565}
\showcasematerialdesignsharp{disabled-by-default}{566}
\showcasematerialdesignsharp{disabled-visible}{567}
\showcasematerialdesignsharp{discount}{568}
\showcasematerialdesignsharp{disc-full}{569}
\showcasematerialdesignsharp{display-settings}{570}
\showcasematerialdesignsharp{diversity-1}{571}
\showcasematerialdesignsharp{diversity-2}{572}
\showcasematerialdesignsharp{diversity-3}{573}
\showcasematerialdesignsharp{dns}{574}
\showcasematerialdesignsharp{dock}{575}
\showcasematerialdesignsharp{document-scanner}{576}
\showcasematerialdesignsharp{domain}{577}
\showcasematerialdesignsharp{domain-add}{578}
\showcasematerialdesignsharp{domain-disabled}{579}
\showcasematerialdesignsharp{domain-verification}{580}
\showcasematerialdesignsharp{done}{581}
\showcasematerialdesignsharp{done-all}{582}
\showcasematerialdesignsharp{done-outline}{583}
\showcasematerialdesignsharp{donut-large}{584}
\showcasematerialdesignsharp{donut-small}{585}
\showcasematerialdesignsharp{doorbell}{586}
\showcasematerialdesignsharp{door-back}{587}
\showcasematerialdesignsharp{door-front}{588}
\showcasematerialdesignsharp{door-sliding}{589}
\showcasematerialdesignsharp{double-arrow}{590}
\showcasematerialdesignsharp{downhill-skiing}{591}
\showcasematerialdesignsharp{download}{592}
\showcasematerialdesignsharp{downloading}{593}
\showcasematerialdesignsharp{download-done}{594}
\showcasematerialdesignsharp{download-for-offline}{595}
\showcasematerialdesignsharp{do-disturb}{596}
\showcasematerialdesignsharp{do-disturb-alt}{597}
\showcasematerialdesignsharp{do-disturb-off}{598}
\showcasematerialdesignsharp{do-disturb-on}{599}
\showcasematerialdesignsharp{do-not-disturb}{600}
\showcasematerialdesignsharp{do-not-disturb-alt}{601}
\showcasematerialdesignsharp{do-not-disturb-off}{602}
\showcasematerialdesignsharp{do-not-disturb-on}{603}
\showcasematerialdesignsharp{do-not-disturb-on-total-silence}{604}
\showcasematerialdesignsharp{do-not-step}{605}
\showcasematerialdesignsharp{do-not-touch}{606}
\showcasematerialdesignsharp{drafts}{607}
\showcasematerialdesignsharp{drag-handle}{608}
\showcasematerialdesignsharp{drag-indicator}{609}
\showcasematerialdesignsharp{draw}{610}
\showcasematerialdesignsharp{drive-eta}{611}
\showcasematerialdesignsharp{drive-file-move}{612}
\showcasematerialdesignsharp{drive-file-move-rtl}{613}
\showcasematerialdesignsharp{drive-file-rename-outline}{614}
\showcasematerialdesignsharp{drive-folder-upload}{615}
\showcasematerialdesignsharp{dry}{616}
\showcasematerialdesignsharp{dry-cleaning}{617}
\showcasematerialdesignsharp{duo}{618}
\showcasematerialdesignsharp{dvr}{619}
\showcasematerialdesignsharp{dynamic-feed}{620}
\showcasematerialdesignsharp{dynamic-form}{621}
\showcasematerialdesignsharp{earbuds}{622}
\showcasematerialdesignsharp{earbuds-battery}{623}
\showcasematerialdesignsharp{east}{624}
\showcasematerialdesignsharp{edgesensor-high}{625}
\showcasematerialdesignsharp{edgesensor-low}{626}
\showcasematerialdesignsharp{edit}{627}
\showcasematerialdesignsharp{edit-attributes}{628}
\showcasematerialdesignsharp{edit-calendar}{629}
\showcasematerialdesignsharp{edit-location}{630}
\showcasematerialdesignsharp{edit-location-alt}{631}
\showcasematerialdesignsharp{edit-note}{632}
\showcasematerialdesignsharp{edit-notifications}{633}
\showcasematerialdesignsharp{edit-off}{634}
\showcasematerialdesignsharp{edit-road}{635}
\showcasematerialdesignsharp{egg}{636}
\showcasematerialdesignsharp{egg-alt}{637}
\showcasematerialdesignsharp{eject}{638}
\showcasematerialdesignsharp{elderly}{639}
\showcasematerialdesignsharp{elderly-woman}{640}
\showcasematerialdesignsharp{electrical-services}{641}
\showcasematerialdesignsharp{electric-bike}{642}
\showcasematerialdesignsharp{electric-bolt}{643}
\showcasematerialdesignsharp{electric-car}{644}
\showcasematerialdesignsharp{electric-meter}{645}
\showcasematerialdesignsharp{electric-moped}{646}
\showcasematerialdesignsharp{electric-rickshaw}{647}
\showcasematerialdesignsharp{electric-scooter}{648}
\showcasematerialdesignsharp{elevator}{649}
\showcasematerialdesignsharp{email}{650}
\showcasematerialdesignsharp{emergency}{651}
\showcasematerialdesignsharp{emergency-recording}{652}
\showcasematerialdesignsharp{emergency-share}{653}
\showcasematerialdesignsharp{emoji-emotions}{654}
\showcasematerialdesignsharp{emoji-events}{655}
\showcasematerialdesignsharp{emoji-food-beverage}{656}
\showcasematerialdesignsharp{emoji-nature}{657}
\showcasematerialdesignsharp{emoji-objects}{658}
\showcasematerialdesignsharp{emoji-people}{659}
\showcasematerialdesignsharp{emoji-symbols}{660}
\showcasematerialdesignsharp{emoji-transportation}{661}
\showcasematerialdesignsharp{energy-savings-leaf}{662}
\showcasematerialdesignsharp{engineering}{663}
\showcasematerialdesignsharp{enhanced-encryption}{664}
\showcasematerialdesignsharp{equalizer}{665}
\showcasematerialdesignsharp{error}{666}
\showcasematerialdesignsharp{error-outline}{667}
\showcasematerialdesignsharp{escalator}{668}
\showcasematerialdesignsharp{escalator-warning}{669}
\showcasematerialdesignsharp{euro}{670}
\showcasematerialdesignsharp{euro-symbol}{671}
\showcasematerialdesignsharp{event}{672}
\showcasematerialdesignsharp{event-available}{673}
\showcasematerialdesignsharp{event-busy}{674}
\showcasematerialdesignsharp{event-note}{675}
\showcasematerialdesignsharp{event-repeat}{676}
\showcasematerialdesignsharp{event-seat}{677}
\showcasematerialdesignsharp{ev-station}{678}
\showcasematerialdesignsharp{exit-to-app}{679}
\showcasematerialdesignsharp{expand}{680}
\showcasematerialdesignsharp{expand-circle-down}{681}
\showcasematerialdesignsharp{expand-less}{682}
\showcasematerialdesignsharp{expand-more}{683}
\showcasematerialdesignsharp{explicit}{684}
\showcasematerialdesignsharp{explore}{685}
\showcasematerialdesignsharp{explore-off}{686}
\showcasematerialdesignsharp{exposure}{687}
\showcasematerialdesignsharp{exposure-neg-1}{688}
\showcasematerialdesignsharp{exposure-neg-2}{689}
\showcasematerialdesignsharp{exposure-plus-1}{690}
\showcasematerialdesignsharp{exposure-plus-2}{691}
\showcasematerialdesignsharp{exposure-zero}{692}
\showcasematerialdesignsharp{extension}{693}
\showcasematerialdesignsharp{extension-off}{694}
\showcasematerialdesignsharp{e-mobiledata}{695}
\showcasematerialdesignsharp{face}{696}
\showcasematerialdesignsharp{face-2}{697}
\showcasematerialdesignsharp{face-3}{698}
\showcasematerialdesignsharp{face-4}{699}
\showcasematerialdesignsharp{face-5}{700}
\showcasematerialdesignsharp{face-6}{701}
\showcasematerialdesignsharp{face-retouching-natural}{702}
\showcasematerialdesignsharp{face-retouching-off}{703}
\showcasematerialdesignsharp{factory}{704}
\showcasematerialdesignsharp{fact-check}{705}
\showcasematerialdesignsharp{family-restroom}{706}
\showcasematerialdesignsharp{fastfood}{707}
\showcasematerialdesignsharp{fast-forward}{708}
\showcasematerialdesignsharp{fast-rewind}{709}
\showcasematerialdesignsharp{favorite}{710}
\showcasematerialdesignsharp{favorite-border}{711}
\showcasematerialdesignsharp{fax}{712}
\showcasematerialdesignsharp{featured-play-list}{713}
\showcasematerialdesignsharp{featured-video}{714}
\showcasematerialdesignsharp{feed}{715}
\showcasematerialdesignsharp{feedback}{716}
\showcasematerialdesignsharp{female}{717}
\showcasematerialdesignsharp{fence}{718}
\showcasematerialdesignsharp{festival}{719}
\showcasematerialdesignsharp{fiber-dvr}{720}
\showcasematerialdesignsharp{fiber-manual-record}{721}
\showcasematerialdesignsharp{fiber-new}{722}
\showcasematerialdesignsharp{fiber-pin}{723}
\showcasematerialdesignsharp{fiber-smart-record}{724}
\showcasematerialdesignsharp{file-copy}{725}
\showcasematerialdesignsharp{file-download}{726}
\showcasematerialdesignsharp{file-download-done}{727}
\showcasematerialdesignsharp{file-download-off}{728}
\showcasematerialdesignsharp{file-open}{729}
\showcasematerialdesignsharp{file-present}{730}
\showcasematerialdesignsharp{file-upload}{731}
\showcasematerialdesignsharp{filter}{732}
\showcasematerialdesignsharp{filter-1}{733}
\showcasematerialdesignsharp{filter-2}{734}
\showcasematerialdesignsharp{filter-3}{735}
\showcasematerialdesignsharp{filter-4}{736}
\showcasematerialdesignsharp{filter-5}{737}
\showcasematerialdesignsharp{filter-6}{738}
\showcasematerialdesignsharp{filter-7}{739}
\showcasematerialdesignsharp{filter-8}{740}
\showcasematerialdesignsharp{filter-9}{741}
\showcasematerialdesignsharp{filter-9-plus}{742}
\showcasematerialdesignsharp{filter-alt}{743}
\showcasematerialdesignsharp{filter-alt-off}{744}
\showcasematerialdesignsharp{filter-b-and-w}{745}
\showcasematerialdesignsharp{filter-center-focus}{746}
\showcasematerialdesignsharp{filter-drama}{747}
\showcasematerialdesignsharp{filter-frames}{748}
\showcasematerialdesignsharp{filter-hdr}{749}
\showcasematerialdesignsharp{filter-list}{750}
\showcasematerialdesignsharp{filter-list-off}{751}
\showcasematerialdesignsharp{filter-none}{752}
\showcasematerialdesignsharp{filter-tilt-shift}{753}
\showcasematerialdesignsharp{filter-vintage}{754}
\showcasematerialdesignsharp{find-in-page}{755}
\showcasematerialdesignsharp{find-replace}{756}
\showcasematerialdesignsharp{fingerprint}{757}
\showcasematerialdesignsharp{fireplace}{758}
\showcasematerialdesignsharp{fire-extinguisher}{759}
\showcasematerialdesignsharp{fire-hydrant-alt}{760}
\showcasematerialdesignsharp{fire-truck}{761}
\showcasematerialdesignsharp{first-page}{762}
\showcasematerialdesignsharp{fitbit}{763}
\showcasematerialdesignsharp{fitness-center}{764}
\showcasematerialdesignsharp{fit-screen}{765}
\showcasematerialdesignsharp{flag}{766}
\showcasematerialdesignsharp{flag-circle}{767}
\showcasematerialdesignsharp{flaky}{768}
\showcasematerialdesignsharp{flare}{769}
\showcasematerialdesignsharp{flashlight-off}{770}
\showcasematerialdesignsharp{flashlight-on}{771}
\showcasematerialdesignsharp{flash-auto}{772}
\showcasematerialdesignsharp{flash-off}{773}
\showcasematerialdesignsharp{flash-on}{774}
\showcasematerialdesignsharp{flatware}{775}
\showcasematerialdesignsharp{flight}{776}
\showcasematerialdesignsharp{flight-class}{777}
\showcasematerialdesignsharp{flight-land}{778}
\showcasematerialdesignsharp{flight-takeoff}{779}
\showcasematerialdesignsharp{flip}{780}
\showcasematerialdesignsharp{flip-camera-android}{781}
\showcasematerialdesignsharp{flip-camera-ios}{782}
\showcasematerialdesignsharp{flip-to-back}{783}
\showcasematerialdesignsharp{flip-to-front}{784}
\showcasematerialdesignsharp{flood}{785}
\showcasematerialdesignsharp{fluorescent}{786}
\showcasematerialdesignsharp{flutter-dash}{787}
\showcasematerialdesignsharp{fmd-bad}{788}
\showcasematerialdesignsharp{fmd-good}{789}
\showcasematerialdesignsharp{folder}{790}
\showcasematerialdesignsharp{folder-copy}{791}
\showcasematerialdesignsharp{folder-delete}{792}
\showcasematerialdesignsharp{folder-off}{793}
\showcasematerialdesignsharp{folder-open}{794}
\showcasematerialdesignsharp{folder-shared}{795}
\showcasematerialdesignsharp{folder-special}{796}
\showcasematerialdesignsharp{folder-zip}{797}
\showcasematerialdesignsharp{follow-the-signs}{798}
\showcasematerialdesignsharp{font-download}{799}
\showcasematerialdesignsharp{font-download-off}{800}
\showcasematerialdesignsharp{food-bank}{801}
\showcasematerialdesignsharp{forest}{802}
\showcasematerialdesignsharp{fork-left}{803}
\showcasematerialdesignsharp{fork-right}{804}
\showcasematerialdesignsharp{format-align-center}{805}
\showcasematerialdesignsharp{format-align-justify}{806}
\showcasematerialdesignsharp{format-align-left}{807}
\showcasematerialdesignsharp{format-align-right}{808}
\showcasematerialdesignsharp{format-bold}{809}
\showcasematerialdesignsharp{format-clear}{810}
\showcasematerialdesignsharp{format-color-fill}{811}
\showcasematerialdesignsharp{format-color-reset}{812}
\showcasematerialdesignsharp{format-color-text}{813}
\showcasematerialdesignsharp{format-indent-decrease}{814}
\showcasematerialdesignsharp{format-indent-increase}{815}
\showcasematerialdesignsharp{format-italic}{816}
\showcasematerialdesignsharp{format-line-spacing}{817}
\showcasematerialdesignsharp{format-list-bulleted}{818}
\showcasematerialdesignsharp{format-list-numbered}{819}
\showcasematerialdesignsharp{format-list-numbered-rtl}{820}
\showcasematerialdesignsharp{format-overline}{821}
\showcasematerialdesignsharp{format-paint}{822}
\showcasematerialdesignsharp{format-quote}{823}
\showcasematerialdesignsharp{format-shapes}{824}
\showcasematerialdesignsharp{format-size}{825}
\showcasematerialdesignsharp{format-strikethrough}{826}
\showcasematerialdesignsharp{format-textdirection-l-to-r}{827}
\showcasematerialdesignsharp{format-textdirection-r-to-l}{828}
\showcasematerialdesignsharp{format-underlined}{829}
\showcasematerialdesignsharp{fort}{830}
\showcasematerialdesignsharp{forum}{831}
\showcasematerialdesignsharp{forward}{832}
\showcasematerialdesignsharp{forward-5}{833}
\showcasematerialdesignsharp{forward-10}{834}
\showcasematerialdesignsharp{forward-30}{835}
\showcasematerialdesignsharp{forward-to-inbox}{836}
\showcasematerialdesignsharp{foundation}{837}
\showcasematerialdesignsharp{free-breakfast}{838}
\showcasematerialdesignsharp{free-cancellation}{839}
\showcasematerialdesignsharp{front-hand}{840}
\showcasematerialdesignsharp{fullscreen}{841}
\showcasematerialdesignsharp{fullscreen-exit}{842}
\showcasematerialdesignsharp{functions}{843}
\showcasematerialdesignsharp{gamepad}{844}
\showcasematerialdesignsharp{games}{845}
\showcasematerialdesignsharp{garage}{846}
\showcasematerialdesignsharp{gas-meter}{847}
\showcasematerialdesignsharp{gavel}{848}
\showcasematerialdesignsharp{generating-tokens}{849}
\showcasematerialdesignsharp{gesture}{850}
\showcasematerialdesignsharp{get-app}{851}
\showcasematerialdesignsharp{gif}{852}
\showcasematerialdesignsharp{gif-box}{853}
\showcasematerialdesignsharp{girl}{854}
\showcasematerialdesignsharp{gite}{855}
\showcasematerialdesignsharp{golf-course}{856}
\showcasematerialdesignsharp{gpp-bad}{857}
\showcasematerialdesignsharp{gpp-good}{858}
\showcasematerialdesignsharp{gpp-maybe}{859}
\showcasematerialdesignsharp{gps-fixed}{860}
\showcasematerialdesignsharp{gps-not-fixed}{861}
\showcasematerialdesignsharp{gps-off}{862}
\showcasematerialdesignsharp{grade}{863}
\showcasematerialdesignsharp{gradient}{864}
\showcasematerialdesignsharp{grading}{865}
\showcasematerialdesignsharp{grain}{866}
\showcasematerialdesignsharp{graphic-eq}{867}
\showcasematerialdesignsharp{grass}{868}
\showcasematerialdesignsharp{grid-3x3}{869}
\showcasematerialdesignsharp{grid-4x4}{870}
\showcasematerialdesignsharp{grid-goldenratio}{871}
\showcasematerialdesignsharp{grid-off}{872}
\showcasematerialdesignsharp{grid-on}{873}
\showcasematerialdesignsharp{grid-view}{874}
\showcasematerialdesignsharp{group}{875}
\showcasematerialdesignsharp{groups}{876}
\showcasematerialdesignsharp{groups-2}{877}
\showcasematerialdesignsharp{groups-3}{878}
\showcasematerialdesignsharp{group-add}{879}
\showcasematerialdesignsharp{group-off}{880}
\showcasematerialdesignsharp{group-remove}{881}
\showcasematerialdesignsharp{group-work}{882}
\showcasematerialdesignsharp{g-mobiledata}{883}
\showcasematerialdesignsharp{g-translate}{884}
\showcasematerialdesignsharp{hail}{885}
\showcasematerialdesignsharp{handshake}{886}
\showcasematerialdesignsharp{handyman}{887}
\showcasematerialdesignsharp{hardware}{888}
\showcasematerialdesignsharp{hd}{889}
\showcasematerialdesignsharp{hdr-auto}{890}
\showcasematerialdesignsharp{hdr-auto-select}{891}
\showcasematerialdesignsharp{hdr-enhanced-select}{892}
\showcasematerialdesignsharp{hdr-off}{893}
\showcasematerialdesignsharp{hdr-off-select}{894}
\showcasematerialdesignsharp{hdr-on}{895}
\showcasematerialdesignsharp{hdr-on-select}{896}
\showcasematerialdesignsharp{hdr-plus}{897}
\showcasematerialdesignsharp{hdr-strong}{898}
\showcasematerialdesignsharp{hdr-weak}{899}
\showcasematerialdesignsharp{headphones}{900}
\showcasematerialdesignsharp{headphones-battery}{901}
\showcasematerialdesignsharp{headset}{902}
\showcasematerialdesignsharp{headset-mic}{903}
\showcasematerialdesignsharp{headset-off}{904}
\showcasematerialdesignsharp{healing}{905}
\showcasematerialdesignsharp{health-and-safety}{906}
\showcasematerialdesignsharp{hearing}{907}
\showcasematerialdesignsharp{hearing-disabled}{908}
\showcasematerialdesignsharp{heart-broken}{909}
\showcasematerialdesignsharp{heat-pump}{910}
\showcasematerialdesignsharp{height}{911}
\showcasematerialdesignsharp{help}{912}
\showcasematerialdesignsharp{help-center}{913}
\showcasematerialdesignsharp{help-outline}{914}
\showcasematerialdesignsharp{hevc}{915}
\showcasematerialdesignsharp{hexagon}{916}
\showcasematerialdesignsharp{hide-image}{917}
\showcasematerialdesignsharp{hide-source}{918}
\showcasematerialdesignsharp{highlight}{919}
\showcasematerialdesignsharp{highlight-alt}{920}
\showcasematerialdesignsharp{highlight-off}{921}
\showcasematerialdesignsharp{high-quality}{922}
\showcasematerialdesignsharp{hiking}{923}
\showcasematerialdesignsharp{history}{924}
\showcasematerialdesignsharp{history-edu}{925}
\showcasematerialdesignsharp{history-toggle-off}{926}
\showcasematerialdesignsharp{hive}{927}
\showcasematerialdesignsharp{hls}{928}
\showcasematerialdesignsharp{hls-off}{929}
\showcasematerialdesignsharp{holiday-village}{930}
\showcasematerialdesignsharp{home}{931}
\showcasematerialdesignsharp{home-max}{932}
\showcasematerialdesignsharp{home-mini}{933}
\showcasematerialdesignsharp{home-repair-service}{934}
\showcasematerialdesignsharp{home-work}{935}
\showcasematerialdesignsharp{horizontal-distribute}{936}
\showcasematerialdesignsharp{horizontal-rule}{937}
\showcasematerialdesignsharp{horizontal-split}{938}
\showcasematerialdesignsharp{hotel}{939}
\showcasematerialdesignsharp{hotel-class}{940}
\showcasematerialdesignsharp{hot-tub}{941}
\showcasematerialdesignsharp{hourglass-bottom}{942}
\showcasematerialdesignsharp{hourglass-disabled}{943}
\showcasematerialdesignsharp{hourglass-empty}{944}
\showcasematerialdesignsharp{hourglass-full}{945}
\showcasematerialdesignsharp{hourglass-top}{946}
\showcasematerialdesignsharp{house}{947}
\showcasematerialdesignsharp{houseboat}{948}
\showcasematerialdesignsharp{house-siding}{949}
\showcasematerialdesignsharp{how-to-reg}{950}
\showcasematerialdesignsharp{how-to-vote}{951}
\showcasematerialdesignsharp{html}{952}
\showcasematerialdesignsharp{http}{953}
\showcasematerialdesignsharp{https}{954}
\showcasematerialdesignsharp{hub}{955}
\showcasematerialdesignsharp{hvac}{956}
\showcasematerialdesignsharp{h-mobiledata}{957}
\showcasematerialdesignsharp{h-plus-mobiledata}{958}
\showcasematerialdesignsharp{icecream}{959}
\showcasematerialdesignsharp{ice-skating}{960}
\showcasematerialdesignsharp{image}{961}
\showcasematerialdesignsharp{imagesearch-roller}{962}
\showcasematerialdesignsharp{image-aspect-ratio}{963}
\showcasematerialdesignsharp{image-not-supported}{964}
\showcasematerialdesignsharp{image-search}{965}
\showcasematerialdesignsharp{important-devices}{966}
\showcasematerialdesignsharp{import-contacts}{967}
\showcasematerialdesignsharp{import-export}{968}
\showcasematerialdesignsharp{inbox}{969}
\showcasematerialdesignsharp{incomplete-circle}{970}
\showcasematerialdesignsharp{indeterminate-check-box}{971}
\showcasematerialdesignsharp{info}{972}
\showcasematerialdesignsharp{input}{973}
\showcasematerialdesignsharp{insert-chart}{974}
\showcasematerialdesignsharp{insert-chart-outlined}{975}
\showcasematerialdesignsharp{insert-comment}{976}
\showcasematerialdesignsharp{insert-drive-file}{977}
\showcasematerialdesignsharp{insert-emoticon}{978}
\showcasematerialdesignsharp{insert-invitation}{979}
\showcasematerialdesignsharp{insert-link}{980}
\showcasematerialdesignsharp{insert-page-break}{981}
\showcasematerialdesignsharp{insert-photo}{982}
\showcasematerialdesignsharp{insights}{983}
\showcasematerialdesignsharp{install-desktop}{984}
\showcasematerialdesignsharp{install-mobile}{985}
\showcasematerialdesignsharp{integration-instructions}{986}
\showcasematerialdesignsharp{interests}{987}
\showcasematerialdesignsharp{interpreter-mode}{988}
\showcasematerialdesignsharp{inventory}{989}
\showcasematerialdesignsharp{inventory-2}{990}
\showcasematerialdesignsharp{invert-colors}{991}
\showcasematerialdesignsharp{invert-colors-off}{992}
\showcasematerialdesignsharp{ios-share}{993}
\showcasematerialdesignsharp{iron}{994}
\showcasematerialdesignsharp{iso}{995}
\showcasematerialdesignsharp{javascript}{996}
\showcasematerialdesignsharp{join-full}{997}
\showcasematerialdesignsharp{join-inner}{998}
\showcasematerialdesignsharp{join-left}{999}
\showcasematerialdesignsharp{join-right}{1000}
\showcasematerialdesignsharp{kayaking}{1001}
\showcasematerialdesignsharp{kebab-dining}{1002}
\showcasematerialdesignsharp{key}{1003}
\showcasematerialdesignsharp{keyboard}{1004}
\showcasematerialdesignsharp{keyboard-alt}{1005}
\showcasematerialdesignsharp{keyboard-arrow-down}{1006}
\showcasematerialdesignsharp{keyboard-arrow-left}{1007}
\showcasematerialdesignsharp{keyboard-arrow-right}{1008}
\showcasematerialdesignsharp{keyboard-arrow-up}{1009}
\showcasematerialdesignsharp{keyboard-backspace}{1010}
\showcasematerialdesignsharp{keyboard-capslock}{1011}
\showcasematerialdesignsharp{keyboard-command-key}{1012}
\showcasematerialdesignsharp{keyboard-control-key}{1013}
\showcasematerialdesignsharp{keyboard-double-arrow-down}{1014}
\showcasematerialdesignsharp{keyboard-double-arrow-left}{1015}
\showcasematerialdesignsharp{keyboard-double-arrow-right}{1016}
\showcasematerialdesignsharp{keyboard-double-arrow-up}{1017}
\showcasematerialdesignsharp{keyboard-hide}{1018}
\showcasematerialdesignsharp{keyboard-option-key}{1019}
\showcasematerialdesignsharp{keyboard-return}{1020}
\showcasematerialdesignsharp{keyboard-tab}{1021}
\showcasematerialdesignsharp{keyboard-voice}{1022}
\showcasematerialdesignsharp{key-off}{1023}
\showcasematerialdesignsharp{king-bed}{1024}
\showcasematerialdesignsharp{kitchen}{1025}
\showcasematerialdesignsharp{kitesurfing}{1026}
\showcasematerialdesignsharp{label}{1027}
\showcasematerialdesignsharp{label-important}{1028}
\showcasematerialdesignsharp{label-off}{1029}
\showcasematerialdesignsharp{lan}{1030}
\showcasematerialdesignsharp{landscape}{1031}
\showcasematerialdesignsharp{landslide}{1032}
\showcasematerialdesignsharp{language}{1033}
\showcasematerialdesignsharp{laptop}{1034}
\showcasematerialdesignsharp{laptop-chromebook}{1035}
\showcasematerialdesignsharp{laptop-mac}{1036}
\showcasematerialdesignsharp{laptop-windows}{1037}
\showcasematerialdesignsharp{last-page}{1038}
\showcasematerialdesignsharp{launch}{1039}
\showcasematerialdesignsharp{layers}{1040}
\showcasematerialdesignsharp{layers-clear}{1041}
\showcasematerialdesignsharp{leaderboard}{1042}
\showcasematerialdesignsharp{leak-add}{1043}
\showcasematerialdesignsharp{leak-remove}{1044}
\showcasematerialdesignsharp{legend-toggle}{1045}
\showcasematerialdesignsharp{lens}{1046}
\showcasematerialdesignsharp{lens-blur}{1047}
\showcasematerialdesignsharp{library-add}{1048}
\showcasematerialdesignsharp{library-add-check}{1049}
\showcasematerialdesignsharp{library-books}{1050}
\showcasematerialdesignsharp{library-music}{1051}
\showcasematerialdesignsharp{light}{1052}
\showcasematerialdesignsharp{lightbulb}{1053}
\showcasematerialdesignsharp{lightbulb-circle}{1054}
\showcasematerialdesignsharp{light-mode}{1055}
\showcasematerialdesignsharp{linear-scale}{1056}
\showcasematerialdesignsharp{line-axis}{1057}
\showcasematerialdesignsharp{line-style}{1058}
\showcasematerialdesignsharp{line-weight}{1059}
\showcasematerialdesignsharp{link}{1060}
\showcasematerialdesignsharp{linked-camera}{1061}
\showcasematerialdesignsharp{link-off}{1062}
\showcasematerialdesignsharp{liquor}{1063}
\showcasematerialdesignsharp{list}{1064}
\showcasematerialdesignsharp{list-alt}{1065}
\showcasematerialdesignsharp{live-help}{1066}
\showcasematerialdesignsharp{live-tv}{1067}
\showcasematerialdesignsharp{living}{1068}
\showcasematerialdesignsharp{local-activity}{1069}
\showcasematerialdesignsharp{local-airport}{1070}
\showcasematerialdesignsharp{local-atm}{1071}
\showcasematerialdesignsharp{local-bar}{1072}
\showcasematerialdesignsharp{local-cafe}{1073}
\showcasematerialdesignsharp{local-car-wash}{1074}
\showcasematerialdesignsharp{local-convenience-store}{1075}
\showcasematerialdesignsharp{local-dining}{1076}
\showcasematerialdesignsharp{local-drink}{1077}
\showcasematerialdesignsharp{local-fire-department}{1078}
\showcasematerialdesignsharp{local-florist}{1079}
\showcasematerialdesignsharp{local-gas-station}{1080}
\showcasematerialdesignsharp{local-grocery-store}{1081}
\showcasematerialdesignsharp{local-hospital}{1082}
\showcasematerialdesignsharp{local-hotel}{1083}
\showcasematerialdesignsharp{local-laundry-service}{1084}
\showcasematerialdesignsharp{local-library}{1085}
\showcasematerialdesignsharp{local-mall}{1086}
\showcasematerialdesignsharp{local-movies}{1087}
\showcasematerialdesignsharp{local-offer}{1088}
\showcasematerialdesignsharp{local-parking}{1089}
\showcasematerialdesignsharp{local-pharmacy}{1090}
\showcasematerialdesignsharp{local-phone}{1091}
\showcasematerialdesignsharp{local-pizza}{1092}
\showcasematerialdesignsharp{local-play}{1093}
\showcasematerialdesignsharp{local-police}{1094}
\showcasematerialdesignsharp{local-post-office}{1095}
\showcasematerialdesignsharp{local-printshop}{1096}
\showcasematerialdesignsharp{local-see}{1097}
\showcasematerialdesignsharp{local-shipping}{1098}
\showcasematerialdesignsharp{local-taxi}{1099}
\showcasematerialdesignsharp{location-city}{1100}
\showcasematerialdesignsharp{location-disabled}{1101}
\showcasematerialdesignsharp{location-off}{1102}
\showcasematerialdesignsharp{location-on}{1103}
\showcasematerialdesignsharp{location-searching}{1104}
\showcasematerialdesignsharp{lock}{1105}
\showcasematerialdesignsharp{lock-clock}{1106}
\showcasematerialdesignsharp{lock-open}{1107}
\showcasematerialdesignsharp{lock-person}{1108}
\showcasematerialdesignsharp{lock-reset}{1109}
\showcasematerialdesignsharp{login}{1110}
\showcasematerialdesignsharp{logout}{1111}
\showcasematerialdesignsharp{logo-dev}{1112}
\showcasematerialdesignsharp{looks}{1113}
\showcasematerialdesignsharp{looks-3}{1114}
\showcasematerialdesignsharp{looks-4}{1115}
\showcasematerialdesignsharp{looks-5}{1116}
\showcasematerialdesignsharp{looks-6}{1117}
\showcasematerialdesignsharp{looks-one}{1118}
\showcasematerialdesignsharp{looks-two}{1119}
\showcasematerialdesignsharp{loop}{1120}
\showcasematerialdesignsharp{loupe}{1121}
\showcasematerialdesignsharp{low-priority}{1122}
\showcasematerialdesignsharp{loyalty}{1123}
\showcasematerialdesignsharp{lte-mobiledata}{1124}
\showcasematerialdesignsharp{lte-plus-mobiledata}{1125}
\showcasematerialdesignsharp{luggage}{1126}
\showcasematerialdesignsharp{lunch-dining}{1127}
\showcasematerialdesignsharp{lyrics}{1128}
\showcasematerialdesignsharp{macro-off}{1129}
\showcasematerialdesignsharp{mail}{1130}
\showcasematerialdesignsharp{mail-lock}{1131}
\showcasematerialdesignsharp{mail-outline}{1132}
\showcasematerialdesignsharp{male}{1133}
\showcasematerialdesignsharp{man}{1134}
\showcasematerialdesignsharp{manage-accounts}{1135}
\showcasematerialdesignsharp{manage-history}{1136}
\showcasematerialdesignsharp{manage-search}{1137}
\showcasematerialdesignsharp{man-2}{1138}
\showcasematerialdesignsharp{man-3}{1139}
\showcasematerialdesignsharp{man-4}{1140}
\showcasematerialdesignsharp{map}{1141}
\showcasematerialdesignsharp{maps-home-work}{1142}
\showcasematerialdesignsharp{maps-ugc}{1143}
\showcasematerialdesignsharp{margin}{1144}
\showcasematerialdesignsharp{markunread}{1145}
\showcasematerialdesignsharp{markunread-mailbox}{1146}
\showcasematerialdesignsharp{mark-as-unread}{1147}
\showcasematerialdesignsharp{mark-chat-read}{1148}
\showcasematerialdesignsharp{mark-chat-unread}{1149}
\showcasematerialdesignsharp{mark-email-read}{1150}
\showcasematerialdesignsharp{mark-email-unread}{1151}
\showcasematerialdesignsharp{mark-unread-chat-alt}{1152}
\showcasematerialdesignsharp{masks}{1153}
\showcasematerialdesignsharp{maximize}{1154}
\showcasematerialdesignsharp{mediation}{1155}
\showcasematerialdesignsharp{media-bluetooth-off}{1156}
\showcasematerialdesignsharp{media-bluetooth-on}{1157}
\showcasematerialdesignsharp{medical-information}{1158}
\showcasematerialdesignsharp{medical-services}{1159}
\showcasematerialdesignsharp{medication}{1160}
\showcasematerialdesignsharp{medication-liquid}{1161}
\showcasematerialdesignsharp{meeting-room}{1162}
\showcasematerialdesignsharp{memory}{1163}
\showcasematerialdesignsharp{menu}{1164}
\showcasematerialdesignsharp{menu-book}{1165}
\showcasematerialdesignsharp{menu-open}{1166}
\showcasematerialdesignsharp{merge}{1167}
\showcasematerialdesignsharp{merge-type}{1168}
\showcasematerialdesignsharp{message}{1169}
\showcasematerialdesignsharp{mic}{1170}
\showcasematerialdesignsharp{microwave}{1171}
\showcasematerialdesignsharp{mic-external-off}{1172}
\showcasematerialdesignsharp{mic-external-on}{1173}
\showcasematerialdesignsharp{mic-none}{1174}
\showcasematerialdesignsharp{mic-off}{1175}
\showcasematerialdesignsharp{military-tech}{1176}
\showcasematerialdesignsharp{minimize}{1177}
\showcasematerialdesignsharp{minor-crash}{1178}
\showcasematerialdesignsharp{miscellaneous-services}{1179}
\showcasematerialdesignsharp{missed-video-call}{1180}
\showcasematerialdesignsharp{mms}{1181}
\showcasematerialdesignsharp{mobiledata-off}{1182}
\showcasematerialdesignsharp{mobile-friendly}{1183}
\showcasematerialdesignsharp{mobile-off}{1184}
\showcasematerialdesignsharp{mobile-screen-share}{1185}
\showcasematerialdesignsharp{mode}{1186}
\showcasematerialdesignsharp{model-training}{1187}
\showcasematerialdesignsharp{mode-comment}{1188}
\showcasematerialdesignsharp{mode-edit}{1189}
\showcasematerialdesignsharp{mode-edit-outline}{1190}
\showcasematerialdesignsharp{mode-fan-off}{1191}
\showcasematerialdesignsharp{mode-night}{1192}
\showcasematerialdesignsharp{mode-of-travel}{1193}
\showcasematerialdesignsharp{mode-standby}{1194}
\showcasematerialdesignsharp{monetization-on}{1195}
\showcasematerialdesignsharp{money}{1196}
\showcasematerialdesignsharp{money-off}{1197}
\showcasematerialdesignsharp{money-off-csred}{1198}
\showcasematerialdesignsharp{monitor}{1199}
\showcasematerialdesignsharp{monitor-heart}{1200}
\showcasematerialdesignsharp{monitor-weight}{1201}
\showcasematerialdesignsharp{monochrome-photos}{1202}
\showcasematerialdesignsharp{mood}{1203}
\showcasematerialdesignsharp{mood-bad}{1204}
\showcasematerialdesignsharp{moped}{1205}
\showcasematerialdesignsharp{more}{1206}
\showcasematerialdesignsharp{more-horiz}{1207}
\showcasematerialdesignsharp{more-time}{1208}
\showcasematerialdesignsharp{more-vert}{1209}
\showcasematerialdesignsharp{mosque}{1210}
\showcasematerialdesignsharp{motion-photos-auto}{1211}
\showcasematerialdesignsharp{motion-photos-off}{1212}
\showcasematerialdesignsharp{motion-photos-on}{1213}
\showcasematerialdesignsharp{motion-photos-pause}{1214}
\showcasematerialdesignsharp{motion-photos-paused}{1215}
\showcasematerialdesignsharp{mouse}{1216}
\showcasematerialdesignsharp{move-down}{1217}
\showcasematerialdesignsharp{move-to-inbox}{1218}
\showcasematerialdesignsharp{move-up}{1219}
\showcasematerialdesignsharp{movie}{1220}
\showcasematerialdesignsharp{movie-creation}{1221}
\showcasematerialdesignsharp{movie-filter}{1222}
\showcasematerialdesignsharp{moving}{1223}
\showcasematerialdesignsharp{mp}{1224}
\showcasematerialdesignsharp{multiline-chart}{1225}
\showcasematerialdesignsharp{multiple-stop}{1226}
\showcasematerialdesignsharp{museum}{1227}
\showcasematerialdesignsharp{music-note}{1228}
\showcasematerialdesignsharp{music-off}{1229}
\showcasematerialdesignsharp{music-video}{1230}
\showcasematerialdesignsharp{my-location}{1231}
\showcasematerialdesignsharp{nat}{1232}
\showcasematerialdesignsharp{nature}{1233}
\showcasematerialdesignsharp{nature-people}{1234}
\showcasematerialdesignsharp{navigate-before}{1235}
\showcasematerialdesignsharp{navigate-next}{1236}
\showcasematerialdesignsharp{navigation}{1237}
\showcasematerialdesignsharp{nearby-error}{1238}
\showcasematerialdesignsharp{nearby-off}{1239}
\showcasematerialdesignsharp{near-me}{1240}
\showcasematerialdesignsharp{near-me-disabled}{1241}
\showcasematerialdesignsharp{nest-cam-wired-stand}{1242}
\showcasematerialdesignsharp{network-cell}{1243}
\showcasematerialdesignsharp{network-check}{1244}
\showcasematerialdesignsharp{network-locked}{1245}
\showcasematerialdesignsharp{network-ping}{1246}
\showcasematerialdesignsharp{network-wifi}{1247}
\showcasematerialdesignsharp{network-wifi-1-bar}{1248}
\showcasematerialdesignsharp{network-wifi-2-bar}{1249}
\showcasematerialdesignsharp{network-wifi-3-bar}{1250}
\showcasematerialdesignsharp{newspaper}{1251}
\showcasematerialdesignsharp{new-label}{1252}
\showcasematerialdesignsharp{new-releases}{1253}
\showcasematerialdesignsharp{next-plan}{1254}
\showcasematerialdesignsharp{next-week}{1255}
\showcasematerialdesignsharp{nfc}{1256}
\showcasematerialdesignsharp{nightlife}{1257}
\showcasematerialdesignsharp{nightlight}{1258}
\showcasematerialdesignsharp{nightlight-round}{1259}
\showcasematerialdesignsharp{nights-stay}{1260}
\showcasematerialdesignsharp{night-shelter}{1261}
\showcasematerialdesignsharp{noise-aware}{1262}
\showcasematerialdesignsharp{noise-control-off}{1263}
\showcasematerialdesignsharp{nordic-walking}{1264}
\showcasematerialdesignsharp{north}{1265}
\showcasematerialdesignsharp{north-east}{1266}
\showcasematerialdesignsharp{north-west}{1267}
\showcasematerialdesignsharp{note}{1268}
\showcasematerialdesignsharp{notes}{1269}
\showcasematerialdesignsharp{note-add}{1270}
\showcasematerialdesignsharp{note-alt}{1271}
\showcasematerialdesignsharp{notifications}{1272}
\showcasematerialdesignsharp{notifications-active}{1273}
\showcasematerialdesignsharp{notifications-none}{1274}
\showcasematerialdesignsharp{notifications-off}{1275}
\showcasematerialdesignsharp{notifications-paused}{1276}
\showcasematerialdesignsharp{notification-add}{1277}
\showcasematerialdesignsharp{notification-important}{1278}
\showcasematerialdesignsharp{not-accessible}{1279}
\showcasematerialdesignsharp{not-interested}{1280}
\showcasematerialdesignsharp{not-listed-location}{1281}
\showcasematerialdesignsharp{not-started}{1282}
\showcasematerialdesignsharp{no-accounts}{1283}
\showcasematerialdesignsharp{no-adult-content}{1284}
\showcasematerialdesignsharp{no-backpack}{1285}
\showcasematerialdesignsharp{no-cell}{1286}
\showcasematerialdesignsharp{no-crash}{1287}
\showcasematerialdesignsharp{no-drinks}{1288}
\showcasematerialdesignsharp{no-encryption}{1289}
\showcasematerialdesignsharp{no-encryption-gmailerrorred}{1290}
\showcasematerialdesignsharp{no-flash}{1291}
\showcasematerialdesignsharp{no-food}{1292}
\showcasematerialdesignsharp{no-luggage}{1293}
\showcasematerialdesignsharp{no-meals}{1294}
\showcasematerialdesignsharp{no-meeting-room}{1295}
\showcasematerialdesignsharp{no-photography}{1296}
\showcasematerialdesignsharp{no-sim}{1297}
\showcasematerialdesignsharp{no-stroller}{1298}
\showcasematerialdesignsharp{no-transfer}{1299}
\showcasematerialdesignsharp{numbers}{1300}
\showcasematerialdesignsharp{offline-bolt}{1301}
\showcasematerialdesignsharp{offline-pin}{1302}
\showcasematerialdesignsharp{offline-share}{1303}
\showcasematerialdesignsharp{oil-barrel}{1304}
\showcasematerialdesignsharp{ondemand-video}{1305}
\showcasematerialdesignsharp{online-prediction}{1306}
\showcasematerialdesignsharp{on-device-training}{1307}
\showcasematerialdesignsharp{opacity}{1308}
\showcasematerialdesignsharp{open-in-browser}{1309}
\showcasematerialdesignsharp{open-in-full}{1310}
\showcasematerialdesignsharp{open-in-new}{1311}
\showcasematerialdesignsharp{open-in-new-off}{1312}
\showcasematerialdesignsharp{open-with}{1313}
\showcasematerialdesignsharp{other-houses}{1314}
\showcasematerialdesignsharp{outbound}{1315}
\showcasematerialdesignsharp{outbox}{1316}
\showcasematerialdesignsharp{outdoor-grill}{1317}
\showcasematerialdesignsharp{outlet}{1318}
\showcasematerialdesignsharp{outlined-flag}{1319}
\showcasematerialdesignsharp{output}{1320}
\showcasematerialdesignsharp{padding}{1321}
\showcasematerialdesignsharp{pages}{1322}
\showcasematerialdesignsharp{pageview}{1323}
\showcasematerialdesignsharp{paid}{1324}
\showcasematerialdesignsharp{palette}{1325}
\showcasematerialdesignsharp{panorama}{1326}
\showcasematerialdesignsharp{panorama-fish-eye}{1327}
\showcasematerialdesignsharp{panorama-horizontal}{1328}
\showcasematerialdesignsharp{panorama-horizontal-select}{1329}
\showcasematerialdesignsharp{panorama-photosphere}{1330}
\showcasematerialdesignsharp{panorama-photosphere-select}{1331}
\showcasematerialdesignsharp{panorama-vertical}{1332}
\showcasematerialdesignsharp{panorama-vertical-select}{1333}
\showcasematerialdesignsharp{panorama-wide-angle}{1334}
\showcasematerialdesignsharp{panorama-wide-angle-select}{1335}
\showcasematerialdesignsharp{pan-tool}{1336}
\showcasematerialdesignsharp{pan-tool-alt}{1337}
\showcasematerialdesignsharp{paragliding}{1338}
\showcasematerialdesignsharp{park}{1339}
\showcasematerialdesignsharp{party-mode}{1340}
\showcasematerialdesignsharp{password}{1341}
\showcasematerialdesignsharp{pattern}{1342}
\showcasematerialdesignsharp{pause}{1343}
\showcasematerialdesignsharp{pause-circle}{1344}
\showcasematerialdesignsharp{pause-circle-filled}{1345}
\showcasematerialdesignsharp{pause-circle-outline}{1346}
\showcasematerialdesignsharp{pause-presentation}{1347}
\showcasematerialdesignsharp{payment}{1348}
\showcasematerialdesignsharp{payments}{1349}
\showcasematerialdesignsharp{pedal-bike}{1350}
\showcasematerialdesignsharp{pending}{1351}
\showcasematerialdesignsharp{pending-actions}{1352}
\showcasematerialdesignsharp{pentagon}{1353}
\showcasematerialdesignsharp{people}{1354}
\showcasematerialdesignsharp{people-alt}{1355}
\showcasematerialdesignsharp{people-outline}{1356}
\showcasematerialdesignsharp{percent}{1357}
\showcasematerialdesignsharp{perm-camera-mic}{1358}
\showcasematerialdesignsharp{perm-contact-calendar}{1359}
\showcasematerialdesignsharp{perm-data-setting}{1360}
\showcasematerialdesignsharp{perm-device-information}{1361}
\showcasematerialdesignsharp{perm-identity}{1362}
\showcasematerialdesignsharp{perm-media}{1363}
\showcasematerialdesignsharp{perm-phone-msg}{1364}
\showcasematerialdesignsharp{perm-scan-wifi}{1365}
\showcasematerialdesignsharp{person}{1366}
\showcasematerialdesignsharp{personal-injury}{1367}
\showcasematerialdesignsharp{personal-video}{1368}
\showcasematerialdesignsharp{person-2}{1369}
\showcasematerialdesignsharp{person-3}{1370}
\showcasematerialdesignsharp{person-4}{1371}
\showcasematerialdesignsharp{person-add}{1372}
\showcasematerialdesignsharp{person-add-alt}{1373}
\showcasematerialdesignsharp{person-add-alt-1}{1374}
\showcasematerialdesignsharp{person-add-disabled}{1375}
\showcasematerialdesignsharp{person-off}{1376}
\showcasematerialdesignsharp{person-outline}{1377}
\showcasematerialdesignsharp{person-pin}{1378}
\showcasematerialdesignsharp{person-pin-circle}{1379}
\showcasematerialdesignsharp{person-remove}{1380}
\showcasematerialdesignsharp{person-remove-alt-1}{1381}
\showcasematerialdesignsharp{person-search}{1382}
\showcasematerialdesignsharp{pest-control}{1383}
\showcasematerialdesignsharp{pest-control-rodent}{1384}
\showcasematerialdesignsharp{pets}{1385}
\showcasematerialdesignsharp{phishing}{1386}
\showcasematerialdesignsharp{phone}{1387}
\showcasematerialdesignsharp{phonelink}{1388}
\showcasematerialdesignsharp{phonelink-erase}{1389}
\showcasematerialdesignsharp{phonelink-lock}{1390}
\showcasematerialdesignsharp{phonelink-off}{1391}
\showcasematerialdesignsharp{phonelink-ring}{1392}
\showcasematerialdesignsharp{phonelink-setup}{1393}
\showcasematerialdesignsharp{phone-android}{1394}
\showcasematerialdesignsharp{phone-bluetooth-speaker}{1395}
\showcasematerialdesignsharp{phone-callback}{1396}
\showcasematerialdesignsharp{phone-disabled}{1397}
\showcasematerialdesignsharp{phone-enabled}{1398}
\showcasematerialdesignsharp{phone-forwarded}{1399}
\showcasematerialdesignsharp{phone-iphone}{1400}
\showcasematerialdesignsharp{phone-locked}{1401}
\showcasematerialdesignsharp{phone-missed}{1402}
\showcasematerialdesignsharp{phone-paused}{1403}
\showcasematerialdesignsharp{photo}{1404}
\showcasematerialdesignsharp{photo-album}{1405}
\showcasematerialdesignsharp{photo-camera}{1406}
\showcasematerialdesignsharp{photo-camera-back}{1407}
\showcasematerialdesignsharp{photo-camera-front}{1408}
\showcasematerialdesignsharp{photo-filter}{1409}
\showcasematerialdesignsharp{photo-library}{1410}
\showcasematerialdesignsharp{photo-size-select-actual}{1411}
\showcasematerialdesignsharp{photo-size-select-large}{1412}
\showcasematerialdesignsharp{photo-size-select-small}{1413}
\showcasematerialdesignsharp{php}{1414}
\showcasematerialdesignsharp{piano}{1415}
\showcasematerialdesignsharp{piano-off}{1416}
\showcasematerialdesignsharp{picture-as-pdf}{1417}
\showcasematerialdesignsharp{picture-in-picture}{1418}
\showcasematerialdesignsharp{picture-in-picture-alt}{1419}
\showcasematerialdesignsharp{pie-chart}{1420}
\showcasematerialdesignsharp{pie-chart-outline}{1421}
\showcasematerialdesignsharp{pin}{1422}
\showcasematerialdesignsharp{pinch}{1423}
\showcasematerialdesignsharp{pin-drop}{1424}
\showcasematerialdesignsharp{pin-end}{1425}
\showcasematerialdesignsharp{pin-invoke}{1426}
\showcasematerialdesignsharp{pivot-table-chart}{1427}
\showcasematerialdesignsharp{pix}{1428}
\showcasematerialdesignsharp{place}{1429}
\showcasematerialdesignsharp{plagiarism}{1430}
\showcasematerialdesignsharp{playlist-add}{1431}
\showcasematerialdesignsharp{playlist-add-check}{1432}
\showcasematerialdesignsharp{playlist-add-check-circle}{1433}
\showcasematerialdesignsharp{playlist-add-circle}{1434}
\showcasematerialdesignsharp{playlist-play}{1435}
\showcasematerialdesignsharp{playlist-remove}{1436}
\showcasematerialdesignsharp{play-arrow}{1437}
\showcasematerialdesignsharp{play-circle}{1438}
\showcasematerialdesignsharp{play-circle-filled}{1439}
\showcasematerialdesignsharp{play-circle-outline}{1440}
\showcasematerialdesignsharp{play-disabled}{1441}
\showcasematerialdesignsharp{play-for-work}{1442}
\showcasematerialdesignsharp{play-lesson}{1443}
\showcasematerialdesignsharp{plumbing}{1444}
\showcasematerialdesignsharp{plus-one}{1445}
\showcasematerialdesignsharp{podcasts}{1446}
\showcasematerialdesignsharp{point-of-sale}{1447}
\showcasematerialdesignsharp{policy}{1448}
\showcasematerialdesignsharp{poll}{1449}
\showcasematerialdesignsharp{polyline}{1450}
\showcasematerialdesignsharp{polymer}{1451}
\showcasematerialdesignsharp{pool}{1452}
\showcasematerialdesignsharp{portable-wifi-off}{1453}
\showcasematerialdesignsharp{portrait}{1454}
\showcasematerialdesignsharp{post-add}{1455}
\showcasematerialdesignsharp{power}{1456}
\showcasematerialdesignsharp{power-input}{1457}
\showcasematerialdesignsharp{power-off}{1458}
\showcasematerialdesignsharp{power-settings-new}{1459}
\showcasematerialdesignsharp{precision-manufacturing}{1460}
\showcasematerialdesignsharp{pregnant-woman}{1461}
\showcasematerialdesignsharp{present-to-all}{1462}
\showcasematerialdesignsharp{preview}{1463}
\showcasematerialdesignsharp{price-change}{1464}
\showcasematerialdesignsharp{price-check}{1465}
\showcasematerialdesignsharp{print}{1466}
\showcasematerialdesignsharp{print-disabled}{1467}
\showcasematerialdesignsharp{priority-high}{1468}
\showcasematerialdesignsharp{privacy-tip}{1469}
\showcasematerialdesignsharp{private-connectivity}{1470}
\showcasematerialdesignsharp{production-quantity-limits}{1471}
\showcasematerialdesignsharp{propane}{1472}
\showcasematerialdesignsharp{propane-tank}{1473}
\showcasematerialdesignsharp{psychology}{1474}
\showcasematerialdesignsharp{psychology-alt}{1475}
\showcasematerialdesignsharp{public}{1476}
\showcasematerialdesignsharp{public-off}{1477}
\showcasematerialdesignsharp{publish}{1478}
\showcasematerialdesignsharp{published-with-changes}{1479}
\showcasematerialdesignsharp{punch-clock}{1480}
\showcasematerialdesignsharp{push-pin}{1481}
\showcasematerialdesignsharp{qr-code}{1482}
\showcasematerialdesignsharp{qr-code-2}{1483}
\showcasematerialdesignsharp{qr-code-scanner}{1484}
\showcasematerialdesignsharp{query-builder}{1485}
\showcasematerialdesignsharp{query-stats}{1486}
\showcasematerialdesignsharp{question-answer}{1487}
\showcasematerialdesignsharp{question-mark}{1488}
\showcasematerialdesignsharp{queue}{1489}
\showcasematerialdesignsharp{queue-music}{1490}
\showcasematerialdesignsharp{queue-play-next}{1491}
\showcasematerialdesignsharp{quickreply}{1492}
\showcasematerialdesignsharp{quiz}{1493}
\showcasematerialdesignsharp{radar}{1494}
\showcasematerialdesignsharp{radio}{1495}
\showcasematerialdesignsharp{radio-button-checked}{1496}
\showcasematerialdesignsharp{radio-button-unchecked}{1497}
\showcasematerialdesignsharp{railway-alert}{1498}
\showcasematerialdesignsharp{ramen-dining}{1499}
\showcasematerialdesignsharp{ramp-left}{1500}
\showcasematerialdesignsharp{ramp-right}{1501}
\showcasematerialdesignsharp{rate-review}{1502}
\showcasematerialdesignsharp{raw-off}{1503}
\showcasematerialdesignsharp{raw-on}{1504}
\showcasematerialdesignsharp{read-more}{1505}
\showcasematerialdesignsharp{real-estate-agent}{1506}
\showcasematerialdesignsharp{receipt}{1507}
\showcasematerialdesignsharp{receipt-long}{1508}
\showcasematerialdesignsharp{recent-actors}{1509}
\showcasematerialdesignsharp{recommend}{1510}
\showcasematerialdesignsharp{record-voice-over}{1511}
\showcasematerialdesignsharp{rectangle}{1512}
\showcasematerialdesignsharp{recycling}{1513}
\showcasematerialdesignsharp{redeem}{1514}
\showcasematerialdesignsharp{redo}{1515}
\showcasematerialdesignsharp{reduce-capacity}{1516}
\showcasematerialdesignsharp{refresh}{1517}
\showcasematerialdesignsharp{remember-me}{1518}
\showcasematerialdesignsharp{remove}{1519}
\showcasematerialdesignsharp{remove-circle}{1520}
\showcasematerialdesignsharp{remove-circle-outline}{1521}
\showcasematerialdesignsharp{remove-done}{1522}
\showcasematerialdesignsharp{remove-from-queue}{1523}
\showcasematerialdesignsharp{remove-moderator}{1524}
\showcasematerialdesignsharp{remove-red-eye}{1525}
\showcasematerialdesignsharp{remove-road}{1526}
\showcasematerialdesignsharp{remove-shopping-cart}{1527}
\showcasematerialdesignsharp{reorder}{1528}
\showcasematerialdesignsharp{repartition}{1529}
\showcasematerialdesignsharp{repeat}{1530}
\showcasematerialdesignsharp{repeat-on}{1531}
\showcasematerialdesignsharp{repeat-one}{1532}
\showcasematerialdesignsharp{repeat-one-on}{1533}
\showcasematerialdesignsharp{replay}{1534}
\showcasematerialdesignsharp{replay-5}{1535}
\showcasematerialdesignsharp{replay-10}{1536}
\showcasematerialdesignsharp{replay-30}{1537}
\showcasematerialdesignsharp{replay-circle-filled}{1538}
\showcasematerialdesignsharp{reply}{1539}
\showcasematerialdesignsharp{reply-all}{1540}
\showcasematerialdesignsharp{report}{1541}
\showcasematerialdesignsharp{report-gmailerrorred}{1542}
\showcasematerialdesignsharp{report-off}{1543}
\showcasematerialdesignsharp{report-problem}{1544}
\showcasematerialdesignsharp{request-page}{1545}
\showcasematerialdesignsharp{request-quote}{1546}
\showcasematerialdesignsharp{reset-tv}{1547}
\showcasematerialdesignsharp{restart-alt}{1548}
\showcasematerialdesignsharp{restaurant}{1549}
\showcasematerialdesignsharp{restaurant-menu}{1550}
\showcasematerialdesignsharp{restore}{1551}
\showcasematerialdesignsharp{restore-from-trash}{1552}
\showcasematerialdesignsharp{restore-page}{1553}
\showcasematerialdesignsharp{reviews}{1554}
\showcasematerialdesignsharp{rice-bowl}{1555}
\showcasematerialdesignsharp{ring-volume}{1556}
\showcasematerialdesignsharp{rocket}{1557}
\showcasematerialdesignsharp{rocket-launch}{1558}
\showcasematerialdesignsharp{roller-shades}{1559}
\showcasematerialdesignsharp{roller-shades-closed}{1560}
\showcasematerialdesignsharp{roller-skating}{1561}
\showcasematerialdesignsharp{roofing}{1562}
\showcasematerialdesignsharp{room}{1563}
\showcasematerialdesignsharp{room-preferences}{1564}
\showcasematerialdesignsharp{room-service}{1565}
\showcasematerialdesignsharp{rotate-90-degrees-ccw}{1566}
\showcasematerialdesignsharp{rotate-90-degrees-cw}{1567}
\showcasematerialdesignsharp{rotate-left}{1568}
\showcasematerialdesignsharp{rotate-right}{1569}
\showcasematerialdesignsharp{roundabout-left}{1570}
\showcasematerialdesignsharp{roundabout-right}{1571}
\showcasematerialdesignsharp{rounded-corner}{1572}
\showcasematerialdesignsharp{route}{1573}
\showcasematerialdesignsharp{router}{1574}
\showcasematerialdesignsharp{rowing}{1575}
\showcasematerialdesignsharp{rss-feed}{1576}
\showcasematerialdesignsharp{rsvp}{1577}
\showcasematerialdesignsharp{rtt}{1578}
\showcasematerialdesignsharp{rule}{1579}
\showcasematerialdesignsharp{rule-folder}{1580}
\showcasematerialdesignsharp{running-with-errors}{1581}
\showcasematerialdesignsharp{run-circle}{1582}
\showcasematerialdesignsharp{rv-hookup}{1583}
\showcasematerialdesignsharp{r-mobiledata}{1584}
\showcasematerialdesignsharp{safety-check}{1585}
\showcasematerialdesignsharp{safety-divider}{1586}
\showcasematerialdesignsharp{sailing}{1587}
\showcasematerialdesignsharp{sanitizer}{1588}
\showcasematerialdesignsharp{satellite}{1589}
\showcasematerialdesignsharp{satellite-alt}{1590}
\showcasematerialdesignsharp{save}{1591}
\showcasematerialdesignsharp{saved-search}{1592}
\showcasematerialdesignsharp{save-alt}{1593}
\showcasematerialdesignsharp{save-as}{1594}
\showcasematerialdesignsharp{savings}{1595}
\showcasematerialdesignsharp{scale}{1596}
\showcasematerialdesignsharp{scanner}{1597}
\showcasematerialdesignsharp{scatter-plot}{1598}
\showcasematerialdesignsharp{schedule}{1599}
\showcasematerialdesignsharp{schedule-send}{1600}
\showcasematerialdesignsharp{schema}{1601}
\showcasematerialdesignsharp{school}{1602}
\showcasematerialdesignsharp{science}{1603}
\showcasematerialdesignsharp{score}{1604}
\showcasematerialdesignsharp{scoreboard}{1605}
\showcasematerialdesignsharp{screenshot}{1606}
\showcasematerialdesignsharp{screenshot-monitor}{1607}
\showcasematerialdesignsharp{screen-lock-landscape}{1608}
\showcasematerialdesignsharp{screen-lock-portrait}{1609}
\showcasematerialdesignsharp{screen-lock-rotation}{1610}
\showcasematerialdesignsharp{screen-rotation}{1611}
\showcasematerialdesignsharp{screen-rotation-alt}{1612}
\showcasematerialdesignsharp{screen-search-desktop}{1613}
\showcasematerialdesignsharp{screen-share}{1614}
\showcasematerialdesignsharp{scuba-diving}{1615}
\showcasematerialdesignsharp{sd}{1616}
\showcasematerialdesignsharp{sd-card}{1617}
\showcasematerialdesignsharp{sd-card-alert}{1618}
\showcasematerialdesignsharp{sd-storage}{1619}
\showcasematerialdesignsharp{search}{1620}
\showcasematerialdesignsharp{search-off}{1621}
\showcasematerialdesignsharp{security}{1622}
\showcasematerialdesignsharp{security-update}{1623}
\showcasematerialdesignsharp{security-update-good}{1624}
\showcasematerialdesignsharp{security-update-warning}{1625}
\showcasematerialdesignsharp{segment}{1626}
\showcasematerialdesignsharp{select-all}{1627}
\showcasematerialdesignsharp{self-improvement}{1628}
\showcasematerialdesignsharp{sell}{1629}
\showcasematerialdesignsharp{send}{1630}
\showcasematerialdesignsharp{send-and-archive}{1631}
\showcasematerialdesignsharp{send-time-extension}{1632}
\showcasematerialdesignsharp{send-to-mobile}{1633}
\showcasematerialdesignsharp{sensors}{1634}
\showcasematerialdesignsharp{sensors-off}{1635}
\showcasematerialdesignsharp{sensor-door}{1636}
\showcasematerialdesignsharp{sensor-occupied}{1637}
\showcasematerialdesignsharp{sensor-window}{1638}
\showcasematerialdesignsharp{sentiment-dissatisfied}{1639}
\showcasematerialdesignsharp{sentiment-neutral}{1640}
\showcasematerialdesignsharp{sentiment-satisfied}{1641}
\showcasematerialdesignsharp{sentiment-satisfied-alt}{1642}
\showcasematerialdesignsharp{sentiment-very-dissatisfied}{1643}
\showcasematerialdesignsharp{sentiment-very-satisfied}{1644}
\showcasematerialdesignsharp{settings}{1645}
\showcasematerialdesignsharp{settings-accessibility}{1646}
\showcasematerialdesignsharp{settings-applications}{1647}
\showcasematerialdesignsharp{settings-backup-restore}{1648}
\showcasematerialdesignsharp{settings-bluetooth}{1649}
\showcasematerialdesignsharp{settings-brightness}{1650}
\showcasematerialdesignsharp{settings-cell}{1651}
\showcasematerialdesignsharp{settings-ethernet}{1652}
\showcasematerialdesignsharp{settings-input-antenna}{1653}
\showcasematerialdesignsharp{settings-input-component}{1654}
\showcasematerialdesignsharp{settings-input-composite}{1655}
\showcasematerialdesignsharp{settings-input-hdmi}{1656}
\showcasematerialdesignsharp{settings-input-svideo}{1657}
\showcasematerialdesignsharp{settings-overscan}{1658}
\showcasematerialdesignsharp{settings-phone}{1659}
\showcasematerialdesignsharp{settings-power}{1660}
\showcasematerialdesignsharp{settings-remote}{1661}
\showcasematerialdesignsharp{settings-suggest}{1662}
\showcasematerialdesignsharp{settings-system-daydream}{1663}
\showcasematerialdesignsharp{settings-voice}{1664}
\showcasematerialdesignsharp{set-meal}{1665}
\showcasematerialdesignsharp{severe-cold}{1666}
\showcasematerialdesignsharp{shape-line}{1667}
\showcasematerialdesignsharp{share}{1668}
\showcasematerialdesignsharp{share-location}{1669}
\showcasematerialdesignsharp{shield}{1670}
\showcasematerialdesignsharp{shield-moon}{1671}
\showcasematerialdesignsharp{shop}{1672}
\showcasematerialdesignsharp{shopping-bag}{1673}
\showcasematerialdesignsharp{shopping-basket}{1674}
\showcasematerialdesignsharp{shopping-cart}{1675}
\showcasematerialdesignsharp{shopping-cart-checkout}{1676}
\showcasematerialdesignsharp{shop-2}{1677}
\showcasematerialdesignsharp{shop-two}{1678}
\showcasematerialdesignsharp{shortcut}{1679}
\showcasematerialdesignsharp{short-text}{1680}
\showcasematerialdesignsharp{shower}{1681}
\showcasematerialdesignsharp{show-chart}{1682}
\showcasematerialdesignsharp{shuffle}{1683}
\showcasematerialdesignsharp{shuffle-on}{1684}
\showcasematerialdesignsharp{shutter-speed}{1685}
\showcasematerialdesignsharp{sick}{1686}
\showcasematerialdesignsharp{signal-cellular-0-bar}{1687}
\showcasematerialdesignsharp{signal-cellular-4-bar}{1688}
\showcasematerialdesignsharp{signal-cellular-alt}{1689}
\showcasematerialdesignsharp{signal-cellular-alt-1-bar}{1690}
\showcasematerialdesignsharp{signal-cellular-alt-2-bar}{1691}
\showcasematerialdesignsharp{signal-cellular-connected-no-internet-0-bar}{1692}
\showcasematerialdesignsharp{signal-cellular-connected-no-internet-4-bar}{1693}
\showcasematerialdesignsharp{signal-cellular-nodata}{1694}
\showcasematerialdesignsharp{signal-cellular-no-sim}{1695}
\showcasematerialdesignsharp{signal-cellular-null}{1696}
\showcasematerialdesignsharp{signal-cellular-off}{1697}
\showcasematerialdesignsharp{signal-wifi-0-bar}{1698}
\showcasematerialdesignsharp{signal-wifi-4-bar}{1699}
\showcasematerialdesignsharp{signal-wifi-4-bar-lock}{1700}
\showcasematerialdesignsharp{signal-wifi-bad}{1701}
\showcasematerialdesignsharp{signal-wifi-connected-no-internet-4}{1702}
\showcasematerialdesignsharp{signal-wifi-off}{1703}
\showcasematerialdesignsharp{signal-wifi-statusbar-4-bar}{1704}
\showcasematerialdesignsharp{signal-wifi-statusbar-connected-no-internet-4}{1705}
\showcasematerialdesignsharp{signal-wifi-statusbar-null}{1706}
\showcasematerialdesignsharp{signpost}{1707}
\showcasematerialdesignsharp{sign-language}{1708}
\showcasematerialdesignsharp{sim-card}{1709}
\showcasematerialdesignsharp{sim-card-alert}{1710}
\showcasematerialdesignsharp{sim-card-download}{1711}
\showcasematerialdesignsharp{single-bed}{1712}
\showcasematerialdesignsharp{sip}{1713}
\showcasematerialdesignsharp{skateboarding}{1714}
\showcasematerialdesignsharp{skip-next}{1715}
\showcasematerialdesignsharp{skip-previous}{1716}
\showcasematerialdesignsharp{sledding}{1717}
\showcasematerialdesignsharp{slideshow}{1718}
\showcasematerialdesignsharp{slow-motion-video}{1719}
\showcasematerialdesignsharp{smartphone}{1720}
\showcasematerialdesignsharp{smart-button}{1721}
\showcasematerialdesignsharp{smart-display}{1722}
\showcasematerialdesignsharp{smart-screen}{1723}
\showcasematerialdesignsharp{smart-toy}{1724}
\showcasematerialdesignsharp{smoke-free}{1725}
\showcasematerialdesignsharp{smoking-rooms}{1726}
\showcasematerialdesignsharp{sms}{1727}
\showcasematerialdesignsharp{sms-failed}{1728}
\showcasematerialdesignsharp{snippet-folder}{1729}
\showcasematerialdesignsharp{snooze}{1730}
\showcasematerialdesignsharp{snowboarding}{1731}
\showcasematerialdesignsharp{snowmobile}{1732}
\showcasematerialdesignsharp{snowshoeing}{1733}
\showcasematerialdesignsharp{soap}{1734}
\showcasematerialdesignsharp{social-distance}{1735}
\showcasematerialdesignsharp{solar-power}{1736}
\showcasematerialdesignsharp{sort}{1737}
\showcasematerialdesignsharp{sort-by-alpha}{1738}
\showcasematerialdesignsharp{sos}{1739}
\showcasematerialdesignsharp{soup-kitchen}{1740}
\showcasematerialdesignsharp{source}{1741}
\showcasematerialdesignsharp{south}{1742}
\showcasematerialdesignsharp{south-america}{1743}
\showcasematerialdesignsharp{south-east}{1744}
\showcasematerialdesignsharp{south-west}{1745}
\showcasematerialdesignsharp{spa}{1746}
\showcasematerialdesignsharp{space-bar}{1747}
\showcasematerialdesignsharp{space-dashboard}{1748}
\showcasematerialdesignsharp{spatial-audio}{1749}
\showcasematerialdesignsharp{spatial-audio-off}{1750}
\showcasematerialdesignsharp{spatial-tracking}{1751}
\showcasematerialdesignsharp{speaker}{1752}
\showcasematerialdesignsharp{speaker-group}{1753}
\showcasematerialdesignsharp{speaker-notes}{1754}
\showcasematerialdesignsharp{speaker-notes-off}{1755}
\showcasematerialdesignsharp{speaker-phone}{1756}
\showcasematerialdesignsharp{speed}{1757}
\showcasematerialdesignsharp{spellcheck}{1758}
\showcasematerialdesignsharp{splitscreen}{1759}
\showcasematerialdesignsharp{spoke}{1760}
\showcasematerialdesignsharp{sports}{1761}
\showcasematerialdesignsharp{sports-bar}{1762}
\showcasematerialdesignsharp{sports-baseball}{1763}
\showcasematerialdesignsharp{sports-basketball}{1764}
\showcasematerialdesignsharp{sports-cricket}{1765}
\showcasematerialdesignsharp{sports-esports}{1766}
\showcasematerialdesignsharp{sports-football}{1767}
\showcasematerialdesignsharp{sports-golf}{1768}
\showcasematerialdesignsharp{sports-gymnastics}{1769}
\showcasematerialdesignsharp{sports-handball}{1770}
\showcasematerialdesignsharp{sports-hockey}{1771}
\showcasematerialdesignsharp{sports-kabaddi}{1772}
\showcasematerialdesignsharp{sports-martial-arts}{1773}
\showcasematerialdesignsharp{sports-mma}{1774}
\showcasematerialdesignsharp{sports-motorsports}{1775}
\showcasematerialdesignsharp{sports-rugby}{1776}
\showcasematerialdesignsharp{sports-score}{1777}
\showcasematerialdesignsharp{sports-soccer}{1778}
\showcasematerialdesignsharp{sports-tennis}{1779}
\showcasematerialdesignsharp{sports-volleyball}{1780}
\showcasematerialdesignsharp{square}{1781}
\showcasematerialdesignsharp{square-foot}{1782}
\showcasematerialdesignsharp{ssid-chart}{1783}
\showcasematerialdesignsharp{stacked-bar-chart}{1784}
\showcasematerialdesignsharp{stacked-line-chart}{1785}
\showcasematerialdesignsharp{stadium}{1786}
\showcasematerialdesignsharp{stairs}{1787}
\showcasematerialdesignsharp{star}{1788}
\showcasematerialdesignsharp{stars}{1789}
\showcasematerialdesignsharp{start}{1790}
\showcasematerialdesignsharp{star-border}{1791}
\showcasematerialdesignsharp{star-border-purple500}{1792}
\showcasematerialdesignsharp{star-half}{1793}
\showcasematerialdesignsharp{star-outline}{1794}
\showcasematerialdesignsharp{star-purple500}{1795}
\showcasematerialdesignsharp{star-rate}{1796}
\showcasematerialdesignsharp{stay-current-landscape}{1797}
\showcasematerialdesignsharp{stay-current-portrait}{1798}
\showcasematerialdesignsharp{stay-primary-landscape}{1799}
\showcasematerialdesignsharp{stay-primary-portrait}{1800}
\showcasematerialdesignsharp{sticky-note-2}{1801}
\showcasematerialdesignsharp{stop}{1802}
\showcasematerialdesignsharp{stop-circle}{1803}
\showcasematerialdesignsharp{stop-screen-share}{1804}
\showcasematerialdesignsharp{storage}{1805}
\showcasematerialdesignsharp{store}{1806}
\showcasematerialdesignsharp{storefront}{1807}
\showcasematerialdesignsharp{store-mall-directory}{1808}
\showcasematerialdesignsharp{storm}{1809}
\showcasematerialdesignsharp{straight}{1810}
\showcasematerialdesignsharp{straighten}{1811}
\showcasematerialdesignsharp{stream}{1812}
\showcasematerialdesignsharp{streetview}{1813}
\showcasematerialdesignsharp{strikethrough-s}{1814}
\showcasematerialdesignsharp{stroller}{1815}
\showcasematerialdesignsharp{style}{1816}
\showcasematerialdesignsharp{subdirectory-arrow-left}{1817}
\showcasematerialdesignsharp{subdirectory-arrow-right}{1818}
\showcasematerialdesignsharp{subject}{1819}
\showcasematerialdesignsharp{subscript}{1820}
\showcasematerialdesignsharp{subscriptions}{1821}
\showcasematerialdesignsharp{subtitles}{1822}
\showcasematerialdesignsharp{subtitles-off}{1823}
\showcasematerialdesignsharp{subway}{1824}
\showcasematerialdesignsharp{summarize}{1825}
\showcasematerialdesignsharp{superscript}{1826}
\showcasematerialdesignsharp{supervised-user-circle}{1827}
\showcasematerialdesignsharp{supervisor-account}{1828}
\showcasematerialdesignsharp{support}{1829}
\showcasematerialdesignsharp{support-agent}{1830}
\showcasematerialdesignsharp{surfing}{1831}
\showcasematerialdesignsharp{surround-sound}{1832}
\showcasematerialdesignsharp{swap-calls}{1833}
\showcasematerialdesignsharp{swap-horiz}{1834}
\showcasematerialdesignsharp{swap-horizontal-circle}{1835}
\showcasematerialdesignsharp{swap-vert}{1836}
\showcasematerialdesignsharp{swap-vertical-circle}{1837}
\showcasematerialdesignsharp{swipe}{1838}
\showcasematerialdesignsharp{swipe-down}{1839}
\showcasematerialdesignsharp{swipe-down-alt}{1840}
\showcasematerialdesignsharp{swipe-left}{1841}
\showcasematerialdesignsharp{swipe-left-alt}{1842}
\showcasematerialdesignsharp{swipe-right}{1843}
\showcasematerialdesignsharp{swipe-right-alt}{1844}
\showcasematerialdesignsharp{swipe-up}{1845}
\showcasematerialdesignsharp{swipe-up-alt}{1846}
\showcasematerialdesignsharp{swipe-vertical}{1847}
\showcasematerialdesignsharp{switch-access-shortcut}{1848}
\showcasematerialdesignsharp{switch-access-shortcut-add}{1849}
\showcasematerialdesignsharp{switch-account}{1850}
\showcasematerialdesignsharp{switch-camera}{1851}
\showcasematerialdesignsharp{switch-left}{1852}
\showcasematerialdesignsharp{switch-right}{1853}
\showcasematerialdesignsharp{switch-video}{1854}
\showcasematerialdesignsharp{synagogue}{1855}
\showcasematerialdesignsharp{sync}{1856}
\showcasematerialdesignsharp{sync-alt}{1857}
\showcasematerialdesignsharp{sync-disabled}{1858}
\showcasematerialdesignsharp{sync-lock}{1859}
\showcasematerialdesignsharp{sync-problem}{1860}
\showcasematerialdesignsharp{system-security-update}{1861}
\showcasematerialdesignsharp{system-security-update-good}{1862}
\showcasematerialdesignsharp{system-security-update-warning}{1863}
\showcasematerialdesignsharp{system-update}{1864}
\showcasematerialdesignsharp{system-update-alt}{1865}
\showcasematerialdesignsharp{tab}{1866}
\showcasematerialdesignsharp{tablet}{1867}
\showcasematerialdesignsharp{tablet-android}{1868}
\showcasematerialdesignsharp{tablet-mac}{1869}
\showcasematerialdesignsharp{table-bar}{1870}
\showcasematerialdesignsharp{table-chart}{1871}
\showcasematerialdesignsharp{table-restaurant}{1872}
\showcasematerialdesignsharp{table-rows}{1873}
\showcasematerialdesignsharp{table-view}{1874}
\showcasematerialdesignsharp{tab-unselected}{1875}
\showcasematerialdesignsharp{tag}{1876}
\showcasematerialdesignsharp{tag-faces}{1877}
\showcasematerialdesignsharp{takeout-dining}{1878}
\showcasematerialdesignsharp{tapas}{1879}
\showcasematerialdesignsharp{tap-and-play}{1880}
\showcasematerialdesignsharp{task}{1881}
\showcasematerialdesignsharp{task-alt}{1882}
\showcasematerialdesignsharp{taxi-alert}{1883}
\showcasematerialdesignsharp{temple-buddhist}{1884}
\showcasematerialdesignsharp{temple-hindu}{1885}
\showcasematerialdesignsharp{terminal}{1886}
\showcasematerialdesignsharp{terrain}{1887}
\showcasematerialdesignsharp{textsms}{1888}
\showcasematerialdesignsharp{texture}{1889}
\showcasematerialdesignsharp{text-decrease}{1890}
\showcasematerialdesignsharp{text-fields}{1891}
\showcasematerialdesignsharp{text-format}{1892}
\showcasematerialdesignsharp{text-increase}{1893}
\showcasematerialdesignsharp{text-rotate-up}{1894}
\showcasematerialdesignsharp{text-rotate-vertical}{1895}
\showcasematerialdesignsharp{text-rotation-angledown}{1896}
\showcasematerialdesignsharp{text-rotation-angleup}{1897}
\showcasematerialdesignsharp{text-rotation-down}{1898}
\showcasematerialdesignsharp{text-rotation-none}{1899}
\showcasematerialdesignsharp{text-snippet}{1900}
\showcasematerialdesignsharp{theaters}{1901}
\showcasematerialdesignsharp{theater-comedy}{1902}
\showcasematerialdesignsharp{thermostat}{1903}
\showcasematerialdesignsharp{thermostat-auto}{1904}
\showcasematerialdesignsharp{thumbs-up-down}{1905}
\showcasematerialdesignsharp{thumb-down}{1906}
\showcasematerialdesignsharp{thumb-down-alt}{1907}
\showcasematerialdesignsharp{thumb-down-off-alt}{1908}
\showcasematerialdesignsharp{thumb-up}{1909}
\showcasematerialdesignsharp{thumb-up-alt}{1910}
\showcasematerialdesignsharp{thumb-up-off-alt}{1911}
\showcasematerialdesignsharp{thunderstorm}{1912}
\showcasematerialdesignsharp{timelapse}{1913}
\showcasematerialdesignsharp{timeline}{1914}
\showcasematerialdesignsharp{timer}{1915}
\showcasematerialdesignsharp{timer-3}{1916}
\showcasematerialdesignsharp{timer-3-select}{1917}
\showcasematerialdesignsharp{timer-10}{1918}
\showcasematerialdesignsharp{timer-10-select}{1919}
\showcasematerialdesignsharp{timer-off}{1920}
\showcasematerialdesignsharp{time-to-leave}{1921}
\showcasematerialdesignsharp{tips-and-updates}{1922}
\showcasematerialdesignsharp{tire-repair}{1923}
\showcasematerialdesignsharp{title}{1924}
\showcasematerialdesignsharp{toc}{1925}
\showcasematerialdesignsharp{today}{1926}
\showcasematerialdesignsharp{toggle-off}{1927}
\showcasematerialdesignsharp{toggle-on}{1928}
\showcasematerialdesignsharp{token}{1929}
\showcasematerialdesignsharp{toll}{1930}
\showcasematerialdesignsharp{tonality}{1931}
\showcasematerialdesignsharp{topic}{1932}
\showcasematerialdesignsharp{tornado}{1933}
\showcasematerialdesignsharp{touch-app}{1934}
\showcasematerialdesignsharp{tour}{1935}
\showcasematerialdesignsharp{toys}{1936}
\showcasematerialdesignsharp{track-changes}{1937}
\showcasematerialdesignsharp{traffic}{1938}
\showcasematerialdesignsharp{train}{1939}
\showcasematerialdesignsharp{tram}{1940}
\showcasematerialdesignsharp{transcribe}{1941}
\showcasematerialdesignsharp{transfer-within-a-station}{1942}
\showcasematerialdesignsharp{transform}{1943}
\showcasematerialdesignsharp{transgender}{1944}
\showcasematerialdesignsharp{transit-enterexit}{1945}
\showcasematerialdesignsharp{translate}{1946}
\showcasematerialdesignsharp{travel-explore}{1947}
\showcasematerialdesignsharp{trending-down}{1948}
\showcasematerialdesignsharp{trending-flat}{1949}
\showcasematerialdesignsharp{trending-up}{1950}
\showcasematerialdesignsharp{trip-origin}{1951}
\showcasematerialdesignsharp{troubleshoot}{1952}
\showcasematerialdesignsharp{try}{1953}
\showcasematerialdesignsharp{tsunami}{1954}
\showcasematerialdesignsharp{tty}{1955}
\showcasematerialdesignsharp{tune}{1956}
\showcasematerialdesignsharp{tungsten}{1957}
\showcasematerialdesignsharp{turned-in}{1958}
\showcasematerialdesignsharp{turned-in-not}{1959}
\showcasematerialdesignsharp{turn-left}{1960}
\showcasematerialdesignsharp{turn-right}{1961}
\showcasematerialdesignsharp{turn-sharp-left}{1962}
\showcasematerialdesignsharp{turn-sharp-right}{1963}
\showcasematerialdesignsharp{turn-slight-left}{1964}
\showcasematerialdesignsharp{turn-slight-right}{1965}
\showcasematerialdesignsharp{tv}{1966}
\showcasematerialdesignsharp{tv-off}{1967}
\showcasematerialdesignsharp{two-wheeler}{1968}
\showcasematerialdesignsharp{type-specimen}{1969}
\showcasematerialdesignsharp{umbrella}{1970}
\showcasematerialdesignsharp{unarchive}{1971}
\showcasematerialdesignsharp{undo}{1972}
\showcasematerialdesignsharp{unfold-less}{1973}
\showcasematerialdesignsharp{unfold-less-double}{1974}
\showcasematerialdesignsharp{unfold-more}{1975}
\showcasematerialdesignsharp{unfold-more-double}{1976}
\showcasematerialdesignsharp{unpublished}{1977}
\showcasematerialdesignsharp{unsubscribe}{1978}
\showcasematerialdesignsharp{upcoming}{1979}
\showcasematerialdesignsharp{update}{1980}
\showcasematerialdesignsharp{update-disabled}{1981}
\showcasematerialdesignsharp{upgrade}{1982}
\showcasematerialdesignsharp{upload}{1983}
\showcasematerialdesignsharp{upload-file}{1984}
\showcasematerialdesignsharp{usb}{1985}
\showcasematerialdesignsharp{usb-off}{1986}
\showcasematerialdesignsharp{u-turn-left}{1987}
\showcasematerialdesignsharp{u-turn-right}{1988}
\showcasematerialdesignsharp{vaccines}{1989}
\showcasematerialdesignsharp{vape-free}{1990}
\showcasematerialdesignsharp{vaping-rooms}{1991}
\showcasematerialdesignsharp{verified}{1992}
\showcasematerialdesignsharp{verified-user}{1993}
\showcasematerialdesignsharp{vertical-align-bottom}{1994}
\showcasematerialdesignsharp{vertical-align-center}{1995}
\showcasematerialdesignsharp{vertical-align-top}{1996}
\showcasematerialdesignsharp{vertical-distribute}{1997}
\showcasematerialdesignsharp{vertical-shades}{1998}
\showcasematerialdesignsharp{vertical-shades-closed}{1999}
\showcasematerialdesignsharp{vertical-split}{2000}
\showcasematerialdesignsharp{vibration}{2001}
\showcasematerialdesignsharp{videocam}{2002}
\showcasematerialdesignsharp{videocam-off}{2003}
\showcasematerialdesignsharp{videogame-asset}{2004}
\showcasematerialdesignsharp{videogame-asset-off}{2005}
\showcasematerialdesignsharp{video-call}{2006}
\showcasematerialdesignsharp{video-camera-back}{2007}
\showcasematerialdesignsharp{video-camera-front}{2008}
\showcasematerialdesignsharp{video-chat}{2009}
\showcasematerialdesignsharp{video-file}{2010}
\showcasematerialdesignsharp{video-label}{2011}
\showcasematerialdesignsharp{video-library}{2012}
\showcasematerialdesignsharp{video-settings}{2013}
\showcasematerialdesignsharp{video-stable}{2014}
\showcasematerialdesignsharp{view-agenda}{2015}
\showcasematerialdesignsharp{view-array}{2016}
\showcasematerialdesignsharp{view-carousel}{2017}
\showcasematerialdesignsharp{view-column}{2018}
\showcasematerialdesignsharp{view-comfy}{2019}
\showcasematerialdesignsharp{view-comfy-alt}{2020}
\showcasematerialdesignsharp{view-compact}{2021}
\showcasematerialdesignsharp{view-compact-alt}{2022}
\showcasematerialdesignsharp{view-cozy}{2023}
\showcasematerialdesignsharp{view-day}{2024}
\showcasematerialdesignsharp{view-headline}{2025}
\showcasematerialdesignsharp{view-in-ar}{2026}
\showcasematerialdesignsharp{view-kanban}{2027}
\showcasematerialdesignsharp{view-list}{2028}
\showcasematerialdesignsharp{view-module}{2029}
\showcasematerialdesignsharp{view-quilt}{2030}
\showcasematerialdesignsharp{view-sidebar}{2031}
\showcasematerialdesignsharp{view-stream}{2032}
\showcasematerialdesignsharp{view-timeline}{2033}
\showcasematerialdesignsharp{view-week}{2034}
\showcasematerialdesignsharp{vignette}{2035}
\showcasematerialdesignsharp{villa}{2036}
\showcasematerialdesignsharp{visibility}{2037}
\showcasematerialdesignsharp{visibility-off}{2038}
\showcasematerialdesignsharp{voicemail}{2039}
\showcasematerialdesignsharp{voice-chat}{2040}
\showcasematerialdesignsharp{voice-over-off}{2041}
\showcasematerialdesignsharp{volcano}{2042}
\showcasematerialdesignsharp{volume-down}{2043}
\showcasematerialdesignsharp{volume-mute}{2044}
\showcasematerialdesignsharp{volume-off}{2045}
\showcasematerialdesignsharp{volume-up}{2046}
\showcasematerialdesignsharp{volunteer-activism}{2047}
\showcasematerialdesignsharp{vpn-key}{2048}
\showcasematerialdesignsharp{vpn-key-off}{2049}
\showcasematerialdesignsharp{vpn-lock}{2050}
\showcasematerialdesignsharp{vrpano}{2051}
\showcasematerialdesignsharp{wallet}{2052}
\showcasematerialdesignsharp{wallpaper}{2053}
\showcasematerialdesignsharp{warehouse}{2054}
\showcasematerialdesignsharp{warning}{2055}
\showcasematerialdesignsharp{warning-amber}{2056}
\showcasematerialdesignsharp{wash}{2057}
\showcasematerialdesignsharp{watch}{2058}
\showcasematerialdesignsharp{watch-later}{2059}
\showcasematerialdesignsharp{watch-off}{2060}
\showcasematerialdesignsharp{water}{2061}
\showcasematerialdesignsharp{waterfall-chart}{2062}
\showcasematerialdesignsharp{water-damage}{2063}
\showcasematerialdesignsharp{water-drop}{2064}
\showcasematerialdesignsharp{waves}{2065}
\showcasematerialdesignsharp{waving-hand}{2066}
\showcasematerialdesignsharp{wb-auto}{2067}
\showcasematerialdesignsharp{wb-cloudy}{2068}
\showcasematerialdesignsharp{wb-incandescent}{2069}
\showcasematerialdesignsharp{wb-iridescent}{2070}
\showcasematerialdesignsharp{wb-shade}{2071}
\showcasematerialdesignsharp{wb-sunny}{2072}
\showcasematerialdesignsharp{wb-twilight}{2073}
\showcasematerialdesignsharp{wc}{2074}
\showcasematerialdesignsharp{web}{2075}
\showcasematerialdesignsharp{webhook}{2076}
\showcasematerialdesignsharp{web-asset}{2077}
\showcasematerialdesignsharp{web-asset-off}{2078}
\showcasematerialdesignsharp{web-stories}{2079}
\showcasematerialdesignsharp{weekend}{2080}
\showcasematerialdesignsharp{west}{2081}
\showcasematerialdesignsharp{whatshot}{2082}
\showcasematerialdesignsharp{wheelchair-pickup}{2083}
\showcasematerialdesignsharp{where-to-vote}{2084}
\showcasematerialdesignsharp{widgets}{2085}
\showcasematerialdesignsharp{width-full}{2086}
\showcasematerialdesignsharp{width-normal}{2087}
\showcasematerialdesignsharp{width-wide}{2088}
\showcasematerialdesignsharp{wifi}{2089}
\showcasematerialdesignsharp{wifi-1-bar}{2090}
\showcasematerialdesignsharp{wifi-2-bar}{2091}
\showcasematerialdesignsharp{wifi-calling}{2092}
\showcasematerialdesignsharp{wifi-calling-3}{2093}
\showcasematerialdesignsharp{wifi-channel}{2094}
\showcasematerialdesignsharp{wifi-find}{2095}
\showcasematerialdesignsharp{wifi-lock}{2096}
\showcasematerialdesignsharp{wifi-off}{2097}
\showcasematerialdesignsharp{wifi-password}{2098}
\showcasematerialdesignsharp{wifi-protected-setup}{2099}
\showcasematerialdesignsharp{wifi-tethering}{2100}
\showcasematerialdesignsharp{wifi-tethering-error}{2101}
\showcasematerialdesignsharp{wifi-tethering-off}{2102}
\showcasematerialdesignsharp{window}{2103}
\showcasematerialdesignsharp{wind-power}{2104}
\showcasematerialdesignsharp{wine-bar}{2105}
\showcasematerialdesignsharp{woman}{2106}
\showcasematerialdesignsharp{woman-2}{2107}
\showcasematerialdesignsharp{work}{2108}
\showcasematerialdesignsharp{workspaces}{2109}
\showcasematerialdesignsharp{workspace-premium}{2110}
\showcasematerialdesignsharp{work-history}{2111}
\showcasematerialdesignsharp{work-off}{2112}
\showcasematerialdesignsharp{work-outline}{2113}
\showcasematerialdesignsharp{wrap-text}{2114}
\showcasematerialdesignsharp{wrong-location}{2115}
\showcasematerialdesignsharp{wysiwyg}{2116}
\showcasematerialdesignsharp{yard}{2117}
\showcasematerialdesignsharp{youtube-searched-for}{2118}
\showcasematerialdesignsharp{zoom-in}{2119}
\showcasematerialdesignsharp{zoom-in-map}{2120}
\showcasematerialdesignsharp{zoom-out}{2121}
\showcasematerialdesignsharp{zoom-out-map}{2122}
\end{materialdesigncase}

\subsection{Two-Tone version (2122 icons)}

The \textsf{PDF} file is \texttt{materialdesign-twotone-all.pdf}

\begin{materialdesigncase}
\showcasematerialdesigntwotone{1k}{1}
\showcasematerialdesigntwotone{1k-plus}{2}
\showcasematerialdesigntwotone{1x-mobiledata}{3}
\showcasematerialdesigntwotone{2k}{4}
\showcasematerialdesigntwotone{2k-plus}{5}
\showcasematerialdesigntwotone{2mp}{6}
\showcasematerialdesigntwotone{3d-rotation}{7}
\showcasematerialdesigntwotone{3g-mobiledata}{8}
\showcasematerialdesigntwotone{3k}{9}
\showcasematerialdesigntwotone{3k-plus}{10}
\showcasematerialdesigntwotone{3mp}{11}
\showcasematerialdesigntwotone{3p}{12}
\showcasematerialdesigntwotone{4g-mobiledata}{13}
\showcasematerialdesigntwotone{4g-plus-mobiledata}{14}
\showcasematerialdesigntwotone{4k}{15}
\showcasematerialdesigntwotone{4k-plus}{16}
\showcasematerialdesigntwotone{4mp}{17}
\showcasematerialdesigntwotone{5g}{18}
\showcasematerialdesigntwotone{5k}{19}
\showcasematerialdesigntwotone{5k-plus}{20}
\showcasematerialdesigntwotone{5mp}{21}
\showcasematerialdesigntwotone{6k}{22}
\showcasematerialdesigntwotone{6k-plus}{23}
\showcasematerialdesigntwotone{6mp}{24}
\showcasematerialdesigntwotone{6-ft-apart}{25}
\showcasematerialdesigntwotone{7k}{26}
\showcasematerialdesigntwotone{7k-plus}{27}
\showcasematerialdesigntwotone{7mp}{28}
\showcasematerialdesigntwotone{8k}{29}
\showcasematerialdesigntwotone{8k-plus}{30}
\showcasematerialdesigntwotone{8mp}{31}
\showcasematerialdesigntwotone{9k}{32}
\showcasematerialdesigntwotone{9k-plus}{33}
\showcasematerialdesigntwotone{9mp}{34}
\showcasematerialdesigntwotone{10k}{35}
\showcasematerialdesigntwotone{10mp}{36}
\showcasematerialdesigntwotone{11mp}{37}
\showcasematerialdesigntwotone{12mp}{38}
\showcasematerialdesigntwotone{13mp}{39}
\showcasematerialdesigntwotone{14mp}{40}
\showcasematerialdesigntwotone{15mp}{41}
\showcasematerialdesigntwotone{16mp}{42}
\showcasematerialdesigntwotone{17mp}{43}
\showcasematerialdesigntwotone{18mp}{44}
\showcasematerialdesigntwotone{18-up-rating}{45}
\showcasematerialdesigntwotone{19mp}{46}
\showcasematerialdesigntwotone{20mp}{47}
\showcasematerialdesigntwotone{21mp}{48}
\showcasematerialdesigntwotone{22mp}{49}
\showcasematerialdesigntwotone{23mp}{50}
\showcasematerialdesigntwotone{24mp}{51}
\showcasematerialdesigntwotone{30fps}{52}
\showcasematerialdesigntwotone{30fps-select}{53}
\showcasematerialdesigntwotone{60fps}{54}
\showcasematerialdesigntwotone{60fps-select}{55}
\showcasematerialdesigntwotone{123}{56}
\showcasematerialdesigntwotone{360}{57}
\showcasematerialdesigntwotone{abc}{58}
\showcasematerialdesigntwotone{accessibility}{59}
\showcasematerialdesigntwotone{accessibility-new}{60}
\showcasematerialdesigntwotone{accessible}{61}
\showcasematerialdesigntwotone{accessible-forward}{62}
\showcasematerialdesigntwotone{access-alarm}{63}
\showcasematerialdesigntwotone{access-alarms}{64}
\showcasematerialdesigntwotone{access-time}{65}
\showcasematerialdesigntwotone{access-time-filled}{66}
\showcasematerialdesigntwotone{account-balance}{67}
\showcasematerialdesigntwotone{account-balance-wallet}{68}
\showcasematerialdesigntwotone{account-box}{69}
\showcasematerialdesigntwotone{account-circle}{70}
\showcasematerialdesigntwotone{account-tree}{71}
\showcasematerialdesigntwotone{ac-unit}{72}
\showcasematerialdesigntwotone{adb}{73}
\showcasematerialdesigntwotone{add}{74}
\showcasematerialdesigntwotone{addchart}{75}
\showcasematerialdesigntwotone{add-alarm}{76}
\showcasematerialdesigntwotone{add-alert}{77}
\showcasematerialdesigntwotone{add-a-photo}{78}
\showcasematerialdesigntwotone{add-box}{79}
\showcasematerialdesigntwotone{add-business}{80}
\showcasematerialdesigntwotone{add-card}{81}
\showcasematerialdesigntwotone{add-chart}{82}
\showcasematerialdesigntwotone{add-circle}{83}
\showcasematerialdesigntwotone{add-circle-outline}{84}
\showcasematerialdesigntwotone{add-comment}{85}
\showcasematerialdesigntwotone{add-home}{86}
\showcasematerialdesigntwotone{add-home-work}{87}
\showcasematerialdesigntwotone{add-ic-call}{88}
\showcasematerialdesigntwotone{add-link}{89}
\showcasematerialdesigntwotone{add-location}{90}
\showcasematerialdesigntwotone{add-location-alt}{91}
\showcasematerialdesigntwotone{add-moderator}{92}
\showcasematerialdesigntwotone{add-photo-alternate}{93}
\showcasematerialdesigntwotone{add-reaction}{94}
\showcasematerialdesigntwotone{add-road}{95}
\showcasematerialdesigntwotone{add-shopping-cart}{96}
\showcasematerialdesigntwotone{add-task}{97}
\showcasematerialdesigntwotone{add-to-drive}{98}
\showcasematerialdesigntwotone{add-to-home-screen}{99}
\showcasematerialdesigntwotone{add-to-photos}{100}
\showcasematerialdesigntwotone{add-to-queue}{101}
\showcasematerialdesigntwotone{adf-scanner}{102}
\showcasematerialdesigntwotone{adjust}{103}
\showcasematerialdesigntwotone{admin-panel-settings}{104}
\showcasematerialdesigntwotone{ads-click}{105}
\showcasematerialdesigntwotone{ad-units}{106}
\showcasematerialdesigntwotone{agriculture}{107}
\showcasematerialdesigntwotone{air}{108}
\showcasematerialdesigntwotone{airlines}{109}
\showcasematerialdesigntwotone{airline-seat-flat}{110}
\showcasematerialdesigntwotone{airline-seat-flat-angled}{111}
\showcasematerialdesigntwotone{airline-seat-individual-suite}{112}
\showcasematerialdesigntwotone{airline-seat-legroom-extra}{113}
\showcasematerialdesigntwotone{airline-seat-legroom-normal}{114}
\showcasematerialdesigntwotone{airline-seat-legroom-reduced}{115}
\showcasematerialdesigntwotone{airline-seat-recline-extra}{116}
\showcasematerialdesigntwotone{airline-seat-recline-normal}{117}
\showcasematerialdesigntwotone{airline-stops}{118}
\showcasematerialdesigntwotone{airplanemode-active}{119}
\showcasematerialdesigntwotone{airplanemode-inactive}{120}
\showcasematerialdesigntwotone{airplane-ticket}{121}
\showcasematerialdesigntwotone{airplay}{122}
\showcasematerialdesigntwotone{airport-shuttle}{123}
\showcasematerialdesigntwotone{alarm}{124}
\showcasematerialdesigntwotone{alarm-add}{125}
\showcasematerialdesigntwotone{alarm-off}{126}
\showcasematerialdesigntwotone{alarm-on}{127}
\showcasematerialdesigntwotone{album}{128}
\showcasematerialdesigntwotone{align-horizontal-center}{129}
\showcasematerialdesigntwotone{align-horizontal-left}{130}
\showcasematerialdesigntwotone{align-horizontal-right}{131}
\showcasematerialdesigntwotone{align-vertical-bottom}{132}
\showcasematerialdesigntwotone{align-vertical-center}{133}
\showcasematerialdesigntwotone{align-vertical-top}{134}
\showcasematerialdesigntwotone{all-inbox}{135}
\showcasematerialdesigntwotone{all-inclusive}{136}
\showcasematerialdesigntwotone{all-out}{137}
\showcasematerialdesigntwotone{alternate-email}{138}
\showcasematerialdesigntwotone{alt-route}{139}
\showcasematerialdesigntwotone{analytics}{140}
\showcasematerialdesigntwotone{anchor}{141}
\showcasematerialdesigntwotone{android}{142}
\showcasematerialdesigntwotone{animation}{143}
\showcasematerialdesigntwotone{announcement}{144}
\showcasematerialdesigntwotone{aod}{145}
\showcasematerialdesigntwotone{apartment}{146}
\showcasematerialdesigntwotone{api}{147}
\showcasematerialdesigntwotone{approval}{148}
\showcasematerialdesigntwotone{apps}{149}
\showcasematerialdesigntwotone{apps-outage}{150}
\showcasematerialdesigntwotone{app-blocking}{151}
\showcasematerialdesigntwotone{app-registration}{152}
\showcasematerialdesigntwotone{app-settings-alt}{153}
\showcasematerialdesigntwotone{app-shortcut}{154}
\showcasematerialdesigntwotone{architecture}{155}
\showcasematerialdesigntwotone{archive}{156}
\showcasematerialdesigntwotone{area-chart}{157}
\showcasematerialdesigntwotone{arrow-back}{158}
\showcasematerialdesigntwotone{arrow-back-ios}{159}
\showcasematerialdesigntwotone{arrow-back-ios-new}{160}
\showcasematerialdesigntwotone{arrow-circle-down}{161}
\showcasematerialdesigntwotone{arrow-circle-left}{162}
\showcasematerialdesigntwotone{arrow-circle-right}{163}
\showcasematerialdesigntwotone{arrow-circle-up}{164}
\showcasematerialdesigntwotone{arrow-downward}{165}
\showcasematerialdesigntwotone{arrow-drop-down}{166}
\showcasematerialdesigntwotone{arrow-drop-down-circle}{167}
\showcasematerialdesigntwotone{arrow-drop-up}{168}
\showcasematerialdesigntwotone{arrow-forward}{169}
\showcasematerialdesigntwotone{arrow-forward-ios}{170}
\showcasematerialdesigntwotone{arrow-left}{171}
\showcasematerialdesigntwotone{arrow-outward}{172}
\showcasematerialdesigntwotone{arrow-right}{173}
\showcasematerialdesigntwotone{arrow-right-alt}{174}
\showcasematerialdesigntwotone{arrow-upward}{175}
\showcasematerialdesigntwotone{article}{176}
\showcasematerialdesigntwotone{art-track}{177}
\showcasematerialdesigntwotone{aspect-ratio}{178}
\showcasematerialdesigntwotone{assessment}{179}
\showcasematerialdesigntwotone{assignment}{180}
\showcasematerialdesigntwotone{assignment-ind}{181}
\showcasematerialdesigntwotone{assignment-late}{182}
\showcasematerialdesigntwotone{assignment-return}{183}
\showcasematerialdesigntwotone{assignment-returned}{184}
\showcasematerialdesigntwotone{assignment-turned-in}{185}
\showcasematerialdesigntwotone{assistant}{186}
\showcasematerialdesigntwotone{assistant-direction}{187}
\showcasematerialdesigntwotone{assistant-photo}{188}
\showcasematerialdesigntwotone{assist-walker}{189}
\showcasematerialdesigntwotone{assured-workload}{190}
\showcasematerialdesigntwotone{atm}{191}
\showcasematerialdesigntwotone{attachment}{192}
\showcasematerialdesigntwotone{attach-email}{193}
\showcasematerialdesigntwotone{attach-file}{194}
\showcasematerialdesigntwotone{attach-money}{195}
\showcasematerialdesigntwotone{attractions}{196}
\showcasematerialdesigntwotone{attribution}{197}
\showcasematerialdesigntwotone{audiotrack}{198}
\showcasematerialdesigntwotone{audio-file}{199}
\showcasematerialdesigntwotone{autofps-select}{200}
\showcasematerialdesigntwotone{autorenew}{201}
\showcasematerialdesigntwotone{auto-awesome}{202}
\showcasematerialdesigntwotone{auto-awesome-mosaic}{203}
\showcasematerialdesigntwotone{auto-awesome-motion}{204}
\showcasematerialdesigntwotone{auto-delete}{205}
\showcasematerialdesigntwotone{auto-fix-high}{206}
\showcasematerialdesigntwotone{auto-fix-normal}{207}
\showcasematerialdesigntwotone{auto-fix-off}{208}
\showcasematerialdesigntwotone{auto-graph}{209}
\showcasematerialdesigntwotone{auto-mode}{210}
\showcasematerialdesigntwotone{auto-stories}{211}
\showcasematerialdesigntwotone{av-timer}{212}
\showcasematerialdesigntwotone{baby-changing-station}{213}
\showcasematerialdesigntwotone{backpack}{214}
\showcasematerialdesigntwotone{backspace}{215}
\showcasematerialdesigntwotone{backup}{216}
\showcasematerialdesigntwotone{backup-table}{217}
\showcasematerialdesigntwotone{back-hand}{218}
\showcasematerialdesigntwotone{badge}{219}
\showcasematerialdesigntwotone{bakery-dining}{220}
\showcasematerialdesigntwotone{balance}{221}
\showcasematerialdesigntwotone{balcony}{222}
\showcasematerialdesigntwotone{ballot}{223}
\showcasematerialdesigntwotone{bar-chart}{224}
\showcasematerialdesigntwotone{batch-prediction}{225}
\showcasematerialdesigntwotone{bathroom}{226}
\showcasematerialdesigntwotone{bathtub}{227}
\showcasematerialdesigntwotone{battery-0-bar}{228}
\showcasematerialdesigntwotone{battery-1-bar}{229}
\showcasematerialdesigntwotone{battery-2-bar}{230}
\showcasematerialdesigntwotone{battery-3-bar}{231}
\showcasematerialdesigntwotone{battery-4-bar}{232}
\showcasematerialdesigntwotone{battery-5-bar}{233}
\showcasematerialdesigntwotone{battery-6-bar}{234}
\showcasematerialdesigntwotone{battery-alert}{235}
\showcasematerialdesigntwotone{battery-charging-full}{236}
\showcasematerialdesigntwotone{battery-full}{237}
\showcasematerialdesigntwotone{battery-saver}{238}
\showcasematerialdesigntwotone{battery-std}{239}
\showcasematerialdesigntwotone{battery-unknown}{240}
\showcasematerialdesigntwotone{beach-access}{241}
\showcasematerialdesigntwotone{bed}{242}
\showcasematerialdesigntwotone{bedroom-baby}{243}
\showcasematerialdesigntwotone{bedroom-child}{244}
\showcasematerialdesigntwotone{bedroom-parent}{245}
\showcasematerialdesigntwotone{bedtime}{246}
\showcasematerialdesigntwotone{bedtime-off}{247}
\showcasematerialdesigntwotone{beenhere}{248}
\showcasematerialdesigntwotone{bento}{249}
\showcasematerialdesigntwotone{bike-scooter}{250}
\showcasematerialdesigntwotone{biotech}{251}
\showcasematerialdesigntwotone{blender}{252}
\showcasematerialdesigntwotone{blind}{253}
\showcasematerialdesigntwotone{blinds}{254}
\showcasematerialdesigntwotone{blinds-closed}{255}
\showcasematerialdesigntwotone{block}{256}
\showcasematerialdesigntwotone{bloodtype}{257}
\showcasematerialdesigntwotone{bluetooth}{258}
\showcasematerialdesigntwotone{bluetooth-audio}{259}
\showcasematerialdesigntwotone{bluetooth-connected}{260}
\showcasematerialdesigntwotone{bluetooth-disabled}{261}
\showcasematerialdesigntwotone{bluetooth-drive}{262}
\showcasematerialdesigntwotone{bluetooth-searching}{263}
\showcasematerialdesigntwotone{blur-circular}{264}
\showcasematerialdesigntwotone{blur-linear}{265}
\showcasematerialdesigntwotone{blur-off}{266}
\showcasematerialdesigntwotone{blur-on}{267}
\showcasematerialdesigntwotone{bolt}{268}
\showcasematerialdesigntwotone{book}{269}
\showcasematerialdesigntwotone{bookmark}{270}
\showcasematerialdesigntwotone{bookmarks}{271}
\showcasematerialdesigntwotone{bookmark-add}{272}
\showcasematerialdesigntwotone{bookmark-added}{273}
\showcasematerialdesigntwotone{bookmark-border}{274}
\showcasematerialdesigntwotone{bookmark-remove}{275}
\showcasematerialdesigntwotone{book-online}{276}
\showcasematerialdesigntwotone{border-all}{277}
\showcasematerialdesigntwotone{border-bottom}{278}
\showcasematerialdesigntwotone{border-clear}{279}
\showcasematerialdesigntwotone{border-color}{280}
\showcasematerialdesigntwotone{border-horizontal}{281}
\showcasematerialdesigntwotone{border-inner}{282}
\showcasematerialdesigntwotone{border-left}{283}
\showcasematerialdesigntwotone{border-outer}{284}
\showcasematerialdesigntwotone{border-right}{285}
\showcasematerialdesigntwotone{border-style}{286}
\showcasematerialdesigntwotone{border-top}{287}
\showcasematerialdesigntwotone{border-vertical}{288}
\showcasematerialdesigntwotone{boy}{289}
\showcasematerialdesigntwotone{branding-watermark}{290}
\showcasematerialdesigntwotone{breakfast-dining}{291}
\showcasematerialdesigntwotone{brightness-1}{292}
\showcasematerialdesigntwotone{brightness-2}{293}
\showcasematerialdesigntwotone{brightness-3}{294}
\showcasematerialdesigntwotone{brightness-4}{295}
\showcasematerialdesigntwotone{brightness-5}{296}
\showcasematerialdesigntwotone{brightness-6}{297}
\showcasematerialdesigntwotone{brightness-7}{298}
\showcasematerialdesigntwotone{brightness-auto}{299}
\showcasematerialdesigntwotone{brightness-high}{300}
\showcasematerialdesigntwotone{brightness-low}{301}
\showcasematerialdesigntwotone{brightness-medium}{302}
\showcasematerialdesigntwotone{broadcast-on-home}{303}
\showcasematerialdesigntwotone{broadcast-on-personal}{304}
\showcasematerialdesigntwotone{broken-image}{305}
\showcasematerialdesigntwotone{browser-not-supported}{306}
\showcasematerialdesigntwotone{browser-updated}{307}
\showcasematerialdesigntwotone{browse-gallery}{308}
\showcasematerialdesigntwotone{brunch-dining}{309}
\showcasematerialdesigntwotone{brush}{310}
\showcasematerialdesigntwotone{bubble-chart}{311}
\showcasematerialdesigntwotone{bug-report}{312}
\showcasematerialdesigntwotone{build}{313}
\showcasematerialdesigntwotone{build-circle}{314}
\showcasematerialdesigntwotone{bungalow}{315}
\showcasematerialdesigntwotone{burst-mode}{316}
\showcasematerialdesigntwotone{business}{317}
\showcasematerialdesigntwotone{business-center}{318}
\showcasematerialdesigntwotone{bus-alert}{319}
\showcasematerialdesigntwotone{cabin}{320}
\showcasematerialdesigntwotone{cable}{321}
\showcasematerialdesigntwotone{cached}{322}
\showcasematerialdesigntwotone{cake}{323}
\showcasematerialdesigntwotone{calculate}{324}
\showcasematerialdesigntwotone{calendar-month}{325}
\showcasematerialdesigntwotone{calendar-today}{326}
\showcasematerialdesigntwotone{calendar-view-day}{327}
\showcasematerialdesigntwotone{calendar-view-month}{328}
\showcasematerialdesigntwotone{calendar-view-week}{329}
\showcasematerialdesigntwotone{call}{330}
\showcasematerialdesigntwotone{call-end}{331}
\showcasematerialdesigntwotone{call-made}{332}
\showcasematerialdesigntwotone{call-merge}{333}
\showcasematerialdesigntwotone{call-missed}{334}
\showcasematerialdesigntwotone{call-missed-outgoing}{335}
\showcasematerialdesigntwotone{call-received}{336}
\showcasematerialdesigntwotone{call-split}{337}
\showcasematerialdesigntwotone{call-to-action}{338}
\showcasematerialdesigntwotone{camera}{339}
\showcasematerialdesigntwotone{cameraswitch}{340}
\showcasematerialdesigntwotone{camera-alt}{341}
\showcasematerialdesigntwotone{camera-enhance}{342}
\showcasematerialdesigntwotone{camera-front}{343}
\showcasematerialdesigntwotone{camera-indoor}{344}
\showcasematerialdesigntwotone{camera-outdoor}{345}
\showcasematerialdesigntwotone{camera-rear}{346}
\showcasematerialdesigntwotone{camera-roll}{347}
\showcasematerialdesigntwotone{campaign}{348}
\showcasematerialdesigntwotone{cancel}{349}
\showcasematerialdesigntwotone{cancel-presentation}{350}
\showcasematerialdesigntwotone{cancel-schedule-send}{351}
\showcasematerialdesigntwotone{candlestick-chart}{352}
\showcasematerialdesigntwotone{card-giftcard}{353}
\showcasematerialdesigntwotone{card-membership}{354}
\showcasematerialdesigntwotone{card-travel}{355}
\showcasematerialdesigntwotone{carpenter}{356}
\showcasematerialdesigntwotone{car-crash}{357}
\showcasematerialdesigntwotone{car-rental}{358}
\showcasematerialdesigntwotone{car-repair}{359}
\showcasematerialdesigntwotone{cases}{360}
\showcasematerialdesigntwotone{casino}{361}
\showcasematerialdesigntwotone{cast}{362}
\showcasematerialdesigntwotone{castle}{363}
\showcasematerialdesigntwotone{cast-connected}{364}
\showcasematerialdesigntwotone{cast-for-education}{365}
\showcasematerialdesigntwotone{catching-pokemon}{366}
\showcasematerialdesigntwotone{category}{367}
\showcasematerialdesigntwotone{celebration}{368}
\showcasematerialdesigntwotone{cell-tower}{369}
\showcasematerialdesigntwotone{cell-wifi}{370}
\showcasematerialdesigntwotone{center-focus-strong}{371}
\showcasematerialdesigntwotone{center-focus-weak}{372}
\showcasematerialdesigntwotone{chair}{373}
\showcasematerialdesigntwotone{chair-alt}{374}
\showcasematerialdesigntwotone{chalet}{375}
\showcasematerialdesigntwotone{change-circle}{376}
\showcasematerialdesigntwotone{change-history}{377}
\showcasematerialdesigntwotone{charging-station}{378}
\showcasematerialdesigntwotone{chat}{379}
\showcasematerialdesigntwotone{chat-bubble}{380}
\showcasematerialdesigntwotone{chat-bubble-outline}{381}
\showcasematerialdesigntwotone{check}{382}
\showcasematerialdesigntwotone{checklist}{383}
\showcasematerialdesigntwotone{checklist-rtl}{384}
\showcasematerialdesigntwotone{checkroom}{385}
\showcasematerialdesigntwotone{check-box}{386}
\showcasematerialdesigntwotone{check-box-outline-blank}{387}
\showcasematerialdesigntwotone{check-circle}{388}
\showcasematerialdesigntwotone{check-circle-outline}{389}
\showcasematerialdesigntwotone{chevron-left}{390}
\showcasematerialdesigntwotone{chevron-right}{391}
\showcasematerialdesigntwotone{child-care}{392}
\showcasematerialdesigntwotone{child-friendly}{393}
\showcasematerialdesigntwotone{chrome-reader-mode}{394}
\showcasematerialdesigntwotone{church}{395}
\showcasematerialdesigntwotone{circle}{396}
\showcasematerialdesigntwotone{circle-notifications}{397}
\showcasematerialdesigntwotone{class}{398}
\showcasematerialdesigntwotone{cleaning-services}{399}
\showcasematerialdesigntwotone{clean-hands}{400}
\showcasematerialdesigntwotone{clear}{401}
\showcasematerialdesigntwotone{clear-all}{402}
\showcasematerialdesigntwotone{close}{403}
\showcasematerialdesigntwotone{closed-caption}{404}
\showcasematerialdesigntwotone{closed-caption-disabled}{405}
\showcasematerialdesigntwotone{closed-caption-off}{406}
\showcasematerialdesigntwotone{close-fullscreen}{407}
\showcasematerialdesigntwotone{cloud}{408}
\showcasematerialdesigntwotone{cloud-circle}{409}
\showcasematerialdesigntwotone{cloud-done}{410}
\showcasematerialdesigntwotone{cloud-download}{411}
\showcasematerialdesigntwotone{cloud-off}{412}
\showcasematerialdesigntwotone{cloud-queue}{413}
\showcasematerialdesigntwotone{cloud-sync}{414}
\showcasematerialdesigntwotone{cloud-upload}{415}
\showcasematerialdesigntwotone{co2}{416}
\showcasematerialdesigntwotone{code}{417}
\showcasematerialdesigntwotone{code-off}{418}
\showcasematerialdesigntwotone{coffee}{419}
\showcasematerialdesigntwotone{coffee-maker}{420}
\showcasematerialdesigntwotone{collections}{421}
\showcasematerialdesigntwotone{collections-bookmark}{422}
\showcasematerialdesigntwotone{colorize}{423}
\showcasematerialdesigntwotone{color-lens}{424}
\showcasematerialdesigntwotone{comment}{425}
\showcasematerialdesigntwotone{comments-disabled}{426}
\showcasematerialdesigntwotone{comment-bank}{427}
\showcasematerialdesigntwotone{commit}{428}
\showcasematerialdesigntwotone{commute}{429}
\showcasematerialdesigntwotone{compare}{430}
\showcasematerialdesigntwotone{compare-arrows}{431}
\showcasematerialdesigntwotone{compass-calibration}{432}
\showcasematerialdesigntwotone{compost}{433}
\showcasematerialdesigntwotone{compress}{434}
\showcasematerialdesigntwotone{computer}{435}
\showcasematerialdesigntwotone{confirmation-number}{436}
\showcasematerialdesigntwotone{connected-tv}{437}
\showcasematerialdesigntwotone{connecting-airports}{438}
\showcasematerialdesigntwotone{connect-without-contact}{439}
\showcasematerialdesigntwotone{construction}{440}
\showcasematerialdesigntwotone{contactless}{441}
\showcasematerialdesigntwotone{contacts}{442}
\showcasematerialdesigntwotone{contact-emergency}{443}
\showcasematerialdesigntwotone{contact-mail}{444}
\showcasematerialdesigntwotone{contact-page}{445}
\showcasematerialdesigntwotone{contact-phone}{446}
\showcasematerialdesigntwotone{contact-support}{447}
\showcasematerialdesigntwotone{content-copy}{448}
\showcasematerialdesigntwotone{content-cut}{449}
\showcasematerialdesigntwotone{content-paste}{450}
\showcasematerialdesigntwotone{content-paste-go}{451}
\showcasematerialdesigntwotone{content-paste-off}{452}
\showcasematerialdesigntwotone{content-paste-search}{453}
\showcasematerialdesigntwotone{contrast}{454}
\showcasematerialdesigntwotone{control-camera}{455}
\showcasematerialdesigntwotone{control-point}{456}
\showcasematerialdesigntwotone{control-point-duplicate}{457}
\showcasematerialdesigntwotone{cookie}{458}
\showcasematerialdesigntwotone{copyright}{459}
\showcasematerialdesigntwotone{copy-all}{460}
\showcasematerialdesigntwotone{coronavirus}{461}
\showcasematerialdesigntwotone{corporate-fare}{462}
\showcasematerialdesigntwotone{cottage}{463}
\showcasematerialdesigntwotone{countertops}{464}
\showcasematerialdesigntwotone{co-present}{465}
\showcasematerialdesigntwotone{create}{466}
\showcasematerialdesigntwotone{create-new-folder}{467}
\showcasematerialdesigntwotone{credit-card}{468}
\showcasematerialdesigntwotone{credit-card-off}{469}
\showcasematerialdesigntwotone{credit-score}{470}
\showcasematerialdesigntwotone{crib}{471}
\showcasematerialdesigntwotone{crisis-alert}{472}
\showcasematerialdesigntwotone{crop}{473}
\showcasematerialdesigntwotone{crop-3-2}{474}
\showcasematerialdesigntwotone{crop-5-4}{475}
\showcasematerialdesigntwotone{crop-7-5}{476}
\showcasematerialdesigntwotone{crop-16-9}{477}
\showcasematerialdesigntwotone{crop-din}{478}
\showcasematerialdesigntwotone{crop-free}{479}
\showcasematerialdesigntwotone{crop-landscape}{480}
\showcasematerialdesigntwotone{crop-original}{481}
\showcasematerialdesigntwotone{crop-portrait}{482}
\showcasematerialdesigntwotone{crop-rotate}{483}
\showcasematerialdesigntwotone{crop-square}{484}
\showcasematerialdesigntwotone{cruelty-free}{485}
\showcasematerialdesigntwotone{css}{486}
\showcasematerialdesigntwotone{currency-bitcoin}{487}
\showcasematerialdesigntwotone{currency-exchange}{488}
\showcasematerialdesigntwotone{currency-franc}{489}
\showcasematerialdesigntwotone{currency-lira}{490}
\showcasematerialdesigntwotone{currency-pound}{491}
\showcasematerialdesigntwotone{currency-ruble}{492}
\showcasematerialdesigntwotone{currency-rupee}{493}
\showcasematerialdesigntwotone{currency-yen}{494}
\showcasematerialdesigntwotone{currency-yuan}{495}
\showcasematerialdesigntwotone{curtains}{496}
\showcasematerialdesigntwotone{curtains-closed}{497}
\showcasematerialdesigntwotone{cyclone}{498}
\showcasematerialdesigntwotone{dangerous}{499}
\showcasematerialdesigntwotone{dark-mode}{500}
\showcasematerialdesigntwotone{dashboard}{501}
\showcasematerialdesigntwotone{dashboard-customize}{502}
\showcasematerialdesigntwotone{dataset}{503}
\showcasematerialdesigntwotone{dataset-linked}{504}
\showcasematerialdesigntwotone{data-array}{505}
\showcasematerialdesigntwotone{data-exploration}{506}
\showcasematerialdesigntwotone{data-object}{507}
\showcasematerialdesigntwotone{data-saver-off}{508}
\showcasematerialdesigntwotone{data-saver-on}{509}
\showcasematerialdesigntwotone{data-thresholding}{510}
\showcasematerialdesigntwotone{data-usage}{511}
\showcasematerialdesigntwotone{date-range}{512}
\showcasematerialdesigntwotone{deblur}{513}
\showcasematerialdesigntwotone{deck}{514}
\showcasematerialdesigntwotone{dehaze}{515}
\showcasematerialdesigntwotone{delete}{516}
\showcasematerialdesigntwotone{delete-forever}{517}
\showcasematerialdesigntwotone{delete-outline}{518}
\showcasematerialdesigntwotone{delete-sweep}{519}
\showcasematerialdesigntwotone{delivery-dining}{520}
\showcasematerialdesigntwotone{density-large}{521}
\showcasematerialdesigntwotone{density-medium}{522}
\showcasematerialdesigntwotone{density-small}{523}
\showcasematerialdesigntwotone{departure-board}{524}
\showcasematerialdesigntwotone{description}{525}
\showcasematerialdesigntwotone{deselect}{526}
\showcasematerialdesigntwotone{design-services}{527}
\showcasematerialdesigntwotone{desk}{528}
\showcasematerialdesigntwotone{desktop-access-disabled}{529}
\showcasematerialdesigntwotone{desktop-mac}{530}
\showcasematerialdesigntwotone{desktop-windows}{531}
\showcasematerialdesigntwotone{details}{532}
\showcasematerialdesigntwotone{developer-board}{533}
\showcasematerialdesigntwotone{developer-board-off}{534}
\showcasematerialdesigntwotone{developer-mode}{535}
\showcasematerialdesigntwotone{devices}{536}
\showcasematerialdesigntwotone{devices-fold}{537}
\showcasematerialdesigntwotone{devices-other}{538}
\showcasematerialdesigntwotone{device-hub}{539}
\showcasematerialdesigntwotone{device-thermostat}{540}
\showcasematerialdesigntwotone{device-unknown}{541}
\showcasematerialdesigntwotone{dialer-sip}{542}
\showcasematerialdesigntwotone{dialpad}{543}
\showcasematerialdesigntwotone{diamond}{544}
\showcasematerialdesigntwotone{difference}{545}
\showcasematerialdesigntwotone{dining}{546}
\showcasematerialdesigntwotone{dinner-dining}{547}
\showcasematerialdesigntwotone{directions}{548}
\showcasematerialdesigntwotone{directions-bike}{549}
\showcasematerialdesigntwotone{directions-boat}{550}
\showcasematerialdesigntwotone{directions-boat-filled}{551}
\showcasematerialdesigntwotone{directions-bus}{552}
\showcasematerialdesigntwotone{directions-bus-filled}{553}
\showcasematerialdesigntwotone{directions-car}{554}
\showcasematerialdesigntwotone{directions-car-filled}{555}
\showcasematerialdesigntwotone{directions-off}{556}
\showcasematerialdesigntwotone{directions-railway}{557}
\showcasematerialdesigntwotone{directions-railway-filled}{558}
\showcasematerialdesigntwotone{directions-run}{559}
\showcasematerialdesigntwotone{directions-subway}{560}
\showcasematerialdesigntwotone{directions-subway-filled}{561}
\showcasematerialdesigntwotone{directions-transit}{562}
\showcasematerialdesigntwotone{directions-transit-filled}{563}
\showcasematerialdesigntwotone{directions-walk}{564}
\showcasematerialdesigntwotone{dirty-lens}{565}
\showcasematerialdesigntwotone{disabled-by-default}{566}
\showcasematerialdesigntwotone{disabled-visible}{567}
\showcasematerialdesigntwotone{discount}{568}
\showcasematerialdesigntwotone{disc-full}{569}
\showcasematerialdesigntwotone{display-settings}{570}
\showcasematerialdesigntwotone{diversity-1}{571}
\showcasematerialdesigntwotone{diversity-2}{572}
\showcasematerialdesigntwotone{diversity-3}{573}
\showcasematerialdesigntwotone{dns}{574}
\showcasematerialdesigntwotone{dock}{575}
\showcasematerialdesigntwotone{document-scanner}{576}
\showcasematerialdesigntwotone{domain}{577}
\showcasematerialdesigntwotone{domain-add}{578}
\showcasematerialdesigntwotone{domain-disabled}{579}
\showcasematerialdesigntwotone{domain-verification}{580}
\showcasematerialdesigntwotone{done}{581}
\showcasematerialdesigntwotone{done-all}{582}
\showcasematerialdesigntwotone{done-outline}{583}
\showcasematerialdesigntwotone{donut-large}{584}
\showcasematerialdesigntwotone{donut-small}{585}
\showcasematerialdesigntwotone{doorbell}{586}
\showcasematerialdesigntwotone{door-back}{587}
\showcasematerialdesigntwotone{door-front}{588}
\showcasematerialdesigntwotone{door-sliding}{589}
\showcasematerialdesigntwotone{double-arrow}{590}
\showcasematerialdesigntwotone{downhill-skiing}{591}
\showcasematerialdesigntwotone{download}{592}
\showcasematerialdesigntwotone{downloading}{593}
\showcasematerialdesigntwotone{download-done}{594}
\showcasematerialdesigntwotone{download-for-offline}{595}
\showcasematerialdesigntwotone{do-disturb}{596}
\showcasematerialdesigntwotone{do-disturb-alt}{597}
\showcasematerialdesigntwotone{do-disturb-off}{598}
\showcasematerialdesigntwotone{do-disturb-on}{599}
\showcasematerialdesigntwotone{do-not-disturb}{600}
\showcasematerialdesigntwotone{do-not-disturb-alt}{601}
\showcasematerialdesigntwotone{do-not-disturb-off}{602}
\showcasematerialdesigntwotone{do-not-disturb-on}{603}
\showcasematerialdesigntwotone{do-not-disturb-on-total-silence}{604}
\showcasematerialdesigntwotone{do-not-step}{605}
\showcasematerialdesigntwotone{do-not-touch}{606}
\showcasematerialdesigntwotone{drafts}{607}
\showcasematerialdesigntwotone{drag-handle}{608}
\showcasematerialdesigntwotone{drag-indicator}{609}
\showcasematerialdesigntwotone{draw}{610}
\showcasematerialdesigntwotone{drive-eta}{611}
\showcasematerialdesigntwotone{drive-file-move}{612}
\showcasematerialdesigntwotone{drive-file-move-rtl}{613}
\showcasematerialdesigntwotone{drive-file-rename-outline}{614}
\showcasematerialdesigntwotone{drive-folder-upload}{615}
\showcasematerialdesigntwotone{dry}{616}
\showcasematerialdesigntwotone{dry-cleaning}{617}
\showcasematerialdesigntwotone{duo}{618}
\showcasematerialdesigntwotone{dvr}{619}
\showcasematerialdesigntwotone{dynamic-feed}{620}
\showcasematerialdesigntwotone{dynamic-form}{621}
\showcasematerialdesigntwotone{earbuds}{622}
\showcasematerialdesigntwotone{earbuds-battery}{623}
\showcasematerialdesigntwotone{east}{624}
\showcasematerialdesigntwotone{edgesensor-high}{625}
\showcasematerialdesigntwotone{edgesensor-low}{626}
\showcasematerialdesigntwotone{edit}{627}
\showcasematerialdesigntwotone{edit-attributes}{628}
\showcasematerialdesigntwotone{edit-calendar}{629}
\showcasematerialdesigntwotone{edit-location}{630}
\showcasematerialdesigntwotone{edit-location-alt}{631}
\showcasematerialdesigntwotone{edit-note}{632}
\showcasematerialdesigntwotone{edit-notifications}{633}
\showcasematerialdesigntwotone{edit-off}{634}
\showcasematerialdesigntwotone{edit-road}{635}
\showcasematerialdesigntwotone{egg}{636}
\showcasematerialdesigntwotone{egg-alt}{637}
\showcasematerialdesigntwotone{eject}{638}
\showcasematerialdesigntwotone{elderly}{639}
\showcasematerialdesigntwotone{elderly-woman}{640}
\showcasematerialdesigntwotone{electrical-services}{641}
\showcasematerialdesigntwotone{electric-bike}{642}
\showcasematerialdesigntwotone{electric-bolt}{643}
\showcasematerialdesigntwotone{electric-car}{644}
\showcasematerialdesigntwotone{electric-meter}{645}
\showcasematerialdesigntwotone{electric-moped}{646}
\showcasematerialdesigntwotone{electric-rickshaw}{647}
\showcasematerialdesigntwotone{electric-scooter}{648}
\showcasematerialdesigntwotone{elevator}{649}
\showcasematerialdesigntwotone{email}{650}
\showcasematerialdesigntwotone{emergency}{651}
\showcasematerialdesigntwotone{emergency-recording}{652}
\showcasematerialdesigntwotone{emergency-share}{653}
\showcasematerialdesigntwotone{emoji-emotions}{654}
\showcasematerialdesigntwotone{emoji-events}{655}
\showcasematerialdesigntwotone{emoji-food-beverage}{656}
\showcasematerialdesigntwotone{emoji-nature}{657}
\showcasematerialdesigntwotone{emoji-objects}{658}
\showcasematerialdesigntwotone{emoji-people}{659}
\showcasematerialdesigntwotone{emoji-symbols}{660}
\showcasematerialdesigntwotone{emoji-transportation}{661}
\showcasematerialdesigntwotone{energy-savings-leaf}{662}
\showcasematerialdesigntwotone{engineering}{663}
\showcasematerialdesigntwotone{enhanced-encryption}{664}
\showcasematerialdesigntwotone{equalizer}{665}
\showcasematerialdesigntwotone{error}{666}
\showcasematerialdesigntwotone{error-outline}{667}
\showcasematerialdesigntwotone{escalator}{668}
\showcasematerialdesigntwotone{escalator-warning}{669}
\showcasematerialdesigntwotone{euro}{670}
\showcasematerialdesigntwotone{euro-symbol}{671}
\showcasematerialdesigntwotone{event}{672}
\showcasematerialdesigntwotone{event-available}{673}
\showcasematerialdesigntwotone{event-busy}{674}
\showcasematerialdesigntwotone{event-note}{675}
\showcasematerialdesigntwotone{event-repeat}{676}
\showcasematerialdesigntwotone{event-seat}{677}
\showcasematerialdesigntwotone{ev-station}{678}
\showcasematerialdesigntwotone{exit-to-app}{679}
\showcasematerialdesigntwotone{expand}{680}
\showcasematerialdesigntwotone{expand-circle-down}{681}
\showcasematerialdesigntwotone{expand-less}{682}
\showcasematerialdesigntwotone{expand-more}{683}
\showcasematerialdesigntwotone{explicit}{684}
\showcasematerialdesigntwotone{explore}{685}
\showcasematerialdesigntwotone{explore-off}{686}
\showcasematerialdesigntwotone{exposure}{687}
\showcasematerialdesigntwotone{exposure-neg-1}{688}
\showcasematerialdesigntwotone{exposure-neg-2}{689}
\showcasematerialdesigntwotone{exposure-plus-1}{690}
\showcasematerialdesigntwotone{exposure-plus-2}{691}
\showcasematerialdesigntwotone{exposure-zero}{692}
\showcasematerialdesigntwotone{extension}{693}
\showcasematerialdesigntwotone{extension-off}{694}
\showcasematerialdesigntwotone{e-mobiledata}{695}
\showcasematerialdesigntwotone{face}{696}
\showcasematerialdesigntwotone{face-2}{697}
\showcasematerialdesigntwotone{face-3}{698}
\showcasematerialdesigntwotone{face-4}{699}
\showcasematerialdesigntwotone{face-5}{700}
\showcasematerialdesigntwotone{face-6}{701}
\showcasematerialdesigntwotone{face-retouching-natural}{702}
\showcasematerialdesigntwotone{face-retouching-off}{703}
\showcasematerialdesigntwotone{factory}{704}
\showcasematerialdesigntwotone{fact-check}{705}
\showcasematerialdesigntwotone{family-restroom}{706}
\showcasematerialdesigntwotone{fastfood}{707}
\showcasematerialdesigntwotone{fast-forward}{708}
\showcasematerialdesigntwotone{fast-rewind}{709}
\showcasematerialdesigntwotone{favorite}{710}
\showcasematerialdesigntwotone{favorite-border}{711}
\showcasematerialdesigntwotone{fax}{712}
\showcasematerialdesigntwotone{featured-play-list}{713}
\showcasematerialdesigntwotone{featured-video}{714}
\showcasematerialdesigntwotone{feed}{715}
\showcasematerialdesigntwotone{feedback}{716}
\showcasematerialdesigntwotone{female}{717}
\showcasematerialdesigntwotone{fence}{718}
\showcasematerialdesigntwotone{festival}{719}
\showcasematerialdesigntwotone{fiber-dvr}{720}
\showcasematerialdesigntwotone{fiber-manual-record}{721}
\showcasematerialdesigntwotone{fiber-new}{722}
\showcasematerialdesigntwotone{fiber-pin}{723}
\showcasematerialdesigntwotone{fiber-smart-record}{724}
\showcasematerialdesigntwotone{file-copy}{725}
\showcasematerialdesigntwotone{file-download}{726}
\showcasematerialdesigntwotone{file-download-done}{727}
\showcasematerialdesigntwotone{file-download-off}{728}
\showcasematerialdesigntwotone{file-open}{729}
\showcasematerialdesigntwotone{file-present}{730}
\showcasematerialdesigntwotone{file-upload}{731}
\showcasematerialdesigntwotone{filter}{732}
\showcasematerialdesigntwotone{filter-1}{733}
\showcasematerialdesigntwotone{filter-2}{734}
\showcasematerialdesigntwotone{filter-3}{735}
\showcasematerialdesigntwotone{filter-4}{736}
\showcasematerialdesigntwotone{filter-5}{737}
\showcasematerialdesigntwotone{filter-6}{738}
\showcasematerialdesigntwotone{filter-7}{739}
\showcasematerialdesigntwotone{filter-8}{740}
\showcasematerialdesigntwotone{filter-9}{741}
\showcasematerialdesigntwotone{filter-9-plus}{742}
\showcasematerialdesigntwotone{filter-alt}{743}
\showcasematerialdesigntwotone{filter-alt-off}{744}
\showcasematerialdesigntwotone{filter-b-and-w}{745}
\showcasematerialdesigntwotone{filter-center-focus}{746}
\showcasematerialdesigntwotone{filter-drama}{747}
\showcasematerialdesigntwotone{filter-frames}{748}
\showcasematerialdesigntwotone{filter-hdr}{749}
\showcasematerialdesigntwotone{filter-list}{750}
\showcasematerialdesigntwotone{filter-list-off}{751}
\showcasematerialdesigntwotone{filter-none}{752}
\showcasematerialdesigntwotone{filter-tilt-shift}{753}
\showcasematerialdesigntwotone{filter-vintage}{754}
\showcasematerialdesigntwotone{find-in-page}{755}
\showcasematerialdesigntwotone{find-replace}{756}
\showcasematerialdesigntwotone{fingerprint}{757}
\showcasematerialdesigntwotone{fireplace}{758}
\showcasematerialdesigntwotone{fire-extinguisher}{759}
\showcasematerialdesigntwotone{fire-hydrant-alt}{760}
\showcasematerialdesigntwotone{fire-truck}{761}
\showcasematerialdesigntwotone{first-page}{762}
\showcasematerialdesigntwotone{fitbit}{763}
\showcasematerialdesigntwotone{fitness-center}{764}
\showcasematerialdesigntwotone{fit-screen}{765}
\showcasematerialdesigntwotone{flag}{766}
\showcasematerialdesigntwotone{flag-circle}{767}
\showcasematerialdesigntwotone{flaky}{768}
\showcasematerialdesigntwotone{flare}{769}
\showcasematerialdesigntwotone{flashlight-off}{770}
\showcasematerialdesigntwotone{flashlight-on}{771}
\showcasematerialdesigntwotone{flash-auto}{772}
\showcasematerialdesigntwotone{flash-off}{773}
\showcasematerialdesigntwotone{flash-on}{774}
\showcasematerialdesigntwotone{flatware}{775}
\showcasematerialdesigntwotone{flight}{776}
\showcasematerialdesigntwotone{flight-class}{777}
\showcasematerialdesigntwotone{flight-land}{778}
\showcasematerialdesigntwotone{flight-takeoff}{779}
\showcasematerialdesigntwotone{flip}{780}
\showcasematerialdesigntwotone{flip-camera-android}{781}
\showcasematerialdesigntwotone{flip-camera-ios}{782}
\showcasematerialdesigntwotone{flip-to-back}{783}
\showcasematerialdesigntwotone{flip-to-front}{784}
\showcasematerialdesigntwotone{flood}{785}
\showcasematerialdesigntwotone{fluorescent}{786}
\showcasematerialdesigntwotone{flutter-dash}{787}
\showcasematerialdesigntwotone{fmd-bad}{788}
\showcasematerialdesigntwotone{fmd-good}{789}
\showcasematerialdesigntwotone{folder}{790}
\showcasematerialdesigntwotone{folder-copy}{791}
\showcasematerialdesigntwotone{folder-delete}{792}
\showcasematerialdesigntwotone{folder-off}{793}
\showcasematerialdesigntwotone{folder-open}{794}
\showcasematerialdesigntwotone{folder-shared}{795}
\showcasematerialdesigntwotone{folder-special}{796}
\showcasematerialdesigntwotone{folder-zip}{797}
\showcasematerialdesigntwotone{follow-the-signs}{798}
\showcasematerialdesigntwotone{font-download}{799}
\showcasematerialdesigntwotone{font-download-off}{800}
\showcasematerialdesigntwotone{food-bank}{801}
\showcasematerialdesigntwotone{forest}{802}
\showcasematerialdesigntwotone{fork-left}{803}
\showcasematerialdesigntwotone{fork-right}{804}
\showcasematerialdesigntwotone{format-align-center}{805}
\showcasematerialdesigntwotone{format-align-justify}{806}
\showcasematerialdesigntwotone{format-align-left}{807}
\showcasematerialdesigntwotone{format-align-right}{808}
\showcasematerialdesigntwotone{format-bold}{809}
\showcasematerialdesigntwotone{format-clear}{810}
\showcasematerialdesigntwotone{format-color-fill}{811}
\showcasematerialdesigntwotone{format-color-reset}{812}
\showcasematerialdesigntwotone{format-color-text}{813}
\showcasematerialdesigntwotone{format-indent-decrease}{814}
\showcasematerialdesigntwotone{format-indent-increase}{815}
\showcasematerialdesigntwotone{format-italic}{816}
\showcasematerialdesigntwotone{format-line-spacing}{817}
\showcasematerialdesigntwotone{format-list-bulleted}{818}
\showcasematerialdesigntwotone{format-list-numbered}{819}
\showcasematerialdesigntwotone{format-list-numbered-rtl}{820}
\showcasematerialdesigntwotone{format-overline}{821}
\showcasematerialdesigntwotone{format-paint}{822}
\showcasematerialdesigntwotone{format-quote}{823}
\showcasematerialdesigntwotone{format-shapes}{824}
\showcasematerialdesigntwotone{format-size}{825}
\showcasematerialdesigntwotone{format-strikethrough}{826}
\showcasematerialdesigntwotone{format-textdirection-l-to-r}{827}
\showcasematerialdesigntwotone{format-textdirection-r-to-l}{828}
\showcasematerialdesigntwotone{format-underlined}{829}
\showcasematerialdesigntwotone{fort}{830}
\showcasematerialdesigntwotone{forum}{831}
\showcasematerialdesigntwotone{forward}{832}
\showcasematerialdesigntwotone{forward-5}{833}
\showcasematerialdesigntwotone{forward-10}{834}
\showcasematerialdesigntwotone{forward-30}{835}
\showcasematerialdesigntwotone{forward-to-inbox}{836}
\showcasematerialdesigntwotone{foundation}{837}
\showcasematerialdesigntwotone{free-breakfast}{838}
\showcasematerialdesigntwotone{free-cancellation}{839}
\showcasematerialdesigntwotone{front-hand}{840}
\showcasematerialdesigntwotone{fullscreen}{841}
\showcasematerialdesigntwotone{fullscreen-exit}{842}
\showcasematerialdesigntwotone{functions}{843}
\showcasematerialdesigntwotone{gamepad}{844}
\showcasematerialdesigntwotone{games}{845}
\showcasematerialdesigntwotone{garage}{846}
\showcasematerialdesigntwotone{gas-meter}{847}
\showcasematerialdesigntwotone{gavel}{848}
\showcasematerialdesigntwotone{generating-tokens}{849}
\showcasematerialdesigntwotone{gesture}{850}
\showcasematerialdesigntwotone{get-app}{851}
\showcasematerialdesigntwotone{gif}{852}
\showcasematerialdesigntwotone{gif-box}{853}
\showcasematerialdesigntwotone{girl}{854}
\showcasematerialdesigntwotone{gite}{855}
\showcasematerialdesigntwotone{golf-course}{856}
\showcasematerialdesigntwotone{gpp-bad}{857}
\showcasematerialdesigntwotone{gpp-good}{858}
\showcasematerialdesigntwotone{gpp-maybe}{859}
\showcasematerialdesigntwotone{gps-fixed}{860}
\showcasematerialdesigntwotone{gps-not-fixed}{861}
\showcasematerialdesigntwotone{gps-off}{862}
\showcasematerialdesigntwotone{grade}{863}
\showcasematerialdesigntwotone{gradient}{864}
\showcasematerialdesigntwotone{grading}{865}
\showcasematerialdesigntwotone{grain}{866}
\showcasematerialdesigntwotone{graphic-eq}{867}
\showcasematerialdesigntwotone{grass}{868}
\showcasematerialdesigntwotone{grid-3x3}{869}
\showcasematerialdesigntwotone{grid-4x4}{870}
\showcasematerialdesigntwotone{grid-goldenratio}{871}
\showcasematerialdesigntwotone{grid-off}{872}
\showcasematerialdesigntwotone{grid-on}{873}
\showcasematerialdesigntwotone{grid-view}{874}
\showcasematerialdesigntwotone{group}{875}
\showcasematerialdesigntwotone{groups}{876}
\showcasematerialdesigntwotone{groups-2}{877}
\showcasematerialdesigntwotone{groups-3}{878}
\showcasematerialdesigntwotone{group-add}{879}
\showcasematerialdesigntwotone{group-off}{880}
\showcasematerialdesigntwotone{group-remove}{881}
\showcasematerialdesigntwotone{group-work}{882}
\showcasematerialdesigntwotone{g-mobiledata}{883}
\showcasematerialdesigntwotone{g-translate}{884}
\showcasematerialdesigntwotone{hail}{885}
\showcasematerialdesigntwotone{handshake}{886}
\showcasematerialdesigntwotone{handyman}{887}
\showcasematerialdesigntwotone{hardware}{888}
\showcasematerialdesigntwotone{hd}{889}
\showcasematerialdesigntwotone{hdr-auto}{890}
\showcasematerialdesigntwotone{hdr-auto-select}{891}
\showcasematerialdesigntwotone{hdr-enhanced-select}{892}
\showcasematerialdesigntwotone{hdr-off}{893}
\showcasematerialdesigntwotone{hdr-off-select}{894}
\showcasematerialdesigntwotone{hdr-on}{895}
\showcasematerialdesigntwotone{hdr-on-select}{896}
\showcasematerialdesigntwotone{hdr-plus}{897}
\showcasematerialdesigntwotone{hdr-strong}{898}
\showcasematerialdesigntwotone{hdr-weak}{899}
\showcasematerialdesigntwotone{headphones}{900}
\showcasematerialdesigntwotone{headphones-battery}{901}
\showcasematerialdesigntwotone{headset}{902}
\showcasematerialdesigntwotone{headset-mic}{903}
\showcasematerialdesigntwotone{headset-off}{904}
\showcasematerialdesigntwotone{healing}{905}
\showcasematerialdesigntwotone{health-and-safety}{906}
\showcasematerialdesigntwotone{hearing}{907}
\showcasematerialdesigntwotone{hearing-disabled}{908}
\showcasematerialdesigntwotone{heart-broken}{909}
\showcasematerialdesigntwotone{heat-pump}{910}
\showcasematerialdesigntwotone{height}{911}
\showcasematerialdesigntwotone{help}{912}
\showcasematerialdesigntwotone{help-center}{913}
\showcasematerialdesigntwotone{help-outline}{914}
\showcasematerialdesigntwotone{hevc}{915}
\showcasematerialdesigntwotone{hexagon}{916}
\showcasematerialdesigntwotone{hide-image}{917}
\showcasematerialdesigntwotone{hide-source}{918}
\showcasematerialdesigntwotone{highlight}{919}
\showcasematerialdesigntwotone{highlight-alt}{920}
\showcasematerialdesigntwotone{highlight-off}{921}
\showcasematerialdesigntwotone{high-quality}{922}
\showcasematerialdesigntwotone{hiking}{923}
\showcasematerialdesigntwotone{history}{924}
\showcasematerialdesigntwotone{history-edu}{925}
\showcasematerialdesigntwotone{history-toggle-off}{926}
\showcasematerialdesigntwotone{hive}{927}
\showcasematerialdesigntwotone{hls}{928}
\showcasematerialdesigntwotone{hls-off}{929}
\showcasematerialdesigntwotone{holiday-village}{930}
\showcasematerialdesigntwotone{home}{931}
\showcasematerialdesigntwotone{home-max}{932}
\showcasematerialdesigntwotone{home-mini}{933}
\showcasematerialdesigntwotone{home-repair-service}{934}
\showcasematerialdesigntwotone{home-work}{935}
\showcasematerialdesigntwotone{horizontal-distribute}{936}
\showcasematerialdesigntwotone{horizontal-rule}{937}
\showcasematerialdesigntwotone{horizontal-split}{938}
\showcasematerialdesigntwotone{hotel}{939}
\showcasematerialdesigntwotone{hotel-class}{940}
\showcasematerialdesigntwotone{hot-tub}{941}
\showcasematerialdesigntwotone{hourglass-bottom}{942}
\showcasematerialdesigntwotone{hourglass-disabled}{943}
\showcasematerialdesigntwotone{hourglass-empty}{944}
\showcasematerialdesigntwotone{hourglass-full}{945}
\showcasematerialdesigntwotone{hourglass-top}{946}
\showcasematerialdesigntwotone{house}{947}
\showcasematerialdesigntwotone{houseboat}{948}
\showcasematerialdesigntwotone{house-siding}{949}
\showcasematerialdesigntwotone{how-to-reg}{950}
\showcasematerialdesigntwotone{how-to-vote}{951}
\showcasematerialdesigntwotone{html}{952}
\showcasematerialdesigntwotone{http}{953}
\showcasematerialdesigntwotone{https}{954}
\showcasematerialdesigntwotone{hub}{955}
\showcasematerialdesigntwotone{hvac}{956}
\showcasematerialdesigntwotone{h-mobiledata}{957}
\showcasematerialdesigntwotone{h-plus-mobiledata}{958}
\showcasematerialdesigntwotone{icecream}{959}
\showcasematerialdesigntwotone{ice-skating}{960}
\showcasematerialdesigntwotone{image}{961}
\showcasematerialdesigntwotone{imagesearch-roller}{962}
\showcasematerialdesigntwotone{image-aspect-ratio}{963}
\showcasematerialdesigntwotone{image-not-supported}{964}
\showcasematerialdesigntwotone{image-search}{965}
\showcasematerialdesigntwotone{important-devices}{966}
\showcasematerialdesigntwotone{import-contacts}{967}
\showcasematerialdesigntwotone{import-export}{968}
\showcasematerialdesigntwotone{inbox}{969}
\showcasematerialdesigntwotone{incomplete-circle}{970}
\showcasematerialdesigntwotone{indeterminate-check-box}{971}
\showcasematerialdesigntwotone{info}{972}
\showcasematerialdesigntwotone{input}{973}
\showcasematerialdesigntwotone{insert-chart}{974}
\showcasematerialdesigntwotone{insert-chart-outlined}{975}
\showcasematerialdesigntwotone{insert-comment}{976}
\showcasematerialdesigntwotone{insert-drive-file}{977}
\showcasematerialdesigntwotone{insert-emoticon}{978}
\showcasematerialdesigntwotone{insert-invitation}{979}
\showcasematerialdesigntwotone{insert-link}{980}
\showcasematerialdesigntwotone{insert-page-break}{981}
\showcasematerialdesigntwotone{insert-photo}{982}
\showcasematerialdesigntwotone{insights}{983}
\showcasematerialdesigntwotone{install-desktop}{984}
\showcasematerialdesigntwotone{install-mobile}{985}
\showcasematerialdesigntwotone{integration-instructions}{986}
\showcasematerialdesigntwotone{interests}{987}
\showcasematerialdesigntwotone{interpreter-mode}{988}
\showcasematerialdesigntwotone{inventory}{989}
\showcasematerialdesigntwotone{inventory-2}{990}
\showcasematerialdesigntwotone{invert-colors}{991}
\showcasematerialdesigntwotone{invert-colors-off}{992}
\showcasematerialdesigntwotone{ios-share}{993}
\showcasematerialdesigntwotone{iron}{994}
\showcasematerialdesigntwotone{iso}{995}
\showcasematerialdesigntwotone{javascript}{996}
\showcasematerialdesigntwotone{join-full}{997}
\showcasematerialdesigntwotone{join-inner}{998}
\showcasematerialdesigntwotone{join-left}{999}
\showcasematerialdesigntwotone{join-right}{1000}
\showcasematerialdesigntwotone{kayaking}{1001}
\showcasematerialdesigntwotone{kebab-dining}{1002}
\showcasematerialdesigntwotone{key}{1003}
\showcasematerialdesigntwotone{keyboard}{1004}
\showcasematerialdesigntwotone{keyboard-alt}{1005}
\showcasematerialdesigntwotone{keyboard-arrow-down}{1006}
\showcasematerialdesigntwotone{keyboard-arrow-left}{1007}
\showcasematerialdesigntwotone{keyboard-arrow-right}{1008}
\showcasematerialdesigntwotone{keyboard-arrow-up}{1009}
\showcasematerialdesigntwotone{keyboard-backspace}{1010}
\showcasematerialdesigntwotone{keyboard-capslock}{1011}
\showcasematerialdesigntwotone{keyboard-command-key}{1012}
\showcasematerialdesigntwotone{keyboard-control-key}{1013}
\showcasematerialdesigntwotone{keyboard-double-arrow-down}{1014}
\showcasematerialdesigntwotone{keyboard-double-arrow-left}{1015}
\showcasematerialdesigntwotone{keyboard-double-arrow-right}{1016}
\showcasematerialdesigntwotone{keyboard-double-arrow-up}{1017}
\showcasematerialdesigntwotone{keyboard-hide}{1018}
\showcasematerialdesigntwotone{keyboard-option-key}{1019}
\showcasematerialdesigntwotone{keyboard-return}{1020}
\showcasematerialdesigntwotone{keyboard-tab}{1021}
\showcasematerialdesigntwotone{keyboard-voice}{1022}
\showcasematerialdesigntwotone{key-off}{1023}
\showcasematerialdesigntwotone{king-bed}{1024}
\showcasematerialdesigntwotone{kitchen}{1025}
\showcasematerialdesigntwotone{kitesurfing}{1026}
\showcasematerialdesigntwotone{label}{1027}
\showcasematerialdesigntwotone{label-important}{1028}
\showcasematerialdesigntwotone{label-off}{1029}
\showcasematerialdesigntwotone{lan}{1030}
\showcasematerialdesigntwotone{landscape}{1031}
\showcasematerialdesigntwotone{landslide}{1032}
\showcasematerialdesigntwotone{language}{1033}
\showcasematerialdesigntwotone{laptop}{1034}
\showcasematerialdesigntwotone{laptop-chromebook}{1035}
\showcasematerialdesigntwotone{laptop-mac}{1036}
\showcasematerialdesigntwotone{laptop-windows}{1037}
\showcasematerialdesigntwotone{last-page}{1038}
\showcasematerialdesigntwotone{launch}{1039}
\showcasematerialdesigntwotone{layers}{1040}
\showcasematerialdesigntwotone{layers-clear}{1041}
\showcasematerialdesigntwotone{leaderboard}{1042}
\showcasematerialdesigntwotone{leak-add}{1043}
\showcasematerialdesigntwotone{leak-remove}{1044}
\showcasematerialdesigntwotone{legend-toggle}{1045}
\showcasematerialdesigntwotone{lens}{1046}
\showcasematerialdesigntwotone{lens-blur}{1047}
\showcasematerialdesigntwotone{library-add}{1048}
\showcasematerialdesigntwotone{library-add-check}{1049}
\showcasematerialdesigntwotone{library-books}{1050}
\showcasematerialdesigntwotone{library-music}{1051}
\showcasematerialdesigntwotone{light}{1052}
\showcasematerialdesigntwotone{lightbulb}{1053}
\showcasematerialdesigntwotone{lightbulb-circle}{1054}
\showcasematerialdesigntwotone{light-mode}{1055}
\showcasematerialdesigntwotone{linear-scale}{1056}
\showcasematerialdesigntwotone{line-axis}{1057}
\showcasematerialdesigntwotone{line-style}{1058}
\showcasematerialdesigntwotone{line-weight}{1059}
\showcasematerialdesigntwotone{link}{1060}
\showcasematerialdesigntwotone{linked-camera}{1061}
\showcasematerialdesigntwotone{link-off}{1062}
\showcasematerialdesigntwotone{liquor}{1063}
\showcasematerialdesigntwotone{list}{1064}
\showcasematerialdesigntwotone{list-alt}{1065}
\showcasematerialdesigntwotone{live-help}{1066}
\showcasematerialdesigntwotone{live-tv}{1067}
\showcasematerialdesigntwotone{living}{1068}
\showcasematerialdesigntwotone{local-activity}{1069}
\showcasematerialdesigntwotone{local-airport}{1070}
\showcasematerialdesigntwotone{local-atm}{1071}
\showcasematerialdesigntwotone{local-bar}{1072}
\showcasematerialdesigntwotone{local-cafe}{1073}
\showcasematerialdesigntwotone{local-car-wash}{1074}
\showcasematerialdesigntwotone{local-convenience-store}{1075}
\showcasematerialdesigntwotone{local-dining}{1076}
\showcasematerialdesigntwotone{local-drink}{1077}
\showcasematerialdesigntwotone{local-fire-department}{1078}
\showcasematerialdesigntwotone{local-florist}{1079}
\showcasematerialdesigntwotone{local-gas-station}{1080}
\showcasematerialdesigntwotone{local-grocery-store}{1081}
\showcasematerialdesigntwotone{local-hospital}{1082}
\showcasematerialdesigntwotone{local-hotel}{1083}
\showcasematerialdesigntwotone{local-laundry-service}{1084}
\showcasematerialdesigntwotone{local-library}{1085}
\showcasematerialdesigntwotone{local-mall}{1086}
\showcasematerialdesigntwotone{local-movies}{1087}
\showcasematerialdesigntwotone{local-offer}{1088}
\showcasematerialdesigntwotone{local-parking}{1089}
\showcasematerialdesigntwotone{local-pharmacy}{1090}
\showcasematerialdesigntwotone{local-phone}{1091}
\showcasematerialdesigntwotone{local-pizza}{1092}
\showcasematerialdesigntwotone{local-play}{1093}
\showcasematerialdesigntwotone{local-police}{1094}
\showcasematerialdesigntwotone{local-post-office}{1095}
\showcasematerialdesigntwotone{local-printshop}{1096}
\showcasematerialdesigntwotone{local-see}{1097}
\showcasematerialdesigntwotone{local-shipping}{1098}
\showcasematerialdesigntwotone{local-taxi}{1099}
\showcasematerialdesigntwotone{location-city}{1100}
\showcasematerialdesigntwotone{location-disabled}{1101}
\showcasematerialdesigntwotone{location-off}{1102}
\showcasematerialdesigntwotone{location-on}{1103}
\showcasematerialdesigntwotone{location-searching}{1104}
\showcasematerialdesigntwotone{lock}{1105}
\showcasematerialdesigntwotone{lock-clock}{1106}
\showcasematerialdesigntwotone{lock-open}{1107}
\showcasematerialdesigntwotone{lock-person}{1108}
\showcasematerialdesigntwotone{lock-reset}{1109}
\showcasematerialdesigntwotone{login}{1110}
\showcasematerialdesigntwotone{logout}{1111}
\showcasematerialdesigntwotone{logo-dev}{1112}
\showcasematerialdesigntwotone{looks}{1113}
\showcasematerialdesigntwotone{looks-3}{1114}
\showcasematerialdesigntwotone{looks-4}{1115}
\showcasematerialdesigntwotone{looks-5}{1116}
\showcasematerialdesigntwotone{looks-6}{1117}
\showcasematerialdesigntwotone{looks-one}{1118}
\showcasematerialdesigntwotone{looks-two}{1119}
\showcasematerialdesigntwotone{loop}{1120}
\showcasematerialdesigntwotone{loupe}{1121}
\showcasematerialdesigntwotone{low-priority}{1122}
\showcasematerialdesigntwotone{loyalty}{1123}
\showcasematerialdesigntwotone{lte-mobiledata}{1124}
\showcasematerialdesigntwotone{lte-plus-mobiledata}{1125}
\showcasematerialdesigntwotone{luggage}{1126}
\showcasematerialdesigntwotone{lunch-dining}{1127}
\showcasematerialdesigntwotone{lyrics}{1128}
\showcasematerialdesigntwotone{macro-off}{1129}
\showcasematerialdesigntwotone{mail}{1130}
\showcasematerialdesigntwotone{mail-lock}{1131}
\showcasematerialdesigntwotone{mail-outline}{1132}
\showcasematerialdesigntwotone{male}{1133}
\showcasematerialdesigntwotone{man}{1134}
\showcasematerialdesigntwotone{manage-accounts}{1135}
\showcasematerialdesigntwotone{manage-history}{1136}
\showcasematerialdesigntwotone{manage-search}{1137}
\showcasematerialdesigntwotone{man-2}{1138}
\showcasematerialdesigntwotone{man-3}{1139}
\showcasematerialdesigntwotone{man-4}{1140}
\showcasematerialdesigntwotone{map}{1141}
\showcasematerialdesigntwotone{maps-home-work}{1142}
\showcasematerialdesigntwotone{maps-ugc}{1143}
\showcasematerialdesigntwotone{margin}{1144}
\showcasematerialdesigntwotone{markunread}{1145}
\showcasematerialdesigntwotone{markunread-mailbox}{1146}
\showcasematerialdesigntwotone{mark-as-unread}{1147}
\showcasematerialdesigntwotone{mark-chat-read}{1148}
\showcasematerialdesigntwotone{mark-chat-unread}{1149}
\showcasematerialdesigntwotone{mark-email-read}{1150}
\showcasematerialdesigntwotone{mark-email-unread}{1151}
\showcasematerialdesigntwotone{mark-unread-chat-alt}{1152}
\showcasematerialdesigntwotone{masks}{1153}
\showcasematerialdesigntwotone{maximize}{1154}
\showcasematerialdesigntwotone{mediation}{1155}
\showcasematerialdesigntwotone{media-bluetooth-off}{1156}
\showcasematerialdesigntwotone{media-bluetooth-on}{1157}
\showcasematerialdesigntwotone{medical-information}{1158}
\showcasematerialdesigntwotone{medical-services}{1159}
\showcasematerialdesigntwotone{medication}{1160}
\showcasematerialdesigntwotone{medication-liquid}{1161}
\showcasematerialdesigntwotone{meeting-room}{1162}
\showcasematerialdesigntwotone{memory}{1163}
\showcasematerialdesigntwotone{menu}{1164}
\showcasematerialdesigntwotone{menu-book}{1165}
\showcasematerialdesigntwotone{menu-open}{1166}
\showcasematerialdesigntwotone{merge}{1167}
\showcasematerialdesigntwotone{merge-type}{1168}
\showcasematerialdesigntwotone{message}{1169}
\showcasematerialdesigntwotone{mic}{1170}
\showcasematerialdesigntwotone{microwave}{1171}
\showcasematerialdesigntwotone{mic-external-off}{1172}
\showcasematerialdesigntwotone{mic-external-on}{1173}
\showcasematerialdesigntwotone{mic-none}{1174}
\showcasematerialdesigntwotone{mic-off}{1175}
\showcasematerialdesigntwotone{military-tech}{1176}
\showcasematerialdesigntwotone{minimize}{1177}
\showcasematerialdesigntwotone{minor-crash}{1178}
\showcasematerialdesigntwotone{miscellaneous-services}{1179}
\showcasematerialdesigntwotone{missed-video-call}{1180}
\showcasematerialdesigntwotone{mms}{1181}
\showcasematerialdesigntwotone{mobiledata-off}{1182}
\showcasematerialdesigntwotone{mobile-friendly}{1183}
\showcasematerialdesigntwotone{mobile-off}{1184}
\showcasematerialdesigntwotone{mobile-screen-share}{1185}
\showcasematerialdesigntwotone{mode}{1186}
\showcasematerialdesigntwotone{model-training}{1187}
\showcasematerialdesigntwotone{mode-comment}{1188}
\showcasematerialdesigntwotone{mode-edit}{1189}
\showcasematerialdesigntwotone{mode-edit-outline}{1190}
\showcasematerialdesigntwotone{mode-fan-off}{1191}
\showcasematerialdesigntwotone{mode-night}{1192}
\showcasematerialdesigntwotone{mode-of-travel}{1193}
\showcasematerialdesigntwotone{mode-standby}{1194}
\showcasematerialdesigntwotone{monetization-on}{1195}
\showcasematerialdesigntwotone{money}{1196}
\showcasematerialdesigntwotone{money-off}{1197}
\showcasematerialdesigntwotone{money-off-csred}{1198}
\showcasematerialdesigntwotone{monitor}{1199}
\showcasematerialdesigntwotone{monitor-heart}{1200}
\showcasematerialdesigntwotone{monitor-weight}{1201}
\showcasematerialdesigntwotone{monochrome-photos}{1202}
\showcasematerialdesigntwotone{mood}{1203}
\showcasematerialdesigntwotone{mood-bad}{1204}
\showcasematerialdesigntwotone{moped}{1205}
\showcasematerialdesigntwotone{more}{1206}
\showcasematerialdesigntwotone{more-horiz}{1207}
\showcasematerialdesigntwotone{more-time}{1208}
\showcasematerialdesigntwotone{more-vert}{1209}
\showcasematerialdesigntwotone{mosque}{1210}
\showcasematerialdesigntwotone{motion-photos-auto}{1211}
\showcasematerialdesigntwotone{motion-photos-off}{1212}
\showcasematerialdesigntwotone{motion-photos-on}{1213}
\showcasematerialdesigntwotone{motion-photos-pause}{1214}
\showcasematerialdesigntwotone{motion-photos-paused}{1215}
\showcasematerialdesigntwotone{mouse}{1216}
\showcasematerialdesigntwotone{move-down}{1217}
\showcasematerialdesigntwotone{move-to-inbox}{1218}
\showcasematerialdesigntwotone{move-up}{1219}
\showcasematerialdesigntwotone{movie}{1220}
\showcasematerialdesigntwotone{movie-creation}{1221}
\showcasematerialdesigntwotone{movie-filter}{1222}
\showcasematerialdesigntwotone{moving}{1223}
\showcasematerialdesigntwotone{mp}{1224}
\showcasematerialdesigntwotone{multiline-chart}{1225}
\showcasematerialdesigntwotone{multiple-stop}{1226}
\showcasematerialdesigntwotone{museum}{1227}
\showcasematerialdesigntwotone{music-note}{1228}
\showcasematerialdesigntwotone{music-off}{1229}
\showcasematerialdesigntwotone{music-video}{1230}
\showcasematerialdesigntwotone{my-location}{1231}
\showcasematerialdesigntwotone{nat}{1232}
\showcasematerialdesigntwotone{nature}{1233}
\showcasematerialdesigntwotone{nature-people}{1234}
\showcasematerialdesigntwotone{navigate-before}{1235}
\showcasematerialdesigntwotone{navigate-next}{1236}
\showcasematerialdesigntwotone{navigation}{1237}
\showcasematerialdesigntwotone{nearby-error}{1238}
\showcasematerialdesigntwotone{nearby-off}{1239}
\showcasematerialdesigntwotone{near-me}{1240}
\showcasematerialdesigntwotone{near-me-disabled}{1241}
\showcasematerialdesigntwotone{nest-cam-wired-stand}{1242}
\showcasematerialdesigntwotone{network-cell}{1243}
\showcasematerialdesigntwotone{network-check}{1244}
\showcasematerialdesigntwotone{network-locked}{1245}
\showcasematerialdesigntwotone{network-ping}{1246}
\showcasematerialdesigntwotone{network-wifi}{1247}
\showcasematerialdesigntwotone{network-wifi-1-bar}{1248}
\showcasematerialdesigntwotone{network-wifi-2-bar}{1249}
\showcasematerialdesigntwotone{network-wifi-3-bar}{1250}
\showcasematerialdesigntwotone{newspaper}{1251}
\showcasematerialdesigntwotone{new-label}{1252}
\showcasematerialdesigntwotone{new-releases}{1253}
\showcasematerialdesigntwotone{next-plan}{1254}
\showcasematerialdesigntwotone{next-week}{1255}
\showcasematerialdesigntwotone{nfc}{1256}
\showcasematerialdesigntwotone{nightlife}{1257}
\showcasematerialdesigntwotone{nightlight}{1258}
\showcasematerialdesigntwotone{nightlight-round}{1259}
\showcasematerialdesigntwotone{nights-stay}{1260}
\showcasematerialdesigntwotone{night-shelter}{1261}
\showcasematerialdesigntwotone{noise-aware}{1262}
\showcasematerialdesigntwotone{noise-control-off}{1263}
\showcasematerialdesigntwotone{nordic-walking}{1264}
\showcasematerialdesigntwotone{north}{1265}
\showcasematerialdesigntwotone{north-east}{1266}
\showcasematerialdesigntwotone{north-west}{1267}
\showcasematerialdesigntwotone{note}{1268}
\showcasematerialdesigntwotone{notes}{1269}
\showcasematerialdesigntwotone{note-add}{1270}
\showcasematerialdesigntwotone{note-alt}{1271}
\showcasematerialdesigntwotone{notifications}{1272}
\showcasematerialdesigntwotone{notifications-active}{1273}
\showcasematerialdesigntwotone{notifications-none}{1274}
\showcasematerialdesigntwotone{notifications-off}{1275}
\showcasematerialdesigntwotone{notifications-paused}{1276}
\showcasematerialdesigntwotone{notification-add}{1277}
\showcasematerialdesigntwotone{notification-important}{1278}
\showcasematerialdesigntwotone{not-accessible}{1279}
\showcasematerialdesigntwotone{not-interested}{1280}
\showcasematerialdesigntwotone{not-listed-location}{1281}
\showcasematerialdesigntwotone{not-started}{1282}
\showcasematerialdesigntwotone{no-accounts}{1283}
\showcasematerialdesigntwotone{no-adult-content}{1284}
\showcasematerialdesigntwotone{no-backpack}{1285}
\showcasematerialdesigntwotone{no-cell}{1286}
\showcasematerialdesigntwotone{no-crash}{1287}
\showcasematerialdesigntwotone{no-drinks}{1288}
\showcasematerialdesigntwotone{no-encryption}{1289}
\showcasematerialdesigntwotone{no-encryption-gmailerrorred}{1290}
\showcasematerialdesigntwotone{no-flash}{1291}
\showcasematerialdesigntwotone{no-food}{1292}
\showcasematerialdesigntwotone{no-luggage}{1293}
\showcasematerialdesigntwotone{no-meals}{1294}
\showcasematerialdesigntwotone{no-meeting-room}{1295}
\showcasematerialdesigntwotone{no-photography}{1296}
\showcasematerialdesigntwotone{no-sim}{1297}
\showcasematerialdesigntwotone{no-stroller}{1298}
\showcasematerialdesigntwotone{no-transfer}{1299}
\showcasematerialdesigntwotone{numbers}{1300}
\showcasematerialdesigntwotone{offline-bolt}{1301}
\showcasematerialdesigntwotone{offline-pin}{1302}
\showcasematerialdesigntwotone{offline-share}{1303}
\showcasematerialdesigntwotone{oil-barrel}{1304}
\showcasematerialdesigntwotone{ondemand-video}{1305}
\showcasematerialdesigntwotone{online-prediction}{1306}
\showcasematerialdesigntwotone{on-device-training}{1307}
\showcasematerialdesigntwotone{opacity}{1308}
\showcasematerialdesigntwotone{open-in-browser}{1309}
\showcasematerialdesigntwotone{open-in-full}{1310}
\showcasematerialdesigntwotone{open-in-new}{1311}
\showcasematerialdesigntwotone{open-in-new-off}{1312}
\showcasematerialdesigntwotone{open-with}{1313}
\showcasematerialdesigntwotone{other-houses}{1314}
\showcasematerialdesigntwotone{outbound}{1315}
\showcasematerialdesigntwotone{outbox}{1316}
\showcasematerialdesigntwotone{outdoor-grill}{1317}
\showcasematerialdesigntwotone{outlet}{1318}
\showcasematerialdesigntwotone{outlined-flag}{1319}
\showcasematerialdesigntwotone{output}{1320}
\showcasematerialdesigntwotone{padding}{1321}
\showcasematerialdesigntwotone{pages}{1322}
\showcasematerialdesigntwotone{pageview}{1323}
\showcasematerialdesigntwotone{paid}{1324}
\showcasematerialdesigntwotone{palette}{1325}
\showcasematerialdesigntwotone{panorama}{1326}
\showcasematerialdesigntwotone{panorama-fish-eye}{1327}
\showcasematerialdesigntwotone{panorama-horizontal}{1328}
\showcasematerialdesigntwotone{panorama-horizontal-select}{1329}
\showcasematerialdesigntwotone{panorama-photosphere}{1330}
\showcasematerialdesigntwotone{panorama-photosphere-select}{1331}
\showcasematerialdesigntwotone{panorama-vertical}{1332}
\showcasematerialdesigntwotone{panorama-vertical-select}{1333}
\showcasematerialdesigntwotone{panorama-wide-angle}{1334}
\showcasematerialdesigntwotone{panorama-wide-angle-select}{1335}
\showcasematerialdesigntwotone{pan-tool}{1336}
\showcasematerialdesigntwotone{pan-tool-alt}{1337}
\showcasematerialdesigntwotone{paragliding}{1338}
\showcasematerialdesigntwotone{park}{1339}
\showcasematerialdesigntwotone{party-mode}{1340}
\showcasematerialdesigntwotone{password}{1341}
\showcasematerialdesigntwotone{pattern}{1342}
\showcasematerialdesigntwotone{pause}{1343}
\showcasematerialdesigntwotone{pause-circle}{1344}
\showcasematerialdesigntwotone{pause-circle-filled}{1345}
\showcasematerialdesigntwotone{pause-circle-outline}{1346}
\showcasematerialdesigntwotone{pause-presentation}{1347}
\showcasematerialdesigntwotone{payment}{1348}
\showcasematerialdesigntwotone{payments}{1349}
\showcasematerialdesigntwotone{pedal-bike}{1350}
\showcasematerialdesigntwotone{pending}{1351}
\showcasematerialdesigntwotone{pending-actions}{1352}
\showcasematerialdesigntwotone{pentagon}{1353}
\showcasematerialdesigntwotone{people}{1354}
\showcasematerialdesigntwotone{people-alt}{1355}
\showcasematerialdesigntwotone{people-outline}{1356}
\showcasematerialdesigntwotone{percent}{1357}
\showcasematerialdesigntwotone{perm-camera-mic}{1358}
\showcasematerialdesigntwotone{perm-contact-calendar}{1359}
\showcasematerialdesigntwotone{perm-data-setting}{1360}
\showcasematerialdesigntwotone{perm-device-information}{1361}
\showcasematerialdesigntwotone{perm-identity}{1362}
\showcasematerialdesigntwotone{perm-media}{1363}
\showcasematerialdesigntwotone{perm-phone-msg}{1364}
\showcasematerialdesigntwotone{perm-scan-wifi}{1365}
\showcasematerialdesigntwotone{person}{1366}
\showcasematerialdesigntwotone{personal-injury}{1367}
\showcasematerialdesigntwotone{personal-video}{1368}
\showcasematerialdesigntwotone{person-2}{1369}
\showcasematerialdesigntwotone{person-3}{1370}
\showcasematerialdesigntwotone{person-4}{1371}
\showcasematerialdesigntwotone{person-add}{1372}
\showcasematerialdesigntwotone{person-add-alt}{1373}
\showcasematerialdesigntwotone{person-add-alt-1}{1374}
\showcasematerialdesigntwotone{person-add-disabled}{1375}
\showcasematerialdesigntwotone{person-off}{1376}
\showcasematerialdesigntwotone{person-outline}{1377}
\showcasematerialdesigntwotone{person-pin}{1378}
\showcasematerialdesigntwotone{person-pin-circle}{1379}
\showcasematerialdesigntwotone{person-remove}{1380}
\showcasematerialdesigntwotone{person-remove-alt-1}{1381}
\showcasematerialdesigntwotone{person-search}{1382}
\showcasematerialdesigntwotone{pest-control}{1383}
\showcasematerialdesigntwotone{pest-control-rodent}{1384}
\showcasematerialdesigntwotone{pets}{1385}
\showcasematerialdesigntwotone{phishing}{1386}
\showcasematerialdesigntwotone{phone}{1387}
\showcasematerialdesigntwotone{phonelink}{1388}
\showcasematerialdesigntwotone{phonelink-erase}{1389}
\showcasematerialdesigntwotone{phonelink-lock}{1390}
\showcasematerialdesigntwotone{phonelink-off}{1391}
\showcasematerialdesigntwotone{phonelink-ring}{1392}
\showcasematerialdesigntwotone{phonelink-setup}{1393}
\showcasematerialdesigntwotone{phone-android}{1394}
\showcasematerialdesigntwotone{phone-bluetooth-speaker}{1395}
\showcasematerialdesigntwotone{phone-callback}{1396}
\showcasematerialdesigntwotone{phone-disabled}{1397}
\showcasematerialdesigntwotone{phone-enabled}{1398}
\showcasematerialdesigntwotone{phone-forwarded}{1399}
\showcasematerialdesigntwotone{phone-iphone}{1400}
\showcasematerialdesigntwotone{phone-locked}{1401}
\showcasematerialdesigntwotone{phone-missed}{1402}
\showcasematerialdesigntwotone{phone-paused}{1403}
\showcasematerialdesigntwotone{photo}{1404}
\showcasematerialdesigntwotone{photo-album}{1405}
\showcasematerialdesigntwotone{photo-camera}{1406}
\showcasematerialdesigntwotone{photo-camera-back}{1407}
\showcasematerialdesigntwotone{photo-camera-front}{1408}
\showcasematerialdesigntwotone{photo-filter}{1409}
\showcasematerialdesigntwotone{photo-library}{1410}
\showcasematerialdesigntwotone{photo-size-select-actual}{1411}
\showcasematerialdesigntwotone{photo-size-select-large}{1412}
\showcasematerialdesigntwotone{photo-size-select-small}{1413}
\showcasematerialdesigntwotone{php}{1414}
\showcasematerialdesigntwotone{piano}{1415}
\showcasematerialdesigntwotone{piano-off}{1416}
\showcasematerialdesigntwotone{picture-as-pdf}{1417}
\showcasematerialdesigntwotone{picture-in-picture}{1418}
\showcasematerialdesigntwotone{picture-in-picture-alt}{1419}
\showcasematerialdesigntwotone{pie-chart}{1420}
\showcasematerialdesigntwotone{pie-chart-outline}{1421}
\showcasematerialdesigntwotone{pin}{1422}
\showcasematerialdesigntwotone{pinch}{1423}
\showcasematerialdesigntwotone{pin-drop}{1424}
\showcasematerialdesigntwotone{pin-end}{1425}
\showcasematerialdesigntwotone{pin-invoke}{1426}
\showcasematerialdesigntwotone{pivot-table-chart}{1427}
\showcasematerialdesigntwotone{pix}{1428}
\showcasematerialdesigntwotone{place}{1429}
\showcasematerialdesigntwotone{plagiarism}{1430}
\showcasematerialdesigntwotone{playlist-add}{1431}
\showcasematerialdesigntwotone{playlist-add-check}{1432}
\showcasematerialdesigntwotone{playlist-add-check-circle}{1433}
\showcasematerialdesigntwotone{playlist-add-circle}{1434}
\showcasematerialdesigntwotone{playlist-play}{1435}
\showcasematerialdesigntwotone{playlist-remove}{1436}
\showcasematerialdesigntwotone{play-arrow}{1437}
\showcasematerialdesigntwotone{play-circle}{1438}
\showcasematerialdesigntwotone{play-circle-filled}{1439}
\showcasematerialdesigntwotone{play-circle-outline}{1440}
\showcasematerialdesigntwotone{play-disabled}{1441}
\showcasematerialdesigntwotone{play-for-work}{1442}
\showcasematerialdesigntwotone{play-lesson}{1443}
\showcasematerialdesigntwotone{plumbing}{1444}
\showcasematerialdesigntwotone{plus-one}{1445}
\showcasematerialdesigntwotone{podcasts}{1446}
\showcasematerialdesigntwotone{point-of-sale}{1447}
\showcasematerialdesigntwotone{policy}{1448}
\showcasematerialdesigntwotone{poll}{1449}
\showcasematerialdesigntwotone{polyline}{1450}
\showcasematerialdesigntwotone{polymer}{1451}
\showcasematerialdesigntwotone{pool}{1452}
\showcasematerialdesigntwotone{portable-wifi-off}{1453}
\showcasematerialdesigntwotone{portrait}{1454}
\showcasematerialdesigntwotone{post-add}{1455}
\showcasematerialdesigntwotone{power}{1456}
\showcasematerialdesigntwotone{power-input}{1457}
\showcasematerialdesigntwotone{power-off}{1458}
\showcasematerialdesigntwotone{power-settings-new}{1459}
\showcasematerialdesigntwotone{precision-manufacturing}{1460}
\showcasematerialdesigntwotone{pregnant-woman}{1461}
\showcasematerialdesigntwotone{present-to-all}{1462}
\showcasematerialdesigntwotone{preview}{1463}
\showcasematerialdesigntwotone{price-change}{1464}
\showcasematerialdesigntwotone{price-check}{1465}
\showcasematerialdesigntwotone{print}{1466}
\showcasematerialdesigntwotone{print-disabled}{1467}
\showcasematerialdesigntwotone{priority-high}{1468}
\showcasematerialdesigntwotone{privacy-tip}{1469}
\showcasematerialdesigntwotone{private-connectivity}{1470}
\showcasematerialdesigntwotone{production-quantity-limits}{1471}
\showcasematerialdesigntwotone{propane}{1472}
\showcasematerialdesigntwotone{propane-tank}{1473}
\showcasematerialdesigntwotone{psychology}{1474}
\showcasematerialdesigntwotone{psychology-alt}{1475}
\showcasematerialdesigntwotone{public}{1476}
\showcasematerialdesigntwotone{public-off}{1477}
\showcasematerialdesigntwotone{publish}{1478}
\showcasematerialdesigntwotone{published-with-changes}{1479}
\showcasematerialdesigntwotone{punch-clock}{1480}
\showcasematerialdesigntwotone{push-pin}{1481}
\showcasematerialdesigntwotone{qr-code}{1482}
\showcasematerialdesigntwotone{qr-code-2}{1483}
\showcasematerialdesigntwotone{qr-code-scanner}{1484}
\showcasematerialdesigntwotone{query-builder}{1485}
\showcasematerialdesigntwotone{query-stats}{1486}
\showcasematerialdesigntwotone{question-answer}{1487}
\showcasematerialdesigntwotone{question-mark}{1488}
\showcasematerialdesigntwotone{queue}{1489}
\showcasematerialdesigntwotone{queue-music}{1490}
\showcasematerialdesigntwotone{queue-play-next}{1491}
\showcasematerialdesigntwotone{quickreply}{1492}
\showcasematerialdesigntwotone{quiz}{1493}
\showcasematerialdesigntwotone{radar}{1494}
\showcasematerialdesigntwotone{radio}{1495}
\showcasematerialdesigntwotone{radio-button-checked}{1496}
\showcasematerialdesigntwotone{radio-button-unchecked}{1497}
\showcasematerialdesigntwotone{railway-alert}{1498}
\showcasematerialdesigntwotone{ramen-dining}{1499}
\showcasematerialdesigntwotone{ramp-left}{1500}
\showcasematerialdesigntwotone{ramp-right}{1501}
\showcasematerialdesigntwotone{rate-review}{1502}
\showcasematerialdesigntwotone{raw-off}{1503}
\showcasematerialdesigntwotone{raw-on}{1504}
\showcasematerialdesigntwotone{read-more}{1505}
\showcasematerialdesigntwotone{real-estate-agent}{1506}
\showcasematerialdesigntwotone{receipt}{1507}
\showcasematerialdesigntwotone{receipt-long}{1508}
\showcasematerialdesigntwotone{recent-actors}{1509}
\showcasematerialdesigntwotone{recommend}{1510}
\showcasematerialdesigntwotone{record-voice-over}{1511}
\showcasematerialdesigntwotone{rectangle}{1512}
\showcasematerialdesigntwotone{recycling}{1513}
\showcasematerialdesigntwotone{redeem}{1514}
\showcasematerialdesigntwotone{redo}{1515}
\showcasematerialdesigntwotone{reduce-capacity}{1516}
\showcasematerialdesigntwotone{refresh}{1517}
\showcasematerialdesigntwotone{remember-me}{1518}
\showcasematerialdesigntwotone{remove}{1519}
\showcasematerialdesigntwotone{remove-circle}{1520}
\showcasematerialdesigntwotone{remove-circle-outline}{1521}
\showcasematerialdesigntwotone{remove-done}{1522}
\showcasematerialdesigntwotone{remove-from-queue}{1523}
\showcasematerialdesigntwotone{remove-moderator}{1524}
\showcasematerialdesigntwotone{remove-red-eye}{1525}
\showcasematerialdesigntwotone{remove-road}{1526}
\showcasematerialdesigntwotone{remove-shopping-cart}{1527}
\showcasematerialdesigntwotone{reorder}{1528}
\showcasematerialdesigntwotone{repartition}{1529}
\showcasematerialdesigntwotone{repeat}{1530}
\showcasematerialdesigntwotone{repeat-on}{1531}
\showcasematerialdesigntwotone{repeat-one}{1532}
\showcasematerialdesigntwotone{repeat-one-on}{1533}
\showcasematerialdesigntwotone{replay}{1534}
\showcasematerialdesigntwotone{replay-5}{1535}
\showcasematerialdesigntwotone{replay-10}{1536}
\showcasematerialdesigntwotone{replay-30}{1537}
\showcasematerialdesigntwotone{replay-circle-filled}{1538}
\showcasematerialdesigntwotone{reply}{1539}
\showcasematerialdesigntwotone{reply-all}{1540}
\showcasematerialdesigntwotone{report}{1541}
\showcasematerialdesigntwotone{report-gmailerrorred}{1542}
\showcasematerialdesigntwotone{report-off}{1543}
\showcasematerialdesigntwotone{report-problem}{1544}
\showcasematerialdesigntwotone{request-page}{1545}
\showcasematerialdesigntwotone{request-quote}{1546}
\showcasematerialdesigntwotone{reset-tv}{1547}
\showcasematerialdesigntwotone{restart-alt}{1548}
\showcasematerialdesigntwotone{restaurant}{1549}
\showcasematerialdesigntwotone{restaurant-menu}{1550}
\showcasematerialdesigntwotone{restore}{1551}
\showcasematerialdesigntwotone{restore-from-trash}{1552}
\showcasematerialdesigntwotone{restore-page}{1553}
\showcasematerialdesigntwotone{reviews}{1554}
\showcasematerialdesigntwotone{rice-bowl}{1555}
\showcasematerialdesigntwotone{ring-volume}{1556}
\showcasematerialdesigntwotone{rocket}{1557}
\showcasematerialdesigntwotone{rocket-launch}{1558}
\showcasematerialdesigntwotone{roller-shades}{1559}
\showcasematerialdesigntwotone{roller-shades-closed}{1560}
\showcasematerialdesigntwotone{roller-skating}{1561}
\showcasematerialdesigntwotone{roofing}{1562}
\showcasematerialdesigntwotone{room}{1563}
\showcasematerialdesigntwotone{room-preferences}{1564}
\showcasematerialdesigntwotone{room-service}{1565}
\showcasematerialdesigntwotone{rotate-90-degrees-ccw}{1566}
\showcasematerialdesigntwotone{rotate-90-degrees-cw}{1567}
\showcasematerialdesigntwotone{rotate-left}{1568}
\showcasematerialdesigntwotone{rotate-right}{1569}
\showcasematerialdesigntwotone{roundabout-left}{1570}
\showcasematerialdesigntwotone{roundabout-right}{1571}
\showcasematerialdesigntwotone{rounded-corner}{1572}
\showcasematerialdesigntwotone{route}{1573}
\showcasematerialdesigntwotone{router}{1574}
\showcasematerialdesigntwotone{rowing}{1575}
\showcasematerialdesigntwotone{rss-feed}{1576}
\showcasematerialdesigntwotone{rsvp}{1577}
\showcasematerialdesigntwotone{rtt}{1578}
\showcasematerialdesigntwotone{rule}{1579}
\showcasematerialdesigntwotone{rule-folder}{1580}
\showcasematerialdesigntwotone{running-with-errors}{1581}
\showcasematerialdesigntwotone{run-circle}{1582}
\showcasematerialdesigntwotone{rv-hookup}{1583}
\showcasematerialdesigntwotone{r-mobiledata}{1584}
\showcasematerialdesigntwotone{safety-check}{1585}
\showcasematerialdesigntwotone{safety-divider}{1586}
\showcasematerialdesigntwotone{sailing}{1587}
\showcasematerialdesigntwotone{sanitizer}{1588}
\showcasematerialdesigntwotone{satellite}{1589}
\showcasematerialdesigntwotone{satellite-alt}{1590}
\showcasematerialdesigntwotone{save}{1591}
\showcasematerialdesigntwotone{saved-search}{1592}
\showcasematerialdesigntwotone{save-alt}{1593}
\showcasematerialdesigntwotone{save-as}{1594}
\showcasematerialdesigntwotone{savings}{1595}
\showcasematerialdesigntwotone{scale}{1596}
\showcasematerialdesigntwotone{scanner}{1597}
\showcasematerialdesigntwotone{scatter-plot}{1598}
\showcasematerialdesigntwotone{schedule}{1599}
\showcasematerialdesigntwotone{schedule-send}{1600}
\showcasematerialdesigntwotone{schema}{1601}
\showcasematerialdesigntwotone{school}{1602}
\showcasematerialdesigntwotone{science}{1603}
\showcasematerialdesigntwotone{score}{1604}
\showcasematerialdesigntwotone{scoreboard}{1605}
\showcasematerialdesigntwotone{screenshot}{1606}
\showcasematerialdesigntwotone{screenshot-monitor}{1607}
\showcasematerialdesigntwotone{screen-lock-landscape}{1608}
\showcasematerialdesigntwotone{screen-lock-portrait}{1609}
\showcasematerialdesigntwotone{screen-lock-rotation}{1610}
\showcasematerialdesigntwotone{screen-rotation}{1611}
\showcasematerialdesigntwotone{screen-rotation-alt}{1612}
\showcasematerialdesigntwotone{screen-search-desktop}{1613}
\showcasematerialdesigntwotone{screen-share}{1614}
\showcasematerialdesigntwotone{scuba-diving}{1615}
\showcasematerialdesigntwotone{sd}{1616}
\showcasematerialdesigntwotone{sd-card}{1617}
\showcasematerialdesigntwotone{sd-card-alert}{1618}
\showcasematerialdesigntwotone{sd-storage}{1619}
\showcasematerialdesigntwotone{search}{1620}
\showcasematerialdesigntwotone{search-off}{1621}
\showcasematerialdesigntwotone{security}{1622}
\showcasematerialdesigntwotone{security-update}{1623}
\showcasematerialdesigntwotone{security-update-good}{1624}
\showcasematerialdesigntwotone{security-update-warning}{1625}
\showcasematerialdesigntwotone{segment}{1626}
\showcasematerialdesigntwotone{select-all}{1627}
\showcasematerialdesigntwotone{self-improvement}{1628}
\showcasematerialdesigntwotone{sell}{1629}
\showcasematerialdesigntwotone{send}{1630}
\showcasematerialdesigntwotone{send-and-archive}{1631}
\showcasematerialdesigntwotone{send-time-extension}{1632}
\showcasematerialdesigntwotone{send-to-mobile}{1633}
\showcasematerialdesigntwotone{sensors}{1634}
\showcasematerialdesigntwotone{sensors-off}{1635}
\showcasematerialdesigntwotone{sensor-door}{1636}
\showcasematerialdesigntwotone{sensor-occupied}{1637}
\showcasematerialdesigntwotone{sensor-window}{1638}
\showcasematerialdesigntwotone{sentiment-dissatisfied}{1639}
\showcasematerialdesigntwotone{sentiment-neutral}{1640}
\showcasematerialdesigntwotone{sentiment-satisfied}{1641}
\showcasematerialdesigntwotone{sentiment-satisfied-alt}{1642}
\showcasematerialdesigntwotone{sentiment-very-dissatisfied}{1643}
\showcasematerialdesigntwotone{sentiment-very-satisfied}{1644}
\showcasematerialdesigntwotone{settings}{1645}
\showcasematerialdesigntwotone{settings-accessibility}{1646}
\showcasematerialdesigntwotone{settings-applications}{1647}
\showcasematerialdesigntwotone{settings-backup-restore}{1648}
\showcasematerialdesigntwotone{settings-bluetooth}{1649}
\showcasematerialdesigntwotone{settings-brightness}{1650}
\showcasematerialdesigntwotone{settings-cell}{1651}
\showcasematerialdesigntwotone{settings-ethernet}{1652}
\showcasematerialdesigntwotone{settings-input-antenna}{1653}
\showcasematerialdesigntwotone{settings-input-component}{1654}
\showcasematerialdesigntwotone{settings-input-composite}{1655}
\showcasematerialdesigntwotone{settings-input-hdmi}{1656}
\showcasematerialdesigntwotone{settings-input-svideo}{1657}
\showcasematerialdesigntwotone{settings-overscan}{1658}
\showcasematerialdesigntwotone{settings-phone}{1659}
\showcasematerialdesigntwotone{settings-power}{1660}
\showcasematerialdesigntwotone{settings-remote}{1661}
\showcasematerialdesigntwotone{settings-suggest}{1662}
\showcasematerialdesigntwotone{settings-system-daydream}{1663}
\showcasematerialdesigntwotone{settings-voice}{1664}
\showcasematerialdesigntwotone{set-meal}{1665}
\showcasematerialdesigntwotone{severe-cold}{1666}
\showcasematerialdesigntwotone{shape-line}{1667}
\showcasematerialdesigntwotone{share}{1668}
\showcasematerialdesigntwotone{share-location}{1669}
\showcasematerialdesigntwotone{shield}{1670}
\showcasematerialdesigntwotone{shield-moon}{1671}
\showcasematerialdesigntwotone{shop}{1672}
\showcasematerialdesigntwotone{shopping-bag}{1673}
\showcasematerialdesigntwotone{shopping-basket}{1674}
\showcasematerialdesigntwotone{shopping-cart}{1675}
\showcasematerialdesigntwotone{shopping-cart-checkout}{1676}
\showcasematerialdesigntwotone{shop-2}{1677}
\showcasematerialdesigntwotone{shop-two}{1678}
\showcasematerialdesigntwotone{shortcut}{1679}
\showcasematerialdesigntwotone{short-text}{1680}
\showcasematerialdesigntwotone{shower}{1681}
\showcasematerialdesigntwotone{show-chart}{1682}
\showcasematerialdesigntwotone{shuffle}{1683}
\showcasematerialdesigntwotone{shuffle-on}{1684}
\showcasematerialdesigntwotone{shutter-speed}{1685}
\showcasematerialdesigntwotone{sick}{1686}
\showcasematerialdesigntwotone{signal-cellular-0-bar}{1687}
\showcasematerialdesigntwotone{signal-cellular-4-bar}{1688}
\showcasematerialdesigntwotone{signal-cellular-alt}{1689}
\showcasematerialdesigntwotone{signal-cellular-alt-1-bar}{1690}
\showcasematerialdesigntwotone{signal-cellular-alt-2-bar}{1691}
\showcasematerialdesigntwotone{signal-cellular-connected-no-internet-0-bar}{1692}
\showcasematerialdesigntwotone{signal-cellular-connected-no-internet-4-bar}{1693}
\showcasematerialdesigntwotone{signal-cellular-nodata}{1694}
\showcasematerialdesigntwotone{signal-cellular-no-sim}{1695}
\showcasematerialdesigntwotone{signal-cellular-null}{1696}
\showcasematerialdesigntwotone{signal-cellular-off}{1697}
\showcasematerialdesigntwotone{signal-wifi-0-bar}{1698}
\showcasematerialdesigntwotone{signal-wifi-4-bar}{1699}
\showcasematerialdesigntwotone{signal-wifi-4-bar-lock}{1700}
\showcasematerialdesigntwotone{signal-wifi-bad}{1701}
\showcasematerialdesigntwotone{signal-wifi-connected-no-internet-4}{1702}
\showcasematerialdesigntwotone{signal-wifi-off}{1703}
\showcasematerialdesigntwotone{signal-wifi-statusbar-4-bar}{1704}
\showcasematerialdesigntwotone{signal-wifi-statusbar-connected-no-internet-4}{1705}
\showcasematerialdesigntwotone{signal-wifi-statusbar-null}{1706}
\showcasematerialdesigntwotone{signpost}{1707}
\showcasematerialdesigntwotone{sign-language}{1708}
\showcasematerialdesigntwotone{sim-card}{1709}
\showcasematerialdesigntwotone{sim-card-alert}{1710}
\showcasematerialdesigntwotone{sim-card-download}{1711}
\showcasematerialdesigntwotone{single-bed}{1712}
\showcasematerialdesigntwotone{sip}{1713}
\showcasematerialdesigntwotone{skateboarding}{1714}
\showcasematerialdesigntwotone{skip-next}{1715}
\showcasematerialdesigntwotone{skip-previous}{1716}
\showcasematerialdesigntwotone{sledding}{1717}
\showcasematerialdesigntwotone{slideshow}{1718}
\showcasematerialdesigntwotone{slow-motion-video}{1719}
\showcasematerialdesigntwotone{smartphone}{1720}
\showcasematerialdesigntwotone{smart-button}{1721}
\showcasematerialdesigntwotone{smart-display}{1722}
\showcasematerialdesigntwotone{smart-screen}{1723}
\showcasematerialdesigntwotone{smart-toy}{1724}
\showcasematerialdesigntwotone{smoke-free}{1725}
\showcasematerialdesigntwotone{smoking-rooms}{1726}
\showcasematerialdesigntwotone{sms}{1727}
\showcasematerialdesigntwotone{sms-failed}{1728}
\showcasematerialdesigntwotone{snippet-folder}{1729}
\showcasematerialdesigntwotone{snooze}{1730}
\showcasematerialdesigntwotone{snowboarding}{1731}
\showcasematerialdesigntwotone{snowmobile}{1732}
\showcasematerialdesigntwotone{snowshoeing}{1733}
\showcasematerialdesigntwotone{soap}{1734}
\showcasematerialdesigntwotone{social-distance}{1735}
\showcasematerialdesigntwotone{solar-power}{1736}
\showcasematerialdesigntwotone{sort}{1737}
\showcasematerialdesigntwotone{sort-by-alpha}{1738}
\showcasematerialdesigntwotone{sos}{1739}
\showcasematerialdesigntwotone{soup-kitchen}{1740}
\showcasematerialdesigntwotone{source}{1741}
\showcasematerialdesigntwotone{south}{1742}
\showcasematerialdesigntwotone{south-america}{1743}
\showcasematerialdesigntwotone{south-east}{1744}
\showcasematerialdesigntwotone{south-west}{1745}
\showcasematerialdesigntwotone{spa}{1746}
\showcasematerialdesigntwotone{space-bar}{1747}
\showcasematerialdesigntwotone{space-dashboard}{1748}
\showcasematerialdesigntwotone{spatial-audio}{1749}
\showcasematerialdesigntwotone{spatial-audio-off}{1750}
\showcasematerialdesigntwotone{spatial-tracking}{1751}
\showcasematerialdesigntwotone{speaker}{1752}
\showcasematerialdesigntwotone{speaker-group}{1753}
\showcasematerialdesigntwotone{speaker-notes}{1754}
\showcasematerialdesigntwotone{speaker-notes-off}{1755}
\showcasematerialdesigntwotone{speaker-phone}{1756}
\showcasematerialdesigntwotone{speed}{1757}
\showcasematerialdesigntwotone{spellcheck}{1758}
\showcasematerialdesigntwotone{splitscreen}{1759}
\showcasematerialdesigntwotone{spoke}{1760}
\showcasematerialdesigntwotone{sports}{1761}
\showcasematerialdesigntwotone{sports-bar}{1762}
\showcasematerialdesigntwotone{sports-baseball}{1763}
\showcasematerialdesigntwotone{sports-basketball}{1764}
\showcasematerialdesigntwotone{sports-cricket}{1765}
\showcasematerialdesigntwotone{sports-esports}{1766}
\showcasematerialdesigntwotone{sports-football}{1767}
\showcasematerialdesigntwotone{sports-golf}{1768}
\showcasematerialdesigntwotone{sports-gymnastics}{1769}
\showcasematerialdesigntwotone{sports-handball}{1770}
\showcasematerialdesigntwotone{sports-hockey}{1771}
\showcasematerialdesigntwotone{sports-kabaddi}{1772}
\showcasematerialdesigntwotone{sports-martial-arts}{1773}
\showcasematerialdesigntwotone{sports-mma}{1774}
\showcasematerialdesigntwotone{sports-motorsports}{1775}
\showcasematerialdesigntwotone{sports-rugby}{1776}
\showcasematerialdesigntwotone{sports-score}{1777}
\showcasematerialdesigntwotone{sports-soccer}{1778}
\showcasematerialdesigntwotone{sports-tennis}{1779}
\showcasematerialdesigntwotone{sports-volleyball}{1780}
\showcasematerialdesigntwotone{square}{1781}
\showcasematerialdesigntwotone{square-foot}{1782}
\showcasematerialdesigntwotone{ssid-chart}{1783}
\showcasematerialdesigntwotone{stacked-bar-chart}{1784}
\showcasematerialdesigntwotone{stacked-line-chart}{1785}
\showcasematerialdesigntwotone{stadium}{1786}
\showcasematerialdesigntwotone{stairs}{1787}
\showcasematerialdesigntwotone{star}{1788}
\showcasematerialdesigntwotone{stars}{1789}
\showcasematerialdesigntwotone{start}{1790}
\showcasematerialdesigntwotone{star-border}{1791}
\showcasematerialdesigntwotone{star-border-purple500}{1792}
\showcasematerialdesigntwotone{star-half}{1793}
\showcasematerialdesigntwotone{star-outline}{1794}
\showcasematerialdesigntwotone{star-purple500}{1795}
\showcasematerialdesigntwotone{star-rate}{1796}
\showcasematerialdesigntwotone{stay-current-landscape}{1797}
\showcasematerialdesigntwotone{stay-current-portrait}{1798}
\showcasematerialdesigntwotone{stay-primary-landscape}{1799}
\showcasematerialdesigntwotone{stay-primary-portrait}{1800}
\showcasematerialdesigntwotone{sticky-note-2}{1801}
\showcasematerialdesigntwotone{stop}{1802}
\showcasematerialdesigntwotone{stop-circle}{1803}
\showcasematerialdesigntwotone{stop-screen-share}{1804}
\showcasematerialdesigntwotone{storage}{1805}
\showcasematerialdesigntwotone{store}{1806}
\showcasematerialdesigntwotone{storefront}{1807}
\showcasematerialdesigntwotone{store-mall-directory}{1808}
\showcasematerialdesigntwotone{storm}{1809}
\showcasematerialdesigntwotone{straight}{1810}
\showcasematerialdesigntwotone{straighten}{1811}
\showcasematerialdesigntwotone{stream}{1812}
\showcasematerialdesigntwotone{streetview}{1813}
\showcasematerialdesigntwotone{strikethrough-s}{1814}
\showcasematerialdesigntwotone{stroller}{1815}
\showcasematerialdesigntwotone{style}{1816}
\showcasematerialdesigntwotone{subdirectory-arrow-left}{1817}
\showcasematerialdesigntwotone{subdirectory-arrow-right}{1818}
\showcasematerialdesigntwotone{subject}{1819}
\showcasematerialdesigntwotone{subscript}{1820}
\showcasematerialdesigntwotone{subscriptions}{1821}
\showcasematerialdesigntwotone{subtitles}{1822}
\showcasematerialdesigntwotone{subtitles-off}{1823}
\showcasematerialdesigntwotone{subway}{1824}
\showcasematerialdesigntwotone{summarize}{1825}
\showcasematerialdesigntwotone{superscript}{1826}
\showcasematerialdesigntwotone{supervised-user-circle}{1827}
\showcasematerialdesigntwotone{supervisor-account}{1828}
\showcasematerialdesigntwotone{support}{1829}
\showcasematerialdesigntwotone{support-agent}{1830}
\showcasematerialdesigntwotone{surfing}{1831}
\showcasematerialdesigntwotone{surround-sound}{1832}
\showcasematerialdesigntwotone{swap-calls}{1833}
\showcasematerialdesigntwotone{swap-horiz}{1834}
\showcasematerialdesigntwotone{swap-horizontal-circle}{1835}
\showcasematerialdesigntwotone{swap-vert}{1836}
\showcasematerialdesigntwotone{swap-vertical-circle}{1837}
\showcasematerialdesigntwotone{swipe}{1838}
\showcasematerialdesigntwotone{swipe-down}{1839}
\showcasematerialdesigntwotone{swipe-down-alt}{1840}
\showcasematerialdesigntwotone{swipe-left}{1841}
\showcasematerialdesigntwotone{swipe-left-alt}{1842}
\showcasematerialdesigntwotone{swipe-right}{1843}
\showcasematerialdesigntwotone{swipe-right-alt}{1844}
\showcasematerialdesigntwotone{swipe-up}{1845}
\showcasematerialdesigntwotone{swipe-up-alt}{1846}
\showcasematerialdesigntwotone{swipe-vertical}{1847}
\showcasematerialdesigntwotone{switch-access-shortcut}{1848}
\showcasematerialdesigntwotone{switch-access-shortcut-add}{1849}
\showcasematerialdesigntwotone{switch-account}{1850}
\showcasematerialdesigntwotone{switch-camera}{1851}
\showcasematerialdesigntwotone{switch-left}{1852}
\showcasematerialdesigntwotone{switch-right}{1853}
\showcasematerialdesigntwotone{switch-video}{1854}
\showcasematerialdesigntwotone{synagogue}{1855}
\showcasematerialdesigntwotone{sync}{1856}
\showcasematerialdesigntwotone{sync-alt}{1857}
\showcasematerialdesigntwotone{sync-disabled}{1858}
\showcasematerialdesigntwotone{sync-lock}{1859}
\showcasematerialdesigntwotone{sync-problem}{1860}
\showcasematerialdesigntwotone{system-security-update}{1861}
\showcasematerialdesigntwotone{system-security-update-good}{1862}
\showcasematerialdesigntwotone{system-security-update-warning}{1863}
\showcasematerialdesigntwotone{system-update}{1864}
\showcasematerialdesigntwotone{system-update-alt}{1865}
\showcasematerialdesigntwotone{tab}{1866}
\showcasematerialdesigntwotone{tablet}{1867}
\showcasematerialdesigntwotone{tablet-android}{1868}
\showcasematerialdesigntwotone{tablet-mac}{1869}
\showcasematerialdesigntwotone{table-bar}{1870}
\showcasematerialdesigntwotone{table-chart}{1871}
\showcasematerialdesigntwotone{table-restaurant}{1872}
\showcasematerialdesigntwotone{table-rows}{1873}
\showcasematerialdesigntwotone{table-view}{1874}
\showcasematerialdesigntwotone{tab-unselected}{1875}
\showcasematerialdesigntwotone{tag}{1876}
\showcasematerialdesigntwotone{tag-faces}{1877}
\showcasematerialdesigntwotone{takeout-dining}{1878}
\showcasematerialdesigntwotone{tapas}{1879}
\showcasematerialdesigntwotone{tap-and-play}{1880}
\showcasematerialdesigntwotone{task}{1881}
\showcasematerialdesigntwotone{task-alt}{1882}
\showcasematerialdesigntwotone{taxi-alert}{1883}
\showcasematerialdesigntwotone{temple-buddhist}{1884}
\showcasematerialdesigntwotone{temple-hindu}{1885}
\showcasematerialdesigntwotone{terminal}{1886}
\showcasematerialdesigntwotone{terrain}{1887}
\showcasematerialdesigntwotone{textsms}{1888}
\showcasematerialdesigntwotone{texture}{1889}
\showcasematerialdesigntwotone{text-decrease}{1890}
\showcasematerialdesigntwotone{text-fields}{1891}
\showcasematerialdesigntwotone{text-format}{1892}
\showcasematerialdesigntwotone{text-increase}{1893}
\showcasematerialdesigntwotone{text-rotate-up}{1894}
\showcasematerialdesigntwotone{text-rotate-vertical}{1895}
\showcasematerialdesigntwotone{text-rotation-angledown}{1896}
\showcasematerialdesigntwotone{text-rotation-angleup}{1897}
\showcasematerialdesigntwotone{text-rotation-down}{1898}
\showcasematerialdesigntwotone{text-rotation-none}{1899}
\showcasematerialdesigntwotone{text-snippet}{1900}
\showcasematerialdesigntwotone{theaters}{1901}
\showcasematerialdesigntwotone{theater-comedy}{1902}
\showcasematerialdesigntwotone{thermostat}{1903}
\showcasematerialdesigntwotone{thermostat-auto}{1904}
\showcasematerialdesigntwotone{thumbs-up-down}{1905}
\showcasematerialdesigntwotone{thumb-down}{1906}
\showcasematerialdesigntwotone{thumb-down-alt}{1907}
\showcasematerialdesigntwotone{thumb-down-off-alt}{1908}
\showcasematerialdesigntwotone{thumb-up}{1909}
\showcasematerialdesigntwotone{thumb-up-alt}{1910}
\showcasematerialdesigntwotone{thumb-up-off-alt}{1911}
\showcasematerialdesigntwotone{thunderstorm}{1912}
\showcasematerialdesigntwotone{timelapse}{1913}
\showcasematerialdesigntwotone{timeline}{1914}
\showcasematerialdesigntwotone{timer}{1915}
\showcasematerialdesigntwotone{timer-3}{1916}
\showcasematerialdesigntwotone{timer-3-select}{1917}
\showcasematerialdesigntwotone{timer-10}{1918}
\showcasematerialdesigntwotone{timer-10-select}{1919}
\showcasematerialdesigntwotone{timer-off}{1920}
\showcasematerialdesigntwotone{time-to-leave}{1921}
\showcasematerialdesigntwotone{tips-and-updates}{1922}
\showcasematerialdesigntwotone{tire-repair}{1923}
\showcasematerialdesigntwotone{title}{1924}
\showcasematerialdesigntwotone{toc}{1925}
\showcasematerialdesigntwotone{today}{1926}
\showcasematerialdesigntwotone{toggle-off}{1927}
\showcasematerialdesigntwotone{toggle-on}{1928}
\showcasematerialdesigntwotone{token}{1929}
\showcasematerialdesigntwotone{toll}{1930}
\showcasematerialdesigntwotone{tonality}{1931}
\showcasematerialdesigntwotone{topic}{1932}
\showcasematerialdesigntwotone{tornado}{1933}
\showcasematerialdesigntwotone{touch-app}{1934}
\showcasematerialdesigntwotone{tour}{1935}
\showcasematerialdesigntwotone{toys}{1936}
\showcasematerialdesigntwotone{track-changes}{1937}
\showcasematerialdesigntwotone{traffic}{1938}
\showcasematerialdesigntwotone{train}{1939}
\showcasematerialdesigntwotone{tram}{1940}
\showcasematerialdesigntwotone{transcribe}{1941}
\showcasematerialdesigntwotone{transfer-within-a-station}{1942}
\showcasematerialdesigntwotone{transform}{1943}
\showcasematerialdesigntwotone{transgender}{1944}
\showcasematerialdesigntwotone{transit-enterexit}{1945}
\showcasematerialdesigntwotone{translate}{1946}
\showcasematerialdesigntwotone{travel-explore}{1947}
\showcasematerialdesigntwotone{trending-down}{1948}
\showcasematerialdesigntwotone{trending-flat}{1949}
\showcasematerialdesigntwotone{trending-up}{1950}
\showcasematerialdesigntwotone{trip-origin}{1951}
\showcasematerialdesigntwotone{troubleshoot}{1952}
\showcasematerialdesigntwotone{try}{1953}
\showcasematerialdesigntwotone{tsunami}{1954}
\showcasematerialdesigntwotone{tty}{1955}
\showcasematerialdesigntwotone{tune}{1956}
\showcasematerialdesigntwotone{tungsten}{1957}
\showcasematerialdesigntwotone{turned-in}{1958}
\showcasematerialdesigntwotone{turned-in-not}{1959}
\showcasematerialdesigntwotone{turn-left}{1960}
\showcasematerialdesigntwotone{turn-right}{1961}
\showcasematerialdesigntwotone{turn-sharp-left}{1962}
\showcasematerialdesigntwotone{turn-sharp-right}{1963}
\showcasematerialdesigntwotone{turn-slight-left}{1964}
\showcasematerialdesigntwotone{turn-slight-right}{1965}
\showcasematerialdesigntwotone{tv}{1966}
\showcasematerialdesigntwotone{tv-off}{1967}
\showcasematerialdesigntwotone{two-wheeler}{1968}
\showcasematerialdesigntwotone{type-specimen}{1969}
\showcasematerialdesigntwotone{umbrella}{1970}
\showcasematerialdesigntwotone{unarchive}{1971}
\showcasematerialdesigntwotone{undo}{1972}
\showcasematerialdesigntwotone{unfold-less}{1973}
\showcasematerialdesigntwotone{unfold-less-double}{1974}
\showcasematerialdesigntwotone{unfold-more}{1975}
\showcasematerialdesigntwotone{unfold-more-double}{1976}
\showcasematerialdesigntwotone{unpublished}{1977}
\showcasematerialdesigntwotone{unsubscribe}{1978}
\showcasematerialdesigntwotone{upcoming}{1979}
\showcasematerialdesigntwotone{update}{1980}
\showcasematerialdesigntwotone{update-disabled}{1981}
\showcasematerialdesigntwotone{upgrade}{1982}
\showcasematerialdesigntwotone{upload}{1983}
\showcasematerialdesigntwotone{upload-file}{1984}
\showcasematerialdesigntwotone{usb}{1985}
\showcasematerialdesigntwotone{usb-off}{1986}
\showcasematerialdesigntwotone{u-turn-left}{1987}
\showcasematerialdesigntwotone{u-turn-right}{1988}
\showcasematerialdesigntwotone{vaccines}{1989}
\showcasematerialdesigntwotone{vape-free}{1990}
\showcasematerialdesigntwotone{vaping-rooms}{1991}
\showcasematerialdesigntwotone{verified}{1992}
\showcasematerialdesigntwotone{verified-user}{1993}
\showcasematerialdesigntwotone{vertical-align-bottom}{1994}
\showcasematerialdesigntwotone{vertical-align-center}{1995}
\showcasematerialdesigntwotone{vertical-align-top}{1996}
\showcasematerialdesigntwotone{vertical-distribute}{1997}
\showcasematerialdesigntwotone{vertical-shades}{1998}
\showcasematerialdesigntwotone{vertical-shades-closed}{1999}
\showcasematerialdesigntwotone{vertical-split}{2000}
\showcasematerialdesigntwotone{vibration}{2001}
\showcasematerialdesigntwotone{videocam}{2002}
\showcasematerialdesigntwotone{videocam-off}{2003}
\showcasematerialdesigntwotone{videogame-asset}{2004}
\showcasematerialdesigntwotone{videogame-asset-off}{2005}
\showcasematerialdesigntwotone{video-call}{2006}
\showcasematerialdesigntwotone{video-camera-back}{2007}
\showcasematerialdesigntwotone{video-camera-front}{2008}
\showcasematerialdesigntwotone{video-chat}{2009}
\showcasematerialdesigntwotone{video-file}{2010}
\showcasematerialdesigntwotone{video-label}{2011}
\showcasematerialdesigntwotone{video-library}{2012}
\showcasematerialdesigntwotone{video-settings}{2013}
\showcasematerialdesigntwotone{video-stable}{2014}
\showcasematerialdesigntwotone{view-agenda}{2015}
\showcasematerialdesigntwotone{view-array}{2016}
\showcasematerialdesigntwotone{view-carousel}{2017}
\showcasematerialdesigntwotone{view-column}{2018}
\showcasematerialdesigntwotone{view-comfy}{2019}
\showcasematerialdesigntwotone{view-comfy-alt}{2020}
\showcasematerialdesigntwotone{view-compact}{2021}
\showcasematerialdesigntwotone{view-compact-alt}{2022}
\showcasematerialdesigntwotone{view-cozy}{2023}
\showcasematerialdesigntwotone{view-day}{2024}
\showcasematerialdesigntwotone{view-headline}{2025}
\showcasematerialdesigntwotone{view-in-ar}{2026}
\showcasematerialdesigntwotone{view-kanban}{2027}
\showcasematerialdesigntwotone{view-list}{2028}
\showcasematerialdesigntwotone{view-module}{2029}
\showcasematerialdesigntwotone{view-quilt}{2030}
\showcasematerialdesigntwotone{view-sidebar}{2031}
\showcasematerialdesigntwotone{view-stream}{2032}
\showcasematerialdesigntwotone{view-timeline}{2033}
\showcasematerialdesigntwotone{view-week}{2034}
\showcasematerialdesigntwotone{vignette}{2035}
\showcasematerialdesigntwotone{villa}{2036}
\showcasematerialdesigntwotone{visibility}{2037}
\showcasematerialdesigntwotone{visibility-off}{2038}
\showcasematerialdesigntwotone{voicemail}{2039}
\showcasematerialdesigntwotone{voice-chat}{2040}
\showcasematerialdesigntwotone{voice-over-off}{2041}
\showcasematerialdesigntwotone{volcano}{2042}
\showcasematerialdesigntwotone{volume-down}{2043}
\showcasematerialdesigntwotone{volume-mute}{2044}
\showcasematerialdesigntwotone{volume-off}{2045}
\showcasematerialdesigntwotone{volume-up}{2046}
\showcasematerialdesigntwotone{volunteer-activism}{2047}
\showcasematerialdesigntwotone{vpn-key}{2048}
\showcasematerialdesigntwotone{vpn-key-off}{2049}
\showcasematerialdesigntwotone{vpn-lock}{2050}
\showcasematerialdesigntwotone{vrpano}{2051}
\showcasematerialdesigntwotone{wallet}{2052}
\showcasematerialdesigntwotone{wallpaper}{2053}
\showcasematerialdesigntwotone{warehouse}{2054}
\showcasematerialdesigntwotone{warning}{2055}
\showcasematerialdesigntwotone{warning-amber}{2056}
\showcasematerialdesigntwotone{wash}{2057}
\showcasematerialdesigntwotone{watch}{2058}
\showcasematerialdesigntwotone{watch-later}{2059}
\showcasematerialdesigntwotone{watch-off}{2060}
\showcasematerialdesigntwotone{water}{2061}
\showcasematerialdesigntwotone{waterfall-chart}{2062}
\showcasematerialdesigntwotone{water-damage}{2063}
\showcasematerialdesigntwotone{water-drop}{2064}
\showcasematerialdesigntwotone{waves}{2065}
\showcasematerialdesigntwotone{waving-hand}{2066}
\showcasematerialdesigntwotone{wb-auto}{2067}
\showcasematerialdesigntwotone{wb-cloudy}{2068}
\showcasematerialdesigntwotone{wb-incandescent}{2069}
\showcasematerialdesigntwotone{wb-iridescent}{2070}
\showcasematerialdesigntwotone{wb-shade}{2071}
\showcasematerialdesigntwotone{wb-sunny}{2072}
\showcasematerialdesigntwotone{wb-twilight}{2073}
\showcasematerialdesigntwotone{wc}{2074}
\showcasematerialdesigntwotone{web}{2075}
\showcasematerialdesigntwotone{webhook}{2076}
\showcasematerialdesigntwotone{web-asset}{2077}
\showcasematerialdesigntwotone{web-asset-off}{2078}
\showcasematerialdesigntwotone{web-stories}{2079}
\showcasematerialdesigntwotone{weekend}{2080}
\showcasematerialdesigntwotone{west}{2081}
\showcasematerialdesigntwotone{whatshot}{2082}
\showcasematerialdesigntwotone{wheelchair-pickup}{2083}
\showcasematerialdesigntwotone{where-to-vote}{2084}
\showcasematerialdesigntwotone{widgets}{2085}
\showcasematerialdesigntwotone{width-full}{2086}
\showcasematerialdesigntwotone{width-normal}{2087}
\showcasematerialdesigntwotone{width-wide}{2088}
\showcasematerialdesigntwotone{wifi}{2089}
\showcasematerialdesigntwotone{wifi-1-bar}{2090}
\showcasematerialdesigntwotone{wifi-2-bar}{2091}
\showcasematerialdesigntwotone{wifi-calling}{2092}
\showcasematerialdesigntwotone{wifi-calling-3}{2093}
\showcasematerialdesigntwotone{wifi-channel}{2094}
\showcasematerialdesigntwotone{wifi-find}{2095}
\showcasematerialdesigntwotone{wifi-lock}{2096}
\showcasematerialdesigntwotone{wifi-off}{2097}
\showcasematerialdesigntwotone{wifi-password}{2098}
\showcasematerialdesigntwotone{wifi-protected-setup}{2099}
\showcasematerialdesigntwotone{wifi-tethering}{2100}
\showcasematerialdesigntwotone{wifi-tethering-error}{2101}
\showcasematerialdesigntwotone{wifi-tethering-off}{2102}
\showcasematerialdesigntwotone{window}{2103}
\showcasematerialdesigntwotone{wind-power}{2104}
\showcasematerialdesigntwotone{wine-bar}{2105}
\showcasematerialdesigntwotone{woman}{2106}
\showcasematerialdesigntwotone{woman-2}{2107}
\showcasematerialdesigntwotone{work}{2108}
\showcasematerialdesigntwotone{workspaces}{2109}
\showcasematerialdesigntwotone{workspace-premium}{2110}
\showcasematerialdesigntwotone{work-history}{2111}
\showcasematerialdesigntwotone{work-off}{2112}
\showcasematerialdesigntwotone{work-outline}{2113}
\showcasematerialdesigntwotone{wrap-text}{2114}
\showcasematerialdesigntwotone{wrong-location}{2115}
\showcasematerialdesigntwotone{wysiwyg}{2116}
\showcasematerialdesigntwotone{yard}{2117}
\showcasematerialdesigntwotone{youtube-searched-for}{2118}
\showcasematerialdesigntwotone{zoom-in}{2119}
\showcasematerialdesigntwotone{zoom-in-map}{2120}
\showcasematerialdesigntwotone{zoom-out}{2121}
\showcasematerialdesigntwotone{zoom-out-map}{2122}
\end{materialdesigncase}

\end{document}